\section{The \'etale plus construction of the relative Picard functor}
\label{pa-chap:PicardEtale}

For a curve \(C\) over \(k\) and an affine test scheme \(\Spec A\), the \'etale-sheafified
relative Picard group is computed by a single \emph{plus construction}: a class is represented
by a relative Picard class on an \'etale cover of \(\Spec A\) satisfying a descent condition,
and two representatives are identified when they agree on a common refinement. This chapter
constructs that plus, on affine test schemes indexed by a \(k\)-algebra \(A\); we recall from
elsewhere in this project the ambient functor it sheafifies. For an object \(T\) of
\(\Over(\Spec k)\) write \(\mathrm{relPic}(C,T)\) for the relative Picard group of \(C\) at
\(T\) --- the quotient of the Picard group of \(C \otimes T\) by the subgroup of classes pulled
back from \(T\) along the second projection --- a commutative group, with restriction
homomorphisms \(g^{*} \colon \mathrm{relPic}(C,T) \to \mathrm{relPic}(C,T')\) for a morphism
\(g \colon T' \to T\) of \(\Over(\Spec k)\), satisfying \((\id_T)^{*} = \id\) and
\((g \circ h)^{*} = h^{*} \circ g^{*}\). For a commutative \(k\)-algebra \(A\), write \(\Spec A\)
also for the object of \(\Over(\Spec k)\) it determines via the structure morphism
\(\Spec A \to \Spec k\) induced by the \(k\)-algebra structure of \(A\); a homomorphism of
\(k\)-algebras \(A \to A'\) induces, via \(\Spec\), a morphism \(\Spec A' \to \Spec A\) of
\(\Over(\Spec k)\). We abbreviate \(\mathrm{relPic}(C,A) := \mathrm{relPic}(C,\Spec A)\).

The plus construction is the Grothendieck--Kleiman construction introduced by
Grothendieck \cite{grothendieck-fga} and described in detail by Kleiman
\cite{kleiman-picard}.

\section{The relative Picard group as a quotient}

The group \(\mathrm{relPic}(C,T)\) of the preamble is realised concretely as a quotient of
the \v Cech Picard group of the product: this section records the quotient presentation and
its exact coset relation, the calculus every later rigidification and degree statement
cascades through.

\begin{definition}[The relative Picard group]
  \label{pa-def:relPic}
  \lean{AlgebraicGeometry.picFromBase, AlgebraicGeometry.mem_picFromBase_iff,
    AlgebraicGeometry.relPic, AlgebraicGeometry.relPicMk,
    AlgebraicGeometry.relPicMk_surjective, AlgebraicGeometry.relPicMk_eq_relPicMk_iff}
  \leanok
  \source{kleiman-picard}

  \dcref{custom}
  Let \(C\) be an object of \(\Over(\Spec k)\) and \(T\) a test object. Write
  \(H_T \le \Pic(C \otimes T)\) for the subgroup of \v Cech Picard classes \emph{pulled back
  from the base}: the range of \(\mathrm{pr}_2^{*} = \mathrm{snd}^{*} \colon \Pic(T) \to
  \Pic(C \otimes T)\), so \(L \in H_T\) iff \(L = \mathrm{snd}^{*} N\) for some
  \(N \in \Pic(T)\). The \emph{relative Picard group} is the quotient
  \[
    \mathrm{relPic}(C,T) \;:=\; \Pic(C \otimes T) \,/\, H_T ,
  \]
  a commutative group, with surjective projection \(\pi_T \colon \Pic(C \otimes T) \to
  \mathrm{relPic}(C,T)\). The \emph{exact coset relation} holds: \(\pi_T L = \pi_T L'\) iff
  \(L / L' \in H_T\) — a genuine range condition, involving no choice.
\end{definition}
\begin{proof}
  \leanok
  \(H_T\) is the range of a group homomorphism, hence a subgroup, and \(\Pic(C \otimes T)\)
  is commutative, so \(H_T\) is normal and the quotient is a commutative group with a
  surjective projection homomorphism. The coset relation is the defining property of a
  quotient group: two elements have equal images exactly when their ratio lies in the
  subgroup.
\end{proof}

\begin{definition}[Restriction of relative Picard classes]
  \label{pa-def:relPicMap}
  \lean{AlgebraicGeometry.picFromBase_le_comap, AlgebraicGeometry.relPicMap,
    AlgebraicGeometry.relPicMap_mk, AlgebraicGeometry.relPicMap_id,
    AlgebraicGeometry.relPicMap_comp, AlgebraicGeometry.relPicFunctor}
  \leanok
  \dcref{custom}

  \uses{pa-def:relPic}
  For a morphism \(g \colon T' \to T\) of test objects, pullback of \v Cech classes along
  \(C \lhd g \colon C \otimes T' \to C \otimes T\) descends to the quotients: a homomorphism
  \(g^{*} \colon \mathrm{relPic}(C,T) \to \mathrm{relPic}(C,T')\) with
  \(g^{*}(\pi_T L) = \pi_{T'}\bigl((C \lhd g)^{*} L\bigr)\). It satisfies
  \((\id_T)^{*} = \id\) and \((h \circ g)^{*} = g'^{*}\)-composition contravariantly, so
  \(T \mapsto \mathrm{relPic}(C,T)\) is a presheaf of commutative groups on
  \(\Over(\Spec k)\) — the relative Picard functor of Grothendieck--Kleiman.
\end{definition}
\begin{proof}
  \leanok
  Well-definedness on cosets: for \(L = \mathrm{snd}_T^{*} N \in H_T\),
  \[
    (C \lhd g)^{*}\,\mathrm{snd}_T^{*} N
      = \bigl(\mathrm{snd}_T \circ (C \lhd g)\bigr)^{*} N
      = (g \circ \mathrm{snd}_{T'})^{*} N
      = \mathrm{snd}_{T'}^{*}(g^{*} N) \in H_{T'},
  \]
  using \(\mathrm{snd}_T \circ (C \lhd g) = g \circ \mathrm{snd}_{T'}\) (the whiskering
  identity of the cartesian product). So the composite
  \(\Pic(C \otimes T) \to \mathrm{relPic}(C,T')\) kills \(H_T\) and factors through the
  quotient. The functor laws follow from
  \(C \lhd \id = \id\) and \(C \lhd (h \circ g) = (C \lhd h) \circ (C \lhd g)\) together
  with functoriality of pullback of \v Cech classes, checked on representatives via the
  computation rule.
\end{proof}

\begin{lemma}[Naturality in the curve]
  \label{pa-lem:relPicMapCurve}
  \lean{AlgebraicGeometry.picFromBase_le_comap_whiskerRight,
    AlgebraicGeometry.relPicMapCurveApp, AlgebraicGeometry.relPicMapCurveApp_mk,
    AlgebraicGeometry.relPicMapCurve, AlgebraicGeometry.relPicMapCurve_app_mk,
    AlgebraicGeometry.relPicMapCurve_id, AlgebraicGeometry.relPicMapCurve_comp,
    AlgebraicGeometry.relPicMapCurveApp_id, AlgebraicGeometry.relPicMapCurveApp_comp,
    AlgebraicGeometry.relPicMap_relPicMapCurveApp,
    AlgebraicGeometry.relPicAlgMap_relPicMapCurveApp}
  \leanok
  \dcref{custom}

  \uses{pa-def:relPic, pa-def:relPicMap, pa-def:relPicAlgMap}
  A morphism of curves \(g \colon C' \to C\) over \(\Spec k\) induces, by pullback along
  \(g \rhd T\) on representatives, a natural transformation of relative Picard functors
  \(\mathrm{relPic}(C,-) \to \mathrm{relPic}(C',-)\); the assignment is contravariantly
  functorial in \(g\) on the nose. In particular, at every test object the transports
  satisfy the functor laws, and they commute with restriction along any test morphism ---
  in particular along \(\Spec f\) for a \(k\)-algebra map \(f\) of affine tests.
\end{lemma}
\begin{proof}
  \leanok
  Well-definedness on cosets: \((g \rhd T)^{*}\,\mathrm{snd}_{C,T}^{*} N =
  \mathrm{snd}_{C',T}^{*} N\) since \(\mathrm{snd}_{C,T} \circ (g \rhd T) =
  \mathrm{snd}_{C',T}\). Naturality against a test morphism \(t \colon T' \to T\) is the
  exchange identity \((C' \lhd t) \circ (g \rhd T') = (g \rhd T) \circ (C \lhd t)\) of
  whiskerings, applied to representatives. The functor laws follow from
  \(\id \rhd T = \id\) and \((g' \circ g) \rhd T = (g \rhd T) \circ (g' \rhd T)\).
\end{proof}

\begin{lemma}[Constant curve morphisms kill relative Picard classes]
  \label{pa-lem:relPicMapCurveConstant}
  \lean{AlgebraicGeometry.cechPicMap_toUnit_whiskerRight_mem_picFromBase,
    AlgebraicGeometry.relPicMapCurveApp_toUnit_comp}
  \leanok
  \dcref{custom}

  \uses{pa-def:relPic, pa-lem:relPicMapCurve}
  Let \(g \colon C' \to C\) be a morphism of curve bundles over \(\Spec k\) that factors
  through the terminal object \(\Spec k\) of \(\Over(\Spec k)\). Then for every test
  object \(T\) the induced map \(\mathrm{relPic}(C,T) \to \mathrm{relPic}(C',T)\) is
  trivial.
\end{lemma}
\begin{proof}
  \leanok
  Write \(g = e \circ u\) with \(u \colon C' \to \Spec k\) the unique morphism to the
  terminal object and \(e \colon \Spec k \to C\). Pullback along \(g \rhd T\) then factors
  through pullback along \(u \rhd T \colon C' \otimes T \to \Spec k \otimes T\), so it
  suffices to show that every Picard class on \(\Spec k \otimes T\) pulls back along
  \(u \rhd T\) into the subgroup \(H_T\) of classes pulled back from \(T\). The second
  projection \(\mathrm{snd} \colon \Spec k \otimes T \to T\) is the left unitor, hence an
  isomorphism, so an arbitrary class \(X\) on \(\Spec k \otimes T\) is
  \(\mathrm{snd}^{*} N\) for \(N := (\mathrm{snd}^{-1})^{*} X\). Since
  \(\mathrm{snd} \circ (u \rhd T) = \mathrm{snd}_{C',T}\), we get
  \((u \rhd T)^{*} X = (u \rhd T)^{*}\,\mathrm{snd}^{*} N = \mathrm{snd}_{C',T}^{*} N \in
  H_T\). Every pulled-back class therefore lies in \(H_T\) and dies in the quotient
  \(\mathrm{relPic}(C',T)\).
\end{proof}

At the \v Cech level the constant-morphism collapse is sharper than
Lemma~\ref{pa-lem:relPicMapCurveConstant}: a morphism of bundles whose set-range is a
single closed point kills every \v Cech Picard class outright --- before passing to any
quotient --- because it factors through the spectrum of a residue field, whose \v Cech
Picard group is trivial.

\begin{theorem}[Morphisms factoring through a residue field kill \v Cech classes]
  \label{pa-thm:cechPicMap_eq_one_of_factorization}
  \lean{AlgebraicGeometry.cechPicMap_eq_one_of_factorization}
  \leanok
  \dcref{custom}

  Let \(K\) be a field, \(D, E\) objects of \(\Over(\Spec K)\), and \(h : D \to E\) a
  morphism of \(\Over(\Spec K)\) whose underlying morphism factors as
  \(D \to \Spec \kappa(x) \to E\) through the canonical morphism from the spectrum of
  the residue field of a point \(x\) of \(E\). Then for every \v Cech Picard class
  \(L \in \Pic(E)\),
  \[
    h^{*}L \;=\; 1 .
  \]
\end{theorem}
\begin{proof}
  \leanok
  Write the factorization as \(h = \iota \circ q\) with
  \(q : D \to \Spec \kappa(x)\) and \(\iota : \Spec \kappa(x) \to E\). Pullback of
  \v Cech Picard classes is functorial, so \(h^{*}L = q^{*}(\iota^{*}L)\). The
  underlying space of \(\Spec \kappa(x)\) is a single point (the spectrum of a field),
  and the \v Cech Picard group of a scheme whose space has at most one point is
  trivial: any two members of a pointed cover coincide as opens --- each is determined
  by whether it is empty, and members contain their base points --- so every unit
  \v Cech \(1\)-cocycle lives on a cover by a single open, where it is a coboundary.
  Hence \(\iota^{*}L = 1\), and \(q^{*}\), a homomorphism, gives
  \(h^{*}L = q^{*}(1) = 1\).
\end{proof}

\begin{corollary}[Morphisms with one-point range kill \v Cech classes]
  \label{pa-cor:cechPicMap_eq_one_of_range_eq_singleton}
  \lean{AlgebraicGeometry.cechPicMap_eq_one_of_range_eq_singleton}
  \leanok
  \dcref{custom}

  \uses{pa-thm:cechPicMap_eq_one_of_factorization, pa-thm:hom_factors_fromSpecResidueField}
  Let \(K\) be a field, \(D, E\) objects of \(\Over(\Spec K)\) with \(D\) reduced, and
  \(h : D \to E\) a morphism of \(\Over(\Spec K)\) whose underlying morphism is
  quasi-compact and whose set-range is a single closed point \(x\) of \(E\). Then
  \(h^{*}L = 1\) for every \v Cech Picard class \(L \in \Pic(E)\).
\end{corollary}
\begin{proof}
  \leanok
  By the factorization through the residue field of a closed point
  (\ref{pa-thm:hom_factors_fromSpecResidueField}), the underlying morphism of \(h\)
  factors as \(D \to \Spec \kappa(x) \to E\). Apply
  \ref{pa-thm:cechPicMap_eq_one_of_factorization}.
\end{proof}

\section{Presented \'etale covers of an affine scheme}

Throughout this section \(A\) is a commutative ring (no field or \(k\)-algebra structure is
needed here). An \'etale cover of \(\Spec A\) in the sense consumed by the plus construction is
a faithfully flat \'etale \(A\)-algebra; to keep the index of covers of a fixed base
small --- which is what makes the plus construction below a legitimate colimit, rather than one
indexed by a proper class of isomorphism types --- we realize such algebras through an explicit
finite-variable presentation, at the cost of carrying a witness of the presentation as data.

\begin{definition}[Presented \'etale cover]
  \label{pa-def:etale_cover}
  \lean{Algebra.EtaleCover, Algebra.EtaleCover.Carrier}
  \leanok
  \dcref{custom}

  A \emph{presented \'etale cover} \(E\) of \(A\) consists of a natural number \(n\), an ideal
  \(I\) of the polynomial ring \(A[x_1,\dots,x_n]\), a proof that the quotient algebra
  \(B_E := A[x_1,\dots,x_n]/I\) is \'etale over \(A\), and a proof that the induced map of
  spectra \(\Spec B_E \to \Spec A\) is surjective. We call \(B_E\) the \emph{carrier} of \(E\).
  Since \'etale algebras are flat, surjectivity of \(\Spec B_E \to \Spec A\) together with
  flatness makes \(B_E\) faithfully flat over \(A\).
\end{definition}

\begin{definition}[Presented covers from an \'etale algebra]
  \label{pa-def:etale_cover_of}
  \lean{Algebra.EtaleCover.of}
  \leanok
  \dcref{custom}

  \uses{pa-def:etale_cover}
  Let \(B\) be an \'etale \(A\)-algebra such that \(\Spec B \to \Spec A\) is surjective. Since
  \'etale algebras are of finite presentation, \(B\) admits a presentation as a quotient of a
  polynomial ring in finitely many variables over \(A\); choosing one such presentation
  produces a presented \'etale cover \(E\) of \(A\), denoted \(E = \mathrm{of}(B)\).
\end{definition}

\begin{lemma}[Carrier of a presented cover from an algebra]
  \label{pa-lem:etale_cover_ofEquiv}
  \lean{Algebra.EtaleCover.ofEquiv}
  \leanok
  \dcref{custom}

  \uses{pa-def:etale_cover_of}
  In the situation of Definition~\ref{pa-def:etale_cover_of}, the carrier \(B_E\) of
  \(E = \mathrm{of}(B)\) is isomorphic to \(B\) as an \(A\)-algebra.
\end{lemma}
\begin{proof}
  \leanok
  Immediate from the construction: \(E\) is built from a chosen presentation
  \(A[x_1,\dots,x_n] \twoheadrightarrow B\), and the carrier of the resulting cover is the
  quotient of \(A[x_1,\dots,x_n]\) by the kernel of that presentation, which is isomorphic to
  \(B\) by the first isomorphism theorem for algebras.
\end{proof}

\begin{definition}[Refinement]
  \label{pa-def:etale_cover_refines}
  \lean{Algebra.EtaleCover.Refines}
  \leanok
  \dcref{custom}

  \uses{pa-def:etale_cover}
  For presented \'etale covers \(E, E'\) of \(A\), we say \(E'\) \emph{refines} \(E\) if there
  exists a homomorphism of \(A\)-algebras \(B_E \to B_{E'}\) (equivalently, a morphism
  \(\Spec B_{E'} \to \Spec B_E\) over \(\Spec A\)). This is a mere existence statement: no
  particular refinement map is remembered as data.
\end{definition}

\begin{lemma}[Refinement is a preorder]
  \label{pa-lem:etale_cover_refines_preorder}
  \lean{Algebra.EtaleCover.refines_refl, Algebra.EtaleCover.Refines.trans}
  \leanok
  \dcref{custom}

  \uses{pa-def:etale_cover_refines}
  Every presented \'etale cover \(E\) of \(A\) refines itself, and if \(E''\) refines \(E'\)
  and \(E'\) refines \(E\) then \(E''\) refines \(E\).
\end{lemma}
\begin{proof}
  \leanok
  Reflexivity: the identity homomorphism of \(B_E\) witnesses that \(E\) refines itself.
  Transitivity: composing a homomorphism \(B_{E'} \to B_{E''}\) witnessing that \(E''\) refines
  \(E'\) with a homomorphism \(B_E \to B_{E'}\) witnessing that \(E'\) refines \(E\) gives a
  homomorphism \(B_E \to B_{E''}\) witnessing that \(E''\) refines \(E\).
\end{proof}

\begin{definition}[Common refinement]
  \label{pa-def:etale_cover_prod}
  \lean{Algebra.EtaleCover.prod, Algebra.EtaleCover.prodEquiv}
  \leanok
  \dcref{custom}

  \uses{pa-def:etale_cover_of}
  For presented \'etale covers \(E, E'\) of \(A\), the \emph{common refinement} \(E \times E'\)
  is the presented \'etale cover of \(A\) with carrier \(B_E \otimes_A B_{E'}\) (\'etale over
  \(A\) because it is \'etale over the \'etale algebra \(B_E\), which is \'etale over \(A\); and
  faithfully flat by transitivity of faithfully flat extensions), obtained by applying
  Definition~\ref{pa-def:etale_cover_of} to \(B_E \otimes_A B_{E'}\).
\end{definition}

\begin{definition}[Coprojections of the common refinement]
  \label{pa-def:etale_cover_prod_coproj}
  \lean{Algebra.EtaleCover.prodInl, Algebra.EtaleCover.prodInr}
  \leanok
  \dcref{custom}

  \uses{pa-def:etale_cover_prod}
  The two canonical homomorphisms of \(A\)-algebras
  \(B_E \to B_{E \times E'} \cong B_E \otimes_A B_{E'}\), \(b \mapsto b \otimes 1\), and
  \(B_{E'} \to B_{E \times E'}\), \(b' \mapsto 1 \otimes b'\), transported along the carrier
  identification of Definition~\ref{pa-def:etale_cover_prod}.
\end{definition}

\begin{lemma}[The common refinement refines both factors]
  \label{pa-lem:etale_cover_prod_refines}
  \lean{Algebra.EtaleCover.prod_refines_left, Algebra.EtaleCover.prod_refines_right}
  \leanok
  \dcref{custom}

  \uses{pa-def:etale_cover_prod_coproj, pa-def:etale_cover_refines}
  \(E \times E'\) refines both \(E\) and \(E'\).
\end{lemma}
\begin{proof}
  \leanok
  The two coprojections of Definition~\ref{pa-def:etale_cover_prod_coproj} are homomorphisms of
  \(A\)-algebras out of \(B_E\), respectively \(B_{E'}\), witnessing the two refinements.
\end{proof}

\begin{definition}[Universal map out of the common refinement]
  \label{pa-def:etale_cover_prodLift}
  \lean{Algebra.EtaleCover.prodLift}
  \leanok
  \dcref{custom}

  \uses{pa-def:etale_cover_prod}
  Given homomorphisms of \(A\)-algebras \(f \colon B_E \to B_F\) and \(g \colon B_{E'} \to
  B_F\), the induced homomorphism of \(A\)-algebras \(B_{E \times E'} \to B_F\), obtained from
  the universal property of the tensor product \(B_E \otimes_A B_{E'}\) applied to \(f\) and
  \(g\) (whose images commute in the commutative ring \(B_F\)), transported along the carrier
  identification of Definition~\ref{pa-def:etale_cover_prod}.
\end{definition}

\begin{lemma}[The universal map is a common section]
  \label{pa-lem:etale_cover_prodLift_comp}
  \lean{Algebra.EtaleCover.prodLift_comp_prodInl, Algebra.EtaleCover.prodLift_comp_prodInr}
  \leanok
  \dcref{custom}

  \uses{pa-def:etale_cover_prodLift, pa-def:etale_cover_prod_coproj}
  In the situation of Definition~\ref{pa-def:etale_cover_prodLift}, the induced map composed with
  each coprojection recovers the corresponding map: the composite \(B_E \to B_{E \times E'} \to
  B_F\) equals \(f\), and the composite \(B_{E'} \to B_{E \times E'} \to B_F\) equals \(g\).
\end{lemma}
\begin{proof}
  \leanok
  This is the defining property of the map induced by the universal property of the tensor
  product: on elementary tensors, \(b \otimes 1 \mapsto f(b) g(1) = f(b)\) and
  \(1 \otimes b' \mapsto f(1) g(b') = g(b')\).
\end{proof}

\begin{theorem}[Directedness]
  \label{pa-thm:etale_cover_exists_refines_prod}
  \lean{Algebra.EtaleCover.exists_refines_prod}
  \leanok
  \dcref{custom}

  \uses{pa-def:etale_cover_refines}
  Any two presented \'etale covers \(E, E'\) of \(A\) admit a common refinement: there is a
  presented \'etale cover \(E''\) of \(A\) which refines both \(E\) and \(E'\).
\end{theorem}
\begin{proof}
  \leanok
  \uses{pa-lem:etale_cover_prod_refines}
  Take \(E'' = E \times E'\); it refines both factors by Lemma~\ref{pa-lem:etale_cover_prod_refines}.
\end{proof}

\begin{definition}[The trivial cover]
  \label{pa-def:etale_cover_self}
  \lean{Algebra.EtaleCover.self, Algebra.EtaleCover.selfEquiv}
  \leanok
  \dcref{custom}

  \uses{pa-def:etale_cover_of}
  The presented \'etale cover of \(A\) by itself, obtained by applying
  Definition~\ref{pa-def:etale_cover_of} to \(B = A\) (with structure morphism \(\Spec A \to
  \Spec A\) the identity, trivially surjective). Its carrier is isomorphic to \(A\).
\end{definition}

\begin{definition}[Canonical map from the trivial cover]
  \label{pa-def:etale_cover_fromSelf}
  \lean{Algebra.EtaleCover.fromSelf}
  \leanok
  \dcref{custom}

  \uses{pa-def:etale_cover_self}
  For a presented \'etale cover \(E\) of \(A\), the canonical homomorphism of \(A\)-algebras
  from the carrier of the trivial cover to \(B_E\): under the identification of the carrier of
  the trivial cover with \(A\) (Definition~\ref{pa-def:etale_cover_self}), the structure
  homomorphism \(A \to B_E\).
\end{definition}

\begin{lemma}[Every cover refines the trivial cover]
  \label{pa-lem:etale_cover_refines_self}
  \lean{Algebra.EtaleCover.refines_self}
  \leanok
  \dcref{custom}

  \uses{pa-def:etale_cover_fromSelf, pa-def:etale_cover_refines}
  Every presented \'etale cover \(E\) of \(A\) refines the trivial cover.
\end{lemma}
\begin{proof}
  \leanok
  The canonical map of Definition~\ref{pa-def:etale_cover_fromSelf} witnesses the refinement.
\end{proof}

\begin{definition}[Base change of a cover]
  \label{pa-def:etale_cover_baseChange}
  \lean{Algebra.EtaleCover.baseChange, Algebra.EtaleCover.baseChangeEquiv}
  \leanok
  \dcref{custom}

  \uses{pa-def:etale_cover_of}
  Let \(A'\) be a commutative \(A\)-algebra and \(E\) a presented \'etale cover of \(A\). The
  \emph{base change} \(E_{A'}\) is the presented \'etale cover of \(A'\) with carrier
  \(A' \otimes_A B_E\) (\'etale over \(A'\) and faithfully flat over \(A'\), both by base
  change of the corresponding properties of \(B_E\) over \(A\)), obtained by applying
  Definition~\ref{pa-def:etale_cover_of} to \(A' \otimes_A B_E\).
\end{definition}

\begin{definition}[Canonical inclusion into the base change]
  \label{pa-def:etale_cover_baseChangeInclude}
  \lean{Algebra.EtaleCover.baseChangeInclude}
  \leanok
  \dcref{custom}

  \uses{pa-def:etale_cover_baseChange}
  The canonical homomorphism of \(A\)-algebras \(\kappa_E \colon B_E \to B_{E_{A'}}\), which
  under the carrier identification \(B_{E_{A'}} \cong A' \otimes_A B_E\) of
  Definition~\ref{pa-def:etale_cover_baseChange} is \(b \mapsto 1 \otimes b\).
\end{definition}

\begin{definition}[Base change of a refinement map]
  \label{pa-def:etale_cover_baseChangeMap}
  \lean{Algebra.EtaleCover.baseChangeMap}
  \leanok
  \dcref{custom}

  \uses{pa-def:etale_cover_baseChange}
  For presented \'etale covers \(E, F\) of \(A\) and a homomorphism of \(A\)-algebras
  \(h \colon B_E \to B_F\), the induced homomorphism of \(A'\)-algebras \(B_{E_{A'}} \to
  B_{F_{A'}}\) which, under the carrier identifications of
  Definition~\ref{pa-def:etale_cover_baseChange}, is \(\id_{A'} \otimes h\).
\end{definition}

\begin{lemma}[Naturality of the canonical inclusion]
  \label{pa-lem:etale_cover_baseChangeMap_comp_baseChangeInclude}
  \lean{Algebra.EtaleCover.baseChangeMap_comp_baseChangeInclude}
  \leanok
  \dcref{custom}

  \uses{pa-def:etale_cover_baseChangeMap, pa-def:etale_cover_baseChangeInclude}
  For \(h \colon B_E \to B_F\) as in Definition~\ref{pa-def:etale_cover_baseChangeMap}, the square
  of \(A\)-algebra homomorphisms commutes: the composite \(B_E \to B_{E_{A'}} \to B_{F_{A'}}\)
  (canonical inclusion followed by the base-changed map) equals the composite
  \(B_E \to B_F \to B_{F_{A'}}\) (\(h\) followed by the canonical inclusion).
\end{lemma}
\begin{proof}
  \leanok
  Both composites send \(b \in B_E\) to \(1 \otimes h(b) \in A' \otimes_A B_F\) under the
  carrier identifications: the first computes \(b \mapsto 1 \otimes b \mapsto
  (\id_{A'}\otimes h)(1 \otimes b) = 1 \otimes h(b)\), and the second computes
  \(b \mapsto h(b) \mapsto 1 \otimes h(b)\).
\end{proof}

\begin{lemma}[Nontrivial carrier]
  \label{pa-lem:etale_cover_nontrivial_carrier}
  \lean{Algebra.EtaleCover.nontrivial_carrier}
  \leanok
  \dcref{custom}

  \uses{pa-def:etale_cover}
  If \(A\) is a nontrivial ring, then the carrier \(B_E\) of any presented \'etale cover \(E\)
  of \(A\) is nontrivial.
\end{lemma}
\begin{proof}
  \leanok
  Since \(A\) is nontrivial, \(\Spec A\) is nonempty; since \(\Spec B_E \to \Spec A\) is
  surjective, \(\Spec B_E\) is nonempty, and a ring with nonempty spectrum is nontrivial.
\end{proof}

\begin{theorem}[Field cofinality]
  \label{pa-thm:etale_cover_field_cofinality}
  \lean{Algebra.EtaleCover.exists_finiteSeparableField_algHom}
  \leanok
  \dcref{custom}

  \uses{pa-def:etale_cover}
  Let \(K\) be a field and \(E\) a presented \'etale cover of \(K\). Then there exist a field
  \(L\), finite and separable over \(K\), and a homomorphism of \(K\)-algebras
  \(B_E \to L\).
\end{theorem}
\begin{proof}
  \leanok
  \uses{pa-lem:etale_cover_nontrivial_carrier}
  By the structure theorem for \'etale algebras over a field, \(B_E\) is isomorphic, as a
  \(K\)-algebra, to a finite product \(\prod_{i \in I} L_i\) of finite separable field
  extensions \(L_i\) of \(K\). By Lemma~\ref{pa-lem:etale_cover_nontrivial_carrier}, \(B_E\) is
  nontrivial (\(K\) being a field is nontrivial), so the product is nonempty: \(I \neq
  \emptyset\), for the empty product is the zero ring. Choose \(i_0 \in I\) and set
  \(L = L_{i_0}\); the composite of the structure isomorphism with the projection onto the
  \(i_0\)-th factor is the required homomorphism \(B_E \to L\).
\end{proof}

A finite separable field extension \(L\) of a field \(K\) is automatically \'etale over \(K\):
it is formally \'etale because it is formally unramified (separability) and formally smooth
(finiteness plus the algebra being a field, hence regular), and it is of finite presentation
because it is a finite-type algebra over a field which is moreover finite as a module.

\begin{definition}[Field extension covers]
  \label{pa-def:etale_cover_ofField}
  \lean{Algebra.EtaleCover.ofField, Algebra.EtaleCover.ofFieldEquiv}
  \leanok
  \dcref{custom}

  \uses{pa-def:etale_cover_of}
  Let \(K\) be a field and \(L\) a finite separable field extension of \(K\). Since \(\Spec K\)
  is a single point, the spectrum map \(\Spec L \to \Spec K\) is trivially surjective, so
  Definition~\ref{pa-def:etale_cover_of} applied to \(L\) produces a presented \'etale cover of
  \(K\) with carrier isomorphic to \(L\).
\end{definition}

\section{The algebra-indexed relative Picard functor}

In this section \(k\) is a field and \(C\) is an object of \(\Over(\Spec k)\).

\begin{definition}[\(\Spec\) of an algebra map, as a morphism of test objects]
  \label{pa-def:overSpecMap}
  \lean{AlgebraicGeometry.Over.overSpecMap}
  \leanok
  \dcref{custom}

  For commutative \(k\)-algebras \(A, B\) and a homomorphism of \(k\)-algebras
  \(f \colon A \to B\), the induced morphism \(\Spec B \to \Spec A\) of \(\Over(\Spec k)\):
  \(\Spec f\), which lies over \(\Spec k\) because \(f\) is a homomorphism of \(k\)-algebras.
\end{definition}

\begin{lemma}[Functoriality of \(\Spec\) on algebra maps]
  \label{pa-lem:overSpecMap_functor}
  \lean{AlgebraicGeometry.Over.overSpecMap_id, AlgebraicGeometry.Over.overSpecMap_comp}
  \leanok
  \dcref{custom}

  \uses{pa-def:overSpecMap}
  \(\Spec\) of the identity homomorphism of \(A\) is the identity morphism of \(\Spec A\); and
  for homomorphisms of \(k\)-algebras \(f \colon A \to B\) and \(g \colon B \to C'\),
  \(\Spec(g \circ f) = \Spec f \circ \Spec g\) as morphisms \(\Spec C' \to \Spec A\).
\end{lemma}
\begin{proof}
  \leanok
  \(\Spec\) is a contravariant functor from commutative rings to schemes, so it sends the
  identity to the identity and a composite \(g \circ f\) to \(\Spec f \circ \Spec g\); the
  morphisms of \(\Over(\Spec k)\) so obtained are equal to the corresponding morphisms of
  schemes on their underlying data, so the identities transport.
\end{proof}

\begin{definition}[Restriction of relative Picard classes along an algebra map]
  \label{pa-def:relPicAlgMap}
  \lean{AlgebraicGeometry.relPicAlgMap}
  \leanok
  \dcref{custom}

  \uses{pa-def:overSpecMap, pa-def:relPicMap}
  For a homomorphism of \(k\)-algebras \(f \colon A \to B\), the induced homomorphism of groups
  \(f^{*} \colon \mathrm{relPic}(C,A) \to \mathrm{relPic}(C,B)\): restriction of relative
  Picard classes (Definition~\ref{pa-def:relPicMap}) along the morphism
  \(\Spec f \colon \Spec B \to \Spec A\) of Definition~\ref{pa-def:overSpecMap}.
\end{definition}

\begin{lemma}[Functoriality in the algebra]
  \label{pa-lem:relPicAlgMap_functor}
  \lean{AlgebraicGeometry.relPicAlgMap_id, AlgebraicGeometry.relPicAlgMap_comp}
  \leanok
  \dcref{custom}

  \uses{pa-def:relPicAlgMap, pa-lem:overSpecMap_functor}
  For every \(k\)-algebra \(A\), \((\id_A)^{*} = \id\) on \(\mathrm{relPic}(C,A)\); and for
  homomorphisms of \(k\)-algebras \(f \colon A \to B\) and \(g \colon B \to C'\),
  \((g \circ f)^{*} = g^{*} \circ f^{*}\) on \(\mathrm{relPic}(C,A)\). Thus
  \(A \mapsto \mathrm{relPic}(C,A)\), \(f \mapsto f^{*}\), is a covariant functor from
  \(k\)-algebras to commutative groups.
\end{lemma}
\begin{proof}
  \leanok
  Restriction of relative Picard classes along \(\Spec\) of a scheme morphism satisfies the
  functor laws \(\id^{*} = \id\) and \((\varphi \circ \psi)^{*} = \psi^{*} \circ \varphi^{*}\)
  in the morphism of \(\Over(\Spec k)\). Combined with
  Lemma~\ref{pa-lem:overSpecMap_functor}, which shows \(\Spec\) of the identity algebra map is the
  identity morphism and \(\Spec\) of a composite \(g \circ f\) is \(\Spec f \circ \Spec g\), the
  two claimed identities for \(f^{*} := \Spec f{}^{*}\) follow.
\end{proof}

\section{The base-field shuffle of the relative Picard group at an affine test}
\label{pa-sec:relPicCrossBase}

In this section \(k \to L\) is a field extension, \(\sigma : \Spec L \to \Spec k\) the
induced morphism, \(C\) an object of \(\Over(\Spec k)\) with base change
\(C_L := \sigma^{*}C \in \Over(\Spec L)\) (\ref{pa-def:baseChange}), and \(A\) an algebra
in a scalar tower \(k \to L \to A\). We write \(\mathrm{relPic}_L(C_L, A)\) for the
relative Picard group of \(C_L\) at the affine test of \(A\) over \(L\), formed over
the base field \(L\) exactly as in Definition~\ref{pa-def:relPic}. The two slice products
\(C \otimes \Spec A\) (over \(k\)) and \(C_L \otimes \Spec A\) (over \(L\)) share the
carrier \(\Spec A\) of their second projections, and their underlying schemes are
identified by the same-carrier comparison \(c\) of
Definition~\ref{pa-def:crossBaseAffineIso}
(Chapter~\ref{pa-chap:Cohomology}). This section descends \(c\) to an isomorphism of the
relative Picard groups, natural in \(A\).

\begin{lemma}[The comparison transports base-pulled classes to base-pulled classes]
  \label{pa-lem:picFromBase_crossBaseAffineIso}
  \lean{AlgebraicGeometry.picFromBase_le_comap_crossBaseAffineIso,
    AlgebraicGeometry.picFromBase_le_comap_crossBaseAffineIso_inv}
  \leanok
  \dcref{custom}

  \uses{pa-def:relPic, pa-def:crossBaseAffineIso, pa-lem:crossBaseAffineIso_snd}
  Write \(\mathrm{snd}_k\), \(\mathrm{snd}_L\) for the second projections of the two
  slice products onto \(\Spec A\). For every \v Cech Picard class \(N\) on \(\Spec A\),
  \[
    c^{*}\bigl(\mathrm{snd}_k^{*}\, N\bigr) = \mathrm{snd}_L^{*}\, N,
    \qquad
    (c^{-1})^{*}\bigl(\mathrm{snd}_L^{*}\, N\bigr) = \mathrm{snd}_k^{*}\, N .
  \]
  In particular pullback along \(c\) carries the subgroup of
  \(\Pic(C \otimes \Spec A)\) of classes pulled back from the base (\ref{pa-def:relPic})
  into the corresponding subgroup of \(\Pic(C_L \otimes \Spec A)\), and pullback along
  \(c^{-1}\) carries the latter into the former.
\end{lemma}
\begin{proof}
  \leanok
  The comparison commutes with the second projections:
  \(\mathrm{snd}_k \circ c = \mathrm{snd}_L\) and
  \(\mathrm{snd}_L \circ c^{-1} = \mathrm{snd}_k\)
  (\ref{pa-lem:crossBaseAffineIso_snd}). Pullback of \v Cech Picard classes is
  contravariantly functorial in the morphism, so
  \(c^{*}(\mathrm{snd}_k^{*} N) = (\mathrm{snd}_k \circ c)^{*} N =
  \mathrm{snd}_L^{*} N\), and likewise along \(c^{-1}\).
\end{proof}

\begin{definition}[The base-field shuffle of relative Picard classes]
  \label{pa-def:relPicCrossBase}
  \lean{AlgebraicGeometry.relPicCrossBase, AlgebraicGeometry.relPicCrossBase_mk,
    AlgebraicGeometry.relPicCrossBaseInv, AlgebraicGeometry.relPicCrossBaseInv_mk}
  \leanok
  \dcref{custom}

  \uses{pa-def:relPic, pa-lem:picFromBase_crossBaseAffineIso, pa-def:crossBaseAffineIso,
    pa-def:baseChange}
  Pullback along \(c\) and along \(c^{-1}\) descend to the quotients of
  Definition~\ref{pa-def:relPic}: there are homomorphisms of groups
  \[
    \mathrm{sh}_A : \mathrm{relPic}(C,A) \longrightarrow \mathrm{relPic}_L(C_L, A),
    \qquad
    \mathrm{sh}'_A : \mathrm{relPic}_L(C_L, A) \longrightarrow \mathrm{relPic}(C,A),
  \]
  with computation rules on cosets
  \(\mathrm{sh}_A(\pi \Lambda) = \pi(c^{*} \Lambda)\) for
  \(\Lambda \in \Pic(C \otimes \Spec A)\) and
  \(\mathrm{sh}'_A(\pi M) = \pi\bigl((c^{-1})^{*} M\bigr)\) for
  \(M \in \Pic(C_L \otimes \Spec A)\).
\end{definition}
\begin{proof}
  \leanok
  Pullback of \v Cech Picard classes along a morphism of schemes is a homomorphism of
  Picard groups. By Lemma~\ref{pa-lem:picFromBase_crossBaseAffineIso} it carries the
  subgroup of base-pulled classes into the subgroup of base-pulled classes, in each
  direction. A homomorphism carrying the distinguished subgroup into the distinguished
  subgroup induces a homomorphism of the quotient groups computing on cosets by the
  stated rule --- the universal property of the quotient (\ref{pa-def:relPic}).
\end{proof}

\begin{theorem}[The base-field shuffle is an isomorphism]
  \label{pa-thm:relPicCrossBaseEquiv}
  \lean{AlgebraicGeometry.relPicCrossBaseEquiv,
    AlgebraicGeometry.relPicCrossBaseEquiv_apply,
    AlgebraicGeometry.relPicCrossBaseEquiv_symm_apply}
  \leanok
  \dcref{custom}

  \uses{pa-def:relPicCrossBase, pa-def:relPic}
  The two descended homomorphisms of Definition~\ref{pa-def:relPicCrossBase} are mutually
  inverse: \(\mathrm{sh}'_A \circ \mathrm{sh}_A = \id\) and
  \(\mathrm{sh}_A \circ \mathrm{sh}'_A = \id\). Hence the base-field shuffle is an
  isomorphism of groups
  \(\mathrm{relPic}(C,A) \cong \mathrm{relPic}_L(C_L, A)\).
\end{theorem}
\begin{proof}
  \leanok
  The projections are surjective (\ref{pa-def:relPic}), so it suffices to check on
  cosets. By the computation rules and contravariant functoriality of pullback,
  \[
    \mathrm{sh}'_A\bigl(\mathrm{sh}_A(\pi \Lambda)\bigr)
      = \pi\bigl((c^{-1})^{*}\, c^{*}\, \Lambda\bigr)
      = \pi\bigl((c \circ c^{-1})^{*}\, \Lambda\bigr)
      = \pi(\Lambda),
  \]
  and symmetrically \(\mathrm{sh}_A(\mathrm{sh}'_A(\pi M)) =
  \pi\bigl((c^{-1} \circ c)^{*} M\bigr) = \pi(M)\). A group homomorphism with a
  two-sided inverse homomorphism is an isomorphism.
\end{proof}

\begin{theorem}[Naturality of the shuffle against restriction along algebra maps]
  \label{pa-thm:relPicCrossBase_relPicAlgMap}
  \lean{AlgebraicGeometry.relPicCrossBase_relPicAlgMap}
  \leanok
  \dcref{custom}

  \uses{pa-def:relPicCrossBase, pa-def:relPicAlgMap, pa-thm:crossBaseAffineIso_naturality,
    pa-def:relPicMap}
  Let \(B\) be a further algebra in a scalar tower \(k \to L \to B\) and
  \(f : A \to B\) a homomorphism of \(L\)-algebras (hence of \(k\)-algebras, through
  the towers). Then restriction and shuffle commute: for every
  \(x \in \mathrm{relPic}(C,A)\),
  \[
    \mathrm{sh}_B\bigl(f^{*}_{k}\, x\bigr) \;=\; f^{*}_{L}\bigl(\mathrm{sh}_A(x)\bigr),
  \]
  where \(f^{*}_{k} : \mathrm{relPic}(C,A) \to \mathrm{relPic}(C,B)\) is restriction
  along \(\Spec f\) over \(k\) (\ref{pa-def:relPicAlgMap}) and
  \(f^{*}_{L} : \mathrm{relPic}_L(C_L,A) \to \mathrm{relPic}_L(C_L,B)\) is the same
  construction over the base field \(L\), at the curve \(C_L\).
\end{theorem}
\begin{proof}
  \leanok
  Check on a coset \(x = \pi(\Lambda)\), \(\Lambda \in \Pic(C \otimes \Spec A)\), the
  projection being surjective (\ref{pa-def:relPic}). By the computation rules of
  restriction (\ref{pa-def:relPicMap}, \ref{pa-def:relPicAlgMap}) and of the shuffle
  (\ref{pa-def:relPicCrossBase}), the left side is
  \(\pi\bigl(c_B^{*}\,\bigl((C \lhd \Spec_k f)_{\mathrm{left}}\bigr)^{*}\Lambda\bigr)\),
  the pullback of \(\Lambda\) along the composite
  \((C \lhd \Spec_k f)_{\mathrm{left}} \circ c_B\), and the right side is
  \(\pi\bigl(\bigl((C_L \lhd \Spec_L f)_{\mathrm{left}}\bigr)^{*} c_A^{*}\,
  \Lambda\bigr)\), the pullback of \(\Lambda\) along
  \(c_A \circ (C_L \lhd \Spec_L f)_{\mathrm{left}}\). The two composites are equal ---
  the naturality square of the same-carrier comparison
  (\ref{pa-thm:crossBaseAffineIso_naturality}) --- so the two pullbacks agree, and so do
  their cosets.
\end{proof}

\section{Descent classes on an \'etale cover}

Throughout this section \(k\) is a field, \(C\) is an object of \(\Over(\Spec k)\), and \(A\) is
a commutative \(k\)-algebra. For a presented \'etale cover \(E\) of \(A\)
(Definition~\ref{pa-def:etale_cover}) with carrier \(B_E\), the self-tensor product
\(B_E \otimes_A B_E\) is the carrier of the common refinement \(E \times E\) of \(E\) with
itself (Definition~\ref{pa-def:etale_cover_prod}): it is the algebra of the \emph{double cover}
\(\Spec(B_E \otimes_A B_E) \to \Spec B_E\), whose two projections to \(\Spec B_E\) record the
overlap data of the cover \(E\).

\begin{definition}[Coprojections of the double cover]
  \label{pa-def:doubleInl_doubleInr}
  \lean{AlgebraicGeometry.doubleInl, AlgebraicGeometry.doubleInr}
  \leanok
  \dcref{custom}

  \uses{pa-def:etale_cover}
  For a presented \'etale cover \(E\) of \(A\), the two homomorphisms of \(k\)-algebras
  \(B_E \to B_E \otimes_A B_E\), \(b \mapsto b \otimes 1\) and \(b \mapsto 1 \otimes b\) ---
  homomorphisms of \(A\)-algebras, in particular of \(k\)-algebras since \(A\) is a
  \(k\)-algebra.
\end{definition}

\begin{definition}[Descent classes]
  \label{pa-def:descentClasses}
  \lean{AlgebraicGeometry.descentClasses, AlgebraicGeometry.mem_descentClasses_iff}
  \leanok
  \source{kleiman-picard}

  \uses{pa-def:doubleInl_doubleInr, pa-def:relPicAlgMap}
  \dcref{custom}
  For a presented \'etale cover \(E\) of \(A\), the subgroup
  \(\mathrm{Desc}(C,E) \subseteq \mathrm{relPic}(C,B_E)\) of classes \(\lambda\) satisfying the
  \emph{descent condition}: the two pullbacks of \(\lambda\) along the coprojections of the
  double cover (Definition~\ref{pa-def:doubleInl_doubleInr}) agree,
  \[
    \mathrm{Desc}(C,E) := \bigl\{\, \lambda \in \mathrm{relPic}(C,B_E) : \mathrm{inl}_E^{*}
    \lambda = \mathrm{inr}_E^{*} \lambda \text{ in } \mathrm{relPic}(C, B_E \otimes_A B_E)
    \,\bigr\},
  \]
  the equalizer of the two homomorphisms \(\mathrm{inl}_E^{*}, \mathrm{inr}_E^{*} \colon
  \mathrm{relPic}(C,B_E) \to \mathrm{relPic}(C, B_E \otimes_A B_E)\) of
  Definition~\ref{pa-def:relPicAlgMap}, hence itself a subgroup.
\end{definition}

\begin{lemma}[Naturality of the descent condition]
  \label{pa-lem:relPicAlgMap_mem_descentClasses}
  \lean{AlgebraicGeometry.relPicAlgMap_mem_descentClasses}
  \leanok
  \dcref{custom}

  \uses{pa-def:descentClasses}
  Let \(E, F\) be presented \'etale covers of \(A\) and \(h \colon B_E \to B_F\) a
  homomorphism of \(A\)-algebras. If \(\lambda \in \mathrm{Desc}(C,E)\), then
  \(h^{*}\lambda \in \mathrm{Desc}(C,F)\).
\end{lemma}
\begin{proof}
  \leanok
  \uses{pa-lem:relPicAlgMap_functor}
  The homomorphism \(h \otimes h \colon B_E \otimes_A B_E \to B_F \otimes_A B_F\) intertwines
  the coprojections: \(\mathrm{inl}_F \circ h = (h \otimes h) \circ \mathrm{inl}_E\) and
  \(\mathrm{inr}_F \circ h = (h \otimes h) \circ \mathrm{inr}_E\), since both sides send
  \(b \mapsto h(b) \otimes 1\), respectively \(b \mapsto 1 \otimes h(b)\). By functoriality
  (Lemma~\ref{pa-lem:relPicAlgMap_functor}),
  \[
    \mathrm{inl}_F^{*}(h^{*}\lambda) = (h \otimes h)^{*}(\mathrm{inl}_E^{*}\lambda)
    = (h \otimes h)^{*}(\mathrm{inr}_E^{*}\lambda) = \mathrm{inr}_F^{*}(h^{*}\lambda),
  \]
  using \(\lambda \in \mathrm{Desc}(C,E)\) in the middle equality. Hence
  \(h^{*}\lambda \in \mathrm{Desc}(C,F)\).
\end{proof}

\begin{definition}[Transport of descent classes]
  \label{pa-def:descentMap}
  \lean{AlgebraicGeometry.descentMap}
  \leanok
  \dcref{custom}

  \uses{pa-lem:relPicAlgMap_mem_descentClasses}
  For \(h \colon B_E \to B_F\) a homomorphism of \(A\)-algebras between the carriers of two
  presented \'etale covers of \(A\), the induced homomorphism of groups \(h_{*} \colon
  \mathrm{Desc}(C,E) \to \mathrm{Desc}(C,F)\), \(\lambda \mapsto h^{*}\lambda\), well-defined by
  Lemma~\ref{pa-lem:relPicAlgMap_mem_descentClasses}.
\end{definition}

\begin{lemma}[Transport along the identity]
  \label{pa-lem:descentMap_id}
  \lean{AlgebraicGeometry.descentMap_id}
  \leanok
  \dcref{custom}

  \uses{pa-def:descentMap}
  For every presented \'etale cover \(E\) of \(A\), transport along the identity homomorphism
  of \(B_E\) is the identity on \(\mathrm{Desc}(C,E)\).
\end{lemma}
\begin{proof}
  \leanok
  \uses{pa-lem:relPicAlgMap_functor}
  Immediate from \((\id_{B_E})^{*} = \id\) (Lemma~\ref{pa-lem:relPicAlgMap_functor}).
\end{proof}

\begin{lemma}[Transport is functorial]
  \label{pa-lem:descentMap_comp}
  \lean{AlgebraicGeometry.descentMap_comp}
  \leanok
  \dcref{custom}

  \uses{pa-def:descentMap}
  For homomorphisms of \(A\)-algebras \(h \colon B_E \to B_F\) and \(h' \colon B_F \to B_G\)
  between carriers of presented \'etale covers of \(A\), transport along \(h' \circ h\) equals
  transport along \(h'\) followed by transport along \(h\): \((h' \circ h)_{*} = h'_{*} \circ
  h_{*}\) on \(\mathrm{Desc}(C,E)\).
\end{lemma}
\begin{proof}
  \leanok
  \uses{pa-lem:relPicAlgMap_functor}
  Immediate from \((h' \circ h)^{*} = h'^{*} \circ h^{*}\) (Lemma~\ref{pa-lem:relPicAlgMap_functor}),
  restricted to \(\mathrm{Desc}(C,E)\).
\end{proof}

\begin{theorem}[Refinement independence, the keystone]
  \label{pa-thm:relPicAlgMap_congr}
  \lean{AlgebraicGeometry.relPicAlgMap_congr}
  \leanok
  \dcref{custom}

  \uses{pa-def:descentClasses}
  Let \(E\) be a presented \'etale cover of \(A\) with carrier \(B = B_E\), let \(R\) be a
  commutative \(k\)-algebra equipped with a compatible \(A\)-algebra structure (so that
  \(k \to A \to R\) is a tower of algebras), and let \(j_1, j_2 \colon B \to R\) be two
  homomorphisms of \(A\)-algebras. If \(\lambda \in \mathrm{Desc}(C,E)\), then
  \(j_1^{*}\lambda = j_2^{*}\lambda\) in \(\mathrm{relPic}(C,R)\).
\end{theorem}
\begin{proof}
  \leanok
  \uses{pa-lem:relPicAlgMap_functor}
  Let \(H \colon B \otimes_A B \to R\) be the homomorphism of \(A\)-algebras induced by the
  universal property of the tensor product from the pair \((j_1, j_2)\) (their images commute
  in the commutative ring \(R\)): \(H(b \otimes b') = j_1(b) j_2(b')\). Then
  \(H \circ \mathrm{inl}_E = j_1\) (since \(H(b \otimes 1) = j_1(b) j_2(1) = j_1(b)\)) and
  \(H \circ \mathrm{inr}_E = j_2\) (since \(H(1 \otimes b) = j_1(1) j_2(b) = j_2(b)\)). By
  functoriality (Lemma~\ref{pa-lem:relPicAlgMap_functor}),
  \[
    j_1^{*}\lambda = H^{*}(\mathrm{inl}_E^{*}\lambda) = H^{*}(\mathrm{inr}_E^{*}\lambda)
    = j_2^{*}\lambda,
  \]
  using \(\lambda \in \mathrm{Desc}(C,E)\), i.e. \(\mathrm{inl}_E^{*}\lambda =
  \mathrm{inr}_E^{*}\lambda\), in the middle equality.
\end{proof}

\begin{theorem}[Refinement-map independence]
  \label{pa-thm:descentMap_congr}
  \lean{AlgebraicGeometry.descentMap_congr}
  \leanok
  \dcref{custom}

  \uses{pa-thm:relPicAlgMap_congr, pa-def:descentMap}
  Let \(E, F\) be presented \'etale covers of \(A\) and \(h_1, h_2 \colon B_E \to B_F\) two
  homomorphisms of \(A\)-algebras. Then \(h_{1,*} = h_{2,*}\) as maps
  \(\mathrm{Desc}(C,E) \to \mathrm{Desc}(C,F)\): transport of a descent class along \(h_1\) and
  along \(h_2\) always agree.
\end{theorem}
\begin{proof}
  \leanok
  Apply Theorem~\ref{pa-thm:relPicAlgMap_congr} with \(R = B_F\) (an \(A\)-algebra which is also a
  \(k\)-algebra compatibly, since \(F\) is a cover of the \(k\)-algebra \(A\)), \(j_1 = h_1\),
  \(j_2 = h_2\).
\end{proof}

\section{The one-step \'etale plus}

\begin{definition}[The refinement equivalence]
  \label{pa-def:picEtAffSetoid}
  \lean{AlgebraicGeometry.picEtAffSetoid}
  \leanok
  \dcref{custom}

  \uses{pa-thm:descentMap_congr, pa-def:etale_cover_prod, pa-def:etale_cover_prod_coproj}
  Consider pairs \((E,\lambda)\) with \(E\) a presented \'etale cover of \(A\) and
  \(\lambda \in \mathrm{Desc}(C,E)\). Declare \((E,\lambda) \sim (F,\mu)\) if there exist a
  presented \'etale cover \(G\) of \(A\) and homomorphisms of \(A\)-algebras
  \(f \colon B_E \to B_G\), \(g \colon B_F \to B_G\) with \(f_{*}\lambda = g_{*}\mu\) in
  \(\mathrm{Desc}(C,G)\). This is an equivalence relation.
\end{definition}
\begin{proof}
  \leanok
  \uses{pa-lem:descentMap_comp}
  \emph{Reflexivity.} \((E,\lambda) \sim (E,\lambda)\) via \(G = E\), \(f = g = \id_{B_E}\).

  \emph{Symmetry.} If \((E,\lambda) \sim (F,\mu)\) via \((G,f,g)\), then \((F,\mu) \sim
  (E,\lambda)\) via \((G,g,f)\).

  \emph{Transitivity.} Suppose \((E,\lambda) \sim (F,\mu)\) via \((G,f,g)\), so \(f_{*}\lambda =
  g_{*}\mu\), and \((F,\mu) \sim (H,\nu)\) via \((G',f',g')\), so \(f'_{*}\mu = g'_{*}\nu\). Use
  the common refinement \(G \times G'\) with the maps \(\mathrm{inl}_{G,G'} \circ f \colon B_E
  \to B_{G \times G'}\) and \(\mathrm{inr}_{G,G'} \circ g' \colon B_H \to B_{G \times G'}\).
  Using functoriality of transport (Lemma~\ref{pa-lem:descentMap_comp}) and refinement-map
  independence (Theorem~\ref{pa-thm:descentMap_congr}, applied to the two maps
  \(\mathrm{inl}_{G,G'} \circ g\) and \(\mathrm{inr}_{G,G'} \circ f'\) out of \(B_F\)),
  \[
    (\mathrm{inl}_{G,G'} \circ f)_{*}\lambda = \mathrm{inl}_{G,G',*}(f_{*}\lambda) =
    \mathrm{inl}_{G,G',*}(g_{*}\mu) = (\mathrm{inl}_{G,G'} \circ g)_{*}\mu
  \]
  \[
    = (\mathrm{inr}_{G,G'} \circ f')_{*}\mu = \mathrm{inr}_{G,G',*}(f'_{*}\mu) =
    \mathrm{inr}_{G,G',*}(g'_{*}\nu) = (\mathrm{inr}_{G,G'} \circ g')_{*}\nu,
  \]
  so \((E,\lambda) \sim (H,\nu)\) via \((G \times G', \mathrm{inl}_{G,G'} \circ f,
  \mathrm{inr}_{G,G'} \circ g')\).
\end{proof}

\begin{definition}[The one-step \'etale plus]
  \label{pa-def:PicEtAff}
  \lean{AlgebraicGeometry.PicEtAff}
  \leanok
  \source{kleiman-picard}

  \uses{pa-def:picEtAffSetoid}
  \dcref{custom}
  The set of pairs \((E,\lambda)\) as in Definition~\ref{pa-def:picEtAffSetoid}, modulo the
  equivalence relation there, is denoted \(\mathrm{PicEtAff}(C,A)\): the one-step \'etale plus
  of the relative Picard functor at \(A\). On affine test schemes this is the \'etale
  sheafification of \(\mathrm{relPic}(C,-)\), one plus construction sufficing because that
  presheaf is separated for the \'etale topology (established elsewhere).
\end{definition}

\begin{definition}[The plus class of a descent class]
  \label{pa-def:PicEtAff_mk}
  \lean{AlgebraicGeometry.PicEtAff.mk}
  \leanok
  \dcref{custom}

  \uses{pa-def:PicEtAff}
  For a presented \'etale cover \(E\) of \(A\) and \(\lambda \in \mathrm{Desc}(C,E)\), write
  \([E,\lambda] \in \mathrm{PicEtAff}(C,A)\) for the class of the pair \((E,\lambda)\).
\end{definition}

\begin{lemma}[Every plus class is so represented]
  \label{pa-lem:PicEtAff_ind}
  \lean{AlgebraicGeometry.PicEtAff.ind}
  \leanok
  \dcref{custom}

  \uses{pa-def:PicEtAff_mk}
  Every element of \(\mathrm{PicEtAff}(C,A)\) is of the form \([E,\lambda]\) for some presented
  \'etale cover \(E\) of \(A\) and \(\lambda \in \mathrm{Desc}(C,E)\).
\end{lemma}
\begin{proof}
  \leanok
  Immediate: \(\mathrm{PicEtAff}(C,A)\) is a quotient of the set of such pairs, and every class
  in a quotient is the class of some representative.
\end{proof}

\begin{theorem}[Equality of plus classes]
  \label{pa-thm:PicEtAff_mk_eq_mk_iff}
  \lean{AlgebraicGeometry.PicEtAff.mk_eq_mk_iff}
  \leanok
  \dcref{custom}

  \uses{pa-def:PicEtAff_mk, pa-def:picEtAffSetoid}
  \([E,\lambda] = [F,\mu]\) in \(\mathrm{PicEtAff}(C,A)\) if and only if there exist a presented
  \'etale cover \(G\) of \(A\) and homomorphisms of \(A\)-algebras \(f \colon B_E \to B_G\),
  \(g \colon B_F \to B_G\) with \(f_{*}\lambda = g_{*}\mu\).
\end{theorem}
\begin{proof}
  \leanok
  This is exactly the defining relation of Definition~\ref{pa-def:picEtAffSetoid}, restated as an
  equality of classes in the quotient.
\end{proof}

\begin{lemma}[Transport to a finer cover preserves the plus class]
  \label{pa-lem:PicEtAff_mk_descentMap}
  \lean{AlgebraicGeometry.PicEtAff.mk_descentMap}
  \leanok
  \dcref{custom}

  \uses{pa-thm:PicEtAff_mk_eq_mk_iff, pa-def:descentMap}
  For \(h \colon B_E \to B_F\) a homomorphism of \(A\)-algebras between carriers of presented
  \'etale covers of \(A\) and \(\lambda \in \mathrm{Desc}(C,E)\), \([F, h_{*}\lambda] =
  [E,\lambda]\).
\end{lemma}
\begin{proof}
  \leanok
  \uses{pa-lem:descentMap_id}
  Take \(G = F\) with \(f = \id_{B_F}\) and \(g = h\) in Theorem~\ref{pa-thm:PicEtAff_mk_eq_mk_iff}:
  \(f_{*}(h_{*}\lambda) = h_{*}\lambda = g_{*}\lambda\) by Lemma~\ref{pa-lem:descentMap_id}.
\end{proof}

\section{The group structure and the unit}

\begin{lemma}[Transport through the universal map collapses]
  \label{pa-lem:PicEtAff_descentMap_prodLift}
  \lean{AlgebraicGeometry.PicEtAff.descentMap_prodLift_inl,
    AlgebraicGeometry.PicEtAff.descentMap_prodLift_inr}
  \leanok
  \dcref{custom}

  \uses{pa-lem:etale_cover_prodLift_comp, pa-lem:descentMap_comp}
  Let \(E, F, H\) be presented \'etale covers of \(A\), \(f \colon B_E \to B_H\) and
  \(g \colon B_F \to B_H\) homomorphisms of \(A\)-algebras, and let
  \(\langle f,g\rangle \colon B_{E \times F} \to B_H\) be the induced map of
  Definition~\ref{pa-def:etale_cover_prodLift}. Then for \(\lambda \in \mathrm{Desc}(C,E)\) and
  \(\mu \in \mathrm{Desc}(C,F)\):
  \[
    \langle f,g\rangle_{*}(\mathrm{inl}_{E,F,*}\lambda) = f_{*}\lambda, \qquad
    \langle f,g\rangle_{*}(\mathrm{inr}_{E,F,*}\mu) = g_{*}\mu.
  \]
\end{lemma}
\begin{proof}
  \leanok
  By functoriality of transport (Lemma~\ref{pa-lem:descentMap_comp}), the left-hand sides equal
  \((\langle f,g\rangle \circ \mathrm{inl}_{E,F})_{*}\lambda\), respectively
  \((\langle f,g\rangle \circ \mathrm{inr}_{E,F})_{*}\mu\); by
  Lemma~\ref{pa-lem:etale_cover_prodLift_comp}, \(\langle f,g\rangle \circ \mathrm{inl}_{E,F} = f\)
  and \(\langle f,g\rangle \circ \mathrm{inr}_{E,F} = g\).
\end{proof}

\begin{theorem}[Multiplication]
  \label{pa-thm:PicEtAff_mk_mul_mk}
  \lean{AlgebraicGeometry.PicEtAff.instMul, AlgebraicGeometry.PicEtAff.mk_mul_mk}
  \leanok
  \dcref{custom}

  \uses{pa-lem:PicEtAff_descentMap_prodLift, pa-thm:PicEtAff_mk_eq_mk_iff, pa-thm:descentMap_congr,
    pa-def:etale_cover_prod_coproj}
  \(\mathrm{PicEtAff}(C,A)\) carries a multiplication, defined on representatives by transport
  to a common refinement product cover: for presented \'etale covers \(E, F\) of \(A\) and
  \(\lambda \in \mathrm{Desc}(C,E)\), \(\mu \in \mathrm{Desc}(C,F)\),
  \[
    [E,\lambda] \cdot [F,\mu] = \bigl[\,E \times F,\ \mathrm{inl}_{E,F,*}\lambda \cdot
    \mathrm{inr}_{E,F,*}\mu\,\bigr].
  \]
\end{theorem}
\begin{proof}
  \leanok
  It remains to check that the formula is independent of the chosen representatives, i.e. that
  if \((E,\lambda) \sim (E',\lambda')\) via \((H,f,f')\) (so \(f_{*}\lambda = f'_{*}\lambda'\))
  and \((F,\mu) \sim (F',\mu')\) via \((G,g,g')\) (so \(g_{*}\mu = g'_{*}\mu'\)), then the two
  resulting products agree. Use the common refinement \(H \times G\) of \(E \times F\) and
  \(E' \times F'\), with the maps
  \(\langle \mathrm{inl}_{H,G} \circ f,\ \mathrm{inr}_{H,G} \circ g\rangle \colon B_{E \times F}
  \to B_{H \times G}\) and
  \(\langle \mathrm{inl}_{H,G} \circ f',\ \mathrm{inr}_{H,G} \circ g'\rangle \colon
  B_{E' \times F'} \to B_{H \times G}\)
  (Definition~\ref{pa-def:etale_cover_prodLift}). Since transport is a homomorphism of groups
  (Definition~\ref{pa-def:descentMap}), transporting the product
  \(\mathrm{inl}_{E,F,*}\lambda \cdot \mathrm{inr}_{E,F,*}\mu\) along the first map yields, by
  Lemma~\ref{pa-lem:PicEtAff_descentMap_prodLift} applied twice,
  \(\mathrm{inl}_{H,G,*}(f_{*}\lambda) \cdot \mathrm{inr}_{H,G,*}(g_{*}\mu)\); transporting the
  primed product along the second map yields
  \(\mathrm{inl}_{H,G,*}(f'_{*}\lambda') \cdot \mathrm{inr}_{H,G,*}(g'_{*}\mu')\); these agree
  term by term since \(f_{*}\lambda = f'_{*}\lambda'\) and \(g_{*}\mu = g'_{*}\mu'\). Hence the
  two products transport to the same class on \(H \times G\), so by
  Theorem~\ref{pa-thm:PicEtAff_mk_eq_mk_iff} they define the same element of
  \(\mathrm{PicEtAff}(C,A)\).
\end{proof}

\begin{theorem}[Multiplication of classes on the same cover]
  \label{pa-thm:PicEtAff_mk_mul_mk_same}
  \lean{AlgebraicGeometry.PicEtAff.mk_mul_mk_same}
  \leanok
  \dcref{custom}

  \uses{pa-thm:PicEtAff_mk_mul_mk, pa-thm:descentMap_congr, pa-lem:PicEtAff_mk_descentMap}
  For a presented \'etale cover \(E\) of \(A\) and \(\lambda, \mu \in \mathrm{Desc}(C,E)\),
  \([E,\lambda] \cdot [E,\mu] = [E, \lambda \mu]\).
\end{theorem}
\begin{proof}
  \leanok
  By Theorem~\ref{pa-thm:PicEtAff_mk_mul_mk}, \([E,\lambda][E,\mu] = [E \times E,
  \mathrm{inl}_{E,E,*}\lambda \cdot \mathrm{inr}_{E,E,*}\mu]\). By refinement-map independence
  (Theorem~\ref{pa-thm:descentMap_congr}, applied to the two maps \(\mathrm{inr}_{E,E}\) and
  \(\mathrm{inl}_{E,E}\) out of \(B_E\)), \(\mathrm{inr}_{E,E,*}\mu = \mathrm{inl}_{E,E,*}\mu\).
  Since transport is a homomorphism of groups,
  \(\mathrm{inl}_{E,E,*}\lambda \cdot \mathrm{inl}_{E,E,*}\mu = \mathrm{inl}_{E,E,*}(\lambda
  \mu)\). Hence \([E,\lambda][E,\mu] = [E \times E, \mathrm{inl}_{E,E,*}(\lambda\mu)] =
  [E,\lambda\mu]\), the last equality by Lemma~\ref{pa-lem:PicEtAff_mk_descentMap}.
\end{proof}

\begin{definition}[Identity]
  \label{pa-def:PicEtAff_one}
  \lean{AlgebraicGeometry.PicEtAff.instOne, AlgebraicGeometry.PicEtAff.one_def}
  \leanok
  \dcref{custom}

  \uses{pa-def:etale_cover_self, pa-def:PicEtAff_mk}
  The identity element of \(\mathrm{PicEtAff}(C,A)\) is \(1 := [\mathrm{self}(A), 1]\), the
  class of the identity descent datum on the trivial cover.
\end{definition}

\begin{theorem}[The identity is represented on every cover]
  \label{pa-thm:PicEtAff_mk_one}
  \lean{AlgebraicGeometry.PicEtAff.mk_one}
  \leanok
  \dcref{custom}

  \uses{pa-def:PicEtAff_one, pa-def:etale_cover_fromSelf, pa-thm:PicEtAff_mk_eq_mk_iff}
  For every presented \'etale cover \(E\) of \(A\), \([E,1] = 1\) in \(\mathrm{PicEtAff}(C,A)\).
\end{theorem}
\begin{proof}
  \leanok
  Apply Theorem~\ref{pa-thm:PicEtAff_mk_eq_mk_iff} with \(G = E\), \(f = \id_{B_E}\), and \(g\)
  the canonical map from the trivial cover (Definition~\ref{pa-def:etale_cover_fromSelf}): since
  transport is a homomorphism of groups, both \(f_{*}1\) and \(g_{*}1\) equal \(1\), regardless
  of \(f\) and \(g\), so they agree.
\end{proof}

\begin{definition}[Inverse]
  \label{pa-def:PicEtAff_inv}
  \lean{AlgebraicGeometry.PicEtAff.instInv, AlgebraicGeometry.PicEtAff.mk_inv}
  \leanok
  \dcref{custom}

  \uses{pa-def:PicEtAff_mk, pa-def:picEtAffSetoid}
  \(\mathrm{PicEtAff}(C,A)\) carries an inversion, defined on representatives by
  \([E,\lambda]^{-1} = [E,\lambda^{-1}]\); well-defined because if \((E,\lambda) \sim (F,\mu)\)
  via \((G,f,g)\), so \(f_{*}\lambda = g_{*}\mu\), then applying the group homomorphism
  \(f_{*}\), respectively \(g_{*}\), to inverses gives \(f_{*}(\lambda^{-1}) =
  (f_{*}\lambda)^{-1} = (g_{*}\mu)^{-1} = g_{*}(\mu^{-1})\), so \((E,\lambda^{-1}) \sim
  (F,\mu^{-1})\) via the same \((G,f,g)\).
\end{definition}

\begin{theorem}[\(\mathrm{PicEtAff}(C,A)\) is a commutative group]
  \label{pa-thm:PicEtAff_commGroup}
  \lean{AlgebraicGeometry.PicEtAff.instCommGroup}
  \leanok
  \dcref{custom}

  \uses{pa-thm:PicEtAff_mk_mul_mk_same, pa-thm:PicEtAff_mk_one, pa-def:PicEtAff_inv,
    pa-lem:PicEtAff_ind, pa-lem:PicEtAff_mk_descentMap, pa-def:etale_cover_prod_coproj}
  With the multiplication of Theorem~\ref{pa-thm:PicEtAff_mk_mul_mk}, the identity of
  Definition~\ref{pa-def:PicEtAff_one}, and the inversion of Definition~\ref{pa-def:PicEtAff_inv},
  \(\mathrm{PicEtAff}(C,A)\) is a commutative group.
\end{theorem}
\begin{proof}
  \leanok
  By Lemma~\ref{pa-lem:PicEtAff_ind} every element is of the form \([E,\lambda]\), so it suffices
  to check the axioms on such representatives, re-expressing every class involved as a class on
  one common cover by Lemma~\ref{pa-lem:PicEtAff_mk_descentMap} (transporting a representative up
  along a chosen refinement map does not change its plus class) and then applying
  Theorem~\ref{pa-thm:PicEtAff_mk_mul_mk_same} to compute products of classes living on the same
  cover.

  \emph{Associativity.} For \(a = [E,\lambda]\), \(b = [F,\mu]\), \(c = [G,\nu]\), rewrite
  \(a\), \(b\), \(c\) as classes on the common cover \((E \times F) \times G\), transporting
  \(\lambda\) along \(\mathrm{inl}_{E \times F, G} \circ \mathrm{inl}_{E,F}\), \(\mu\) along
  \(\mathrm{inl}_{E \times F, G} \circ \mathrm{inr}_{E,F}\), and \(\nu\) along
  \(\mathrm{inr}_{E \times F, G}\) (Lemma~\ref{pa-lem:PicEtAff_mk_descentMap}); the three
  transported representatives \(\lambda', \mu', \nu' \in \mathrm{Desc}(C, (E \times F) \times
  G)\) now satisfy, by Theorem~\ref{pa-thm:PicEtAff_mk_mul_mk_same}, \((ab)c = [(E \times F)
  \times G, (\lambda'\mu')\nu']\) and \(a(bc) = [(E \times F) \times G, \lambda'(\mu'\nu')]\);
  these agree by associativity of multiplication in the group \(\mathrm{Desc}(C,(E \times F)
  \times G)\).

  \emph{Identity.} For \(a = [E,\lambda]\), \(1 \cdot a = [\mathrm{self}(A) \times E,
  \mathrm{inl}_{*}1 \cdot \mathrm{inr}_{*}\lambda] = [\mathrm{self}(A) \times E,
  \mathrm{inr}_{*}\lambda] = [E,\lambda] = a\), using that transport preserves the identity and
  Lemma~\ref{pa-lem:PicEtAff_mk_descentMap}; similarly \(a \cdot 1 = a\).

  \emph{Inverses.} For \(a = [E,\lambda]\), \(a^{-1} a = [E,\lambda^{-1}][E,\lambda] =
  [E,\lambda^{-1}\lambda] = [E,1] = 1\), by Theorem~\ref{pa-thm:PicEtAff_mk_mul_mk_same} and
  Theorem~\ref{pa-thm:PicEtAff_mk_one}.

  \emph{Commutativity.} For \(a = [E,\lambda]\), \(b = [F,\mu]\), transport both to the common
  cover \(E \times F\) via \(\mathrm{inl}_{E,F}\), respectively \(\mathrm{inr}_{E,F}\)
  (Lemma~\ref{pa-lem:PicEtAff_mk_descentMap}): \(a = [E \times F, \mathrm{inl}_{*}\lambda]\), \(b =
  [E \times F, \mathrm{inr}_{*}\mu]\); by Theorem~\ref{pa-thm:PicEtAff_mk_mul_mk_same}, \(ab = [E
  \times F, \mathrm{inl}_{*}\lambda \cdot \mathrm{inr}_{*}\mu]\) and \(ba = [E \times F,
  \mathrm{inr}_{*}\mu \cdot \mathrm{inl}_{*}\lambda]\), which agree because \(\mathrm{Desc}(C,E
  \times F)\), a subgroup of the commutative group \(\mathrm{relPic}(C, B_E \otimes_A B_F)\),
  is commutative.
\end{proof}

\begin{lemma}[Tautological descent datum]
  \label{pa-lem:PicEtAff_tautological_mem_descentClasses}
  \lean{AlgebraicGeometry.PicEtAff.tautological_mem_descentClasses}
  \leanok
  \dcref{custom}

  \uses{pa-def:descentClasses, pa-def:doubleInl_doubleInr}
  Let \(E\) be a presented \'etale cover of \(A\) and \(\iota_E \colon A \to B_E\) the
  structure homomorphism. For every \(x \in \mathrm{relPic}(C,A)\), the class
  \(\iota_E^{*}x \in \mathrm{relPic}(C,B_E)\) satisfies the descent condition:
  \(\iota_E^{*}x \in \mathrm{Desc}(C,E)\).
\end{lemma}
\begin{proof}
  \leanok
  The two composites \(\mathrm{inl}_E \circ \iota_E\) and \(\mathrm{inr}_E \circ \iota_E\), both
  maps \(A \to B_E \otimes_A B_E\), coincide: each sends \(a \mapsto \iota_E(a) \otimes 1 = 1
  \otimes \iota_E(a)\), the two descriptions of the structure homomorphism of the tensor
  product \(B_E \otimes_A B_E\) over \(A\). By functoriality (Lemma~\ref{pa-lem:relPicAlgMap_functor}),
  \(\mathrm{inl}_E^{*}(\iota_E^{*}x) = (\mathrm{inl}_E \circ \iota_E)^{*}x = (\mathrm{inr}_E
  \circ \iota_E)^{*}x = \mathrm{inr}_E^{*}(\iota_E^{*}x)\).
\end{proof}

\begin{definition}[The unit of the plus construction]
  \label{pa-def:PicEtAff_unit}
  \lean{AlgebraicGeometry.PicEtAff.unit}
  \leanok
  \source{kleiman-picard}

  \uses{pa-lem:PicEtAff_tautological_mem_descentClasses, pa-def:etale_cover_self,
    pa-thm:PicEtAff_mk_one, pa-thm:PicEtAff_mk_mul_mk_same}
  \dcref{custom}
  The homomorphism of groups \(\eta_A \colon \mathrm{relPic}(C,A) \to
  \mathrm{PicEtAff}(C,A)\) sending \(x\) to the class \([\mathrm{self}(A), \iota^{*}x]\) of \(x\)
  pulled back to the trivial cover with its tautological descent datum
  (Lemma~\ref{pa-lem:PicEtAff_tautological_mem_descentClasses}).
\end{definition}
\begin{proof}
  \leanok
  \(\eta_A\) preserves the identity: the tautological descent datum of \(1 \in
  \mathrm{relPic}(C,A)\) is \(1 \in \mathrm{Desc}(C,\mathrm{self}(A))\) (transport preserves
  the identity), so \(\eta_A(1) = [\mathrm{self}(A),1] = 1\) by
  Theorem~\ref{pa-thm:PicEtAff_mk_one}. It preserves products: the tautological descent datum of
  \(xy\) is the product of the tautological descent data of \(x\) and \(y\) (transport is a
  homomorphism), so \(\eta_A(xy) = [\mathrm{self}(A), \iota^{*}x \cdot \iota^{*}y] =
  [\mathrm{self}(A),\iota^{*}x][\mathrm{self}(A),\iota^{*}y] = \eta_A(x)\eta_A(y)\) by
  Theorem~\ref{pa-thm:PicEtAff_mk_mul_mk_same}.
\end{proof}

\section{Functoriality in the test algebra}

Throughout this section \(k\) is a field, \(C\) is an object of \(\Over(\Spec k)\), \(A\) is a
commutative \(k\)-algebra, and \(A'\) is a commutative \(k\)-algebra additionally equipped with
an \(A\)-algebra structure compatible with the tower \(k \to A \to A'\) (following the
homomorphism-along-instances idiom, this exhibits a homomorphism of \(k\)-algebras \(A \to
A'\), extracted from the tower as \(\tau \colon A \to A'\)). For a presented \'etale cover \(E\)
of \(A\), write \(E_{A'}\) for its base change (Definition~\ref{pa-def:etale_cover_baseChange}) and
\(\kappa_E \colon B_E \to B_{E_{A'}}\) for the canonical inclusion
(Definition~\ref{pa-def:etale_cover_baseChangeInclude}).

\begin{lemma}[Descent survives base change of the cover]
  \label{pa-lem:baseChangeInclude_mem_descentClasses}
  \lean{AlgebraicGeometry.baseChangeInclude_mem_descentClasses}
  \leanok
  \dcref{custom}

  \uses{pa-def:etale_cover_baseChangeInclude, pa-thm:relPicAlgMap_congr, pa-def:descentClasses}
  Let \(E\) be a presented \'etale cover of \(A\) and \(\lambda \in \mathrm{Desc}(C,E)\). Then
  \(\kappa_E^{*}\lambda \in \mathrm{Desc}(C, E_{A'})\).
\end{lemma}
\begin{proof}
  \leanok
  The two coprojections of the double cover of \(E_{A'}\)
  (Definition~\ref{pa-def:doubleInl_doubleInr}), composed with \(\kappa_E\), are two
  homomorphisms of \(A\)-algebras \(B_E \to B_{E_{A'}} \otimes_{A'} B_{E_{A'}}\) (both factor
  through the \(A\)-algebra structure of the target inherited from \(A \to A'\)): namely
  \(\mathrm{inl} \circ \kappa_E\) and \(\mathrm{inr} \circ \kappa_E\), sending \(b \mapsto
  (1 \otimes b) \otimes 1\), respectively \(b \mapsto 1 \otimes (1 \otimes b)\), under the
  carrier identifications. By the keystone (Theorem~\ref{pa-thm:relPicAlgMap_congr}), applied with
  \(R = B_{E_{A'}} \otimes_{A'} B_{E_{A'}}\) and \(j_1 = \mathrm{inl} \circ \kappa_E\), \(j_2 =
  \mathrm{inr} \circ \kappa_E\), the class \(\lambda \in \mathrm{Desc}(C,E)\) satisfies
  \(j_1^{*}\lambda = j_2^{*}\lambda\); by functoriality this reads \(\mathrm{inl}^{*}
  (\kappa_E^{*}\lambda) = \mathrm{inr}^{*}(\kappa_E^{*}\lambda)\), the descent condition for
  \(\kappa_E^{*}\lambda\) on \(E_{A'}\).
\end{proof}

\begin{definition}[Transport along base change of the cover]
  \label{pa-def:descentBaseChange}
  \lean{AlgebraicGeometry.descentBaseChange}
  \leanok
  \dcref{custom}

  \uses{pa-lem:baseChangeInclude_mem_descentClasses}
  The induced homomorphism of groups \(\mathrm{Desc}(C,E) \to \mathrm{Desc}(C,E_{A'})\),
  \(\lambda \mapsto \kappa_E^{*}\lambda\).
\end{definition}

\begin{lemma}[Naturality of base-change transport]
  \label{pa-lem:descentMap_baseChangeMap}
  \lean{AlgebraicGeometry.descentMap_baseChangeMap}
  \leanok
  \dcref{custom}

  \uses{pa-def:descentBaseChange, pa-lem:etale_cover_baseChangeMap_comp_baseChangeInclude,
    pa-def:descentMap}
  Let \(E, F\) be presented \'etale covers of \(A\) and \(h \colon B_E \to B_F\) a
  homomorphism of \(A\)-algebras. Then base-change transport commutes with refinement
  transport: for \(\lambda \in \mathrm{Desc}(C,E)\),
  \[
    \mathrm{baseChangeMap}(h)_{*}\bigl(\kappa_{E}^{*}\lambda\bigr) = \kappa_{F}^{*}(h_{*}\lambda)
  \]
  in \(\mathrm{Desc}(C,F_{A'})\), where \(\mathrm{baseChangeMap}(h) \colon B_{E_{A'}} \to
  B_{F_{A'}}\) is the base change of \(h\) (Definition~\ref{pa-def:etale_cover_baseChangeMap}).
\end{lemma}
\begin{proof}
  \leanok
  By Lemma~\ref{pa-lem:etale_cover_baseChangeMap_comp_baseChangeInclude},
  \(\mathrm{baseChangeMap}(h) \circ \kappa_E = \kappa_F \circ h\) as homomorphisms of
  \(A\)-algebras \(B_E \to B_{F_{A'}}\). Both sides of the claimed identity equal, by
  functoriality of transport, the transport of \(\lambda\) along this common composite.
\end{proof}

\begin{definition}[Restriction along a base extension]
  \label{pa-def:PicEtAff_map}
  \lean{AlgebraicGeometry.PicEtAff.map}
  \leanok
  \dcref{custom}

  \uses{pa-def:descentBaseChange, pa-lem:descentMap_baseChangeMap, pa-thm:PicEtAff_mk_mul_mk,
    pa-thm:descentMap_congr}
  The homomorphism of groups \(\tau^{*} \colon \mathrm{PicEtAff}(C,A) \to
  \mathrm{PicEtAff}(C,A')\), \([E,\lambda] \mapsto [E_{A'}, \kappa_E^{*}\lambda]\).
\end{definition}
\begin{proof}
  \leanok
  \emph{Well-defined on plus classes.} If \((E,\lambda) \sim (F,\mu)\) via \((G,f,g)\)
  (so \(f_{*}\lambda = g_{*}\mu\)), then base-changing gives \((E_{A'},\kappa_E^{*}\lambda) \sim
  (F_{A'},\kappa_F^{*}\mu)\) via \((G_{A'}, \mathrm{baseChangeMap}(f),
  \mathrm{baseChangeMap}(g))\): by Lemma~\ref{pa-lem:descentMap_baseChangeMap},
  \(\mathrm{baseChangeMap}(f)_{*}(\kappa_E^{*}\lambda) = \kappa_G^{*}(f_{*}\lambda) =
  \kappa_G^{*}(g_{*}\mu) = \mathrm{baseChangeMap}(g)_{*}(\kappa_F^{*}\mu)\).

  \emph{Multiplicativity.} For \(a = [E,\lambda]\), \(b = [F,\mu]\), first transport \(a\) and
  \(b\) to the common refinement \(E \times F\) (Lemma~\ref{pa-lem:PicEtAff_mk_descentMap}):
  \(a = [E \times F, \mathrm{inl}_{*}\lambda]\), \(b = [E \times F, \mathrm{inr}_{*}\mu]\), so
  \(\tau^{*}a = [(E \times F)_{A'}, \kappa_{E \times F}^{*}(\mathrm{inl}_{*}\lambda)]\) and
  \(\tau^{*}b = [(E \times F)_{A'}, \kappa_{E \times F}^{*}(\mathrm{inr}_{*}\mu)]\); by
  Theorem~\ref{pa-thm:PicEtAff_mk_mul_mk_same},
  \[
    \tau^{*}a \cdot \tau^{*}b = \bigl[(E \times F)_{A'},\ \kappa_{E \times F}^{*}
    (\mathrm{inl}_{*}\lambda) \cdot \kappa_{E \times F}^{*}(\mathrm{inr}_{*}\mu)\bigr].
  \]
  On the other hand, \(ab = [E \times F, \mathrm{inl}_{*}\lambda \cdot \mathrm{inr}_{*}\mu]\)
  (Theorem~\ref{pa-thm:PicEtAff_mk_mul_mk}), so, since \(\kappa_{E \times F}^{*}\) is a
  homomorphism of groups, \(\tau^{*}(ab) = [(E \times F)_{A'}, \kappa_{E \times F}^{*}
  (\mathrm{inl}_{*}\lambda) \cdot \kappa_{E \times F}^{*}(\mathrm{inr}_{*}\mu)]\) as well. Hence
  \(\tau^{*}(ab) = \tau^{*}a \cdot \tau^{*}b\).
\end{proof}

\begin{theorem}[Naturality of the unit]
  \label{pa-thm:PicEtAff_map_unit}
  \lean{AlgebraicGeometry.PicEtAff.map_unit}
  \leanok
  \dcref{custom}

  \uses{pa-def:PicEtAff_map, pa-def:PicEtAff_unit}
  For \(x \in \mathrm{relPic}(C,A)\), \(\tau^{*}(\eta_A(x)) = \eta_{A'}(\tau^{*}x)\) in
  \(\mathrm{PicEtAff}(C,A')\), where on the right \(\tau^{*} \colon \mathrm{relPic}(C,A) \to
  \mathrm{relPic}(C,A')\) is restriction along the algebra map \(\tau\)
  (Definition~\ref{pa-def:relPicAlgMap}).
\end{theorem}
\begin{proof}
  \leanok
  \uses{pa-lem:PicEtAff_mk_descentMap, pa-lem:PicEtAff_tautological_mem_descentClasses}
  Let \(\sigma \colon B_{\mathrm{self}(A')} \to B_{\mathrm{self}(A)_{A'}}\) be the canonical map
  from the trivial cover of \(A'\) (Definition~\ref{pa-def:etale_cover_fromSelf}, for the cover
  \(\mathrm{self}(A)_{A'}\) of \(A'\)). The two composite \(A\)-algebra homomorphisms
  \(A \to B_{\mathrm{self}(A)_{A'}}\) given by \(\kappa_{\mathrm{self}(A)} \circ \iota_{
  \mathrm{self}(A)}\) and by \(\sigma \circ \iota_{\mathrm{self}(A')} \circ \tau\) both equal
  the structure homomorphism of \(B_{\mathrm{self}(A)_{A'}}\) as an \(A\)-algebra: the first,
  because \(\kappa_{\mathrm{self}(A)}\) and \(\iota_{\mathrm{self}(A)}\) are structure
  homomorphisms; the second, because algebra-map formation is compatible with the tower
  \(A \to A' \to B_{\mathrm{self}(A)_{A'}}\) (the structure map of an \(A\)-algebra obtained by
  restriction of scalars along \(A \to A'\) factors through \(A'\)). Hence the two composites
  coincide, and by functoriality (Lemma~\ref{pa-lem:relPicAlgMap_functor}), for \(x \in
  \mathrm{relPic}(C,A)\),
  \[
    \kappa_{\mathrm{self}(A)}^{*}(\iota_{\mathrm{self}(A)}^{*}x) = \sigma^{*}\bigl(
    \iota_{\mathrm{self}(A')}^{*}(\tau^{*}x)\bigr).
  \]
  The left side represents \(\tau^{*}(\eta_A(x)) = [\mathrm{self}(A)_{A'},
  \kappa_{\mathrm{self}(A)}^{*}(\iota_{\mathrm{self}(A)}^{*}x)]\)
  (Definition~\ref{pa-def:PicEtAff_map} applied to the representative \([\mathrm{self}(A),
  \iota_{\mathrm{self}(A)}^{*}x]\) of \(\eta_A(x)\)). By the displayed equality this equals
  \([\mathrm{self}(A)_{A'}, \sigma^{*}(\iota_{\mathrm{self}(A')}^{*}(\tau^{*}x))]\), which by
  Lemma~\ref{pa-lem:PicEtAff_mk_descentMap} (transport along \(\sigma\)) equals
  \([\mathrm{self}(A'), \iota_{\mathrm{self}(A')}^{*}(\tau^{*}x)] = \eta_{A'}(\tau^{*}x)\).
\end{proof}

\begin{theorem}[Restriction along the trivial extension is the identity]
  \label{pa-thm:PicEtAff_map_id}
  \lean{AlgebraicGeometry.PicEtAff.map_id}
  \leanok
  \dcref{custom}

  \uses{pa-def:PicEtAff_map}
  When \(A' = A\) with \(\tau = \id_A\), the induced homomorphism \(\tau^{*} \colon
  \mathrm{PicEtAff}(C,A) \to \mathrm{PicEtAff}(C,A)\) is the identity.
\end{theorem}
\begin{proof}
  \leanok
  \uses{pa-lem:PicEtAff_mk_descentMap}
  For \([E,\lambda] \in \mathrm{PicEtAff}(C,A)\), base-changing \(E\) along \(A \to A\) gives a
  cover with canonical inclusion \(\kappa_E \colon B_E \to B_{E_A}\) itself an \(A\)-algebra
  refinement map, and by construction transport of \(\lambda\) along \(\kappa_E\) computes the
  same class as \(\kappa_E^{*}\lambda\) in the sense of Definition~\ref{pa-def:descentMap}. Hence
  \(\tau^{*}[E,\lambda] = [E_A, \kappa_E^{*}\lambda] = [E,\lambda]\) by
  Lemma~\ref{pa-lem:PicEtAff_mk_descentMap}, applied to the refinement map \(\kappa_E \colon B_E
  \to B_{E_A}\).
\end{proof}

\begin{theorem}[Restriction is transitive]
  \label{pa-thm:PicEtAff_map_map}
  \lean{AlgebraicGeometry.PicEtAff.map_map}
  \leanok
  \dcref{custom}

  \uses{pa-def:PicEtAff_map, pa-thm:relPicAlgMap_congr, pa-thm:PicEtAff_mk_eq_mk_iff}
  Let \(A''\) be a commutative \(k\)-algebra with compatible algebra structures forming a tower
  \(k \to A \to A' \to A''\). Then, writing \(\tau' \colon A' \to A''\) and \(\tau'' \colon A
  \to A''\) for the induced maps, restriction along \(A \to A' \to A''\) equals restriction
  along \(A \to A''\) directly:
  \(\tau'^{*} \circ \tau^{*} = \tau''^{*}\) as homomorphisms \(\mathrm{PicEtAff}(C,A) \to
  \mathrm{PicEtAff}(C,A'')\).
\end{theorem}
\begin{proof}
  \leanok
  Fix \([E,\lambda] \in \mathrm{PicEtAff}(C,A)\), and write \(E_{A''}\) for the direct base
  change along \(A \to A''\) and \((E_{A'})_{A''}\) for the iterated base change along
  \(A \to A' \to A''\). Both \(\tau'^{*}(\tau^{*}[E,\lambda])\) and \(\tau''^{*}[E,\lambda]\)
  restrict, respectively, to classes on \((E_{A'})_{A''}\) and on \(E_{A''}\); the two
  canonical inclusions of \(B_E\) into these two carriers, mapped further into the common
  refinement \((E_{A'})_{A''} \times E_{A''}\) via its two coprojections
  (Definition~\ref{pa-def:etale_cover_prod_coproj}), are two homomorphisms of \(A\)-algebras out
  of \(B_E\) landing in the same ring. By the keystone (Theorem~\ref{pa-thm:relPicAlgMap_congr}),
  since \(\lambda \in \mathrm{Desc}(C,E)\), these two composite pullbacks of \(\lambda\) agree;
  unwinding this equality through functoriality of transport identifies the transports of
  \(\tau'^{*}(\tau^{*}[E,\lambda])\) and \(\tau''^{*}[E,\lambda]\) to the common refinement
  \((E_{A'})_{A''} \times E_{A''}\), so by Theorem~\ref{pa-thm:PicEtAff_mk_eq_mk_iff},
  \(\tau'^{*}(\tau^{*}[E,\lambda]) = \tau''^{*}[E,\lambda]\).
\end{proof}

\section{The transport core: comparing the plus across curves and base fields}

Throughout this section \(k_D\), \(k_E\) and \(k_T\) are fields, with \(k_T\) an extension
of both \(k_D\) and \(k_E\); \(D\) is an object of \(\Over(\Spec k_D)\) and \(E\) an object
of \(\Over(\Spec k_E)\). A \emph{test algebra} is a commutative \(k_T\)-algebra \(B\),
regarded as a \(k_D\)- and \(k_E\)-algebra through the two embeddings; every construction
performed on test algebras below (\'etale covers and their carriers, products, double
covers, base changes along maps of test algebras) again yields test algebras. The two
consumers of this section are the curve-direction functoriality of the plus (the case
\(k_D = k_E = k_T\), treated at the end) and the base-field shuffle
\(\mathrm{PicEtAff}(C, A) \simeq \mathrm{PicEtAff}(C_L, A)\) (the cases
\(\{k_D, k_E\} = \{k, L\}\), \(k_T = L\)).

\begin{definition}[Transport family]
  \label{pa-def:relPicTransportFamily}
  \lean{AlgebraicGeometry.RelPicTransportFamily}
  \leanok
  \dcref{custom}

  A \emph{transport family} from \(D\) to \(E\) assigns to every test algebra \(B\) a
  morphism
  \[
    h_B \colon D \otimes \Spec B \longrightarrow E \otimes \Spec B
  \]
  of the underlying schemes of the two product bundles (the product on the left taken over
  \(\Spec k_D\), on the right over \(\Spec k_E\)), such that:
  \begin{enumerate}
  \item \emph{(projection compatibility)} \(\mathrm{snd}_E \circ h_B = \mathrm{snd}_D\) as
    morphisms to the shared scheme \(\Spec B\);
  \item \emph{(naturality)} for every homomorphism \(f \colon B \to B'\) of test algebras
    that is simultaneously \(k_D\)-linear and \(k_E\)-linear,
    \((E \lhd \Spec f) \circ h_{B'} = h_B \circ (D \lhd \Spec f)\) as morphisms
    \(D \otimes \Spec B' \to E \otimes \Spec B\).
  \end{enumerate}
\end{definition}

\begin{definition}[Relative Picard transport]
  \label{pa-def:relPicTransport}
  \lean{AlgebraicGeometry.RelPicTransportFamily.picFromBase_le_comap_hom,
    AlgebraicGeometry.RelPicTransportFamily.relPicHom,
    AlgebraicGeometry.RelPicTransportFamily.relPicHom_mk}
  \leanok
  \dcref{custom}

  \uses{pa-def:relPicTransportFamily, pa-def:relPic}
  Let \(h\) be a transport family from \(D\) to \(E\). For a test algebra \(B\), the
  homomorphism of groups
  \(h_B^{*} \colon \mathrm{relPic}(E, \Spec B) \to \mathrm{relPic}(D, \Spec B)\) induced by
  pullback of \v Cech Picard classes along \(h_B\) on representatives.
\end{definition}
\begin{proof}
  \leanok
  Well-definedness on the cosets: if \(X = \mathrm{snd}_E^{*} N\) is pulled back from
  \(\Spec B\), then \(h_B^{*} X = h_B^{*}\, \mathrm{snd}_E^{*} N = \mathrm{snd}_D^{*} N\)
  by projection compatibility, so \(h_B^{*}\) carries the subgroup \(H\) of classes pulled
  back from the test into the corresponding subgroup on the \(D\)-side and descends to the
  quotients.
\end{proof}

\begin{lemma}[Bilateral naturality of the relative Picard transport]
  \label{pa-lem:relPicTransport_naturality}
  \lean{AlgebraicGeometry.RelPicTransportFamily.relPicAlgMap_relPicHom}
  \leanok
  \dcref{custom}

  \uses{pa-def:relPicTransport, pa-def:relPicAlgMap}
  Let \(h\) be a transport family from \(D\) to \(E\) and \(f \colon B \to B'\) a
  homomorphism of test algebras that is both \(k_D\)- and \(k_E\)-linear. Then
  \(f^{*} \circ h_B^{*} = h_{B'}^{*} \circ f^{*}\) as maps
  \(\mathrm{relPic}(E, \Spec B) \to \mathrm{relPic}(D, \Spec B')\), where the outer
  restrictions are taken along \(f\) over \(k_D\) and over \(k_E\) respectively.
\end{lemma}
\begin{proof}
  \leanok
  On a representative \v Cech class the two sides pull back along
  \(h_B \circ (D \lhd \Spec f)\) and \((E \lhd \Spec f) \circ h_{B'}\) respectively, and
  these morphisms are equal by the naturality axiom of the family; conclude by
  functoriality of pullback of \v Cech classes and the computation rules of both
  transports.
\end{proof}

\begin{lemma}[Pointwise-inverse families invert the transport]
  \label{pa-lem:relPicTransport_inverse}
  \lean{AlgebraicGeometry.RelPicTransportFamily.relPicHom_relPicHom}
  \leanok
  \dcref{custom}

  \uses{pa-def:relPicTransport}
  Let \(h\) be a transport family from \(D\) to \(E\) and \(h'\) a transport family from
  \(E\) to \(D\) such that \(h_B \circ h'_B = \id\) on \(E \otimes \Spec B\) for a test
  algebra \(B\). Then \(h'^{*}_B \circ h_B^{*} = \id\) on
  \(\mathrm{relPic}(E, \Spec B)\).
\end{lemma}
\begin{proof}
  \leanok
  On a representative, the round trip pulls back along \(h_B \circ h'_B = \id\), and
  pullback along the identity is the identity.
\end{proof}

\begin{lemma}[Descent classes transport]
  \label{pa-lem:relPicTransport_descent}
  \lean{AlgebraicGeometry.RelPicTransportFamily.relPicHom_mem_descentClasses,
    AlgebraicGeometry.RelPicTransportFamily.descentHom,
    AlgebraicGeometry.RelPicTransportFamily.descentHom_coe,
    AlgebraicGeometry.RelPicTransportFamily.descentHom_descentMap}
  \leanok
  \dcref{custom}

  \uses{pa-def:descentClasses, pa-def:descentMap, pa-lem:relPicTransport_naturality}
  Let \(h\) be a transport family from \(D\) to \(E\), \(A\) a test algebra and \(F\) a
  presented \'etale cover of \(A\) with carrier \(B_F\). Then
  \(h_{B_F}^{*}\bigl(\mathrm{Desc}(E, F)\bigr) \subseteq \mathrm{Desc}(D, F)\), and the
  induced homomorphism of groups
  \(h_F^{*} \colon \mathrm{Desc}(E, F) \to \mathrm{Desc}(D, F)\) commutes with refinement
  transport: \(h_{G}^{*}(r_{*}\lambda) = r_{*}(h_F^{*}\lambda)\) for every refinement map
  \(r \colon B_F \to B_G\) of covers of \(A\).
\end{lemma}
\begin{proof}
  \leanok
  The two coprojections \(B_F \rightrightarrows B_F \otimes_A B_F\) into the double cover
  are homomorphisms of test algebras linear over both base fields, so by
  Lemma~\ref{pa-lem:relPicTransport_naturality} the descent condition
  (equal pullbacks along the two coprojections) is preserved by \(h_{B_F}^{*}\).
  Refinement transport is restriction along the refinement map, which is again linear over
  both base fields, so the commutation is a second instance of
  Lemma~\ref{pa-lem:relPicTransport_naturality}.
\end{proof}

\begin{definition}[Transport of the plus]
  \label{pa-def:PicEtAffTransport}
  \lean{AlgebraicGeometry.RelPicTransportFamily.picEtAffHom,
    AlgebraicGeometry.RelPicTransportFamily.picEtAffHom_mk}
  \leanok
  \dcref{custom}

  \uses{pa-def:PicEtAff, pa-def:PicEtAff_mk, pa-thm:PicEtAff_mk_eq_mk_iff, pa-thm:PicEtAff_mk_mul_mk,
    pa-lem:relPicTransport_descent}
  Let \(h\) be a transport family from \(D\) to \(E\) and \(A\) a test algebra. The
  homomorphism of groups
  \(h^{*} \colon \mathrm{PicEtAff}(E, A) \to \mathrm{PicEtAff}(D, A)\),
  \([F, \lambda] \mapsto [F, h_F^{*}\lambda]\), over the literally shared index of
  presented \'etale covers of \(A\).
\end{definition}
\begin{proof}
  \leanok
  \emph{Well-defined.} If \((F, \lambda) \sim (F', \mu)\) via a common refinement
  \((G, r, r')\), so that \(r_{*}\lambda = r'_{*}\mu\) in \(\mathrm{Desc}(E, G)\), then
  applying \(h_G^{*}\) and commuting it past the two refinement transports
  (Lemma~\ref{pa-lem:relPicTransport_descent}) gives
  \(r_{*}(h_F^{*}\lambda) = r'_{*}(h_{F'}^{*}\mu)\), i.e.\
  \((F, h_F^{*}\lambda) \sim (F', h_{F'}^{*}\mu)\).

  \emph{Multiplicative.} For \(a = [F, \lambda]\) and \(b = [F', \mu]\), the product is
  represented on the product cover \(F \times F'\) by
  \(\mathrm{inl}_{*}\lambda \cdot \mathrm{inr}_{*}\mu\)
  (Theorem~\ref{pa-thm:PicEtAff_mk_mul_mk}); since \(h_{F \times F'}^{*}\) is a homomorphism
  of groups commuting with the two coprojection transports
  (Lemma~\ref{pa-lem:relPicTransport_descent}), its value on this representative is
  \(\mathrm{inl}_{*}(h_F^{*}\lambda) \cdot \mathrm{inr}_{*}(h_{F'}^{*}\mu)\), which
  represents \(h^{*}a \cdot h^{*}b\) (Theorem~\ref{pa-thm:PicEtAff_mk_mul_mk} again).
\end{proof}

\begin{theorem}[The transport is compatible with the unit]
  \label{pa-thm:PicEtAffTransport_unit}
  \lean{AlgebraicGeometry.RelPicTransportFamily.picEtAffHom_unit}
  \leanok
  \dcref{custom}

  \uses{pa-def:PicEtAffTransport, pa-def:PicEtAff_unit}
  For \(x \in \mathrm{relPic}(E, \Spec A)\),
  \(h^{*}(\eta_A(x)) = \eta_A(h_A^{*} x)\) in \(\mathrm{PicEtAff}(D, A)\).
\end{theorem}
\begin{proof}
  \leanok
  \uses{pa-lem:relPicTransport_naturality}
  Both sides are represented on the trivial cover \(\mathrm{self}(A)\); the representatives
  are the pullbacks of \(x\) along the two composites of \(h\) with the structure map
  \(A \to B_{\mathrm{self}(A)}\), which agree by the bilateral naturality square
  (Lemma~\ref{pa-lem:relPicTransport_naturality}) applied to the structure map, a
  homomorphism of test algebras linear over both base fields.
\end{proof}

\begin{theorem}[The transport is compatible with restriction of the test]
  \label{pa-thm:PicEtAffTransport_map}
  \lean{AlgebraicGeometry.RelPicTransportFamily.map_picEtAffHom,
    AlgebraicGeometry.RelPicTransportFamily.mapAlg_picEtAffHom}
  \leanok
  \dcref{custom}

  \uses{pa-def:PicEtAffTransport, pa-def:PicEtAff_map}
  Let \(\tau \colon A \to A'\) be a homomorphism of test algebras linear over both base
  fields. Then \(\tau^{*} \circ h^{*} = h^{*} \circ \tau^{*}\) as maps
  \(\mathrm{PicEtAff}(E, A) \to \mathrm{PicEtAff}(D, A')\).
\end{theorem}
\begin{proof}
  \leanok
  \uses{pa-lem:relPicTransport_naturality, pa-lem:relPicTransport_descent}
  On a representative \([F, \lambda]\), both sides are represented on the base-changed
  cover \(F_{A'}\): one by transporting \(\lambda\) along the canonical inclusion
  \(\kappa_F \colon B_F \to B_{F_{A'}}\) and then applying the family transport, the other
  in the opposite order. The canonical inclusion is a homomorphism of test algebras linear
  over both base fields, so the two representatives agree by
  Lemma~\ref{pa-lem:relPicTransport_naturality}.
\end{proof}

\begin{theorem}[Pointwise-inverse families invert the transport of the plus]
  \label{pa-thm:PicEtAffTransport_inverse}
  \lean{AlgebraicGeometry.RelPicTransportFamily.picEtAffHom_picEtAffHom}
  \leanok
  \dcref{custom}

  \uses{pa-def:PicEtAffTransport, pa-lem:relPicTransport_inverse}
  Let \(h\) be a transport family from \(D\) to \(E\) and \(h'\) a transport family from
  \(E\) to \(D\) with \(h_B \circ h'_B = \id\) for every test algebra \(B\). Then
  \(h'^{*} \circ h^{*} = \id\) on \(\mathrm{PicEtAff}(E, A)\) for every test algebra
  \(A\). In particular a pair of transport families that are pointwise mutually inverse
  induces an isomorphism of the plus groups.
\end{theorem}
\begin{proof}
  \leanok
  On a representative \([F, \lambda]\), the round trip transports \(\lambda\) by
  \(h'^{*}_{B_F} \circ h_{B_F}^{*} = \id\) (Lemma~\ref{pa-lem:relPicTransport_inverse}) on the
  same cover \(F\).
\end{proof}

\section{Functoriality of the plus in the curve}

Throughout this section \(k\) is a field, \(D, E\) are objects of \(\Over(\Spec k)\), and
\(A\) is a commutative \(k\)-algebra.

\begin{definition}[Curve-direction transport of the plus]
  \label{pa-def:PicEtAff_curveMap}
  \lean{AlgebraicGeometry.curveTransportFamily,
    AlgebraicGeometry.curveTransportFamily_hom,
    AlgebraicGeometry.curveTransportFamily_relPicHom,
    AlgebraicGeometry.PicEtAff.curveMap,
    AlgebraicGeometry.PicEtAff.curveMap_mk,
    AlgebraicGeometry.PicEtAff.curveTransportFamily_descentHom_coe,
    AlgebraicGeometry.PicEtAff.curveMap_mk_eq_mk}
  \leanok
  \dcref{custom}

  \uses{pa-def:relPicTransportFamily, pa-def:PicEtAffTransport, pa-lem:relPicMapCurve}
  Let \(g \colon D \to E\) be a morphism of \(\Over(\Spec k)\). The assignment
  \(B \mapsto g \rhd \Spec B\) is a transport family from \(D\) to \(E\) in the sense of
  Definition~\ref{pa-def:relPicTransportFamily} (with \(k_D = k_E = k_T = k\)); its induced
  relative Picard transport is the curve-direction transport of
  Lemma~\ref{pa-lem:relPicMapCurve}, and the induced homomorphism of the plus groups is
  denoted
  \[
    g^{*} \colon \mathrm{PicEtAff}(E, A) \longrightarrow \mathrm{PicEtAff}(D, A),
    \qquad [F, \lambda] \longmapsto [F, (g \rhd \Spec B_F)^{*}\lambda],
  \]
  transporting each descent-class representative on its own cover.
\end{definition}
\begin{proof}
  \leanok
  Projection compatibility is the whiskering identity
  \(\mathrm{snd}_E \circ (g \rhd T) = \mathrm{snd}_D\); naturality against \(\Spec f\) for
  a \(k\)-algebra map \(f\) is the exchange identity
  \((E \lhd \Spec f) \circ (g \rhd \Spec B') = (g \rhd \Spec B) \circ (D \lhd \Spec f)\)
  of whiskerings.
\end{proof}

\begin{theorem}[Laws of the curve-direction transport]
  \label{pa-thm:PicEtAff_curveMap_laws}
  \lean{AlgebraicGeometry.PicEtAff.curveMap_id,
    AlgebraicGeometry.PicEtAff.curveMap_comp,
    AlgebraicGeometry.PicEtAff.curveMap_unit,
    AlgebraicGeometry.PicEtAff.map_curveMap,
    AlgebraicGeometry.PicEtAff.mapAlg_curveMap}
  \leanok
  \dcref{custom}

  \uses{pa-def:PicEtAff_curveMap, pa-def:PicEtAff_unit, pa-def:PicEtAff_map}
  The curve-direction transport of the plus is contravariantly functorial in the curve:
  \(\id_E^{*} = \id\), and \((g' \circ g)^{*} = g^{*} \circ g'^{*}\) for
  \(g \colon D \to E\) and \(g' \colon E \to F\). It is compatible with the unit of the
  plus, \(g^{*} \circ \eta_A = \eta_A \circ (g \rhd \Spec A)^{*}\), and it commutes with
  restriction along any base extension and along any \(k\)-algebra map of affine tests.
\end{theorem}
\begin{proof}
  \leanok
  \uses{pa-lem:relPicMapCurve, pa-thm:PicEtAffTransport_unit, pa-thm:PicEtAffTransport_map}
  The functor laws hold on each descent-class representative by the pointwise functor laws
  of the relative Picard curve transport (Lemma~\ref{pa-lem:relPicMapCurve}), since the cover
  index is untouched. Unit compatibility is Theorem~\ref{pa-thm:PicEtAffTransport_unit} and
  the restriction squares are Theorem~\ref{pa-thm:PicEtAffTransport_map}, both applied to the
  whisker family.
\end{proof}

\section{The base-field shuffle of the plus}
\label{pa-sec:PicEtAffBaseFieldShuffle}

Throughout this section \(k \to L\) is a field extension, \(C\) is an object of
\(\Over(\Spec k)\) with base change \(C_L \in \Over(\Spec L)\) (\ref{pa-def:baseChange}),
and \(A\) is an algebra in a scalar tower \(k \to L \to A\). The presented \'etale
covers of \(A\) (Definition~\ref{pa-def:etale_cover}) do not mention the base field, so the
plus groups \(\mathrm{PicEtAff}(C, A)\) (formed over \(k\)) and
\(\mathrm{PicEtAff}(C_L, A)\) (formed over \(L\)) are indexed by the \emph{same} covers.
This section transports one to the other along the same-carrier comparison
\(c_B \colon C_L \otimes \Spec B \to C \otimes \Spec B\) of
Definition~\ref{pa-def:crossBaseAffineIso} (Chapter~\ref{pa-chap:Cohomology}), by
instantiating the transport core of Definition~\ref{pa-def:relPicTransportFamily} in both
directions.

\begin{definition}[The base-field shuffle of the plus]
  \label{pa-def:PicEtAff_baseFieldShuffle}
  \lean{AlgebraicGeometry.crossBaseAffineIso_inv_naturality,
    AlgebraicGeometry.crossBaseTransportFamily,
    AlgebraicGeometry.crossBaseTransportFamilyInv,
    AlgebraicGeometry.crossBaseTransportFamily_hom,
    AlgebraicGeometry.crossBaseTransportFamilyInv_hom,
    AlgebraicGeometry.PicEtAff.baseFieldShuffle,
    AlgebraicGeometry.PicEtAff.baseFieldShuffle_apply,
    AlgebraicGeometry.PicEtAff.baseFieldShuffle_symm_apply,
    AlgebraicGeometry.PicEtAff.baseFieldShuffle_mk,
    AlgebraicGeometry.PicEtAff.baseFieldShuffle_symm_mk}
  \leanok
  \dcref{custom}

  \uses{pa-def:relPicTransportFamily, pa-def:PicEtAffTransport, pa-thm:PicEtAffTransport_inverse,
    pa-def:crossBaseAffineIso, pa-def:baseChange}
  The same-carrier comparisons \(B \mapsto c_B\) and \(B \mapsto c_B^{-1}\) are transport
  families in the sense of Definition~\ref{pa-def:relPicTransportFamily} — from \(C_L\) to
  \(C\) with \((k_D, k_E, k_T) = (L, k, L)\), and from \(C\) to \(C_L\) with
  \((k_D, k_E, k_T) = (k, L, L)\), respectively. The induced transports of the plus
  (Definition~\ref{pa-def:PicEtAffTransport}) are mutually inverse, and the resulting
  isomorphism of groups is the \emph{base-field shuffle}
  \[
    \mathrm{Sh}_A \colon \mathrm{PicEtAff}(C, A) \xrightarrow{\ \cong\ }
      \mathrm{PicEtAff}(C_L, A),
    \qquad [F, \lambda] \longmapsto [F, c_{B_F}^{*}\lambda],
  \]
  transporting each descent-class representative on its own cover; its inverse is
  \([F, \mu] \mapsto [F, (c_{B_F}^{-1})^{*}\mu]\).
\end{definition}
\begin{proof}
  \leanok
  \uses{pa-lem:crossBaseAffineIso_snd, pa-thm:crossBaseAffineIso_naturality}
  Projection compatibility of both families is the compatibility of the comparison with
  the two second projections onto the shared \(\Spec B\)
  (\ref{pa-lem:crossBaseAffineIso_snd}). Naturality of the forward family against a map
  \(f \colon B \to B'\) of test algebras that is \(L\)-linear (hence \(k\)-linear through
  the towers) is the naturality square of the comparison
  (\ref{pa-thm:crossBaseAffineIso_naturality}); naturality of the backward family follows by
  composing that square with \(c_{B'}^{-1}\) on the left and \(c_B^{-1}\) on the right.
  The two families are pointwise mutually inverse, so the induced transports of the plus
  are mutually inverse isomorphisms by Theorem~\ref{pa-thm:PicEtAffTransport_inverse}.
\end{proof}

\begin{theorem}[The shuffle restricts the relative Picard shuffle cover-wise]
  \label{pa-thm:PicEtAff_baseFieldShuffle_relPic}
  \lean{AlgebraicGeometry.crossBaseTransportFamily_relPicHom,
    AlgebraicGeometry.crossBaseTransportFamilyInv_relPicHom,
    AlgebraicGeometry.PicEtAff.crossBaseTransportFamily_descentHom_coe,
    AlgebraicGeometry.PicEtAff.crossBaseTransportFamilyInv_descentHom_coe}
  \leanok
  \dcref{custom}

  \uses{pa-def:PicEtAff_baseFieldShuffle, pa-def:relPicCrossBase, pa-def:relPicTransport}
  The relative Picard transports induced by the two families of
  Definition~\ref{pa-def:PicEtAff_baseFieldShuffle} are the base-field shuffle
  \(\mathrm{sh}\) of relative Picard classes and its inverse
  (Definition~\ref{pa-def:relPicCrossBase}): for every presented \'etale cover \(F\) of
  \(A\) with carrier \(B_F\) and every descent class
  \(\lambda \in \mathrm{Desc}(C, F)\),
  \(c_{B_F}^{*}\lambda = \mathrm{sh}_{B_F}(\lambda)\) in
  \(\mathrm{relPic}_L(C_L, B_F)\), and symmetrically
  \((c_{B_F}^{-1})^{*}\mu = \mathrm{sh}'_{B_F}(\mu)\) for
  \(\mu \in \mathrm{Desc}(C_L, F)\). In particular every statement about the relative
  Picard shuffle — notably its degree invariance at field tests — applies verbatim to
  the descent-class representatives of the plus shuffle.
\end{theorem}
\begin{proof}
  \leanok
  Both transports are defined as descent of pullback along the same morphism \(c_{B_F}\)
  (resp.\ \(c_{B_F}^{-1}\)) of the same product carriers to the same quotients
  (Definition~\ref{pa-def:relPicTransport} at the comparison family, and
  Definition~\ref{pa-def:relPicCrossBase}); they agree on every coset by their shared
  computation rule.
\end{proof}

\begin{theorem}[Laws of the base-field shuffle]
  \label{pa-thm:PicEtAff_baseFieldShuffle_laws}
  \lean{AlgebraicGeometry.PicEtAff.baseFieldShuffle_mk_eq_mk,
    AlgebraicGeometry.PicEtAff.baseFieldShuffle_symm_mk_eq_mk,
    AlgebraicGeometry.PicEtAff.baseFieldShuffle_unit,
    AlgebraicGeometry.PicEtAff.baseFieldShuffle_symm_unit,
    AlgebraicGeometry.PicEtAff.map_baseFieldShuffle,
    AlgebraicGeometry.PicEtAff.map_baseFieldShuffle_symm,
    AlgebraicGeometry.PicEtAff.mapAlg_baseFieldShuffle,
    AlgebraicGeometry.PicEtAff.mapAlg_baseFieldShuffle_symm}
  \leanok
  \dcref{custom}

  \uses{pa-def:PicEtAff_baseFieldShuffle, pa-def:PicEtAff_unit, pa-def:PicEtAff_map,
    pa-thm:PicEtAff_baseFieldShuffle_relPic}
  The base-field shuffle transports identifications of representatives: if
  \((F, \lambda)\) and \((F', \lambda')\) represent the same class of
  \(\mathrm{PicEtAff}(C, A)\), then \((F, c_{B_F}^{*}\lambda)\) and
  \((F', c_{B_{F'}}^{*}\lambda')\) represent the same class of
  \(\mathrm{PicEtAff}(C_L, A)\), and symmetrically for the inverse. It is compatible
  with the units of the two plus constructions,
  \(\mathrm{Sh}_A \circ \eta_A = \eta_A \circ \mathrm{sh}_A\) and
  \(\mathrm{Sh}_A^{-1} \circ \eta_A = \eta_A \circ \mathrm{sh}'_A\), and it commutes
  with restriction of the test: for every homomorphism \(\tau \colon A \to A'\) of
  algebras in scalar towers \(k \to L \to A\), \(k \to L \to A'\) that is \(L\)-linear
  (hence \(k\)-linear),
  \(\tau^{*} \circ \mathrm{Sh}_A = \mathrm{Sh}_{A'} \circ \tau^{*}\) and
  \(\tau^{*} \circ \mathrm{Sh}_A^{-1} = \mathrm{Sh}_A^{-1} \circ \tau^{*}\), where the
  restrictions are formed over \(L\) on the \(C_L\)-side and over \(k\) on the
  \(C\)-side.
\end{theorem}
\begin{proof}
  \leanok
  \uses{pa-thm:PicEtAffTransport_unit, pa-thm:PicEtAffTransport_map,
    pa-thm:relPicCrossBase_relPicAlgMap}
  Transport of identifications holds because the shuffle is well defined on the plus and
  computes cover-wise on representatives
  (Definition~\ref{pa-def:PicEtAff_baseFieldShuffle}): apply \(\mathrm{Sh}_A\) to the equal
  classes and read off the representatives. Unit compatibility is
  Theorem~\ref{pa-thm:PicEtAffTransport_unit} at the two comparison families, with the
  relative Picard transports identified by
  Theorem~\ref{pa-thm:PicEtAff_baseFieldShuffle_relPic}. The restriction squares are
  Theorem~\ref{pa-thm:PicEtAffTransport_map} at the two families, applied to \(\tau\)
  regarded simultaneously as a \(k\)- and an \(L\)-algebra map; on relative Picard
  representatives they restrict to the naturality square of the relative Picard shuffle
  (\ref{pa-thm:relPicCrossBase_relPicAlgMap}).
\end{proof}

\section{Separatedness I: pullback along the projection is injective}

The plus construction of the preceding sections is a legitimate computation of the \'etale
sheafification only because \emph{one} plus suffices, and that rests on the relative Picard
functor being separated for the \'etale topology. This section establishes the geometric heart
of that separatedness: pullback of Picard classes along the projection
\(\mathrm{pr} \colon C \otimes T \to T\) is \emph{injective}, for every test object \(T\).

Throughout, \(k\) is a field and \(C\) is an object of \(\Over(\Spec k)\) which is proper,
geometrically irreducible and geometrically reduced over \(k\). We recall from elsewhere in
this project the definitional Picard layer on a scheme \(X\): a \emph{pointed cover}
\(\mathcal U\) of \(X\) assigns to each point \(x \in X\) an open \(\mathcal U(x) \ni x\);
pointed covers are ordered by pointwise inclusion (a \emph{refinement}
\(\mathcal W \le \mathcal U\)), with pointwise intersection as meet; a \emph{unit \v Cech
\(1\)-cocycle} \(\gamma\) on \(\mathcal U\) assigns to each pair \(x,y\) a unit
\(\gamma_{xy} \in \struct X(\mathcal U(x) \cap \mathcal U(y))^{\times}\) satisfying
\(\gamma_{xy}\gamma_{yz} = \gamma_{xz}\) on triple intersections; \(\check H^1(\mathcal U)\)
is the group of such cocycles modulo coboundaries \((\delta\alpha)_{xy} =
\alpha_x\alpha_y^{-1}\) for families of units \(\alpha_x \in
\struct X(\mathcal U(x))^{\times}\); restriction along a refinement is a group homomorphism
\(\check H^1(\mathcal U) \to \check H^1(\mathcal W)\); and the Picard group
\(\Pic(X)\) is the colimit of the \(\check H^1(\mathcal U)\) along the (directed) refinement
order, with structure maps written \([\mathcal U, a] \in \Pic(X)\), and with contravariant
pullback \(f^{*} \colon \Pic(Y) \to \Pic(X)\) along a morphism \(f \colon X \to Y\), computed
by pulling back the cover pointwise and the cocycle by \(f^{\sharp}\).

\begin{lemma}[Separation for unit sections]
  \label{pa-lem:units_separation}
  \lean{AlgebraicGeometry.Scheme.unitsRestrict_eq_of_locally_eq}
  \leanok
  \dcref{custom}

  Let \(X\) be a scheme, \(V \subseteq X\) open, and \((W_i)_{i \in I}\) a family of opens
  with \(W_i \subseteq V\) and \(V \subseteq \bigcup_i W_i\). If two units
  \(s, t \in \struct X(V)^{\times}\) satisfy \(s|_{W_i} = t|_{W_i}\) for all \(i\), then
  \(s = t\).
\end{lemma}
\begin{proof}
  \leanok
  Units are equal once their underlying sections are, and the underlying sections agree on a
  cover of \(V\); conclude by the separation half of the sheaf axiom for \(\struct X\).
\end{proof}

\begin{lemma}[Gluing for unit sections]
  \label{pa-lem:units_gluing}
  \lean{AlgebraicGeometry.Scheme.exists_unitsRestrict_eq}
  \leanok
  \dcref{custom}

  Let \(X\) be a scheme, \(V \subseteq X\) open, and \((W_i)_{i \in I}\) a family of opens
  with \(W_i \subseteq V\) and \(V \subseteq \bigcup_i W_i\). Any family of units
  \(u_i \in \struct X(W_i)^{\times}\) with \(u_i|_{W_i \cap W_j} = u_j|_{W_i \cap W_j}\) for
  all \(i,j\) glues to a unit \(s \in \struct X(V)^{\times}\) with \(s|_{W_i} = u_i\).
\end{lemma}
\begin{proof}
  \leanok
  \uses{pa-lem:units_separation}
  The underlying sections \((u_i)\) form a compatible family, as do the inverses
  \((u_i^{-1})\); by the sheaf axiom they glue to sections \(s, s' \in \struct X(V)\)
  respectively. The products \(s s'\) and \(s' s\) restrict to \(1\) on every \(W_i\), hence
  equal \(1\) by the separation half of the sheaf axiom; so \(s\) is a unit, and it restricts
  to \(u_i\) because its underlying section does.
\end{proof}

\begin{lemma}[Restriction of unit sections is transitive]
  \label{pa-lem:unitsRestrict_trans}
  \lean{AlgebraicGeometry.Scheme.unitsRestrict_unitsRestrict}
  \leanok
  \dcref{custom}

  For opens \(W \subseteq V \subseteq U\) of a scheme \(X\) and
  \(u \in \struct X(U)^{\times}\), restricting \(u\) to \(V\) and then to \(W\) is
  restricting \(u\) to \(W\).
\end{lemma}
\begin{proof}
  \leanok
  Functoriality of the structure presheaf, applied to the underlying sections.
\end{proof}

\begin{theorem}[Refinement injectivity in degree one]
  \label{pa-thm:unitsRes_injective}
  \lean{AlgebraicGeometry.Scheme.unitsRes_injective}
  \leanok
  \dcref{custom}

  Let \(X\) be a scheme and \(\mathcal W \le \mathcal U\) pointed covers of \(X\). Then
  restriction \(\check H^1(\mathcal U) \to \check H^1(\mathcal W)\) is injective.
\end{theorem}
\begin{proof}
  \leanok
  \uses{pa-lem:units_separation, pa-lem:units_gluing, pa-lem:unitsRestrict_trans}
  Restriction is a group homomorphism, so it suffices to show a cocycle \(\gamma\) on
  \(\mathcal U\) whose restriction to \(\mathcal W\) is a coboundary, say
  \(\gamma_{zz'}|_{\mathcal W(z) \cap \mathcal W(z')} = \alpha_z \alpha_{z'}^{-1}\) with
  \(\alpha_z \in \struct X(\mathcal W(z))^{\times}\), is itself a coboundary. Fix a point
  \(i\) and consider, for each point \(z\), the \emph{comparison unit}
  \[ s_{i,z} := \alpha_z \cdot \gamma_{iz}^{-1}
     \in \struct X\bigl(\mathcal U(i) \cap \mathcal W(z)\bigr)^{\times}, \]
  which makes sense since \(\mathcal W(z) \subseteq \mathcal U(z)\). On the overlap
  \((\mathcal U(i) \cap \mathcal W(z)) \cap (\mathcal U(i) \cap \mathcal W(z'))\) the cocycle
  identity \(\gamma_{iz}\gamma_{zz'} = \gamma_{iz'}\) at \((i,z,z')\) and the coboundary
  relation \(\alpha_z \gamma_{zz'} = \alpha_{z'}\) give
  \(s_{i,z'} = \alpha_z\gamma_{zz'}\cdot(\gamma_{iz}\gamma_{zz'})^{-1} = s_{i,z}\). Since
  every point \(w \in \mathcal U(i)\) lies in \(\mathcal U(i) \cap \mathcal W(w)\), these
  opens cover \(\mathcal U(i)\), and by gluing (Lemma~\ref{pa-lem:units_gluing}) there is
  \(\beta_i \in \struct X(\mathcal U(i))^{\times}\) with
  \(\beta_i|_{\mathcal U(i) \cap \mathcal W(z)} = s_{i,z}\) for all \(z\). Finally, for any
  \(i, j\), on each member \(\mathcal U(i) \cap \mathcal U(j) \cap \mathcal W(z)\) of a cover
  of \(\mathcal U(i) \cap \mathcal U(j)\) we compute, using the cocycle identity at
  \((i,j,z)\),
  \[ \beta_i \gamma_{ij} = \alpha_z \gamma_{iz}^{-1}\gamma_{ij}
       = \alpha_z \gamma_{jz}^{-1} = \beta_j , \]
  so \(\beta_i \gamma_{ij} = \beta_j\) on \(\mathcal U(i) \cap \mathcal U(j)\) by separation
  (Lemma~\ref{pa-lem:units_separation}); that is, \(\gamma = \delta(\beta^{-1})\) is a
  coboundary.
\end{proof}

\begin{corollary}[Triviality is detected on the representing cover]
  \label{pa-cor:cechPic_mk_eq_one}
  \lean{AlgebraicGeometry.Scheme.CechPic.mk_eq_one_iff}
  \leanok
  \dcref{custom}

  For a pointed cover \(\mathcal U\) of a scheme \(X\) and \(a \in \check H^1(\mathcal U)\),
  the Picard class \([\mathcal U, a] \in \Pic(X)\) is trivial if and only if \(a = 1\).
\end{corollary}
\begin{proof}
  \leanok
  \uses{pa-thm:unitsRes_injective}
  If \([\mathcal U, a] = 1\) then \(a\) and \(1\) restrict to the same class on some common
  refinement \(\mathcal W \le \mathcal U\); by refinement injectivity
  (Theorem~\ref{pa-thm:unitsRes_injective}) already \(a = 1\). The converse is immediate.
\end{proof}

\begin{corollary}[\(\check H^1\) of a cover embeds in the Picard group]
  \label{pa-cor:cechPic_mk_injective}
  \lean{AlgebraicGeometry.Scheme.CechPic.mk_injective}
  \leanok
  \dcref{custom}

  For every pointed cover \(\mathcal U\) of a scheme \(X\), the canonical map
  \(\check H^1(\mathcal U) \to \Pic(X)\), \(a \mapsto [\mathcal U, a]\), is injective.
\end{corollary}
\begin{proof}
  \leanok
  \uses{pa-cor:cechPic_mk_eq_one}
  The map is a group homomorphism, and its kernel is trivial by
  Corollary~\ref{pa-cor:cechPic_mk_eq_one}.
\end{proof}

\begin{lemma}[Everything is flat over a field]
  \label{pa-lem:flat_of_field}
  \lean{AlgebraicGeometry.Flat.of_field}
  \leanok
  \dcref{custom}

  Every morphism of schemes \(X \to \Spec K\), \(K\) a field, is flat.
\end{lemma}
\begin{proof}
  \leanok
  Flatness is Zariski-local on the source, so it suffices to treat
  \(\Spec B \to \Spec K\), i.e.\ a ring homomorphism \(K \to B\); every module over a field
  is free, hence flat.
\end{proof}

\begin{lemma}[Affine refinement of pointed covers]
  \label{pa-lem:pointedCover_affine_refinement}
  \lean{AlgebraicGeometry.Scheme.PointedCover.exists_isAffineOpen_le}
  \leanok
  \dcref{custom}

  Every pointed cover of a scheme admits a refinement all of whose members are affine open.
\end{lemma}
\begin{proof}
  \leanok
  The affine opens form a basis of the topology of a scheme: for each point \(x\) choose an
  affine open \(V(x)\) with \(x \in V(x) \subseteq \mathcal U(x)\).
\end{proof}

\begin{theorem}[Sections over an affine open pull back bijectively]
  \label{pa-thm:isIso_appLE_snd}
  \lean{AlgebraicGeometry.Over.isIso_appLE_snd}
  \leanok
  \dcref{custom}

  Let \(T\) be an object of \(\Over(\Spec k)\) and \(V \subseteq T\) an affine open. Then
  pullback of sections along the second projection
  \(\mathrm{pr} \colon C \otimes T \to T\) is an isomorphism
  \[ \Gamma(T, V) \;\xrightarrow{\ \sim\ }\; \Gamma\bigl(C \otimes T,\
     \mathrm{pr}^{-1}V\bigr). \]
\end{theorem}
\begin{proof}
  \leanok
  \uses{pa-lem:flat_of_field, pa-cor:curve_sections}
  The underlying scheme of \(C \otimes T\) is the fibre product
  \(C \times_{\Spec k} T\). Since \(C\) is proper over the field \(k\), its underlying space
  is quasi-compact and quasi-separated, and \(T \to \Spec k\) is flat
  (Lemma~\ref{pa-lem:flat_of_field}); therefore the canonical square of section rings
  \[ \begin{array}{ccc}
     \Gamma(\Spec k, \struct{}) & \longrightarrow & \Gamma(C, \struct C) \\
     \downarrow & & \downarrow \\
     \Gamma(T, V) & \longrightarrow & \Gamma(C \otimes T, \mathrm{pr}^{-1}V)
     \end{array} \]
  --- with the affine slots at \(\Spec k\) (all of it) and at \(V \subseteq T\), and the
  quasi-compact quasi-separated slot at all of \(C\) --- is a pushout of commutative rings:
  the flat-base-change theorem for sections of fibre products. The top edge is an
  isomorphism, \(\Gamma(C, \struct C) = k\) (Corollary~\ref{pa-cor:curve_sections}); a pushout
  of an isomorphism is an isomorphism, so the bottom edge is one too.
\end{proof}

\begin{corollary}[Sections pull back injectively over every open]
  \label{pa-cor:appLE_snd_injective}
  \lean{AlgebraicGeometry.Over.appLE_snd_injective}
  \leanok
  \dcref{custom}

  Let \(T\) be an object of \(\Over(\Spec k)\) and \(W \subseteq T\) any open. Then pullback
  of sections along the second projection,
  \(\Gamma(T, W) \to \Gamma(C \otimes T, \mathrm{pr}^{-1}W)\), is injective.
\end{corollary}
\begin{proof}
  \leanok
  \uses{pa-thm:isIso_appLE_snd}
  Let \(s, t \in \Gamma(T,W)\) have equal pullbacks. For every affine open \(V \subseteq W\),
  restriction to \(\mathrm{pr}^{-1}V\) followed by the bijection of
  Theorem~\ref{pa-thm:isIso_appLE_snd} shows \(s|_V = t|_V\); the affine opens contained in
  \(W\) cover \(W\), so \(s = t\) by separation for \(\struct T\).
\end{proof}

\begin{corollary}[Unit sections over an affine open pull back bijectively]
  \label{pa-cor:unitsAppLE_snd_bijective}
  \lean{AlgebraicGeometry.Over.unitsAppLE_snd_bijective}
  \leanok
  \dcref{custom}

  Let \(T\) be an object of \(\Over(\Spec k)\) and \(V \subseteq T\) an affine open. Then
  pullback along the second projection is a bijection
  \(\Gamma(T, V)^{\times} \to \Gamma(C \otimes T, \mathrm{pr}^{-1}V)^{\times}\) of unit
  groups.
\end{corollary}
\begin{proof}
  \leanok
  \uses{pa-thm:isIso_appLE_snd, pa-cor:appLE_snd_injective}
  Injectivity on units follows from injectivity on sections
  (Corollary~\ref{pa-cor:appLE_snd_injective}). For surjectivity, a unit
  \(y \in \Gamma(C \otimes T, \mathrm{pr}^{-1}V)^{\times}\) and its inverse are pullbacks of
  sections \(a, b \in \Gamma(T,V)\) by Theorem~\ref{pa-thm:isIso_appLE_snd}; the products
  \(ab\) and \(ba\) pull back to \(1\), hence equal \(1\) by injectivity, so \(a\) is a unit
  pulling back to \(y\).
\end{proof}

\begin{lemma}[The projection is surjective]
  \label{pa-lem:surjective_snd}
  \lean{AlgebraicGeometry.Over.surjective_snd_left}
  \leanok
  \dcref{custom}

  For every object \(T\) of \(\Over(\Spec k)\), the second projection
  \(C \otimes T \to T\) is surjective on points.
\end{lemma}
\begin{proof}
  \leanok
  The space of \(C\) is irreducible (geometric irreducibility over the one-point base), in
  particular nonempty, and \(\Spec k\) is a single point; so \(C \to \Spec k\) is
  surjective. Surjectivity of morphisms of schemes is stable under base change, and the
  second projection is the base change of \(C \to \Spec k\) along \(T \to \Spec k\).
\end{proof}

\begin{theorem}[Pullback along the projection is injective on Picard groups]
  \label{pa-thm:prPullback_injective}
  \lean{AlgebraicGeometry.Over.prPullback_injective}
  \leanok
  \source{kleiman-picard}

  \dcref{custom}
  For every object \(T\) of \(\Over(\Spec k)\), pullback of Picard classes along the second
  projection, \(\mathrm{pr}^{*} \colon \Pic(T) \to \Pic(C \otimes T)\), is injective.
\end{theorem}
\begin{proof}
  \leanok
  \uses{pa-cor:cechPic_mk_eq_one, pa-lem:pointedCover_affine_refinement, pa-thm:unitsRes_injective,
    pa-cor:appLE_snd_injective, pa-cor:unitsAppLE_snd_bijective, pa-lem:surjective_snd,
    pa-lem:unitsRestrict_trans}
  Pullback is a group homomorphism, so let \(\lambda \in \Pic(T)\) satisfy
  \(\mathrm{pr}^{*}\lambda = 1\). By Lemma~\ref{pa-lem:pointedCover_affine_refinement} and
  invariance of the class under refinement we may represent \(\lambda = [\mathcal U, \gamma]\)
  with \(\mathcal U\) a pointed cover of \(T\) all of whose members \(\mathcal U(t)\) are
  affine. The pullback class is represented on the pulled-back pointed cover
  \(z \mapsto \mathrm{pr}^{-1}\mathcal U(\mathrm{pr}(z))\) by the pulled-back cocycle
  \(\mathrm{pr}^{\sharp}\gamma\), and it is trivial; by
  Corollary~\ref{pa-cor:cechPic_mk_eq_one}, \(\mathrm{pr}^{\sharp}\gamma\) is a coboundary
  \emph{on the pulled-back cover itself}: there are units
  \(w_z \in \struct{C \otimes T}(\mathrm{pr}^{-1}\mathcal U(\mathrm{pr} z))^{\times}\) with
  \[ w_z \cdot \mathrm{pr}^{\sharp}(\gamma_{\mathrm{pr}z, \mathrm{pr}z'}) = w_{z'}
     \qquad\text{on } \mathrm{pr}^{-1}\bigl(\mathcal U(\mathrm{pr}z) \cap
       \mathcal U(\mathrm{pr}z')\bigr). \]
  The projection is surjective (Lemma~\ref{pa-lem:surjective_snd}); choose for each point
  \(t \in T\) a point \(\sigma(t) \in C \otimes T\) over it. Since \(\mathcal U(t)\) is
  affine, the unit \(w_{\sigma(t)} \in
  \struct{C \otimes T}(\mathrm{pr}^{-1}\mathcal U(t))^{\times}\) is the pullback of a unique
  unit \(v_t \in \struct T(\mathcal U(t))^{\times}\)
  (Corollary~\ref{pa-cor:unitsAppLE_snd_bijective}). For any \(t, t'\), applying the injectivity
  of pullback of unit sections on \(\mathcal U(t) \cap \mathcal U(t')\)
  (Corollary~\ref{pa-cor:appLE_snd_injective}, applied to the underlying sections of units) to
  the displayed coboundary relation at \((\sigma(t), \sigma(t'))\) yields
  \(v_t \cdot \gamma_{t t'} = v_{t'}\) on \(\mathcal U(t) \cap \mathcal U(t')\). Hence
  \(\gamma\) is a coboundary, \(\lambda = [\mathcal U, \gamma] = 1\).
\end{proof}

\section{The \v Cech--Picard dictionary}

For a commutative ring \(A\), write \(\operatorname{Pic}(A)\) for its \emph{module} Picard
group: the group, under \(\otimes_A\), of isomorphism classes of invertible \(A\)-modules
--- those \(A\)-modules \(N\) admitting an \(A\)-module \(N'\) with \(N \otimes_A N' \cong A\).
This section identifies, for an affine scheme \(X\) with global ring \(A := \Gamma(X,\struct
X)\), the \v Cech Picard group \(\Pic(X)\) recalled at the head of the previous section ---
the colimit of the groups \(\check H^1(\mathcal U)\) of unit \v Cech \(1\)-cocycles on pointed
covers \(\mathcal U\) of \(X\) --- with \(\operatorname{Pic}(A)\).

For \(f \in A\), write \(D(f) \subseteq X\) for its basic open, the largest open on which the
image of \(f\) in every stalk is a unit; \(D(f) \cap D(g) = D(fg)\), and, \(X\) being affine, a
family \((D(f))_{f \in S}\) covers \(X\) exactly when \(S\) generates the unit ideal of \(A\).
A \emph{finite basic refinement} of a pointed cover \(\mathcal U\) is a finite pointed cover
\(\mathcal W\) refining \(\mathcal U\) all of whose members are basic opens \(D(r_i)\),
\(i \in I\), of elements \(r_i \in A\); since \(X\) is quasi-compact, being affine, every
pointed cover admits one. A unit \v Cech \(1\)-cocycle \(\delta\) on such a \(\mathcal W\)
determines a \emph{descent datum} for the finite faithfully flat algebra
\(A \to \prod_{i \in I} \Gamma(X, D(r_i))\), by twisting the tautological gluing data of the
free rank-one modules \(\Gamma(X,D(r_i))\) by the transition units \(\delta_{ij}\); faithfully
flat descent produces from this datum an \(A\)-module, well defined up to isomorphism, the
\emph{module descended from \(\delta\)}.

\begin{theorem}[Trivializing elements span the unit ideal]
  \label{pa-thm:invertible_span_free}
  \lean{Module.Invertible.span_tensor_free_eq_top}
  \leanok
  \dcref{custom}

  Let \(A\) be a commutative ring and \(N\) an invertible \(A\)-module. The ideal generated by
  \[
    \{\, f \in A : S \otimes_A N \text{ is a free } S\text{-module for every localization } S
    \text{ of } A \text{ away from } f \,\}
  \]
  is the unit ideal of \(A\).
\end{theorem}
\begin{proof}
  \leanok
  An invertible module is finitely presented, and for every prime ideal \(\mathfrak p\) of
  \(A\) the localization of \(N\) at \(\mathfrak p\) is a finite flat module over the local
  ring \(A_{\mathfrak p}\), hence free; by finite presentation, freeness at \(\mathfrak p\)
  spreads out to freeness of \(S \otimes_A N\) for \(S\) a localization away from some single
  element \(f \notin \mathfrak p\). If the ideal generated by all such \(f\) were proper, it
  would be contained in some maximal ideal \(\mathfrak m\), contradicting the existence of a
  trivializing element outside \(\mathfrak m\). Hence the ideal is all of \(A\).
\end{proof}

\begin{lemma}[Existence of a trivializing family]
  \label{pa-lem:trivializingFamily_nonempty}
  \lean{AlgebraicGeometry.Scheme.TrivializingFamily.nonempty}
  \leanok
  \dcref{custom}

  \uses{pa-thm:invertible_span_free}
  Let \(X\) be an affine scheme with global ring \(A = \Gamma(X,\struct X)\) and \(N\) an
  invertible \(A\)-module. There exists a \emph{trivializing family} for \(N\): an assignment,
  to each point \(x \in X\), of a section \(s_x \in A\) with \(x \in D(s_x)\), together with a
  trivialization
  \[
    \Gamma(X,D(s_x)) \otimes_A N \;\cong\; \Gamma(X,D(s_x))
  \]
  of the base change of \(N\) to \(D(s_x)\).
\end{lemma}
\begin{proof}
  \leanok
  By Theorem~\ref{pa-thm:invertible_span_free} the elements \(f \in A\) over which \(N\) becomes
  free after localizing away from \(f\) generate the unit ideal, so their basic opens \(D(f)\)
  cover \(X\). For each point \(x \in X\) choose one such \(f\) with \(x \in D(f)\) and set
  \(s_x := f\), together with the resulting trivialization of \(\Gamma(X,D(s_x)) \otimes_A
  N\).
\end{proof}

\begin{definition}[The cocycle of a trivializing family]
  \label{pa-def:trivializingFamily_cocycle}
  \lean{AlgebraicGeometry.Scheme.TrivializingFamily.cocycle}
  \leanok
  \dcref{custom}

  \uses{pa-lem:trivializingFamily_nonempty}
  Let \(F = (s_x, \mathrm{triv}_x)_{x \in X}\) be a trivializing family for \(N\) as in
  Lemma~\ref{pa-lem:trivializingFamily_nonempty}, and let \(\mathcal U_F(x) := D(s_x)\) be the
  associated pointed cover. For points \(x,y \in X\), pushing \(\mathrm{triv}_x\) and
  \(\mathrm{triv}_y\) forward along the restriction maps to \(D(s_xs_y) = \mathcal U_F(x) \cap
  \mathcal U_F(y)\) produces two trivializations of \(N\) over \(D(s_xs_y)\); their
  \emph{transition unit} is the unique unit \(\gamma_{xy} \in \Gamma(X,D(s_xs_y))^{\times}\)
  with
  \[
    \gamma_{xy} \cdot \bigl(\text{pushforward of } \mathrm{triv}_x\bigr)(1 \otimes n) =
    \bigl(\text{pushforward of } \mathrm{triv}_y\bigr)(1 \otimes n) \qquad \text{for all } n \in
    N.
  \]
  The assignment \((x,y) \mapsto \gamma_{xy}\) is a unit \v Cech \(1\)-cocycle on
  \(\mathcal U_F\).
\end{definition}
\begin{proof}
  \leanok
  The diagonal identity \(\gamma_{xx} = 1\) holds because the two pushforwards of
  \(\mathrm{triv}_x\) to \(D(s_x^2)\) coincide, both being the unique \(A\)-algebra
  homomorphism out of the localization \(\Gamma(X,D(s_x))\) composed with \(\mathrm{triv}_x\).
  For the cocycle identity, push the trivializations \(\mathrm{triv}_x, \mathrm{triv}_y,
  \mathrm{triv}_z\) of a triple \(x,y,z\) forward to the triple overlap \(D(s_xs_ys_z)\): since
  the transition unit between two trivializations composes along a chain of further
  pushforwards --- pushing \(\gamma_{xy}\) and \(\gamma_{yz}\) both forward to
  \(D(s_xs_ys_z)\) compares \(\mathrm{triv}_x\) to \(\mathrm{triv}_y\) and \(\mathrm{triv}_y\)
  to \(\mathrm{triv}_z\) there --- their product equals the transition unit comparing
  \(\mathrm{triv}_x\) directly to \(\mathrm{triv}_z\), which restricts to \(\gamma_{xz}\).
\end{proof}

\begin{theorem}[Descent identification for the trivializing cocycle]
  \label{pa-thm:trivializingFamily_pic_cocycle}
  \lean{AlgebraicGeometry.Scheme.TrivializingFamily.pic_cocycle}
  \leanok
  \dcref{custom}

  \uses{pa-def:trivializingFamily_cocycle}
  In the situation of Definition~\ref{pa-def:trivializingFamily_cocycle}, let \(\mathcal W \le
  \mathcal U_F\) be any finite basic refinement of the trivializing cover. Then the class in
  \(\operatorname{Pic}(A)\) of the module descended from the cocycle \(\gamma|_{\mathcal W}\)
  obtained by restricting the trivializing cocycle to \(\mathcal W\) equals the class of \(N\).
\end{theorem}
\begin{proof}
  \leanok
  Write \(\mathcal W(i) = D(r_i)\), \(i \in I\), with \(r_i\) a multiple of \(s_{p(i)}\) for a
  chosen point \(p(i)\) with \(\mathcal W(i) \subseteq \mathcal U_F(p(i))\). Restricting a
  transition unit is compatible with pushing forward its defining trivializations further, so
  the restricted cocycle \(\gamma|_{\mathcal W}\) is exactly the transition cocycle, on
  \(\mathcal W\), of the trivializations of \(N\) obtained by pushing \(\mathrm{triv}_{p(i)}\)
  forward from \(D(s_{p(i)})\) to \(D(r_i)\) for each \(i\). These pushed-forward
  trivializations assemble into one trivialization of \(N\) over the cover algebra
  \(\prod_i \Gamma(X,D(r_i))\), and it intertwines the descent datum twisted by
  \(\gamma|_{\mathcal W}\) with the canonical descent datum of \(N\) base-changed to the cover
  algebra, by the defining property of the transition units on each factor of the twisted
  double-overlap algebra \(\prod_i \Gamma(X,D(r_i)) \otimes_A \prod_i \Gamma(X,D(r_i))\). Since
  the cover algebra is faithfully flat over \(A\) (a finite family of basic opens covering the
  affine \(X\)), uniqueness of descended modules identifies the module descended from
  \(\gamma|_{\mathcal W}\) with \(N\) itself, so their classes in \(\operatorname{Pic}(A)\)
  agree.
\end{proof}

\begin{definition}[The \v Cech--Picard descent homomorphism]
  \label{pa-def:cechPic_toPic}
  \lean{AlgebraicGeometry.Scheme.CechPic.toPic}
  \leanok
  \dcref{custom}

  For an affine scheme \(X\) with global ring \(A = \Gamma(X,\struct X)\), the homomorphism of
  groups
  \[
    \operatorname{toPic} \colon \Pic(X) \to \operatorname{Pic}(A), \qquad [\mathcal U,\gamma]
    \mapsto \bigl(\text{class of the module descended from } \gamma|_{\mathcal W}\bigr),
  \]
  where \(\mathcal W\) is any finite basic refinement of \(\mathcal U\) and
  \(\gamma|_{\mathcal W}\) is the cocycle obtained by restricting \(\gamma\) along
  \(\mathcal W \le \mathcal U\).
\end{definition}
\begin{proof}
  \leanok
  The class of the descended module does not depend on the choice of finite basic refinement
  \(\mathcal W\) of \(\mathcal U\): any two such refinements are dominated by a common one, along
  which the two descended modules become identified. Nor does it depend on the representative
  \((\mathcal U,\gamma)\) of the class \([\mathcal U,\gamma] \in \Pic(X)\): a further common
  refinement of two cohomologous representing cocycles carries cohomologous restricted
  cocycles, hence isomorphic descended modules. So \(\operatorname{toPic}\) is a well-defined
  function; it is a homomorphism because descending along a common refinement is compatible
  with the pointwise product of cocycles and with tensor product of the descended modules.
\end{proof}

\begin{theorem}[Injectivity of the descent homomorphism]
  \label{pa-thm:cechPic_toPic_injective}
  \lean{AlgebraicGeometry.Scheme.CechPic.toPic_injective}
  \leanok
  \dcref{custom}

  \uses{pa-def:cechPic_toPic, pa-cor:cechPic_mk_eq_one}
  For every affine scheme \(X\), the descent homomorphism \(\operatorname{toPic} \colon
  \Pic(X) \to \operatorname{Pic}(A)\) of Definition~\ref{pa-def:cechPic_toPic} is injective.
\end{theorem}
\begin{proof}
  \leanok
  It suffices to show that a class \([\mathcal U,\gamma]\) with
  \(\operatorname{toPic}[\mathcal U,\gamma] = 1\) is itself trivial. Represent
  \([\mathcal U,\gamma]\) along a finite basic refinement \(\mathcal W \le \mathcal U\) as in
  Definition~\ref{pa-def:cechPic_toPic}: since the module descended from \(\gamma|_{\mathcal W}\)
  is isomorphic to \(A\), its twisted descent datum is isomorphic to the trivial descent datum
  on \(\prod_i \Gamma(X,\mathcal W(i))\), so the transition units of \(\gamma\) restricted to
  \(\mathcal W\) differ from those of the trivial cocycle by an index-wise coboundary. Since the
  opens of \(\mathcal W\) refine those of \(\mathcal U\) and, member by member, cover them, this
  coboundary glues, by the sheaf property of \(\struct X\), to a coboundary for \(\gamma\)
  itself on \(\mathcal U\): \(\gamma = 1\) in \(\check H^1(\mathcal U)\). By
  Corollary~\ref{pa-cor:cechPic_mk_eq_one}, \([\mathcal U,\gamma] = 1\) in \(\Pic(X)\).
\end{proof}

\begin{theorem}[Surjectivity of the descent homomorphism]
  \label{pa-thm:cechPic_toPic_surjective}
  \lean{AlgebraicGeometry.Scheme.CechPic.toPic_surjective}
  \leanok
  \dcref{custom}

  \uses{pa-def:cechPic_toPic, pa-lem:trivializingFamily_nonempty, pa-thm:trivializingFamily_pic_cocycle}
  For every affine scheme \(X\), the descent homomorphism \(\operatorname{toPic} \colon
  \Pic(X) \to \operatorname{Pic}(A)\) is surjective.
\end{theorem}
\begin{proof}
  \leanok
  Let \(c \in \operatorname{Pic}(A)\) be represented by an invertible \(A\)-module \(N\). By
  Lemma~\ref{pa-lem:trivializingFamily_nonempty}, \(N\) admits a trivializing family, giving a
  pointed cover \(\mathcal U_F\) together with the unit \v Cech \(1\)-cocycle \(\gamma_F\) of
  Definition~\ref{pa-def:trivializingFamily_cocycle}. Since \(X\) is affine, hence quasi-compact,
  \(\mathcal U_F\) admits a finite basic refinement \(\mathcal W\). By
  Theorem~\ref{pa-thm:trivializingFamily_pic_cocycle}, \(\operatorname{toPic}[\mathcal
  U_F,\gamma_F] = c\). As \(c\) was arbitrary, \(\operatorname{toPic}\) is surjective.
\end{proof}

\begin{theorem}[The \v Cech--Picard dictionary]
  \label{pa-thm:cechPicEquivPic}
  \lean{AlgebraicGeometry.Scheme.cechPicEquivPic}
  \leanok
  \dcref{custom}

  \uses{pa-thm:cechPic_toPic_injective, pa-thm:cechPic_toPic_surjective}
  For every affine scheme \(X\) with global ring \(A = \Gamma(X,\struct X)\), the descent
  homomorphism of Definition~\ref{pa-def:cechPic_toPic} is an isomorphism of groups
  \[
    \operatorname{toPic} \colon \Pic(X) \;\xrightarrow{\ \sim\ }\; \operatorname{Pic}(A).
  \]
  In particular the definitional \v Cech Picard group \(\Pic(X)\) on pointed Zariski covers of
  \(X\) is canonically isomorphic to mathlib's Picard group of invertible \(A\)-modules.
\end{theorem}
\begin{proof}
  \leanok
  \(\operatorname{toPic}\) is a bijective group homomorphism by
  Theorem~\ref{pa-thm:cechPic_toPic_injective} (injectivity) and
  Theorem~\ref{pa-thm:cechPic_toPic_surjective} (surjectivity), hence an isomorphism.
\end{proof}

\section{The cocycle presentation of the dictionary, and its naturality in the base}

The \v Cech--Picard dictionary of the previous section passes through a \emph{cocycle}
presentation of the descended module, algebraically equivalent to but more computable than
the trivializing-family presentation used there: a descent \(1\)-cocycle is a single unit of
\(B \otimes_A B\), for \(B\) the (faithfully flat) cover algebra, satisfying a normalization and
a cocycle identity. This section records that presentation and its main structural property,
\emph{naturality in the base ring}: the descended module, and hence its class in
\(\operatorname{Pic}\), is compatible with base change of \(A\) along a further ring map
\(A \to A'\). This is the algebraic counterpart of the compatibility, in the test object, of
the functors compared by Kleiman's Comparison Theorem for the relative Picard functor.

Throughout, \(A \to B\) is a homomorphism of commutative rings, exhibiting \(B\) as a
commutative \(A\)-algebra. Write \(\mu \colon B \otimes_A B \to B\) for the multiplication
homomorphism, \(x \otimes y \mapsto xy\), and, in the triple tensor product
\(B \otimes_A (B \otimes_A B)\), write
\[
  \partial_{23}(w) := 1 \otimes w, \qquad
  \partial_{12}(x \otimes y) := x \otimes (y \otimes 1), \qquad
  \partial_{13}(x \otimes y) := x \otimes (1 \otimes y)
\]
for the three \(A\)-algebra homomorphisms \(B \otimes_A B \to B \otimes_A (B \otimes_A B)\)
recording the three faces of the \(2\)-simplex of the cover \(B\) of \(A\).

\begin{definition}[Descent \(1\)-cocycle]
  \label{pa-def:module_descent_cocycle}
  \lean{Module.IsDescentCocycle}
  \leanok
  \dcref{custom}

  A \emph{descent \(1\)-cocycle} relative to \(A \to B\) is a unit \(u \in (B \otimes_A
  B)^{\times}\) such that:
  \begin{enumerate}
    \item (normalization) \(\mu(u) = 1\);
    \item (cocycle identity) \(\partial_{23}(u) \cdot \partial_{12}(u) = \partial_{13}(u)\) in
      \(B \otimes_A (B \otimes_A B)\).
  \end{enumerate}
  When \(B = \prod_i A_{f_i}\) is a finite product of localizations covering \(\Spec A\), this
  is exactly a \v Cech \(1\)-cocycle of units on the cover by the basic opens \(D(f_i)\).
\end{definition}

\begin{definition}[The descended module of a cocycle]
  \label{pa-def:module_descended}
  \lean{Module.IsDescentCocycle.descended}
  \leanok
  \dcref{custom}

  \uses{pa-def:module_descent_cocycle}
  Let \(u\) be a descent \(1\)-cocycle relative to \(A \to B\). The \emph{descended} \(A\)-module
  of \(u\) is the \(A\)-submodule
  \[
    \mathrm{descended}(u) := \{\, x \in B : u \cdot (x \otimes 1) = 1 \otimes x \,\} \subseteq B,
  \]
  the equalizer of the two \(A\)-linear maps \(B \to B \otimes_A B\), \(x \mapsto u \cdot (x
  \otimes 1)\) and \(x \mapsto 1 \otimes x\); equivalently, the descended module of the descent
  datum on the \(B\)-module \(B\) whose coaction is \(x \mapsto u \cdot (x \otimes 1)\)
  (counitality is the normalization condition and coassociativity is the cocycle identity).
\end{definition}

\begin{definition}[The Picard class of a cocycle]
  \label{pa-def:module_picClass}
  \lean{Module.IsDescentCocycle.picClass}
  \leanok
  \dcref{custom}

  \uses{pa-def:module_descended}
  Assume moreover that \(B\) is faithfully flat over \(A\), and let \(u\) be a descent
  \(1\)-cocycle relative to \(A \to B\). The \emph{Picard class} of \(u\) is the class
  \(\mathrm{picClass}(u) \in \operatorname{Pic}(A)\) of the descended module
  \(\mathrm{descended}(u)\).
\end{definition}
\begin{proof}
  \leanok
  \(B\), as a module over itself, is free of rank one, hence invertible, and the coaction \(x
  \mapsto u \cdot (x \otimes 1)\) is a descent datum on \(B\) relative to \(A \to B\)
  (Definition~\ref{pa-def:module_descended}) whose descended module is \(\mathrm{descended}(u)\) by
  definition. Since invertibility of modules descends along faithfully flat ring maps together
  with a descent datum, \(\mathrm{descended}(u)\) is an invertible \(A\)-module, so it has a
  well-defined class in \(\operatorname{Pic}(A)\).
\end{proof}

\begin{definition}[Base change of the tensor square and tensor cube of a cover]
  \label{pa-def:module_tensorSqBaseChange}
  \lean{Module.tensorSqBaseChange, Module.tensorCubeBaseChange}
  \leanok
  \dcref{custom}

  Let \(A \to A'\) be a homomorphism of commutative rings and \(B\) a commutative
  \(A\)-algebra. The \emph{base change of the tensor square} of the cover \(B\) along
  \(A \to A'\) is the canonical \(A\)-algebra homomorphism
  \[
    B \otimes_A B \;\longrightarrow\; (A' \otimes_A B) \otimes_{A'} (A' \otimes_A B), \qquad
    x \otimes y \mapsto (1 \otimes x) \otimes (1 \otimes y).
  \]
  Likewise, the \emph{base change of the tensor cube} is the canonical \(A\)-algebra
  homomorphism on triple tensor products
  \[
    B \otimes_A (B \otimes_A B) \;\longrightarrow\; (A' \otimes_A B) \otimes_{A'} \bigl( (A'
    \otimes_A B) \otimes_{A'} (A' \otimes_A B) \bigr), \qquad
    x \otimes (y \otimes z) \mapsto (1 \otimes x) \otimes \bigl( (1 \otimes y) \otimes (1
    \otimes z) \bigr).
  \]
\end{definition}

\begin{theorem}[Base change of a descent cocycle]
  \label{pa-thm:module_descentCocycle_baseChange}
  \lean{Module.IsDescentCocycle.baseChange}
  \leanok
  \dcref{custom}

  \uses{pa-def:module_descent_cocycle, pa-def:module_tensorSqBaseChange}
  Let \(A \to A'\) and \(A \to B\) be homomorphisms of commutative rings and let \(u\) be a
  descent \(1\)-cocycle relative to \(A \to B\). Then the image of \(u\) under the base-change
  homomorphism of Definition~\ref{pa-def:module_tensorSqBaseChange} is a descent \(1\)-cocycle
  relative to \(A' \to A' \otimes_A B\).
\end{theorem}
\begin{proof}
  \leanok
  Write \(\varphi \colon B \otimes_A B \to (A' \otimes_A B) \otimes_{A'} (A' \otimes_A B)\) for
  the base-change homomorphism and \(\psi\) for its analogue on triple tensor products
  (Definition~\ref{pa-def:module_tensorSqBaseChange}). On elementary tensors,
  \[
    \mu'(\varphi(x \otimes y)) = (1 \otimes x)(1 \otimes y) = 1 \otimes xy = \iota(\mu(x \otimes
    y)),
  \]
  where \(\iota \colon B \to A' \otimes_A B\), \(b \mapsto 1 \otimes b\), is the canonical
  inclusion and \(\mu'\) is the multiplication of \(A' \otimes_A B\); and, using
  \(1 \otimes_{A'} 1 = 1\) under the identification \(A' \otimes_A B \cong A' \otimes_{A'} (A'
  \otimes_A B)\),
  \[
    \partial'_{23}(\varphi(x \otimes y)) = \psi(\partial_{23}(x \otimes y)), \qquad
    \partial'_{12}(\varphi(x \otimes y)) = \psi(\partial_{12}(x \otimes y)), \qquad
    \partial'_{13}(\varphi(x \otimes y)) = \psi(\partial_{13}(x \otimes y)),
  \]
  where \(\partial'_{23}, \partial'_{12}, \partial'_{13}\) are the face maps of \(A' \to A'
  \otimes_A B\). Both sides of each displayed identity are \(A\)-bilinear in \(x, y\), so the
  identities extend from elementary tensors to all of \(B \otimes_A B\) by linearity.

  Applying \(\mu'\) to \(\varphi(u)\) and using the first identity together with the
  normalization \(\mu(u) = 1\) gives \(\mu'(\varphi(u)) = \iota(\mu(u)) = \iota(1) = 1\), the
  normalization of \(\varphi(u)\). Applying \(\psi\) to the cocycle identity \(\partial_{23}(u)
  \cdot \partial_{12}(u) = \partial_{13}(u)\) and using the remaining three identities gives
  \[
    \partial'_{23}(\varphi(u)) \cdot \partial'_{12}(\varphi(u)) = \psi\bigl(\partial_{23}(u)
    \cdot \partial_{12}(u)\bigr) = \psi(\partial_{13}(u)) = \partial'_{13}(\varphi(u)),
  \]
  the cocycle identity for \(\varphi(u)\). Hence \(\varphi(u)\) is a descent \(1\)-cocycle
  relative to \(A' \to A' \otimes_A B\).
\end{proof}

\begin{theorem}[Base change of the descended module]
  \label{pa-thm:module_descendedBaseChangeEquiv}
  \lean{Module.descendedBaseChangeEquiv}
  \leanok
  \dcref{custom}

  \uses{pa-thm:module_descentCocycle_baseChange, pa-def:module_descended}
  Let \(A \to A'\) and \(A \to B\) be homomorphisms of commutative rings with \(B\) faithfully
  flat over \(A\), and let \(u\) be a descent \(1\)-cocycle relative to \(A \to B\), with base
  change \(u' := \varphi(u)\) relative to \(A' \to A' \otimes_A B\)
  (Theorem~\ref{pa-thm:module_descentCocycle_baseChange}). Then there is a canonical isomorphism of
  \(A'\)-modules
  \[
    A' \otimes_A \mathrm{descended}(u) \;\cong\; \mathrm{descended}(u').
  \]
\end{theorem}
\begin{proof}
  \leanok
  Write \(B' := A' \otimes_A B\); since faithful flatness is preserved by base change, \(B'\) is
  faithfully flat over \(A'\), so \(\mathrm{descended}(u')\) is the descended module of the
  tautological descent datum on the \(B'\)-module \(B'\) twisted by \(u'\)
  (Definition~\ref{pa-def:module_descended}), and by the uniqueness property of descended modules
  it suffices to produce a \(B'\)-linear isomorphism
  \(B' \otimes_{A'} \bigl(A' \otimes_A \mathrm{descended}(u)\bigr) \cong B'\) intertwining the
  base change of the coaction \(x \mapsto u \cdot (x \otimes 1)\) on \(\mathrm{descended}(u)\)
  with the coaction \(y \mapsto u' \cdot (y \otimes 1)\) on \(B'\).

  Cancelling the iterated base change identifies \(B'\)-modules
  \[
    B' \otimes_{A'} \bigl(A' \otimes_A \mathrm{descended}(u)\bigr) \;\cong\; B' \otimes_A
    \mathrm{descended}(u) \;\cong\; B' \otimes_B \bigl(B \otimes_A \mathrm{descended}(u)\bigr),
  \]
  and tensoring the descent isomorphism \(B \otimes_A \mathrm{descended}(u) \cong B\) attached to
  \(u\) with \(B'\) over \(B\) turns the right side into \(B' \otimes_B B \cong B'\). On
  elementary tensors, the resulting isomorphism sends \(b' \otimes (a' \otimes n)\), for
  \(b' \in B'\), \(a' \in A'\), \(n \in \mathrm{descended}(u)\), to \((a' b') \cdot n\) (viewing
  \(n \in \mathrm{descended}(u) \subseteq B\) inside \(B'\) via \(\iota\)); using the membership
  identity \(u \cdot (n \otimes 1) = 1 \otimes n\), this identification carries the base-changed
  coaction on the left to the coaction \(y \mapsto u' \cdot (y \otimes 1)\) on the right. Hence
  the isomorphism intertwines the two coactions, and uniqueness of descended modules produces
  the canonical isomorphism \(A' \otimes_A \mathrm{descended}(u) \cong \mathrm{descended}(u')\).
\end{proof}

\begin{theorem}[Naturality of the Picard class in the base ring]
  \label{pa-thm:module_picClass_baseChange}
  \lean{Module.IsDescentCocycle.picClass_baseChange}
  \leanok
  \source{kleiman-picard}

  \uses{pa-thm:module_descendedBaseChangeEquiv, pa-def:module_picClass}
  \dcref{custom}
  Let \(A \to A'\) and \(A \to B\) be homomorphisms of commutative rings with \(B\) faithfully
  flat over \(A\) and \(A' \otimes_A B\) faithfully flat over \(A'\), and let \(u\) be a descent
  \(1\)-cocycle relative to \(A \to B\) with base change \(u'\) relative to \(A' \to A'
  \otimes_A B\) (Theorem~\ref{pa-thm:module_descentCocycle_baseChange}). Then
  \[
    \mathrm{picClass}(u') = \bigl[A' \otimes_A \mathrm{descended}(u)\bigr] \quad \text{in }
    \operatorname{Pic}(A'),
  \]
  that is, the Picard class of the base-changed cocycle is the image of \(\mathrm{picClass}(u)\)
  under the base-change homomorphism \(\operatorname{Pic}(A) \to \operatorname{Pic}(A')\),
  \([N] \mapsto [A' \otimes_A N]\).
\end{theorem}
\begin{proof}
  \leanok
  By Theorem~\ref{pa-thm:module_descendedBaseChangeEquiv}, \(A' \otimes_A \mathrm{descended}(u)
  \cong \mathrm{descended}(u')\) as \(A'\)-modules, so these two invertible \(A'\)-modules have
  the same class in \(\operatorname{Pic}(A')\). The left side represents the image of
  \(\mathrm{picClass}(u) = [\mathrm{descended}(u)]\) under the base-change homomorphism
  \([N] \mapsto [A' \otimes_A N]\), and the right side represents \(\mathrm{picClass}(u')\) by
  definition; hence \(\mathrm{picClass}(u') = \bigl[A' \otimes_A \mathrm{descended}(u)\bigr]\)
  is the image of \(\mathrm{picClass}(u)\) under this homomorphism.
\end{proof}

\section{Base change of a finite localization cover cocycle}
\label{pa-sec:awayCover_baseChange}

Fix a homomorphism of commutative rings \(A \to A'\) and a finite family \(f \colon \iota \to
A\).  For each index we fix localization models
\[
  S_i = A_{f_i}, \qquad T_{ij} = A_{f_i f_j}, \qquad W_{ijk} = A_{f_i f_j f_k},
\]
and the analogous \(A'\)-side models \(S'_i, T'_{ij}, W'_{ijk}\) for the pushed family
\(f'_i := \varphi(f_i)\), where \(\varphi \colon A \to A'\) is the structure map; each primed
model also carries the \(A\)-algebra structure of the tower \(A \to A' \to \cdot\).  This
section transports the whole \v Cech picture of the cover \(f\) to the cover \(f'\).

\begin{definition}[Base-change maps between localization models]
  \label{pa-def:awayCover_baseChangeMaps}
  \lean{IsLocalization.AwayCover.mapAway, IsLocalization.AwayCover.mapOverlap,
    IsLocalization.AwayCover.mapTriple}
  \leanok
  \dcref{custom}

  Because \(f_i\) becomes invertible in \(S'_i\) through the tower \(A \to A' \to S'_i\),
  there is a unique \(A\)-algebra homomorphism \(S_i = A_{f_i} \to S'_i\); denote it by
  \(\mathrm{mapAway}_i\).  Likewise, \(f_i f_j\) is invertible in \(T'_{ij}\) and
  \(f_i f_j f_k\) in \(W'_{ijk}\), giving unique \(A\)-algebra homomorphisms
  \[
    \mathrm{mapOverlap}_{ij} \colon T_{ij} \to T'_{ij}, \qquad
    \mathrm{mapTriple}_{ijk} \colon W_{ijk} \to W'_{ijk}.
  \]
\end{definition}

\begin{theorem}[Base change of a cover cocycle]
  \label{pa-thm:awayCover_isCoverCocycle_baseChange}
  \lean{IsLocalization.AwayCover.IsCoverCocycle.baseChange}
  \leanok
  \dcref{custom}

  \uses{pa-def:awayCover_baseChangeMaps}
  Let \(\gamma = (\gamma_{ij})\) with \(\gamma_{ij} \in T_{ij}^{\times}\) be a cover cocycle
  for \(f\): its diagonal restrictions \(\gamma_{ii}\) map to \(1\) in \(T_{ii}\), and on
  triple overlaps the cocycle identity
  \(\partial_{23}\gamma_{jk} \cdot \partial_{12}\gamma_{ij} = \partial_{13}\gamma_{ik}\) holds
  in \(W_{ijk}\).  Then the pushed family
  \(\gamma'_{ij} := \mathrm{mapOverlap}_{ij}(\gamma_{ij}) \in (T'_{ij})^{\times}\) is a cover
  cocycle for \(f'\).
\end{theorem}
\begin{proof}
  \leanok
  Each restriction map in the \v Cech picture — the diagonal identification \(T_{ii} \to
  T_{ij}\) and the three face maps \(\partial_{12}, \partial_{13}, \partial_{23}\) from
  \(T\)-overlaps to \(W\)-triples — fits into a square with its primed analogue and the maps of
  Definition~\ref{pa-def:awayCover_baseChangeMaps}.  Both composites in each such square are
  \(A\)-algebra homomorphisms out of a localization of \(A\), which admits at most one such
  homomorphism to any fixed target; hence the squares commute.  Applying the diagonal square to
  \(\gamma_{ii}\) and using \(\gamma_{ii} \mapsto 1\) shows \(\gamma'_{ii} \mapsto 1\); applying
  the three face squares to the cocycle identity for \(\gamma\) and using multiplicativity of
  \(\mathrm{mapTriple}_{ijk}\) yields the cocycle identity for \(\gamma'\).  Thus \(\gamma'\) is
  a cover cocycle for \(f'\).
\end{proof}

\begin{definition}[Base-change map of cover algebras]
  \label{pa-def:awayCover_piBaseChangeLift}
  \lean{IsLocalization.AwayCover.piBaseChangeLift}
  \leanok
  \dcref{custom}

  \uses{pa-def:awayCover_baseChangeMaps}
  The maps \(\mathrm{mapAway}_i\) assemble into a componentwise \(A\)-algebra homomorphism
  \(\prod_i S_i \to \prod_i S'_i\), and hence, by the universal property of base change, into a
  canonical \(A'\)-algebra homomorphism
  \[
    j \colon A' \otimes_A \Bigl(\prod_i S_i\Bigr) \longrightarrow \prod_i S'_i, \qquad
    a' \otimes s \mapsto a' \cdot \bigl(\mathrm{mapAway}_i(s_i)\bigr)_i,
  \]
  the \emph{base-change map of cover algebras}.
\end{definition}

\begin{theorem}[Base change of the descent unit]
  \label{pa-thm:awayCover_cocycleUnit_baseChange}
  \lean{IsLocalization.AwayCover.cocycleUnit_baseChange}
  \leanok
  \dcref{custom}

  \uses{pa-def:awayCover_baseChangeMaps, pa-def:awayCover_piBaseChangeLift,
    pa-thm:awayCover_isCoverCocycle_baseChange, pa-def:module_tensorSqBaseChange}
  Under the canonical identification of the double product \(\prod_{i,j} T_{ij}\) with the
  tensor square \(\bigl(\prod_i S_i\bigr) \otimes_A \bigl(\prod_i S_i\bigr)\), a cover cocycle
  \(\gamma\) determines a descent unit \(c(\gamma) \in \bigl((\prod_i S_i) \otimes_A (\prod_i
  S_i)\bigr)^{\times}\).  Then the image of \(c(\gamma)\) under the base change of the tensor
  square (Definition~\ref{pa-def:module_tensorSqBaseChange}) along \(A \to A'\), followed by the
  tensor square of the base-change map \(j\)
  (Definition~\ref{pa-def:awayCover_piBaseChangeLift}), is the descent unit \(c(\gamma')\) of the
  pushed cocycle \(\gamma'\) (Theorem~\ref{pa-thm:awayCover_isCoverCocycle_baseChange}):
  \[
    (j \otimes j)\bigl(\mathrm{tensorSqBaseChange}(c(\gamma))\bigr) = c(\gamma')
    \quad\text{in }\Bigl(\bigl(\textstyle\prod_i S'_i\bigr) \otimes_{A'} \bigl(\prod_i
    S'_i\bigr)\Bigr)^{\times}.
  \]
\end{theorem}
\begin{proof}
  \leanok
  Both sides are units, so it suffices to compare their values under the identification of the
  tensor square \(\bigl(\prod_i S'_i\bigr) \otimes_{A'} \bigl(\prod_i S'_i\bigr)\) with the
  double product \(\prod_{i,j} T'_{ij}\).  The composite
  \[
    \Bigl(\prod_{i,j} T_{ij}\Bigr) \longrightarrow \bigl(\textstyle\prod_i S_i\bigr) \otimes_A
    \bigl(\prod_i S_i\bigr) \xrightarrow{\ \mathrm{tensorSqBaseChange}\ } \bigl(\textstyle\prod_i
    S'_i\bigr) \otimes_{A'} \bigl(\prod_i S'_i\bigr) \xrightarrow{\ j \otimes j\ }
    \Bigl(\prod_{i,j} T'_{ij}\Bigr)
  \]
  is an \(A\)-algebra homomorphism out of the product of localizations \(\prod_{i,j} T_{ij}\);
  by the pi-extensionality principle it is determined by its values on the coordinate
  idempotents \(e_{ij}\).  A direct computation on \(e_{ij}\) — unfolding
  \(\mathrm{tensorSqBaseChange}\), the tensor product \(j \otimes j\), and the componentwise
  maps \(\mathrm{mapAway}\) — shows the composite agrees coordinatewise with
  \(\bigl(\mathrm{mapOverlap}_{ij}\bigr)_{i,j}\).  Applying this identity of homomorphisms to the
  unit representing \(\gamma\) gives, coordinatewise, \(\gamma'_{ij} =
  \mathrm{mapOverlap}_{ij}(\gamma_{ij})\), i.e.\ the descent unit \(c(\gamma')\).
\end{proof}

\begin{lemma}[Base change preserves covering families]
  \label{pa-lem:awayCover_span_range_baseChange}
  \lean{IsLocalization.AwayCover.span_range_algebraMap_eq_top}
  \leanok
  \dcref{custom}

  If the family \(f\) is covering, \(\mathrm{span}\,\{f_i\} = A\), then the pushed family
  \(f'_i = \varphi(f_i)\) is covering as well, \(\mathrm{span}\,\{f'_i\} = A'\).
\end{lemma}
\begin{proof}
  \leanok
  The ideal of \(A'\) generated by the pushed family is the extension along \(\varphi \colon A
  \to A'\) of the ideal generated by \(f\), that is \(\mathrm{span}\,\{\varphi(f_i)\}\) is the
  image ideal \(\varphi_{*}\bigl(\mathrm{span}\,\{f_i\}\bigr)\).  By hypothesis
  \(\mathrm{span}\,\{f_i\} = A\) is the unit ideal, and the extension of the unit ideal is again
  the unit ideal \(A'\); hence \(\mathrm{span}\,\{f'_i\} = A'\).
\end{proof}

\begin{lemma}[Base change of a covering product of localizations is faithfully flat]
  \label{pa-lem:awayCover_faithfullyFlat_tensor}
  \lean{IsLocalization.AwayCover.faithfullyFlat_tensor}
  \leanok
  \dcref{custom}

  \uses{pa-lem:awayCover_span_range_baseChange}
  If \(f\) is a covering family, \(\mathrm{span}\,\{f_i\} = A\), then the base change \(A'
  \otimes_A \bigl(\prod_i S_i\bigr)\) of the cover algebra is faithfully flat as an
  \(A'\)-module.
\end{lemma}
\begin{proof}
  \leanok
  For each index the base change \(A' \otimes_A S_i\) of the localization \(S_i = A_{f_i}\) is
  the localization \(A'_{f'_i}\) of \(A'\) at the pushed element \(f'_i = \varphi(f_i)\).  Since
  the pushed family \(f'\) is again covering
  (Lemma~\ref{pa-lem:awayCover_span_range_baseChange}), the finite product \(\prod_i (A' \otimes_A
  S_i)\) of localizations at a covering family is faithfully flat over \(A'\).  Finally base
  change commutes with finite products, \(A' \otimes_A \bigl(\prod_i S_i\bigr) \cong \prod_i (A'
  \otimes_A S_i)\) as \(A'\)-modules, and faithful flatness transports along this isomorphism.
\end{proof}

\begin{theorem}[Naturality of the descent Picard class in the base ring]
  \label{pa-thm:awayCover_pic_baseChange}
  \lean{IsLocalization.AwayCover.pic_baseChange}
  \leanok
  \source{kleiman-picard}

  \uses{pa-thm:awayCover_cocycleUnit_baseChange, pa-thm:awayCover_isCoverCocycle_baseChange,
    pa-thm:module_picClass_baseChange, pa-def:module_picClass,
    pa-lem:awayCover_span_range_baseChange, pa-lem:awayCover_faithfullyFlat_tensor}
  \dcref{custom}
  Assume the localization covers are faithfully flat, \(\prod_i S_i\) over \(A\) and \(\prod_i
  S'_i\) over \(A'\), and that \(f\) is a covering family, \(\mathrm{span}\,\{f_i\} = A\).  Then
  the Picard class of the descended module of the pushed cocycle \(\gamma'\) is the image of the
  Picard class of \(\gamma\) under the base-change homomorphism \(\Pic(A) \to \Pic(A')\),
  \([N] \mapsto [A' \otimes_A N]\):
  \[
    \mathrm{picClass}(c(\gamma')) = \bigl[A' \otimes_A \mathrm{descended}(c(\gamma))\bigr]
    \quad\text{in }\Pic(A').
  \]
\end{theorem}
\begin{proof}
  \leanok
  Since \(\mathrm{span}\,\{f_i\} = A\), the pushed family \(f'\) again spans \(A'\) (the ideal it
  generates is the image of the unit ideal under \(A \to A'\)), so \(\prod_i (A' \otimes_A S_i)
  \cong A' \otimes_A \prod_i S_i\) is faithfully flat over \(A'\), each factor being a
  localization \(A'_{f'_i}\).  By Theorem~\ref{pa-thm:awayCover_cocycleUnit_baseChange} the descent
  unit \(c(\gamma')\) is the image of the base-changed unit
  \(\mathrm{tensorSqBaseChange}(c(\gamma))\) under the algebra map \(j\) of cover algebras
  (Definition~\ref{pa-def:awayCover_piBaseChangeLift}); as the Picard class of a descent unit is
  unchanged by transport along an algebra map of the cover, \(\mathrm{picClass}(c(\gamma'))\)
  equals the Picard class of \(\mathrm{tensorSqBaseChange}(c(\gamma))\), which is the descent
  unit of the base change of the module cocycle \(c(\gamma)\).  By
  Theorem~\ref{pa-thm:module_picClass_baseChange} this last class is \(\bigl[A' \otimes_A
  \mathrm{descended}(c(\gamma))\bigr]\), the image of
  \(\mathrm{picClass}(c(\gamma)) = [\mathrm{descended}(c(\gamma))]\)
  (Definition~\ref{pa-def:module_picClass}) under \([N] \mapsto [A' \otimes_A N]\).
\end{proof}

\section{The projection-descent equivalence for unit sections}
\label{pa-sec:projectionUnits}

Throughout this section \(C\) is a curve — proper, geometrically irreducible and geometrically
reduced — over a field \(k\), and \(T\) is a test object of \(\Over(\Spec k)\).  Write
\(\mathrm{pr} \colon (C \otimes T)_{\mathrm{left}} \to T_{\mathrm{left}}\) for the second
projection.

\begin{definition}[The projection-descent equivalence for unit sections]
  \label{pa-def:unitsSndEquiv}
  \lean{AlgebraicGeometry.Over.unitsSndEquiv}
  \leanok
  \dcref{custom}

  \uses{pa-cor:unitsAppLE_snd_bijective}
  For an affine open \(V \subseteq T_{\mathrm{left}}\), pullback of unit sections along the
  second projection is a group isomorphism
  \[
    \mathrm{unitsSndEquiv} \colon \Gamma(T_{\mathrm{left}}, V)^{\times} \;\xrightarrow{\ \sim\
    }\; \Gamma\bigl((C \otimes T)_{\mathrm{left}}, \mathrm{pr}^{-1} V\bigr)^{\times},
  \]
  the packaging as a multiplicative equivalence of the bijection
  (Corollary~\ref{pa-cor:unitsAppLE_snd_bijective}) that pullback of unit sections over an affine
  open is bijective.  Its inverse is the \emph{descent} of a unit on the product to a unit on
  the base.
\end{definition}
\begin{proof}
  \leanok
  Pullback of unit sections is a group homomorphism, and it is bijective over the affine open
  \(V\) by Corollary~\ref{pa-cor:unitsAppLE_snd_bijective}; a bijective group homomorphism is a
  group isomorphism.
\end{proof}

\begin{lemma}[Uniqueness of the descended unit]
  \label{pa-lem:unitsSndEquiv_symm_eq_of_unitsAppLE}
  \lean{AlgebraicGeometry.Over.unitsSndEquiv_symm_eq_of_unitsAppLE}
  \leanok
  \dcref{custom}

  \uses{pa-def:unitsSndEquiv}
  Let \(V \subseteq T_{\mathrm{left}}\) be an affine open, \(v \in \Gamma(T_{\mathrm{left}},
  V)^{\times}\) and \(w \in \Gamma((C \otimes T)_{\mathrm{left}}, \mathrm{pr}^{-1}
  V)^{\times}\).  If \(v\) pulls back to \(w\) along the projection, then \(v\) is the descent
  of \(w\), that is \(v = \mathrm{unitsSndEquiv}^{-1}(w)\).
\end{lemma}
\begin{proof}
  \leanok
  By Definition~\ref{pa-def:unitsSndEquiv}, \(\mathrm{unitsSndEquiv}(v)\) is the pullback of \(v\),
  which by hypothesis is \(w\); applying the inverse isomorphism gives \(v =
  \mathrm{unitsSndEquiv}^{-1}(w)\).
\end{proof}

\begin{lemma}[Naturality of the descent equivalence in the affine open]
  \label{pa-lem:unitsSndEquiv_unitsRestrict}
  \lean{AlgebraicGeometry.Over.unitsSndEquiv_unitsRestrict,
    AlgebraicGeometry.Over.unitsSndEquiv_symm_unitsRestrict}
  \leanok
  \dcref{custom}

  \uses{pa-def:unitsSndEquiv}
  For affine opens \(V' \subseteq V \subseteq T_{\mathrm{left}}\), pullback of unit sections
  commutes with restriction:
  \[
    \mathrm{unitsSndEquiv}\bigl(v|_{V'}\bigr) = \bigl(\mathrm{unitsSndEquiv}(v)\bigr)
    \big|_{\mathrm{pr}^{-1} V'}, \qquad v \in \Gamma(T_{\mathrm{left}}, V)^{\times},
  \]
  and dually the descent commutes with restriction.
\end{lemma}
\begin{proof}
  \leanok
  Both sides are the value of \(\mathrm{pr}\) on unit sections composed with restriction; the two
  restriction squares of the structure sheaf commute, so the pullback of the restriction equals
  the restriction of the pullback.  Applying \(\mathrm{unitsSndEquiv}^{-1}\) and using the
  round-trips (Definition~\ref{pa-def:unitsSndEquiv}) gives the dual identity.
\end{proof}

\begin{lemma}[The base square of the second projection]
  \label{pa-lem:snd_left_naturality}
  \lean{AlgebraicGeometry.Over.snd_left_naturality,
    AlgebraicGeometry.Over.snd_left_preimage_naturality}
  \leanok
  \dcref{custom}

  For a morphism \(g \colon T' \to T\) of test objects, whiskering \(C \lhd g\) on the product
  intertwines the two projections:
  \[
    (C \lhd g)_{\mathrm{left}} \circ (\mathrm{snd}_{C,T})_{\mathrm{left}}
      = (\mathrm{snd}_{C,T'})_{\mathrm{left}} \circ g_{\mathrm{left}}.
  \]
  Consequently the two legs of this square have equal open preimages: for every open \(V
  \subseteq T_{\mathrm{left}}\), \((C \lhd g)^{-1}\mathrm{pr}^{-1} V = \mathrm{pr}'^{-1}
  g_{\mathrm{left}}^{-1} V\).
\end{lemma}
\begin{proof}
  \leanok
  The identity is the projection-compatibility of the cartesian monoidal structure,
  \((C \lhd g) \circ \mathrm{snd}_{C,T} = \mathrm{snd}_{C,T'} \circ g\), transported through the
  forgetful functor \(\Over(\Spec k) \to \mathrm{Sch}\).  Taking open preimages of \(V\) through
  both composites and using functoriality of preimage gives the second equality.
\end{proof}

\begin{theorem}[Naturality of the descent equivalence in the test object]
  \label{pa-thm:unitsSndEquiv_naturality}
  \lean{AlgebraicGeometry.Over.unitsSndEquiv_naturality}
  \leanok
  \dcref{custom}

  \uses{pa-def:unitsSndEquiv, pa-lem:snd_left_naturality}
  Let \(g \colon T' \to T\) be a morphism of test objects and \(V \subseteq T_{\mathrm{left}}\)
  an affine open whose preimage \(g_{\mathrm{left}}^{-1} V \subseteq T'_{\mathrm{left}}\) is also
  affine.  For \(v \in \Gamma(T_{\mathrm{left}}, V)^{\times}\), descending \(v\) on \(T\) and
  then pulling back along \(C \lhd g\) equals first pulling \(v\) back along \(g\) and then
  descending on \(T'\); both are units over \(\mathrm{pr}'^{-1}(g_{\mathrm{left}}^{-1} V)\).
\end{theorem}
\begin{proof}
  \leanok
  Unfolding both sides as pullbacks of \(v\) along the two legs of the base square of
  Lemma~\ref{pa-lem:snd_left_naturality}, the two composite pullbacks agree because the two legs are
  equal morphisms and the source opens are identified by the preimage identity of that lemma.
  The equality of pullbacks along equal morphisms then closes the goal.
\end{proof}

\section{Descent in stages for unit \(1\)-cocycles}
\label{pa-sec:unitDescentComposite}

Let \(A \to B \to P\) be a tower of commutative rings.  This section relates \(A\)-descent
along the composite \(A \to P\) with \(B\)-descent along \(B \to P\), the cover being factored
through the intermediate ring \(B\).

\begin{definition}[The collapse maps on tensor squares and cubes]
  \label{pa-def:module_tensorCollapse}
  \lean{Module.tensorCollapse, Module.tensorCubeCollapse}
  \leanok
  \dcref{custom}

  The \emph{collapse map} is the canonical \(A\)-algebra homomorphism
  \[
    \pi \colon P \otimes_A P \longrightarrow P \otimes_B P, \qquad x \otimes y \mapsto x \otimes
    y,
  \]
  identifying the two \(A\)-bilinear factors as \(B\)-bilinear.  Its analogue on tensor cubes is
  the \(A\)-algebra homomorphism
  \[
    \pi_3 \colon P \otimes_A (P \otimes_A P) \longrightarrow P \otimes_B (P \otimes_B P), \qquad
    x \otimes (y \otimes z) \mapsto x \otimes (y \otimes z),
  \]
  the receptacle for the collapse of the three face maps.
\end{definition}

\begin{theorem}[Collapse of a descent cocycle]
  \label{pa-thm:module_descentCocycle_collapse}
  \lean{Module.IsDescentCocycle.collapse}
  \leanok
  \dcref{custom}

  \uses{pa-def:module_descent_cocycle, pa-def:module_tensorCollapse}
  Let \(v \in (P \otimes_A P)^{\times}\) be an \(A\)-descent \(1\)-cocycle relative to \(A \to
  P\).  Then its collapse \(\pi(v)\) is a \(B\)-descent \(1\)-cocycle relative to \(B \to P\).
\end{theorem}
\begin{proof}
  \leanok
  On elementary tensors, the multiplication \(\mu_B\) of \(P \otimes_B P\) satisfies \(\mu_B(\pi(x
  \otimes y)) = xy = \mu_A(x \otimes y)\), and each \(B\)-face map factors through the collapse:
  \(\partial^B_{12}(\pi(w)) = \pi_3(\partial^A_{12}(w))\) and likewise for \(\partial_{13},
  \partial_{23}\); both sides of each identity are \(A\)-bilinear, so the identities hold on all
  of \(P \otimes_A P\).  Applying \(\mu_B\) to \(\pi(v)\) and using the normalization
  \(\mu_A(v) = 1\) gives \(\mu_B(\pi(v)) = 1\); applying \(\pi_3\) to the cocycle identity
  \(\partial^A_{23}(v) \cdot \partial^A_{12}(v) = \partial^A_{13}(v)\) and using
  multiplicativity of \(\pi_3\) gives \(\partial^B_{23}(\pi(v)) \cdot \partial^B_{12}(\pi(v)) =
  \partial^B_{13}(\pi(v))\).  Hence \(\pi(v)\) is a \(B\)-descent \(1\)-cocycle.
\end{proof}

\begin{lemma}[The descended module lies in the collapsed descended module]
  \label{pa-lem:module_descended_le_collapse}
  \lean{Module.descended_le_descended_collapse}
  \leanok
  \dcref{custom}

  \uses{pa-def:module_descended, pa-thm:module_descentCocycle_collapse}
  For an \(A\)-descent \(1\)-cocycle \(v\), the \(A\)-descended module of \(v\) is contained,
  as a subset of \(P\), in the \(B\)-descended module of the collapse \(\pi(v)\).
\end{lemma}
\begin{proof}
  \leanok
  An element \(m \in P\) lies in the \(A\)-descended module of \(v\) iff \(v \cdot (m \otimes_A
  1) = 1 \otimes_A m\) in \(P \otimes_A P\).  Applying the collapse \(\pi\) and using
  \(\pi(m \otimes_A 1) = m \otimes_B 1\), \(\pi(1 \otimes_A m) = 1 \otimes_B m\) and
  multiplicativity gives \(\pi(v) \cdot (m \otimes_B 1) = 1 \otimes_B m\), i.e.\ \(m\) lies in
  the \(B\)-descended module of \(\pi(v)\).
\end{proof}

\begin{theorem}[Descent in stages for the descended module]
  \label{pa-thm:module_descendedCollapseEquiv}
  \lean{Module.descendedCollapseEquiv}
  \leanok
  \dcref{custom}

  \uses{pa-thm:module_descentCocycle_collapse, pa-def:module_descended}
  Assume \(A \to B\) and \(B \to P\) are faithfully flat (so \(A \to P\) is faithfully flat as
  well).  For an \(A\)-descent \(1\)-cocycle \(v\), the base change to \(B\) of the
  \(A\)-descended module of \(v\) is canonically isomorphic to the \(B\)-descended module of the
  collapse:
  \[
    B \otimes_A \mathrm{descended}(v) \;\cong\; \mathrm{descended}(\pi(v))
    \qquad (B\text{-linear}),\quad b \otimes m \mapsto b \cdot m.
  \]
\end{theorem}
\begin{proof}
  \leanok
  As \(B \to P\) is faithfully flat, \(\mathrm{descended}(\pi(v))\) is recognized by the
  uniqueness property of descended modules from any \(B\)-module \(N\) equipped with a
  \(P\)-linear isomorphism \(P \otimes_B N \cong P\) intertwining the tautological descent datum
  twisted by \(\pi(v)\).  Take \(N = B \otimes_A \mathrm{descended}(v)\) and the composite
  \[
    P \otimes_B \bigl(B \otimes_A \mathrm{descended}(v)\bigr) \;\cong\; P \otimes_A
    \mathrm{descended}(v) \;\cong\; P,
  \]
  the first isomorphism cancelling the iterated base change, the second the \(A\)-descent
  isomorphism attached to \(v\).  A computation on elementary tensors, using the membership
  identity \(v \cdot (m \otimes_A 1) = 1 \otimes_A m\), shows this composite intertwines the
  base change of the coaction \(x \mapsto v \cdot (x \otimes 1)\) with the coaction \(y \mapsto
  \pi(v) \cdot (y \otimes 1)\).  Uniqueness of descended modules then yields the stated
  isomorphism.
\end{proof}

\begin{theorem}[Naturality of the Picard class under descent in stages]
  \label{pa-thm:module_picClass_collapse}
  \lean{Module.IsDescentCocycle.picClass_collapse}
  \leanok
  \source{kleiman-picard}

  \uses{pa-thm:module_descendedCollapseEquiv, pa-def:module_picClass}
  \dcref{custom}
  Under the hypotheses of Theorem~\ref{pa-thm:module_descendedCollapseEquiv}, the Picard class over
  \(B\) of the collapsed cocycle is the image of the Picard class over \(A\) of \(v\) under the
  base-change homomorphism \(\Pic(A) \to \Pic(B)\):
  \[
    \mathrm{picClass}(\pi(v)) = \bigl[B \otimes_A \mathrm{descended}(v)\bigr] \quad\text{in }
    \Pic(B).
  \]
\end{theorem}
\begin{proof}
  \leanok
  By Theorem~\ref{pa-thm:module_descendedCollapseEquiv}, \(B \otimes_A \mathrm{descended}(v) \cong
  \mathrm{descended}(\pi(v))\) as \(B\)-modules, so the two invertible \(B\)-modules have the
  same class in \(\Pic(B)\).  The right side is \(\mathrm{picClass}(\pi(v))\) by definition
  (Definition~\ref{pa-def:module_picClass}), and the left side is the image of
  \(\mathrm{picClass}(v) = [\mathrm{descended}(v)]\) under \([N] \mapsto [B \otimes_A N]\).
\end{proof}

\section{Naturality of the \v Cech--Picard dictionary in the affine scheme}
\label{pa-sec:cechPic_naturality}

\begin{theorem}[Naturality of the \v Cech--Picard descent homomorphism in the affine scheme]
  \label{pa-thm:cechPic_toPic_map}
  \lean{AlgebraicGeometry.Scheme.CechPic.toPic_map}
  \leanok
  \source{kleiman-picard}

  \uses{pa-def:cechPic_toPic, pa-thm:awayCover_pic_baseChange, pa-thm:prPullback_injective}
  \dcref{2.5,custom}
  For a morphism \(g \colon X \to Y\) of affine schemes, base change of Picard classes along the
  ring map \(g^{\#} \colon \Gamma(Y, \top) \to \Gamma(X, \top)\) intertwines the \v
  Cech--Picard descent homomorphism (Definition~\ref{pa-def:cechPic_toPic}) with the pullback of
  \v Cech classes; that is, for every \(L \in Y.\mathrm{CechPic}\),
  \[
    \mathrm{toPic}_X(g^{*} L) = \Pic(g^{\#})\bigl(\mathrm{toPic}_Y(L)\bigr) \quad\text{in }
    \Pic\bigl(\Gamma(X, \top)\bigr).
  \]
  This is the separatedness half of Kleiman, \emph{The Picard Scheme}, Theorem 2.5(1), in its
  \v Cech incarnation.
\end{theorem}
\begin{proof}
  \leanok
  Represent \(L\) by a unit cocycle \(\gamma\) on a pointed cover \(\mathcal U\) of \(Y\) and
  choose a finite basic refinement \(P\) of \(\mathcal U\).  Then \(\mathrm{toPic}_Y(L)\) is the
  Picard class descended from the cover cocycle of \(\gamma\) on the finite localization cover
  \(P.r \colon P.\iota \to \Gamma(Y, \top)\).  Base-changing this whole \v Cech picture along
  \(g^{\#}\) (Theorem~\ref{pa-thm:awayCover_pic_baseChange}) computes \(\Pic(g^{\#})(\mathrm{toPic}_Y
  L)\) as the Picard class of the pushed cover cocycle on the pushed family \(f'_i = g^{\#}(P.r_i)\),
  whose basic opens \(X.\mathrm{basicOpen}(f'_i) = g^{-1} Y.\mathrm{basicOpen}(P.r_i)\) are the
  pullbacks of \(P\)'s basic opens.

  On the other side, \(\mathrm{toPic}_X(g^{*} L)\) is the Picard class descended from the cover
  cocycle of \(g^{*}\gamma\) on \emph{any} finite basic refinement of the pulled-back cover
  \(\mathcal U^{\mathrm{pb}}\).  As \(g\) need not be surjective, the pushed family \(f'\) is not
  itself a basic refinement of \(\mathcal U^{\mathrm{pb}}\); take the product refinement \(R\)
  indexed by \(P.\iota \times Q.\iota\), for an auxiliary basic refinement \(Q\) of \(\mathcal
  U^{\mathrm{pb}}\), with sections \(g^{\#}(P.r_i) \cdot Q.r_j\).  The pushed cover cocycle,
  refined to \(R\) by the first projection, differs from the cover cocycle of \(g^{*}\gamma\) on
  \(R\) by the index-wise coboundary of the cross evaluations \(\gamma\) pulled back through
  \(g\); two point-selections on a common refinement always differ by such a coboundary.  Since
  the descended Picard class is unchanged by a coboundary, the two Picard classes coincide, which
  is the asserted identity.
\end{proof}

\begin{theorem}[Algebra form of the naturality]
  \label{pa-thm:cechPic_toPic_mapAlgebra}
  \lean{AlgebraicGeometry.Scheme.CechPic.toPic_mapAlgebra}
  \leanok
  \dcref{custom}

  \uses{pa-thm:cechPic_toPic_map}
  With \(\Gamma(X, \top)\) regarded as a \(\Gamma(Y, \top)\)-algebra through \(g^{\#}\), the
  identity of Theorem~\ref{pa-thm:cechPic_toPic_map} reads \(\mathrm{toPic}_X(g^{*} L) =
  \mathrm{mapAlgebra}\bigl(\mathrm{toPic}_Y(L)\bigr)\) for the base-change homomorphism of
  Picard groups attached to the algebra structure.
\end{theorem}
\begin{proof}
  \leanok
  The base-change homomorphism attached to the \(\Gamma(Y, \top)\)-algebra \(\Gamma(X, \top)\)
  is by construction \(\Pic(g^{\#})\); the statement is therefore
  Theorem~\ref{pa-thm:cechPic_toPic_map} rephrased.
\end{proof}

\section{The class-level coherence seed for \'etale separatedness}
\label{pa-sec:etaleSeparatedness}

Let \(k\) be a field and \(A \to B\) a homomorphism of \(k\)-algebras, and let \(C\) be a curve
over \(k\).  Write \(\Spec A, \Spec B\) for the corresponding test objects of \(\Over(\Spec
k)\).  This section records the class-level coherence seed of the (C1) \'etale-separatedness
theorem — the separatedness half of Kleiman, \emph{The Picard Scheme}, Theorem 2.5(1).

\begin{definition}[The two coprojections of a double cover of a test algebra]
  \label{pa-def:tensorInl_tensorInr}
  \lean{AlgebraicGeometry.tensorInl, AlgebraicGeometry.tensorInr}
  \leanok
  \dcref{custom}

  The two \emph{coprojections} of the double cover \(B \otimes_A B\) of \(B\) are the
  \(k\)-algebra homomorphisms
  \[
    \mathrm{tensorInl} \colon B \to B \otimes_A B, \quad b \mapsto b \otimes 1, \qquad
    \mathrm{tensorInr} \colon B \to B \otimes_A B, \quad b \mapsto 1 \otimes b.
  \]
\end{definition}

\begin{lemma}[Diagonal coherence of the two coprojections]
  \label{pa-lem:tensorInl_comp_ofId}
  \lean{AlgebraicGeometry.tensorInl_comp_ofId}
  \leanok
  \dcref{custom}

  \uses{pa-def:tensorInl_tensorInr}
  Precomposed with the base inclusion \(A \to B\), the two coprojections agree:
  \(\mathrm{tensorInl} \circ (A \to B) = \mathrm{tensorInr} \circ (A \to B)\) as \(k\)-algebra
  homomorphisms \(A \to B \otimes_A B\).
\end{lemma}
\begin{proof}
  \leanok
  For \(a \in A\) with image \(\bar a \in B\), both sides send \(a\) to \(\bar a \otimes 1\)
  respectively \(1 \otimes \bar a\); over \(A\) these are equal, \(\bar a \otimes 1 = a \cdot (1
  \otimes 1) = 1 \otimes \bar a\), because tensoring is over \(A\).
\end{proof}

\begin{theorem}[The class-level coherence seed]
  \label{pa-thm:cechPicMap_tensorInl_eq_tensorInr}
  \lean{AlgebraicGeometry.Over.cechPicMap_tensorInl_eq_tensorInr}
  \leanok
  \source{kleiman-picard}

  \uses{pa-def:tensorInl_tensorInr, pa-lem:tensorInl_comp_ofId, pa-thm:prPullback_injective,
    pa-lem:snd_left_naturality, pa-def:overSpecMap, pa-lem:overSpecMap_functor}
  \dcref{custom}
  Let \(L\) be a \v Cech Picard class on the curve-product \(C \otimes \Spec A\) and \(N\) a \v
  Cech Picard class on \(\Spec B\).  Suppose that over the test algebra \(A \to B\) the class
  \(L\) restricts to the pullback of \(N\) from the base:
  \[
    \mathrm{pr}_B^{*} N = (C \lhd \Spec(A \to B))^{*} L.
  \]
  Then the two base changes of \(N\) to \(\Spec(B \otimes_A B)\) along the coprojections
  coincide:
  \(\Spec(\mathrm{tensorInl})^{*} N = \Spec(\mathrm{tensorInr})^{*} N\).
\end{theorem}
\begin{proof}
  \leanok
  For each coprojection \(j \in \{\mathrm{tensorInl}, \mathrm{tensorInr}\}\), pull back
  \(\Spec(j)^{*} N\) along the projection \(\mathrm{pr}_{B \otimes_A B}\).  Using functoriality
  of \v Cech pullback (Lemma~\ref{pa-lem:overSpecMap_functor}), the base square of the projection
  (Lemma~\ref{pa-lem:snd_left_naturality}) and the hypothesis, this equals \((C \lhd \Spec(j
  \circ (A \to B)))^{*} L\).  By the diagonal coherence
  (Lemma~\ref{pa-lem:tensorInl_comp_ofId}), \(\mathrm{tensorInl} \circ (A \to B) = \mathrm{tensorInr}
  \circ (A \to B)\), so the two composites \(\Spec(\mathrm{tensorInl})^{*} N\) and
  \(\Spec(\mathrm{tensorInr})^{*} N\) have the same pullback along \(\mathrm{pr}_{B \otimes_A
  B}\).  Pullback along the curve projection is injective on Picard classes
  (Theorem~\ref{pa-thm:prPullback_injective}), hence the two base changes agree.
\end{proof}

\section{The Amitsur toolkit for \'etale separatedness}
\label{pa-sec:amitsurCochain}

This section collects three reusable groups of statements consumed by the coherent-witness
step of the (C1) \'etale-separatedness assembly, all revolving around the fact that only global
units cross the curve projection \(\mathrm{pr}\).

\begin{theorem}[Global gluing of a compatible unit family]
  \label{pa-thm:exists_global_unit_of_compatible}
  \lean{AlgebraicGeometry.Scheme.exists_global_unit_of_compatible,
    AlgebraicGeometry.Scheme.global_unit_ext}
  \leanok
  \dcref{custom}

  \uses{pa-lem:units_gluing, pa-lem:units_separation}
  Let \(\mathcal U\) be a pointed cover of a scheme \(X\).  A family of units \(\beta_x \in
  \Gamma(X, \mathcal U_x)^{\times}\) whose members agree on pairwise overlaps glues to a single
  global unit \(w \in \Gamma(X, \top)^{\times}\) restricting to each \(\beta_x\); and a global
  unit that restricts to \(1\) on every member of a pointed cover — more precisely, two global
  units agreeing on every member — is uniquely determined.
\end{theorem}
\begin{proof}
  \leanok
  The members \(\mathcal U_x\) cover \(\top\) because \(x \in \mathcal U_x\).  The gluing
  statement is the sheaf gluing for unit sections (Lemma~\ref{pa-lem:units_gluing}) at the
  tautological cover \(V = \top\), \(W = \mathcal U\); the uniqueness statement is the
  separation for unit sections (Lemma~\ref{pa-lem:units_separation}) at the same cover.
\end{proof}

\begin{definition}[The three cofaces on the triple tensor]
  \label{pa-def:tensorFaces}
  \lean{AlgebraicGeometry.tensorFace₁₂, AlgebraicGeometry.tensorFace₁₃,
    AlgebraicGeometry.tensorFace₂₃}
  \leanok
  \dcref{custom}

  Extending the coprojection idiom of Definition~\ref{pa-def:tensorInl_tensorInr}, the three
  \emph{cofaces} on the triple tensor are the \(k\)-algebra homomorphisms \(B \otimes_A B \to B
  \otimes_A (B \otimes_A B)\)
  \[
    \mathrm{tensorFace}_{12}(x \otimes y) = x \otimes (y \otimes 1), \quad
    \mathrm{tensorFace}_{13}(x \otimes y) = x \otimes (1 \otimes y), \quad
    \mathrm{tensorFace}_{23}(w) = 1 \otimes w,
  \]
  the \(k\)-restrictions of the \(A\)-linear descent cofaces.
\end{definition}

\begin{lemma}[Simplicial coincidences of the cofaces]
  \label{pa-lem:tensorFace_comp}
  \lean{AlgebraicGeometry.tensorFace₁₂_comp_tensorInl,
    AlgebraicGeometry.tensorFace₁₂_comp_tensorInr,
    AlgebraicGeometry.tensorFace₁₃_comp_tensorInr}
  \leanok
  \dcref{custom}

  \uses{pa-def:tensorFaces, pa-def:tensorInl_tensorInr}
  The six composites \(\mathrm{tensorFace} \circ \mathrm{coprojection} \colon B \to B \otimes_A
  (B \otimes_A B)\) collapse onto the three insertions \(x \mapsto x \otimes (1 \otimes 1)\),
  \(x \mapsto 1 \otimes (x \otimes 1)\), \(x \mapsto 1 \otimes (1 \otimes x)\).  In particular
  \(\mathrm{tensorFace}_{12} \circ \mathrm{tensorInl} = \mathrm{tensorFace}_{13} \circ
  \mathrm{tensorInl}\), \(\mathrm{tensorFace}_{12} \circ \mathrm{tensorInr} =
  \mathrm{tensorFace}_{23} \circ \mathrm{tensorInl}\), and \(\mathrm{tensorFace}_{13} \circ
  \mathrm{tensorInr} = \mathrm{tensorFace}_{23} \circ \mathrm{tensorInr}\).
\end{lemma}
\begin{proof}
  \leanok
  Each identity is checked on a generator \(b \in B\): both sides insert \(b\) into the same
  tensor position and fill the remaining two positions with \(1\), by the definitions of the
  cofaces and coprojections.
\end{proof}

\begin{definition}[The top-open descent equivalence for global units]
  \label{pa-def:unitsSndTopEquiv}
  \lean{AlgebraicGeometry.Over.unitsSndTopEquiv}
  \leanok
  \dcref{custom}

  \uses{pa-def:unitsSndEquiv}
  Specializing the projection-descent equivalence (Definition~\ref{pa-def:unitsSndEquiv}) to the
  affine open \(V = \top\), pullback of global units along the second projection is a group
  isomorphism
  \[
    \Gamma\bigl((\Spec R)_{\mathrm{left}}, \top\bigr)^{\times} \;\xrightarrow{\ \sim\ }\;
    \Gamma\bigl((C \otimes \Spec R)_{\mathrm{left}}, \top\bigr)^{\times}
  \]
  for a \(k\)-algebra \(R\); its forward map is pullback of global units along \(\mathrm{pr}\).
\end{definition}
\begin{proof}
  \leanok
  Both \((\Spec R)_{\mathrm{left}}\) and \((C \otimes \Spec R)_{\mathrm{left}}\) have affine top
  open, and the projection preimage of \(\top\) is \(\top\); so
  Definition~\ref{pa-def:unitsSndEquiv} applies at \(V = \top\), and its forward map is pullback of
  global units.
\end{proof}

\begin{theorem}[Global-unit pullback is bijective]
  \label{pa-thm:appTop_units_bijective}
  \lean{AlgebraicGeometry.Over.appTop_units_surjective,
    AlgebraicGeometry.Over.appTop_units_injective}
  \leanok
  \dcref{custom}

  \uses{pa-def:unitsSndTopEquiv}
  Pullback of global units along the second projection \(\mathrm{pr} \colon (C \otimes \Spec
  R)_{\mathrm{left}} \to (\Spec R)_{\mathrm{left}}\) is bijective: every global unit on the
  product descends to a unique global unit on the base.
\end{theorem}
\begin{proof}
  \leanok
  This is the bijectivity of the group isomorphism of
  Definition~\ref{pa-def:unitsSndTopEquiv}: surjectivity gives the descent of any global unit as
  the image of the inverse isomorphism, and injectivity is separation of global units by the
  projection.
\end{proof}

\begin{theorem}[Naturality of the top-open descent]
  \label{pa-thm:unitsSndTopEquiv_naturality}
  \lean{AlgebraicGeometry.Over.unitsSndTopEquiv_naturality}
  \leanok
  \dcref{custom}

  \uses{pa-def:unitsSndTopEquiv, pa-thm:unitsSndEquiv_naturality}
  For a \(k\)-algebra homomorphism \(\varphi \colon R \to R'\), descending a global unit on
  \(\Spec R\) and pulling it back along \(C \lhd \Spec\varphi\) equals first pulling back along
  \(\Spec\varphi\) and then descending on \(\Spec R'\).
\end{theorem}
\begin{proof}
  \leanok
  This is the naturality of the projection-descent equivalence in the test object
  (Theorem~\ref{pa-thm:unitsSndEquiv_naturality}) specialized to \(V = \top\) along the morphism
  \(\Spec\varphi\); all the opens occurring there are \(\top\) on the nose, so the general
  naturality square applies directly.
\end{proof}

\section{The pullback calculus for unit cochains}
\label{pa-sec:cochainPullback}

The coherent-witness construction of the next section manipulates unit \v Cech cochains by
pulling them back along morphisms of schemes and combining the results multiplicatively. This
section records the elementary calculus of those operations over abstract schemes, so that the
concrete instances in the construction are single applications.

Throughout, for a morphism \(f \colon X \to Y\), a unit section \(s \in \struct Y(V)^{\times}\)
on an open \(V \subseteq Y\), and an open \(O \subseteq X\) contained in the preimage
\(f^{-1}(V)\), write \(f^{\sharp}_O s \in \struct X(O)^{\times}\) for the pullback of \(s\)
along \(f\), restricted to \(O\). This operation is a group homomorphism on units, is
functorial in the morphism, \((r \circ f)^{\sharp}_O s = f^{\sharp}_O (r^{\sharp} s)\) whenever
the composite is defined, and commutes with shrinking \(O\). When \(u \in \struct Y(\top)^{\times}\)
is a \emph{global} unit, its pullback \(f^{\sharp}_O(u|_\top)\) is the restriction to \(O\) of
the global-section pullback \(f^{\sharp}_\top u \in \struct X(\top)^{\times}\).

\begin{lemma}[Congruence in the morphism]
  \label{pa-lem:unitsAppLE_congr_hom}
  \lean{AlgebraicGeometry.Scheme.Hom.unitsAppLE_section_congr_hom,
    AlgebraicGeometry.Scheme.Hom.unitsAppLE_pair_congr_hom}
  \leanok
  \dcref{custom}

  Let \(g, g' \colon X \to Y\) be two morphisms with \(g = g'\), let \(V\) assign to each point
  of \(Y\) an open of \(Y\), and let \(s\) assign to each point \(y\) a unit section
  \(s_y \in \struct Y(V(y))^{\times}\). Then for every \(x \in X\) and every open
  \(O \subseteq g^{-1}(V(g(x)))\),
  \[
    g^{\sharp}_O\, s_{g(x)} = (g')^{\sharp}_O\, s_{g'(x)} .
  \]
  The same holds for doubly point-indexed families \(V(y,y')\), \(s_{y,y'}\), such as the
  pair evaluations of a unit \v Cech cocycle.
\end{lemma}
\begin{proof}
  \leanok
  Since \(g = g'\), the two sides are identical: the equality of morphisms transports the
  source base point, the open \(V(g(x)) = V(g'(x))\), and the section
  \(s_{g(x)} = s_{g'(x)}\) simultaneously, so the pullbacks agree.
\end{proof}

\begin{lemma}[Pullback of a coboundary relation]
  \label{pa-lem:unitsAppLE_coboundary_rel}
  \lean{AlgebraicGeometry.Scheme.Hom.unitsAppLE_coboundary_rel}
  \leanok
  \dcref{custom}

  Let \(\varphi \colon X \to Y\), \(r_1 \colon Y \to Z_1\), \(r_2 \colon Y \to Z_2\). Suppose a
  unit \(0\)-cochain \(c\) on a pointed cover \(\mathcal C\) of \(Y\) cobounds the comparison of
  two pair-indexed unit families \(E_1\) on \(Z_1\) and \(E_2\) on \(Z_2\), pulled back along
  \(r_1\) and \(r_2\): for all \(y, y' \in Y\), on the overlap
  \(\mathcal C(y) \cap \mathcal C(y')\),
  \[
    c_y \cdot \bigl(r_1^{\sharp} E_1\bigr)_{y y'}
      = \bigl(r_2^{\sharp} E_2\bigr)_{y y'} \cdot c_{y'} .
  \]
  Then for every \(x, x' \in X\) and every open \(O \subseteq \varphi^{-1}(\mathcal C(\varphi x)
  \cap \mathcal C(\varphi x'))\), the pullback of the relation along \(\varphi\) holds in
  composite normal form:
  \[
    \varphi^{\sharp}_O\, c_{\varphi x} \cdot \bigl((r_1\varphi)^{\sharp}_O E_1\bigr)_{x x'}
      = \bigl((r_2\varphi)^{\sharp}_O E_2\bigr)_{x x'} \cdot \varphi^{\sharp}_O\, c_{\varphi x'},
  \]
  where \(r_i\varphi\) denotes the composite \(X \xrightarrow{\varphi} Y \xrightarrow{r_i} Z_i\).
\end{lemma}
\begin{proof}
  \leanok
  Apply the pullback homomorphism \(\varphi^{\sharp}_O\) to both sides of the given relation at
  the points \(\varphi x, \varphi x'\). Being a group homomorphism on units it distributes over
  the products; and by functoriality it merges each pullback,
  \(\varphi^{\sharp}_O(r_i^{\sharp} E_i) = (r_i\varphi)^{\sharp}_O E_i\), while leaving
  \(\varphi^{\sharp}_O c\) unchanged. The two normal forms result.
\end{proof}

\begin{lemma}[The global-unit pullback in restriction form]
  \label{pa-lem:global_unit_pullback}
  \lean{AlgebraicGeometry.Scheme.Hom.unitsAppLE_unitsRestrict_top,
    AlgebraicGeometry.unitsRestrict_top_self, AlgebraicGeometry.unitsAppLE_top_global}
  \leanok
  \dcref{custom}

  Let \(f \colon X \to Y\) and \(u \in \struct Y(\top)^{\times}\) a global unit. The restriction
  of \(u\) along \(\top \le \top\) is \(u\) itself; and for any point-indexed open \(V\) and any
  open \(O \subseteq f^{-1}(V)\), pulling the restricted unit back along \(f\) equals the
  restriction to \(O\) of the global-section pullback of \(u\):
  \[
    f^{\sharp}_O\bigl(u|_V\bigr) = \bigl(f^{\sharp}_\top u\bigr)\big|_O .
  \]
  In particular, at \(V = \top\) the open-indexed and global spellings of pulling back a global
  unit coincide.
\end{lemma}
\begin{proof}
  \leanok
  The restriction along the identity inclusion \(\top \le \top\) is the identity map on
  sections. For the second identity, both sides are the image of \(u\) under one and the same
  ring map \(\struct Y(\top) \to \struct X(O)\) — the composite of the comorphism of \(f\) on
  global sections with restriction to \(O\) — so they agree; the \(V = \top\) case is the same
  statement read backwards.
\end{proof}

\begin{lemma}[Functoriality of the global-unit pullback]
  \label{pa-lem:units_map_appTop_comp}
  \lean{AlgebraicGeometry.units_map_appTop_comp}
  \leanok
  \dcref{custom}

  For morphisms \(f \colon X \to Y\), \(g \colon Y \to Z\) and a global unit
  \(u \in \struct Z(\top)^{\times}\), pulling \(u\) back along \(g\) and then along \(f\) equals
  pulling it back along the composite \(g \circ f\):
  \[
    f^{\sharp}_\top\bigl(g^{\sharp}_\top u\bigr) = (g \circ f)^{\sharp}_\top u .
  \]
\end{lemma}
\begin{proof}
  \leanok
  The comorphism of a composite of morphisms is the composite of the comorphisms in the
  reverse order; applying this on global sections and passing to units gives the claim.
\end{proof}

\begin{lemma}[Pushing a pullback through a product of cochain values]
  \label{pa-lem:unitsAppLE_merge}
  \lean{AlgebraicGeometry.unitsAppLE_mulDiv_comp, AlgebraicGeometry.unitsAppLE_mulRatio_comp}
  \leanok
  \dcref{custom}

  Let \(\varphi \colon X \to Y\) and let \(a, b, c\) be pulled-back cochain values along
  morphisms \(g_1, g_2, g_3 \colon Y \to Z\), each a unit section on a common open \(V\) of
  \(Y\). For any open \(O \subseteq \varphi^{-1}(V)\),
  \[
    \varphi^{\sharp}_O\Bigl(\tfrac{a \cdot b}{c}\Bigr)
      = \frac{(g_1\varphi)^{\sharp}_O a \cdot (g_2\varphi)^{\sharp}_O b}{(g_3\varphi)^{\sharp}_O c},
  \]
  and likewise for a combination of the shape \(t \cdot (a/b)\); pushing the pullback through
  the combination merges each factor into a composite pullback.
\end{lemma}
\begin{proof}
  \leanok
  The pullback \(\varphi^{\sharp}_O\) is a group homomorphism on units, so it distributes over
  the products and quotients; each individual pullback then merges by functoriality,
  \(\varphi^{\sharp}_O(g^{\sharp} x) = (g\varphi)^{\sharp}_O x\).
\end{proof}

\begin{lemma}[Pullback of a glued identity in composite normal form]
  \label{pa-lem:unitsAppLE_glued_pullback}
  \lean{AlgebraicGeometry.unitsAppLE_defect_pullback, AlgebraicGeometry.unitsAppLE_ratio_pullback}
  \leanok
  \dcref{custom}

  \uses{pa-lem:global_unit_pullback, pa-lem:unitsAppLE_merge}
  Let \(\varphi \colon X \to Y\) and morphisms \(m_1, m_2, m_3 \colon Y \to S\) with composites
  \(t_i = m_i \varphi \colon X \to S\). If a global unit \(W\) on \(Y\) restricts on an open
  \(V\) to a combination \(\tfrac{(m_1^{\sharp}\theta)(m_2^{\sharp}\theta)}{m_3^{\sharp}\theta}\)
  of pullbacks of a cochain \(\theta\), then the global-section pullback of \(W\) along
  \(\varphi\) restricts on any \(O \subseteq \varphi^{-1}(V)\) to the corresponding combination
  \(\tfrac{(t_1^{\sharp}\theta)(t_2^{\sharp}\theta)}{t_3^{\sharp}\theta}\) of the composite
  pullbacks. The analogous statement holds for a comparison combination of the shape
  \(t \cdot (a/b)\).
\end{lemma}
\begin{proof}
  \leanok
  By Lemma~\ref{pa-lem:global_unit_pullback} the restriction of the global-section pullback of
  \(W\) to \(O\) equals \(\varphi^{\sharp}_O\) of the \(V\)-restriction of \(W\), which by
  hypothesis is the given combination of \(m_i\)-pullbacks; pushing \(\varphi^{\sharp}_O\)
  through that combination and merging the compositions (Lemma~\ref{pa-lem:unitsAppLE_merge})
  produces the combination of \(t_i = m_i\varphi\) pullbacks. The output composites are the
  named \(t_i\) by hypothesis.
\end{proof}

\section{The Amitsur-coherent \v Cech witness}
\label{pa-sec:coherentWitness}

Let \(k\) be a field, \(A \to B\) a homomorphism of \(k\)-algebras, and \(C\) an object of
\(\Over(\Spec k)\) which is proper, geometrically irreducible and geometrically reduced over
\(k\). Retain the coprojections \(q_1 = \Spec(\mathrm{tensorInl})\),
\(q_2 = \Spec(\mathrm{tensorInr}) \colon \Spec B \to \Spec(B \otimes_A B)\) of
Definition~\ref{pa-def:tensorInl_tensorInr} and the three cofaces
\(f_{12}, f_{13}, f_{23} = \Spec(\mathrm{tensorFace}_{\bullet}) \colon \Spec(B \otimes_A B) \to
\Spec(B \otimes_A (B \otimes_A B))\) of Definition~\ref{pa-def:tensorFaces}. This section
implements the coherent-witness step of the \'etale-separatedness argument of Kleiman,
\emph{The Picard Scheme}: starting from the class equality \(q_1^{*} N = q_2^{*} N\) supplied by
the class-level coherence seed (Theorem~\ref{pa-thm:cechPicMap_tensorInl_eq_tensorInr}), it
produces a \v Cech witness cochain that is moreover \emph{Amitsur-coherent} over the triple
tensor, correcting an arbitrary witness by a single global unit.

\begin{lemma}[The coface coincidences at \(\Spec\) level and on the curve product]
  \label{pa-lem:overSpec_whisker_face}
  \lean{AlgebraicGeometry.Over.overSpecMap_left_face₁₂_inl,
    AlgebraicGeometry.Over.overSpecMap_left_face₁₂_inr,
    AlgebraicGeometry.Over.overSpecMap_left_face₁₃_inr,
    AlgebraicGeometry.Over.whiskerLeft_face₁₂_inl,
    AlgebraicGeometry.Over.whiskerLeft_face₁₂_inr,
    AlgebraicGeometry.Over.whiskerLeft_face₁₃_inr}
  \leanok
  \dcref{custom}

  \uses{pa-lem:tensorFace_comp, pa-def:tensorFaces, pa-def:tensorInl_tensorInr, pa-def:overSpecMap}
  The three coface--coprojection coincidences of Lemma~\ref{pa-lem:tensorFace_comp} persist after
  applying \(\Spec\) and after whiskering with \(C\): as morphisms of the corresponding
  curve products,
  \[
    f_{12} \circ q_1 = f_{13} \circ q_1, \qquad
    f_{12} \circ q_2 = f_{23} \circ q_1, \qquad
    f_{13} \circ q_2 = f_{23} \circ q_2,
  \]
  and the same three identities hold with each map replaced by its whiskering \(C \lhd (-)\).
\end{lemma}
\begin{proof}
  \leanok
  Each coincidence is the image of the corresponding tensor-algebra coincidence of
  Lemma~\ref{pa-lem:tensorFace_comp} under the functor \(\Spec\) (on test objects), which is
  contravariant on algebra maps, respectively under the further whiskering functor
  \(C \lhd (-)\); both functors preserve composition, so the equal composites of algebra maps
  become equal composites of morphisms.
\end{proof}

\begin{lemma}[The diagonal coincidence on the curve product]
  \label{pa-lem:whisker_inl_ofId}
  \lean{AlgebraicGeometry.Over.whiskerLeft_inl_ofId}
  \leanok
  \dcref{custom}

  \uses{pa-lem:tensorInl_comp_ofId, pa-def:tensorInl_tensorInr, pa-def:overSpecMap}
  Composed with the whiskered base inclusion \(C \lhd \Spec(A \to B)\), the two whiskered
  coprojections agree:
  \[
    (C \lhd q_1) \circ \bigl(C \lhd \Spec(A \to B)\bigr)
      = (C \lhd q_2) \circ \bigl(C \lhd \Spec(A \to B)\bigr).
  \]
\end{lemma}
\begin{proof}
  \leanok
  This is the image of the diagonal coherence \(\mathrm{tensorInl} \circ (A \to B) =
  \mathrm{tensorInr} \circ (A \to B)\) of Lemma~\ref{pa-lem:tensorInl_comp_ofId} under \(\Spec\)
  and the whiskering functor \(C \lhd (-)\), both of which preserve composition.
\end{proof}

\begin{definition}[The common refinement of the three coface pullbacks]
  \label{pa-def:amitsurCover}
  \lean{AlgebraicGeometry.amitsurCover}
  \leanok
  \dcref{custom}

  \uses{pa-def:tensorFaces}
  For a pointed cover \(\mathcal V\) of \(\Spec(B \otimes_A B)\), its \emph{Amitsur refinement}
  is the pointwise intersection
  \[
    \mathcal V_{123} := f_{23}^{*}\mathcal V \sqcap f_{12}^{*}\mathcal V \sqcap f_{13}^{*}\mathcal V
  \]
  of the pullbacks of \(\mathcal V\) along the three cofaces, a pointed cover of the triple
  tensor \(\Spec(B \otimes_A (B \otimes_A B))\).
\end{definition}

\begin{definition}[An Amitsur-coherent \v Cech witness]
  \label{pa-def:CoherentCechWitness}
  \lean{AlgebraicGeometry.CoherentCechWitness}
  \leanok
  \source{kleiman-picard}

  \uses{pa-def:amitsurCover, pa-def:tensorInl_tensorInr, pa-def:tensorFaces}
  \dcref{custom}
  Let \(\gamma\) be a unit \v Cech cocycle on a pointed cover \(\mathcal N\) of \(\Spec B\). An
  \emph{Amitsur-coherent \v Cech witness} for \(\gamma\) consists of a pointed cover
  \(\mathcal V\) of \(\Spec(B \otimes_A B)\) refining both coprojection pullbacks
  \(q_1^{*}\mathcal N\) and \(q_2^{*}\mathcal N\), together with a unit \(0\)-cochain
  \(\theta\) on \(\mathcal V\) such that:
  \begin{enumerate}
    \item (\emph{witness}) \(\theta\) cobounds the comparison of the two coprojection pullbacks
      of \(\gamma\): on the overlaps of \(\mathcal V\),
      \(\theta_x \cdot (q_1^{\sharp}\gamma)_{xy} = (q_2^{\sharp}\gamma)_{xy} \cdot \theta_y\);
    \item (\emph{coherent}) on the Amitsur refinement \(\mathcal V_{123}\), the three coface
      pullbacks of \(\theta\) satisfy the Amitsur cocycle relation
      \((f_{23}^{\sharp}\theta) \cdot (f_{12}^{\sharp}\theta) = f_{13}^{\sharp}\theta\).
  \end{enumerate}
  The witness relation records, at cocycle level, the class equality
  \(q_1^{*} N = q_2^{*} N\) for \(N = [\mathcal N, \gamma]\); the coherence relation is the
  triple-overlap cocycle identity that lets \(\theta\) itself descend.
\end{definition}

\begin{lemma}[Global-unit correction of a witness]
  \label{pa-lem:nonempty_of_defect_eq_coboundary}
  \lean{AlgebraicGeometry.CoherentCechWitness.nonempty_of_defect_eq_coboundary}
  \leanok
  \source{kleiman-picard}

  \uses{pa-def:CoherentCechWitness, pa-def:amitsurCover, pa-lem:global_unit_pullback}
  \dcref{custom}
  Let \(\theta_0\) be a witness cochain on a common refinement \(\mathcal V\) of the two
  coprojection pullbacks of \(\mathcal N\) (satisfying the witness relation for \(\gamma\)), and
  let \(\chi\) be a global unit on \(\Spec(B \otimes_A B)\) whose Amitsur coboundary
  \((f_{23}^{\sharp}\chi)(f_{12}^{\sharp}\chi)/(f_{13}^{\sharp}\chi)\) restricts, on
  \(\mathcal V_{123}\), to the Amitsur defect
  \((f_{23}^{\sharp}\theta_0)(f_{12}^{\sharp}\theta_0)/(f_{13}^{\sharp}\theta_0)\) of
  \(\theta_0\). Then the corrected cochain \(\theta := \theta_0/\chi\) is an Amitsur-coherent
  \v Cech witness for \(\gamma\).
\end{lemma}
\begin{proof}
  \leanok
  Take the cover of the witness to be \(\mathcal V\) and the cochain to be
  \(\theta = \theta_0/(\chi|)\), where \(\chi|\) is the restriction of the global unit \(\chi\)
  to the members of \(\mathcal V\). For the witness relation: dividing both sides of the
  witness relation for \(\theta_0\) by the central unit \(\chi\), which restricts the same way
  at \(x\) and \(y\), preserves it, so \(\theta\) is again a witness. For coherence: the Amitsur
  defect of \(\theta\) is the Amitsur defect of \(\theta_0\) divided by the Amitsur combination
  \((f_{23}^{\sharp}\chi)(f_{12}^{\sharp}\chi)/(f_{13}^{\sharp}\chi)\) of the coface restrictions
  of \(\chi\). Each coface pullback of the restricted global unit \(\chi\) is the restriction of
  its global-section pullback (Lemma~\ref{pa-lem:global_unit_pullback}), so this Amitsur
  combination is exactly the restriction of the global unit displayed in the hypothesis; by
  hypothesis it equals the Amitsur defect of \(\theta_0\). The quotient is therefore \(1\),
  i.e.\ \((f_{23}^{\sharp}\theta)(f_{12}^{\sharp}\theta) = f_{13}^{\sharp}\theta\).
\end{proof}

\begin{definition}[The Amitsur defect of a witness cochain]
  \label{pa-def:defectCochain}
  \lean{AlgebraicGeometry.Over.defectCochain}
  \leanok
  \dcref{custom}

  \uses{pa-def:amitsurCover}
  The \emph{Amitsur defect} of a unit \(0\)-cochain \(\theta_0\) on a pointed cover
  \(\mathcal V\) of \(\Spec(B \otimes_A B)\) is the unit \(0\)-cochain
  \[
    \partial\theta_0 :=
      \frac{(f_{23}^{\sharp}\theta_0) \cdot (f_{12}^{\sharp}\theta_0)}{f_{13}^{\sharp}\theta_0}
  \]
  on the Amitsur refinement \(\mathcal V_{123}\). A witness is Amitsur-coherent precisely when
  its defect is the constant unit \(1\).
\end{definition}

\begin{theorem}[The Amitsur defect of a witness glues to a global unit]
  \label{pa-thm:exists_glued_defect}
  \lean{AlgebraicGeometry.Over.exists_glued_defect}
  \leanok
  \dcref{custom}

  \uses{pa-def:defectCochain, pa-lem:unitsAppLE_coboundary_rel, pa-lem:unitsAppLE_congr_hom,
    pa-lem:overSpec_whisker_face, pa-thm:exists_global_unit_of_compatible}
  If \(\theta_0\) is a witness cochain for \(\gamma\) on a common refinement \(\mathcal V\)
  (satisfying the witness relation), then its Amitsur defect \(\partial\theta_0\) has trivial
  \v Cech coboundary and glues to a global unit \(\bar\omega\) on the triple tensor
  \(\Spec(B \otimes_A (B \otimes_A B))\), restricting to \(\partial\theta_0\) on each member of
  \(\mathcal V_{123}\).
\end{theorem}
\begin{proof}
  \leanok
  Pull the witness relation for \(\theta_0\) back along each of the three cofaces
  \(f_{23}, f_{12}, f_{13}\) (Lemma~\ref{pa-lem:unitsAppLE_coboundary_rel}), obtaining on the
  Amitsur refinement three coboundary relations between the coface pullbacks of \(\theta_0\) and
  the pulled-back coprojections of \(\gamma\). By the coface coincidences
  (Lemma~\ref{pa-lem:overSpec_whisker_face}) the six composite maps
  \(f_{\bullet} \circ q_{\bullet}\) collapse onto the three insertions, so — after rewriting the
  point-indexed \(\gamma\)-terms accordingly (Lemma~\ref{pa-lem:unitsAppLE_congr_hom}) — the three
  relations share matching \(\gamma\)-factors. Multiplying the \(f_{23}\)- and
  \(f_{12}\)-relations and dividing by the \(f_{13}\)-relation, the \(\gamma\)-factors cancel
  telescopically, leaving that the defect values \(\partial\theta_0\) at any two points of
  \(\mathcal V_{123}\) agree on their overlap. Thus \(\partial\theta_0\) is a compatible unit
  family, and by global gluing of a compatible family
  (Theorem~\ref{pa-thm:exists_global_unit_of_compatible}) it glues to a global unit \(\bar\omega\).
\end{proof}

\begin{lemma}[Compatibility of a ratio cochain from two coboundary relations]
  \label{pa-lem:ratio_compat}
  \lean{AlgebraicGeometry.Scheme.Hom.unitsAppLE_coboundary_ratio_compat}
  \leanok
  \dcref{custom}

  \uses{pa-lem:unitsAppLE_coboundary_rel, pa-lem:unitsAppLE_congr_hom}
  Let \(\varphi_t \colon X \to Y_c\) and \(\varphi_1, \varphi_2 \colon X \to Y_a\). Suppose a
  unit \(0\)-cochain \(c\) on \(Y_c\) cobounds the comparison of a family \(E_Z\) along two maps
  \(Y_c \to Z\), and a unit \(0\)-cochain \(a\) on \(Y_a\) cobounds the comparison of \(E_Z\)
  (along a map \(Y_a \to Z\)) with a family \(E_W\) (along a map \(Y_a \to W\)). If the
  pulled-back comparison composites coincide pairwise and the two maps \(\varphi_1, \varphi_2\)
  agree after the map to \(W\) (the diagonal coincidence), then the ratio cochain
  \(\varphi_t^{\sharp} c \cdot (\varphi_2^{\sharp} a / \varphi_1^{\sharp} a)\) on \(X\) is
  \v Cech-compatible: its values at any two points agree on a common open.
\end{lemma}
\begin{proof}
  \leanok
  Pull the two coboundary relations back to \(X\): the relation for \(c\) along \(\varphi_t\),
  and the relation for \(a\) along each of \(\varphi_1, \varphi_2\)
  (Lemma~\ref{pa-lem:unitsAppLE_coboundary_rel}). The pairwise coincidence of composites lets us
  identify the \(E_Z\)-terms of the \(c\)-relation with those of the \(a\)-relations, and the
  diagonal coincidence identifies the two \(E_W\)-terms of the \(\varphi_1\)- and
  \(\varphi_2\)-pullbacks (Lemma~\ref{pa-lem:unitsAppLE_congr_hom}). Writing \(Q_1, Q_2\) for the
  common \(E_Z\)-terms and \(L\) for the common \(E_W\)-term, the three relations read
  \(c_x Q_1 = Q_2 c_{x'}\), \(a_i Q_i = L\,a_i'\); a direct computation in the unit group then
  gives \(c_x (a_2/a_1) = c_{x'}(a_2'/a_1')\), which is compatibility of the ratio cochain.
\end{proof}

\begin{definition}[The comparison cover on the curve product]
  \label{pa-def:comparisonCover}
  \lean{AlgebraicGeometry.Over.comparisonCover}
  \leanok
  \dcref{custom}

  \uses{pa-def:tensorInl_tensorInr}
  On the curve product \(X_{B \otimes B} := C \otimes \Spec(B \otimes_A B)\), the
  \emph{comparison cover} of a witness cover \(\mathcal V\) of \(\Spec(B \otimes_A B)\) and an
  upstairs cover \(\mathcal A\) of \(X_B := C \otimes \Spec B\) is the pointwise intersection
  \[
    \mathrm{pr}^{*}\mathcal V \;\sqcap\; \bigl(u_1^{*}\mathcal A \sqcap u_2^{*}\mathcal A\bigr),
  \]
  where \(\mathrm{pr} \colon X_{B \otimes B} \to \Spec(B \otimes_A B)\) is the projection and
  \(u_1, u_2 = C \lhd q_1, C \lhd q_2 \colon X_{B \otimes B} \to X_B\) are the whiskered
  coprojections.
\end{definition}

\begin{definition}[The comparison cochain]
  \label{pa-def:comparisonCochain}
  \lean{AlgebraicGeometry.Over.comparisonCochain}
  \leanok
  \dcref{custom}

  \uses{pa-def:comparisonCover}
  Given a witness cochain \(\theta_0\) on \(\mathcal V\) and an upstairs witness cochain
  \(\alpha\) on \(\mathcal A\), their \emph{comparison cochain} on the comparison cover is the
  unit \(0\)-cochain
  \[
    \mathrm{pr}^{\sharp}\theta_0 \cdot
      \frac{u_2^{\sharp}\alpha}{u_1^{\sharp}\alpha} .
  \]
\end{definition}

\begin{theorem}[The comparison cochain is \v Cech-compatible]
  \label{pa-thm:comparisonCochain_compat}
  \lean{AlgebraicGeometry.Over.comparisonCochain_compat}
  \leanok
  \dcref{custom}

  \uses{pa-def:comparisonCochain, pa-lem:ratio_compat, pa-lem:whisker_inl_ofId, pa-lem:snd_left_naturality}
  Let \(\theta_0\) be a witness cochain for \(\gamma\) downstairs and \(\alpha\) a witness
  cochain upstairs for the lift — that is, \(\alpha\) cobounds the comparison of the projection
  pullback of \(\gamma\) with the pullback along \(C \lhd \Spec(A \to B)\) of a unit cocycle
  \(\lambda\) representing the class \(L\). Then the comparison cochain of \(\theta_0\) against
  \(\alpha\) is \v Cech-compatible, hence has trivial coboundary.
\end{theorem}
\begin{proof}
  \leanok
  Apply the compatibility principle for ratio cochains (Lemma~\ref{pa-lem:ratio_compat}) with
  \(\varphi_t = \mathrm{pr}\) and \(\varphi_1 = u_1\), \(\varphi_2 = u_2\): the witness relation
  for \(\theta_0\) downstairs and the witness relation for \(\alpha\) upstairs are the two
  coboundary relations. The base squares of the projection
  (Lemma~\ref{pa-lem:snd_left_naturality}) identify \(\mathrm{pr}\) followed by each coprojection
  \(q_i\) with \(u_i\) followed by the projection \(\mathrm{pr}_B\), giving the pairwise
  coincidence of the \(\gamma\)-composites; and the diagonal coincidence
  (Lemma~\ref{pa-lem:whisker_inl_ofId}) identifies \(u_1\) and \(u_2\) after the whiskered base
  inclusion, cancelling the \(\lambda\)-terms carrying the class \(L\). The hypotheses of
  Lemma~\ref{pa-lem:ratio_compat} are met, and its conclusion is exactly compatibility of the
  comparison cochain.
\end{proof}

\begin{theorem}[The comparison cochain glues and descends to the correcting unit]
  \label{pa-thm:exists_comparison_unit}
  \lean{AlgebraicGeometry.Over.exists_comparison_unit}
  \leanok
  \dcref{custom}

  \uses{pa-thm:comparisonCochain_compat, pa-thm:exists_global_unit_of_compatible, pa-def:unitsSndTopEquiv}
  In the setting of Theorem~\ref{pa-thm:comparisonCochain_compat}, the comparison cochain glues to
  a global unit on the curve product \(X_{B \otimes B}\), which descends through the projection
  to a global unit \(\chi\) on \(\Spec(B \otimes_A B)\) whose projection pullback restricts to
  the comparison cochain.
\end{theorem}
\begin{proof}
  \leanok
  By Theorem~\ref{pa-thm:comparisonCochain_compat} the comparison cochain is a compatible unit
  family, so it glues to a global unit \(\psi\) on \(X_{B \otimes B}\)
  (Theorem~\ref{pa-thm:exists_global_unit_of_compatible}). Global-unit pullback along the
  projection is a group isomorphism (Definition~\ref{pa-def:unitsSndTopEquiv}); let \(\chi\) be the
  image of \(\psi\) under its inverse. Then the projection pullback of \(\chi\) is \(\psi\),
  which restricts to the comparison cochain on each member of the comparison cover.
\end{proof}

\begin{theorem}[The glued defect is the Amitsur coboundary of the comparison unit]
  \label{pa-thm:glued_defect_eq_amitsur_coboundary}
  \lean{AlgebraicGeometry.Over.glued_defect_eq_amitsur_coboundary}
  \leanok
  \dcref{custom}

  \uses{pa-def:defectCochain, pa-def:comparisonCochain, pa-lem:unitsAppLE_glued_pullback,
    pa-lem:global_unit_pullback, pa-lem:units_map_appTop_comp, pa-lem:overSpec_whisker_face,
    pa-thm:appTop_units_bijective, pa-lem:snd_left_naturality}
  Suppose the glued Amitsur defect \(\bar\omega\) of a witness cochain \(\theta_0\) restricts to
  the Amitsur defect \(\partial\theta_0\), and the comparison unit \(\chi\) pulls back on the
  curve product to the comparison cochain of \(\theta_0\) against an upstairs witness
  \(\alpha\). Then \(\bar\omega\) is the Amitsur coboundary of \(\chi\):
  \[
    \bar\omega
      = \frac{(f_{23}^{\sharp}\chi) \cdot (f_{12}^{\sharp}\chi)}{f_{13}^{\sharp}\chi} .
  \]
\end{theorem}
\begin{proof}
  \leanok
  Both sides are global units on the triple tensor; since global-unit pullback along the curve
  projection is injective (Theorem~\ref{pa-thm:appTop_units_bijective}), it suffices to compare
  their pullbacks to the triple-tensor curve product. Pull \(\bar\omega\) back along the
  projection: by Lemma~\ref{pa-lem:unitsAppLE_glued_pullback} and the base squares of the
  projection (Lemma~\ref{pa-lem:snd_left_naturality}) it restricts to the Amitsur combination of
  the three whiskered-coface pullbacks of \(\theta_0\). Pull each coface pullback of \(\chi\)
  back along the projection: by functoriality of the global-unit pullback
  (Lemma~\ref{pa-lem:global_unit_pullback}, Lemma~\ref{pa-lem:units_map_appTop_comp}) and the base
  squares it becomes the corresponding whiskered-coface pullback of the descended comparison
  unit, which by Lemma~\ref{pa-lem:unitsAppLE_glued_pullback} together with the coface coincidences
  (Lemma~\ref{pa-lem:overSpec_whisker_face}) restricts to the same \(\theta_0\)-term times a
  telescoping combination of correction factors built from \(\alpha\). Forming the Amitsur
  combination of the three, the \(\alpha\)-telescope cancels exactly, leaving precisely the
  Amitsur combination of \(\theta_0\)-terms found for \(\bar\omega\). The two global units thus
  have equal projection pullbacks, hence are equal.
\end{proof}

\begin{theorem}[Existence of the Amitsur-coherent \v Cech witness]
  \label{pa-thm:exists_coherentCechWitness}
  \lean{AlgebraicGeometry.Over.exists_coherentCechWitness}
  \leanok
  \source{kleiman-picard}

  \uses{pa-def:CoherentCechWitness, pa-thm:cechPicMap_tensorInl_eq_tensorInr, pa-thm:exists_glued_defect,
    pa-thm:exists_comparison_unit, pa-thm:glued_defect_eq_amitsur_coboundary,
    pa-lem:nonempty_of_defect_eq_coboundary}
  \dcref{custom}
  Let \(C\) be proper, geometrically irreducible and geometrically reduced over the field
  \(k\), let \(k \to A \to B\) be a tower of algebras, let \(L\) be a \v Cech Picard class on
  \(C \otimes \Spec A\), and let \(\gamma\) be a unit cocycle on a pointed cover \(\mathcal N\)
  of \(\Spec B\) whose class \(N\) pulls back, on the curve product, from \(L\):
  \(\mathrm{pr}_B^{*} N = (C \lhd \Spec(A \to B))^{*} L\). Then there exists an Amitsur-coherent
  \v Cech witness for \(\gamma\) — a witness for the class equality
  \(q_1^{*} N = q_2^{*} N\) over \(\Spec(B \otimes_A B)\) whose coface pullbacks are Amitsur-coherent
  over the triple tensor.
\end{theorem}
\begin{proof}
  \leanok
  By the class-level coherence seed (Theorem~\ref{pa-thm:cechPicMap_tensorInl_eq_tensorInr}) the
  hypothesis gives the equality \(q_1^{*} N = q_2^{*} N\); unwinding it through the
  equality criterion for \v Cech Picard classes yields a common refinement \(\mathcal V\) of the
  two coprojection pullbacks of \(\mathcal N\) and a witness cochain \(\theta_0\) on it
  satisfying the witness relation. Likewise, the hypothesis
  \(\mathrm{pr}_B^{*} N = (C \lhd \Spec(A \to B))^{*} L\) yields an upstairs witness cochain
  \(\alpha\) on a cover \(\mathcal A\) of \(X_B\), cobounding the comparison of the projection
  pullback of \(\gamma\) with the pullback of a representative \(\lambda\) of \(L\).
  The Amitsur defect of \(\theta_0\) glues to a global unit \(\bar\omega\)
  (Theorem~\ref{pa-thm:exists_glued_defect}); the comparison of \(\theta_0\) against
  \(\alpha\) glues and descends to a global unit \(\chi\) on \(\Spec(B \otimes_A B)\)
  (Theorem~\ref{pa-thm:exists_comparison_unit}); and \(\bar\omega\) is identified
  with the Amitsur coboundary of \(\chi\) (Theorem~\ref{pa-thm:glued_defect_eq_amitsur_coboundary}).
  Thus \(\chi\) satisfies the hypothesis of the correction lemma
  (Lemma~\ref{pa-lem:nonempty_of_defect_eq_coboundary}), and the corrected cochain
  \(\theta_0/\chi\) is the required Amitsur-coherent \v Cech witness.
\end{proof}

\section{The simplicial degeneracies of the tensor tower}
\label{pa-sec:specDegeneracy}

Complementing the cofaces of Definition~\ref{pa-def:tensorFaces} and their \(\Spec\)-level
coincidences (Lemma~\ref{pa-lem:overSpec_whisker_face}), the composite descent unit pulls back
along the two \emph{degeneracies} of the tensor tower: the multiplication, whose \(\Spec\)
is the diagonal, and the first-two-positions multiplication, whose \(\Spec\) is a
degeneracy of the triple tensor.

\begin{definition}[The multiplication and the first degeneracy]
  \label{pa-def:tensorMul}
  \lean{AlgebraicGeometry.tensorMul, AlgebraicGeometry.descentMul₁₂,
    AlgebraicGeometry.tensorMul₁₂}
  \leanok
  \dcref{custom}

  \uses{pa-def:tensorInl_tensorInr, pa-def:tensorFaces}
  The \emph{multiplication} is the \(k\)-algebra homomorphism
  \[
    \mathrm{tensorMul} \colon B \otimes_A B \to B, \qquad x \otimes y \mapsto xy ,
  \]
  whose \(\Spec\) is the diagonal \(\Delta \colon \Spec B \to \Spec(B \otimes_A B)\). The
  \emph{first degeneracy} is the \(k\)-algebra homomorphism multiplying the first two tensor
  positions,
  \[
    \mathrm{tensorMul}_{12} \colon B \otimes_A (B \otimes_A B) \to B \otimes_A B, \qquad
    x \otimes (y \otimes z) \mapsto xy \otimes z ,
  \]
  whose \(\Spec\) is the degeneracy \(\sigma \colon \Spec(B \otimes_A B) \to
  \Spec(B \otimes_A (B \otimes_A B))\).
\end{definition}

\begin{lemma}[Degeneracy identities at ring level]
  \label{pa-lem:tensorMul_simplicial}
  \lean{AlgebraicGeometry.tensorMul_comp_tensorInl,
    AlgebraicGeometry.tensorMul_comp_tensorInr,
    AlgebraicGeometry.tensorMul₁₂_comp_tensorFace₂₃,
    AlgebraicGeometry.tensorMul₁₂_comp_tensorFace₁₂,
    AlgebraicGeometry.tensorMul₁₂_comp_tensorFace₁₃}
  \leanok
  \dcref{custom}

  \uses{pa-def:tensorMul, pa-def:tensorInl_tensorInr, pa-def:tensorFaces}
  The degeneracies satisfy the simplicial identities against the coprojections and cofaces:
  \[
    \mathrm{tensorMul} \circ \mathrm{tensorInl} = \mathrm{id}_B
      = \mathrm{tensorMul} \circ \mathrm{tensorInr} ,
  \]
  and, on the triple tensor,
  \[
    \mathrm{tensorMul}_{12} \circ \mathrm{tensorFace}_{23} = \mathrm{id}, \quad
    \mathrm{tensorMul}_{12} \circ \mathrm{tensorFace}_{13} = \mathrm{id}, \quad
    \mathrm{tensorMul}_{12} \circ \mathrm{tensorFace}_{12}
      = \mathrm{tensorInl} \circ \mathrm{tensorMul} .
  \]
\end{lemma}
\begin{proof}
  \leanok
  Each identity is an equality of \(k\)-algebra homomorphisms, checked on a generator. For the
  first pair, \(\mathrm{tensorMul}(\mathrm{tensorInl}\,b) = \mathrm{tensorMul}(b \otimes 1) = b\)
  and \(\mathrm{tensorMul}(\mathrm{tensorInr}\,b) = \mathrm{tensorMul}(1 \otimes b) = b\). For
  the degeneracy against the cofaces, compute on \(x \otimes y\): \(\mathrm{tensorMul}_{12}(1
  \otimes (x \otimes y)) = x \otimes y\); \(\mathrm{tensorMul}_{12}(x \otimes (1 \otimes y)) =
  x \otimes y\); and \(\mathrm{tensorMul}_{12}(x \otimes (y \otimes 1)) = xy \otimes 1 =
  \mathrm{tensorInl}(xy) = \mathrm{tensorInl}(\mathrm{tensorMul}(x \otimes y))\).
\end{proof}

\begin{lemma}[The degeneracy coincidences at \(\Spec\) level]
  \label{pa-lem:overSpecMap_left_degeneracy}
  \lean{AlgebraicGeometry.Over.overSpecMap_left_mul_inl,
    AlgebraicGeometry.Over.overSpecMap_left_mul_inr,
    AlgebraicGeometry.Over.overSpecMap_left_mul₁₂_face₂₃,
    AlgebraicGeometry.Over.overSpecMap_left_mul₁₂_face₁₂,
    AlgebraicGeometry.Over.overSpecMap_left_mul₁₂_face₁₃}
  \leanok
  \dcref{custom}

  \uses{pa-lem:tensorMul_simplicial, pa-def:overSpecMap}
  Applying \(\Spec\) to Lemma~\ref{pa-lem:tensorMul_simplicial} gives the degeneracy coincidences
  of the affine test objects: as morphisms \(\Spec B \to \Spec B\) and
  \(\Spec(B \otimes_A B) \to \Spec(B \otimes_A B)\),
  \[
    \Delta \circ q_1 = \mathrm{id}, \quad \Delta \circ q_2 = \mathrm{id}, \quad
    \sigma \circ f_{23} = \mathrm{id}, \quad \sigma \circ f_{13} = \mathrm{id}, \quad
    \sigma \circ f_{12} = q_1 \circ \Delta .
  \]
\end{lemma}
\begin{proof}
  \leanok
  Each is the image of the corresponding ring-level identity of
  Lemma~\ref{pa-lem:tensorMul_simplicial} under the contravariant functor \(\Spec\) on test
  objects (Definition~\ref{pa-def:overSpecMap}), which turns each equal composite of algebra maps
  into the reversed equal composite of morphisms and sends the identity algebra map to the
  identity morphism.
\end{proof}

\section{The trivialization-corrected witness cochain}
\label{pa-sec:witnessCorrection}

The values \(\theta_x\) of an Amitsur-coherent \v Cech witness
(Definition~\ref{pa-def:CoherentCechWitness}) do not by themselves glue over a basic open of the
tensor square: their coboundary is the comparison of the two coprojection pullbacks of
\(\gamma\). The composite descent unit corrects them by cocycle values of \(\gamma\) anchored
at two distinguished points, after which they glue, satisfy the Amitsur cocycle identity, and
collapse onto the Zariski cover cocycle on the diagonal. This section develops that calculus
over abstract schemes; the concrete tensor-tower instantiation is
Section~\ref{pa-sec:witnessComponents}.

Throughout, \(r_1, r_2 \colon Y \to Z\) are two morphisms (playing the coprojections
\(\Spec(B \otimes_A B) \to \Spec B\)); \(\mathcal C\) is a pointed cover of \(Y\) with a unit
\(0\)-cochain \(c\) (the witness cover and cochain); \(\mathcal N\) is a pointed cover of
\(Z\) with a unit cocycle \(\gamma_Z\) (the trivializing cover of \(\Spec B\)); \(a, b \in Z\)
are two points; and \(D \subseteq Y\) is an open on which \(r_1\) lands in \(\mathcal N_a\) and
\(r_2\) in \(\mathcal N_b\).

\begin{definition}[The corrected witness cochain]
  \label{pa-def:unitsCorrCochain}
  \lean{AlgebraicGeometry.unitsCorrCochain}
  \leanok
  \dcref{custom}

  The \emph{corrected witness value} at a point \(y \in Y\) is the unit on
  \(\mathcal C_y \cap D\)
  \[
    (\mathrm{corr}\,c)_y := r_1^{\sharp}\gamma_Z(a,\, r_1 y) \cdot (c_y)^{-1} \cdot
      r_2^{\sharp}\gamma_Z(r_2 y,\, b) ,
  \]
  the two correction factors being the coprojection pullbacks of the cocycle values of
  \(\gamma_Z\) anchored \(a \to r_1 y\) on the left and \(r_2 y \to b\) on the right, with the
  witness entering inverted. This is the unique orientation whose coboundary cancels the
  coboundary of \(\theta^{-1}\), and which on the diagonal — where the coherent witness is
  \(1\) — telescopes to the cocycle value \(\gamma_Z(a, b)\).
\end{definition}

\begin{theorem}[The corrected values are \v Cech-compatible]
  \label{pa-thm:unitsCorrCochain_compat}
  \lean{AlgebraicGeometry.unitsCorrCochain_compat}
  \leanok
  \dcref{custom}

  \uses{pa-def:unitsCorrCochain}
  Assume the witness relation for \(c\): on the overlaps of \(\mathcal C\),
  \(c_x \cdot r_1^{\sharp}\gamma_Z(r_1 x, r_1 y) = r_2^{\sharp}\gamma_Z(r_2 x, r_2 y) \cdot c_y\).
  Then the corrected values agree on overlaps: the restrictions of \((\mathrm{corr}\,c)_y\) and
  \((\mathrm{corr}\,c)_{y'}\) to \((\mathcal C_y \cap D) \cap (\mathcal C_{y'} \cap D)\)
  coincide.
\end{theorem}
\begin{proof}
  \leanok
  Restrict all data to the overlap \(O := (\mathcal C_y \cap D) \cap (\mathcal C_{y'} \cap D)\).
  Three relations hold on \(O\): the cocycle identity of \(\gamma_Z\) at \((a, r_1 y, r_1 y')\)
  pulled back along \(r_1\), which expresses the left factor at \(y\) times the
  \(r_1\)-comparison term \(r_1^{\sharp}\gamma_Z(r_1 y, r_1 y')\) as the left factor at \(y'\);
  the witness relation for \(c\) at \((y, y')\), relating the two witness values through the
  \(r_1\)- and \(r_2\)-comparison terms; and the cocycle identity at \((r_2 y, r_2 y', b)\)
  pulled back along \(r_2\), which relates the two right factors symmetrically. Writing the
  three as \(A_y Q_1 = A_{y'}\), \(t_y Q_1 = Q_2 t_{y'}\), \(Q_2 C_{y'} = C_y\) in the
  commutative group of units on \(O\) (with \(A\) the left factors, \(t\) the witnesses, \(C\)
  the right factors, \(Q_1, Q_2\) the comparison terms), substituting \(Q_2 = t_y Q_1 t_{y'}^{-1}\)
  and cancelling \(t_y^{-1} t_y\) turns \(A_y t_y^{-1} C_y\) into \(A_{y'} t_{y'}^{-1} C_{y'}\),
  which is the asserted equality of corrected values.
\end{proof}

\begin{theorem}[The corrected values glue to a unit]
  \label{pa-thm:exists_glued_unitsCorrCochain}
  \lean{AlgebraicGeometry.exists_glued_unitsCorrCochain}
  \leanok
  \dcref{custom}

  \uses{pa-def:unitsCorrCochain, pa-thm:unitsCorrCochain_compat, pa-lem:units_gluing}
  Under the witness relation of Theorem~\ref{pa-thm:unitsCorrCochain_compat}, the corrected values
  glue to a single unit \(u \in \Gamma(Y, D)^{\times}\) whose restriction to each
  \(\mathcal C_y \cap D\) is \((\mathrm{corr}\,c)_y\).
\end{theorem}
\begin{proof}
  \leanok
  The opens \(\{\mathcal C_y \cap D\}_{y}\) cover \(D\), since \(y \in \mathcal C_y\). The
  corrected values agree on overlaps by Theorem~\ref{pa-thm:unitsCorrCochain_compat}, so the sheaf
  gluing for unit sections (Lemma~\ref{pa-lem:units_gluing}) produces the glued unit \(u\) on
  \(D\).
\end{proof}

\begin{theorem}[Pullback expansion of a glued corrected unit]
  \label{pa-thm:unitsAppLE_glued_corr}
  \lean{AlgebraicGeometry.Scheme.Hom.unitsAppLE_glued_corr}
  \leanok
  \dcref{custom}

  \uses{pa-def:unitsCorrCochain, pa-thm:exists_glued_unitsCorrCochain}
  Let \(\varphi \colon X' \to Y\) and let \(O \subseteq X'\) be an open inside
  \(\varphi^{-1}(\mathcal C_{\varphi x} \cap D)\). Then the \(\varphi\)-pullback of a glued
  corrected unit \(u\) expands into its three factors, the correction factors pulled back along
  the composites \(\varphi \circ r_1\) and \(\varphi \circ r_2\):
  \[
    \varphi^{\sharp} u =
      (\varphi r_1)^{\sharp}\gamma_Z(a,\, (\varphi r_1)x) \cdot
      \bigl(\varphi^{\sharp} c_{\varphi x}\bigr)^{-1} \cdot
      (\varphi r_2)^{\sharp}\gamma_Z((\varphi r_2)x,\, b) .
  \]
\end{theorem}
\begin{proof}
  \leanok
  By the defining property of \(u\) (Theorem~\ref{pa-thm:exists_glued_unitsCorrCochain}), its
  restriction to \(\mathcal C_{\varphi x} \cap D\) is \((\mathrm{corr}\,c)_{\varphi x}\). Pull
  this equality back along \(\varphi\) into \(O\) and distribute the pullback over the product
  and inverse; each point-indexed factor pulled back along \(\varphi\) then along \(r_i\)
  becomes its pullback along the composite \(\varphi \circ r_i\), which is the displayed normal
  form.
\end{proof}

\begin{theorem}[Coherence kills the witness along a coface section]
  \label{pa-thm:unitsAppLE_eq_one_of_coherent}
  \lean{AlgebraicGeometry.unitsAppLE_eq_one_of_coherent}
  \leanok
  \dcref{custom}

  \uses{pa-def:CoherentCechWitness}
  Let \(g_{23}, g_{12}, g_{13} \colon T \to Y\) satisfy the Amitsur coherence relation for a
  witness cochain \(c\) on a pointed cover \(\mathcal C\) of \(Y\):
  \((g_{23}^{\sharp} c) \cdot (g_{12}^{\sharp} c) = g_{13}^{\sharp} c\) on the common
  refinement of the three coface pullbacks. If \(\sigma \colon Y \to T\) and \(\delta_r \colon
  Y \to Y\) satisfy \(\sigma \circ g_{23} = \mathrm{id}\), \(\sigma \circ g_{12} = \delta_r\)
  and \(\sigma \circ g_{13} = \mathrm{id}\), then the \(\delta_r\)-pullback of \(c\) is
  trivial: \(\delta_r^{\sharp} c = 1\).
\end{theorem}
\begin{proof}
  \leanok
  Pull the coherence relation at the point \(\sigma y\) back along \(\sigma\) onto a suitable
  open. By the composite-section identities \(\sigma \circ g_{23} = \mathrm{id}\),
  \(\sigma \circ g_{12} = \delta_r\), \(\sigma \circ g_{13} = \mathrm{id}\), the three pulled
  back witness terms become \(c\), \(\delta_r^{\sharp} c\) and \(c\) respectively, so the
  relation reads \(c \cdot \delta_r^{\sharp} c = c\). Cancelling the unit \(c\) gives
  \(\delta_r^{\sharp} c = 1\).
\end{proof}

\begin{theorem}[Diagonal triviality of a coherent witness]
  \label{pa-thm:unitsAppLE_diag_eq_one_of_coherent}
  \lean{AlgebraicGeometry.unitsAppLE_diag_eq_one_of_coherent}
  \leanok
  \dcref{custom}

  \uses{pa-thm:unitsAppLE_eq_one_of_coherent}
  Let \(\delta \colon Z \to Y\) and \(\delta_r \colon Y \to Y\) satisfy \(\delta \circ
  \delta_r = \delta\), and suppose the \(\delta_r\)-pullback of \(c\) is trivial. Then the
  \(\delta\)-pullback of \(c\) is trivial: \(\delta^{\sharp} c = 1\).
\end{theorem}
\begin{proof}
  \leanok
  Pull the identity \(\delta_r^{\sharp} c = 1\) (at the point \(\delta v\)) back along
  \(\delta\); by \(\delta \circ \delta_r = \delta\) the composite \(\delta \circ \delta_r\)
  reduces to \(\delta\), so the pullback of \(\delta_r^{\sharp} c\) along \(\delta\) is
  \(\delta^{\sharp} c\), while the pullback of the constant \(1\) is \(1\). Hence
  \(\delta^{\sharp} c = 1\).
\end{proof}

\begin{theorem}[The glued corrected units satisfy the Amitsur cocycle identity]
  \label{pa-thm:glued_corr_amitsur}
  \lean{AlgebraicGeometry.glued_corr_amitsur}
  \leanok
  \dcref{custom}

  \uses{pa-def:unitsCorrCochain, pa-thm:unitsAppLE_glued_corr}
  Let \(g_{23}, g_{12}, g_{13} \colon T \to Y\) satisfy the simplicial coincidences
  \(g_{13} \circ r_1 = g_{12} \circ r_1\), \(g_{23} \circ r_1 = g_{12} \circ r_2\),
  \(g_{23} \circ r_2 = g_{13} \circ r_2\), and let \(c\) be Amitsur-coherent. Let \(u_{23},
  u_{12}, u_{13}\) be the glued corrected units for the point pairs \((b, c')\), \((a, b)\),
  \((a, c')\) on opens \(D_{23}, D_{12}, D_{13}\). Then, on any \(D_T\) inside the three
  preimages,
  \[
    g_{23}^{\sharp} u_{23} \cdot g_{12}^{\sharp} u_{12} = g_{13}^{\sharp} u_{13} .
  \]
\end{theorem}
\begin{proof}
  \leanok
  Check the identity on each member of the cover \(\{(\mathcal C\text{-triple-overlap})_z \cap
  D_T\}\) of \(D_T\). On such a member, expand the three pulled back glued units by
  Theorem~\ref{pa-thm:unitsAppLE_glued_corr}; the six correction factors, pulled back along the
  composites \(g_{\bullet} \circ r_{\bullet}\), collapse by the three simplicial coincidences
  onto matching cocycle values. Grouping, the left correction factor of \(u_{23}\) and the
  right correction factor of \(u_{12}\) are the two \(\gamma_Z\)-values anchored at \(b\); by
  the cocycle identity of \(\gamma_Z\) at \((b, \cdot, b)\) and \(\gamma_Z(b,b) = 1\) they
  cancel, leaving the remaining correction factors as the correction factors of \(u_{13}\). The
  three witness terms combine by the coherence relation
  \((g_{23}^{\sharp} c)(g_{12}^{\sharp} c) = g_{13}^{\sharp} c\). What remains is exactly
  \(g_{13}^{\sharp} u_{13}\); this is the commutative-group telescope of the paired correction
  factors against the coherence relation.
\end{proof}

\begin{theorem}[The glued corrected unit collapses onto the cocycle value on the diagonal]
  \label{pa-thm:glued_corr_collapse}
  \lean{AlgebraicGeometry.glued_corr_collapse}
  \leanok
  \dcref{custom}

  \uses{pa-def:unitsCorrCochain, pa-thm:unitsAppLE_glued_corr}
  Let \(\delta \colon Z \to Y\) be a common section of \(r_1, r_2\) (\(\delta \circ r_1 =
  \mathrm{id} = \delta \circ r_2\)) along which the witness is trivial
  (\(\delta^{\sharp} c = 1\)). Then, on an open \(E\) with \(\delta(E) \subseteq D\) and
  \(E \subseteq \mathcal N_a \cap \mathcal N_b\), the \(\delta\)-pullback of the glued corrected
  unit for the pair \((a, b)\) is the restriction of the cocycle value:
  \(\delta^{\sharp} u = \gamma_Z(a, b)\).
\end{theorem}
\begin{proof}
  \leanok
  Check the identity on each member of the cover \(\{\delta^{-1}\mathcal C_{\delta v} \cap E\}\)
  of \(E\). On such a member, expand \(\delta^{\sharp} u\) by
  Theorem~\ref{pa-thm:unitsAppLE_glued_corr}; since \(\delta \circ r_1 = \mathrm{id} = \delta
  \circ r_2\), the two correction factors become \(\gamma_Z(a, \delta v)\) and
  \(\gamma_Z(\delta v, b)\), and the witness factor is \(1\) by hypothesis. The product of the
  two correction factors telescopes through the cocycle identity of \(\gamma_Z\) at
  \((a, \delta v, b)\) to \(\gamma_Z(a, b)\).
\end{proof}

\section{The glued witness components on the tensor square}
\label{pa-sec:witnessComponents}

We now instantiate the corrected-witness calculus on the tensor towers \(\Spec B \leftarrow
\Spec(B \otimes_A B) \leftarrow \Spec(B \otimes_A (B \otimes_A B))\). Fix an
Amitsur-coherent \v Cech witness \(\theta' \colon \mathrm{CoherentCechWitness}\) for
\(\gamma\) (Definition~\ref{pa-def:CoherentCechWitness}) and a \emph{finite basic refinement}
\(P\) of the trivializing cover \(\mathcal N\): a finite family of sections \(P.r_i \in
\Gamma(\Spec B, \top)\) whose basic opens \(D(P.r_i)\) refine \(\mathcal N\), each with a
distinguished point \(P.\mathrm{pt}_i \in \mathcal N\). Write \(q_1, q_2\) for the
coprojections and \(f_{12}, f_{13}, f_{23}\) for the cofaces.

\begin{definition}[The double and triple basic sections]
  \label{pa-def:pairTripleSection}
  \lean{AlgebraicGeometry.Over.pairSection, AlgebraicGeometry.Over.tripleSection}
  \leanok
  \dcref{custom}

  \uses{pa-def:tensorInl_tensorInr, pa-def:tensorFaces}
  The \emph{double basic section} cutting out the double basic open of \(\Spec(B \otimes_A B)\)
  is
  \[
    \mathrm{pairSection}(i, j) := q_1^{\sharp}(P.r_i) \cdot q_2^{\sharp}(P.r_j)
    \in \Gamma(\Spec(B \otimes_A B), \top) ,
  \]
  and the \emph{triple basic section} cutting out the triple basic open of
  \(\Spec(B \otimes_A (B \otimes_A B))\) is
  \[
    \mathrm{tripleSection}(i, j, l) := (f_{12} q_1)^{\sharp}(P.r_i) \cdot
      \bigl( (f_{12} q_2)^{\sharp}(P.r_j) \cdot (f_{13} q_2)^{\sharp}(P.r_l) \bigr) ,
  \]
  spelled through the three canonical insertions \(f_{12} \circ q_1\), \(f_{12} \circ q_2\),
  \(f_{13} \circ q_2\).
\end{definition}

\begin{definition}[The glued witness component]
  \label{pa-def:witnessBasicSection}
  \lean{AlgebraicGeometry.Over.witnessBasicSection}
  \leanok
  \dcref{custom}

  \uses{pa-def:pairTripleSection, pa-def:unitsCorrCochain, pa-thm:exists_glued_unitsCorrCochain,
    pa-def:CoherentCechWitness}
  The \emph{glued witness component} \(\mathrm{witnessBasicSection}(i, j)\) is the glued unit on
  the double basic open \(D(\mathrm{pairSection}(i, j))\) obtained from
  Theorem~\ref{pa-thm:exists_glued_unitsCorrCochain} applied to the coprojections \(q_1, q_2\), the
  witness cover and cochain of \(\theta'\), the trivializing cocycle \(\gamma\), and the two
  distinguished points \(P.\mathrm{pt}_i, P.\mathrm{pt}_j\); its restriction to each
  \(\mathcal V_x \cap D(\mathrm{pairSection}(i, j))\) is the corrected witness value at \(x\).
\end{definition}

\begin{theorem}[Diagonal triviality of the coherent witness]
  \label{pa-thm:witnessAppLE_diag_eq_one}
  \lean{AlgebraicGeometry.Over.witnessAppLE_diag_eq_one}
  \leanok
  \dcref{custom}

  \uses{pa-thm:unitsAppLE_diag_eq_one_of_coherent, pa-thm:unitsAppLE_eq_one_of_coherent,
    pa-lem:overSpecMap_left_degeneracy, pa-def:CoherentCechWitness}
  The witness cochain \(\theta\) of \(\theta'\) is trivial along the diagonal \(\Delta =
  \Spec(\mathrm{tensorMul})\): the \(\Delta\)-pullback of \(\theta\) is \(1\).
\end{theorem}
\begin{proof}
  \leanok
  Apply the diagonal-triviality transport (Theorem~\ref{pa-thm:unitsAppLE_diag_eq_one_of_coherent})
  with \(\delta = \Delta\) and \(\delta_r = q_1 \circ \Delta\); the required identity
  \(\Delta \circ (q_1 \circ \Delta) = \Delta\) follows from \(\Delta \circ q_1 = \mathrm{id}\)
  (Lemma~\ref{pa-lem:overSpecMap_left_degeneracy}). Its hypothesis — triviality of the
  \((q_1 \circ \Delta)\)-pullback of \(\theta\) — is
  Theorem~\ref{pa-thm:unitsAppLE_eq_one_of_coherent} for the degeneracy \(\sigma\) against the
  three cofaces, whose section identities \(\sigma \circ f_{23} = \mathrm{id}\),
  \(\sigma \circ f_{12} = q_1 \circ \Delta\), \(\sigma \circ f_{13} = \mathrm{id}\) are the
  degeneracy coincidences of Lemma~\ref{pa-lem:overSpecMap_left_degeneracy} and whose coherence
  hypothesis is the Amitsur coherence of \(\theta'\).
\end{proof}

\begin{theorem}[The Amitsur cocycle identity of the glued components]
  \label{pa-thm:witnessBasicSection_amitsur}
  \lean{AlgebraicGeometry.Over.witnessBasicSection_amitsur}
  \leanok
  \dcref{custom}

  \uses{pa-def:witnessBasicSection, pa-thm:glued_corr_amitsur, pa-lem:overSpec_whisker_face,
    pa-def:pairTripleSection, pa-def:CoherentCechWitness}
  Over the triple basic open of \(\Spec(B \otimes_A (B \otimes_A B))\), the coface pullbacks of
  the glued witness components satisfy the Amitsur cocycle identity:
  \[
    f_{23}^{\sharp}\,\mathrm{witnessBasicSection}(j, l) \cdot
      f_{12}^{\sharp}\,\mathrm{witnessBasicSection}(i, j)
      = f_{13}^{\sharp}\,\mathrm{witnessBasicSection}(i, l) .
  \]
\end{theorem}
\begin{proof}
  \leanok
  This is Theorem~\ref{pa-thm:glued_corr_amitsur} applied along the cofaces \(f_{23}, f_{12},
  f_{13}\) and coprojections \(q_1, q_2\), for the point pairs
  \((P.\mathrm{pt}_j, P.\mathrm{pt}_l)\), \((P.\mathrm{pt}_i, P.\mathrm{pt}_j)\),
  \((P.\mathrm{pt}_i, P.\mathrm{pt}_l)\) on the three double basic opens. Its simplicial
  coincidences \(f_{12} q_1 = f_{13} q_1\), \(f_{12} q_2 = f_{23} q_1\), \(f_{13} q_2 = f_{23}
  q_2\) are the coface coincidences of Lemma~\ref{pa-lem:overSpec_whisker_face}, and its coherence
  input is the Amitsur coherence of \(\theta'\); the containments of the triple basic open in
  the three coface preimages hold because \(\mathrm{tripleSection}\) is divisible by each coface
  pullback of the relevant \(\mathrm{pairSection}\).
\end{proof}

\begin{theorem}[The diagonal collapse of the glued components]
  \label{pa-thm:witnessBasicSection_collapse}
  \lean{AlgebraicGeometry.Over.witnessBasicSection_collapse}
  \leanok
  \dcref{custom}

  \uses{pa-def:witnessBasicSection, pa-thm:glued_corr_collapse, pa-thm:witnessAppLE_diag_eq_one,
    pa-lem:overSpecMap_left_degeneracy, pa-def:pairTripleSection}
  The \(\Delta\)-pullback of the glued witness component onto the Zariski overlap
  \(D(P.r_i \cdot P.r_j) \subseteq \Spec B\) is the Zariski cover cocycle value \(\gamma_{ij}\)
  of \(P\):
  \[
    \Delta^{\sharp}\,\mathrm{witnessBasicSection}(i, j) = \gamma_{ij} .
  \]
\end{theorem}
\begin{proof}
  \leanok
  This is Theorem~\ref{pa-thm:glued_corr_collapse} applied with the common section \(\delta =
  \Delta\) of \(q_1, q_2\), whose identities \(\Delta \circ q_1 = \mathrm{id} = \Delta \circ
  q_2\) are the degeneracy coincidences of Lemma~\ref{pa-lem:overSpecMap_left_degeneracy}. Its
  hypothesis of diagonal triviality of the witness is
  Theorem~\ref{pa-thm:witnessAppLE_diag_eq_one}; the point pair is \((P.\mathrm{pt}_i,
  P.\mathrm{pt}_j)\), and the target open \(D(P.r_i \cdot P.r_j)\) lands in \(\mathcal
  N_{P.\mathrm{pt}_i} \cap \mathcal N_{P.\mathrm{pt}_j}\) and in the \(\Delta\)-preimage of the
  double basic open. The resulting cocycle value \(\gamma(P.\mathrm{pt}_i, P.\mathrm{pt}_j)\)
  restricted to the overlap is, by definition, the Zariski cover cocycle \(\gamma_{ij}\) of the
  basic refinement.
\end{proof}

\section{Section rings on basic opens as two-base localizations}
\label{pa-sec:witnessAway}

To pass the three identities of Section~\ref{pa-sec:witnessComponents} to the pure-algebra
assembly of Section~\ref{pa-sec:piAssembly}, the section rings on the basic opens of the tensor
towers are identified with the algebraic tensor localizations. All section rings of
\(\Spec R\) carry the canonical \(R\)-algebra structure whose structure map is the inverse of
the global-sections isomorphism \(\Gamma\Spec\)-iso followed by restriction.

\begin{theorem}[Section rings on basic opens are localizations of the ring]
  \label{pa-thm:isLocalization_away_sections}
  \lean{AlgebraicGeometry.isLocalization_away_sections}
  \leanok
  \dcref{custom}

  For a commutative ring \(R\) and a global section \(g \in \Gamma(\Spec R, \top)\), the
  section ring on the basic open \(D(g)\) is the localization of \(R\) itself away from the
  element \(\Gamma\Spec\text{-}\mathrm{iso}(g) \in R\).
\end{theorem}
\begin{proof}
  \leanok
  The canonical \(\Gamma(\Spec R, \top)\)-algebra structure on \(\Gamma(\Spec R, D(g))\)
  exhibits it as the localization of \(\Gamma(\Spec R, \top)\) away from \(g\). Transport this
  along the ring isomorphism \(R \cong \Gamma(\Spec R, \top)\) given by the global-sections
  isomorphism: the powers of \(g\) correspond to the powers of its image in \(R\), and the two
  algebra structures agree, so \(\Gamma(\Spec R, D(g))\) is the localization of \(R\) away from
  that image.
\end{proof}

\begin{lemma}[The localization elements of the basic sections]
  \label{pa-lem:GammaSpecIso_hom_pairTripleSection}
  \lean{AlgebraicGeometry.Over.ΓSpecIso_hom_pairSection,
    AlgebraicGeometry.Over.ΓSpecIso_hom_tripleSection}
  \leanok
  \dcref{custom}

  \uses{pa-def:pairTripleSection, pa-def:tensorInl_tensorInr, pa-def:tensorFaces}
  Write \(b_i \in B\) for the element cutting out \(D(P.r_i)\). The global-sections
  isomorphism carries the double basic section to the two-base localization element
  \[
    \mathrm{pairSection}(i, j) \mapsto (b_i \otimes 1)(1 \otimes b_j)
    \quad\text{in } B \otimes_A B ,
  \]
  and the triple basic section to
  \[
    \mathrm{tripleSection}(i, j, l) \mapsto
      (b_i \otimes 1)\bigl(1 \otimes \bigl((b_j \otimes 1)(1 \otimes b_l)\bigr)\bigr)
    \quad\text{in } B \otimes_A (B \otimes_A B) .
  \]
\end{lemma}
\begin{proof}
  \leanok
  The global-sections isomorphism intertwines \(q_1^{\sharp}\) with \(\mathrm{tensorInl}\) and
  \(q_2^{\sharp}\) with \(\mathrm{tensorInr}\) (its naturality against
  \(\Spec(\mathrm{tensorInl})\), \(\Spec(\mathrm{tensorInr})\)), so it sends
  \(q_1^{\sharp}(P.r_i)\) to \(b_i \otimes 1\) and \(q_2^{\sharp}(P.r_j)\) to \(1 \otimes b_j\);
  multiplying gives the double element. For the triple element, apply naturality once more
  against the cofaces \(f_{12}, f_{13}\), computing the three insertions
  \(f_{12}^{\sharp} q_1^{\sharp}(P.r_i)\), \(f_{12}^{\sharp} q_2^{\sharp}(P.r_j)\),
  \(f_{13}^{\sharp} q_2^{\sharp}(P.r_l)\) as \(b_i \otimes (1 \otimes 1)\),
  \(1 \otimes (b_j \otimes 1)\), \(1 \otimes (1 \otimes b_l)\), and rearranging the product.
\end{proof}

\begin{definition}[The two-base localization identifications]
  \label{pa-def:pairTripleAwayEquiv}
  \lean{AlgebraicGeometry.Over.pairAwayEquiv, AlgebraicGeometry.Over.tripleAwayEquiv}
  \leanok
  \dcref{custom}

  \uses{pa-def:tensorAwayEquiv, pa-thm:isLocalization_away_sections, pa-thm:isLocalization_away_tensor,
    pa-lem:GammaSpecIso_hom_pairTripleSection, pa-def:pairTripleSection}
  Let \(S_i := \Gamma(\Spec B, D(P.r_i))\), a localization of \(B\) away from \(b_i\)
  (Theorem~\ref{pa-thm:isLocalization_away_sections}). The comparison of tensor localizations
  (Definition~\ref{pa-def:tensorAwayEquiv}) gives the two-base identifications
  \[
    \mathrm{pairAwayEquiv} \colon S_i \otimes_A S_j \;\xrightarrow{\ \sim\ }\;
      \Gamma\bigl(\Spec(B \otimes_A B), D(\mathrm{pairSection}(i, j))\bigr), \qquad
  \]
  \[
    \mathrm{tripleAwayEquiv} \colon S_i \otimes_A (S_j \otimes_A S_l) \;\xrightarrow{\ \sim\ }\;
      \Gamma\bigl(\Spec(B \otimes_A (B \otimes_A B)), D(\mathrm{tripleSection}(i, j, l))\bigr) ,
  \]
  the target section rings being the localizations of \(B \otimes_A B\), resp.\
  \(B \otimes_A (B \otimes_A B)\), at the two-base elements of
  Lemma~\ref{pa-lem:GammaSpecIso_hom_pairTripleSection}. For the triple identification the inner
  factor \(S_j \otimes_A S_l\) is first recognized as a localization of \(B \otimes_A B\)
  (Theorem~\ref{pa-thm:isLocalization_away_tensor}), so that
  \(S_i \otimes_A (S_j \otimes_A S_l)\) is a two-base localization over
  \(B \otimes_A (B \otimes_A B)\).
\end{definition}

\section{Transport of the simplicial maps through the two-base identifications}
\label{pa-sec:witnessTransport}

The restriction maps of the section rings along the diagonal \(\Delta\) and the three cofaces
correspond, under the identifications of Definition~\ref{pa-def:pairTripleAwayEquiv}, to the
index-wise collapse and faces of the pure-algebra assembly. Each is an instance of the
uniqueness of maps out of a localization: both composites are computed on structure-map
images and agree there, and extensionality is invoked over \(B \otimes_A B\) — where the
components are \(\mathrm{Away}\) localizations — never over \(A\).

\begin{lemma}[The diagonal pullback is the index-wise collapse]
  \label{pa-lem:pairAwayEquiv_appLE_diag}
  \lean{AlgebraicGeometry.Over.pairAwayEquiv_appLE_diag}
  \leanok
  \dcref{custom}

  \uses{pa-def:pairTripleAwayEquiv, pa-def:componentCollapse, pa-thm:isLocalization_away_tensor,
    pa-lem:overSpecMap_left_degeneracy}
  Under the identification \(\mathrm{pairAwayEquiv}\), the \(\Delta\)-pullback
  \(\Gamma(D(\mathrm{pairSection}(i, j))) \to \Gamma(D(P.r_i \cdot P.r_j))\) is the index-wise
  collapse \(\kappa_{ij} \colon S_i \otimes_A S_j \to T_{ij}\) onto the Zariski overlap ring
  \(T_{ij} = \Gamma(\Spec B, D(P.r_i \cdot P.r_j))\).
\end{lemma}
\begin{proof}
  \leanok
  Both maps are \(A\)-algebra homomorphisms out of the localization \(S_i \otimes_A S_j\) of
  \(B \otimes_A B\) (Theorem~\ref{pa-thm:isLocalization_away_tensor}), so by uniqueness of maps out
  of a localization it suffices to compare their composites with the structure map
  \(B \otimes_A B \to S_i \otimes_A S_j\). On a structure-map image the \(\Delta\)-pullback
  computes, through the naturality of the global-sections isomorphism against
  \(\Spec(\mathrm{tensorMul})\) and the diagonal coincidence
  (Lemma~\ref{pa-lem:overSpecMap_left_degeneracy}), as the multiplication \(B \otimes_A B \to B\)
  followed by the structure map into \(T_{ij}\); and the index-wise collapse on a structure-map
  image is that same composite (Definition~\ref{pa-def:componentCollapse}). The two agree, hence so
  do the maps.
\end{proof}

\begin{lemma}[The coface pullbacks are the index-wise faces]
  \label{pa-lem:tripleAwayEquiv_faceA}
  \lean{AlgebraicGeometry.Over.tripleAwayEquiv_faceA₂₃,
    AlgebraicGeometry.Over.tripleAwayEquiv_faceA₁₂,
    AlgebraicGeometry.Over.tripleAwayEquiv_faceA₁₃}
  \leanok
  \dcref{custom}

  \uses{pa-def:pairTripleAwayEquiv, pa-def:indexwiseMaps, pa-lem:faceA_comp_tensorMap,
    pa-thm:isLocalization_away_tensor, pa-lem:overSpec_whisker_face, pa-def:pairTripleSection}
  Under the identifications \(\mathrm{pairAwayEquiv}\) and \(\mathrm{tripleAwayEquiv}\), the
  three coface pullbacks \(f_{23}^{\sharp}, f_{12}^{\sharp}, f_{13}^{\sharp}\) between the
  basic-open section rings are the index-wise faces \(\mathrm{face}_{23}, \mathrm{face}_{12},
  \mathrm{face}_{13}\) of Definition~\ref{pa-def:indexwiseMaps}.
\end{lemma}
\begin{proof}
  \leanok
  Fix the coface \(\bullet \in \{12, 13, 23\}\). Both \(\mathrm{tripleAwayEquiv} \circ
  \mathrm{face}_{\bullet}\) and \(f_{\bullet}^{\sharp} \circ \mathrm{pairAwayEquiv}\) are
  \(A\)-algebra homomorphisms out of the localization \(S_j \otimes_A S_l\) (resp.\ \(S_i
  \otimes_A S_j\), \(S_i \otimes_A S_l\)) of \(B \otimes_A B\)
  (Theorem~\ref{pa-thm:isLocalization_away_tensor}), so by uniqueness of maps out of a localization
  it suffices to compare their composites with the structure map \(B \otimes_A B \to S_\bullet
  \otimes_A S_\bullet\). On a structure-map image the coface pullback computes, through the
  naturality of the global-sections isomorphism and the coface coincidences
  (Lemma~\ref{pa-lem:overSpec_whisker_face}), as the descent coface followed by the triple
  structure map; and the index-wise face intertwines the double and triple structure maps over
  that same descent coface (Lemma~\ref{pa-lem:faceA_comp_tensorMap}). The two composites agree.
\end{proof}

\section{The composite descent unit of the \'etale-separatedness assembly}
\label{pa-sec:witnessAssembly}

Assembling the glued witness components into the composite \(A\)-descent unit yields the two
halves of the deliverable and the headline class equation handed onward. Write \(S_i :=
\Gamma(\Spec B, D(P.r_i))\) and \(P_{\!\bullet} := \prod_i S_i\).

\begin{definition}[The witness component]
  \label{pa-def:witnessComponent}
  \lean{AlgebraicGeometry.Over.witnessComponent}
  \leanok
  \dcref{custom}

  \uses{pa-def:witnessBasicSection, pa-def:pairTripleAwayEquiv}
  The \emph{witness component} \(w_{ij} \in (S_i \otimes_A S_j)^{\times}\) is the glued witness
  component \(\mathrm{witnessBasicSection}(i, j)\) transported through the two-base
  identification \(\mathrm{pairAwayEquiv}\) (Definition~\ref{pa-def:pairTripleAwayEquiv}).
\end{definition}

\begin{definition}[The composite descent unit of the assembly]
  \label{pa-def:assemblyUnit}
  \lean{AlgebraicGeometry.Over.assemblyUnit}
  \leanok
  \dcref{custom}

  \uses{pa-def:witnessComponent, pa-def:piAssemblyUnit}
  The \emph{composite descent unit} \(\mathrm{assemblyUnit} \in (P_{\!\bullet} \otimes_A
  P_{\!\bullet})^{\times}\) is the pi-assembly (Definition~\ref{pa-def:piAssemblyUnit}) of the
  family of witness components \((w_{ij})\).
\end{definition}

\begin{theorem}[The witness components collapse onto the Zariski cover cocycle]
  \label{pa-thm:componentCollapse_witnessComponent}
  \lean{AlgebraicGeometry.Over.componentCollapse_witnessComponent}
  \leanok
  \dcref{custom}

  \uses{pa-def:witnessComponent, pa-lem:pairAwayEquiv_appLE_diag, pa-thm:witnessBasicSection_collapse,
    pa-def:componentCollapse}
  The index-wise collapse of the witness component is the Zariski cover cocycle value:
  \(\kappa_{ij}(w_{ij}) = \gamma_{ij}\).
\end{theorem}
\begin{proof}
  \leanok
  By Lemma~\ref{pa-lem:pairAwayEquiv_appLE_diag} the collapse \(\kappa_{ij}\) is, through
  \(\mathrm{pairAwayEquiv}\), the \(\Delta\)-pullback; applied to the witness component it is the
  \(\Delta\)-pullback of \(\mathrm{witnessBasicSection}(i, j)\), which is \(\gamma_{ij}\) by
  Theorem~\ref{pa-thm:witnessBasicSection_collapse}.
\end{proof}

\begin{theorem}[The witness components satisfy the index-wise cocycle identity]
  \label{pa-thm:faceA_cocycle_witnessComponent}
  \lean{AlgebraicGeometry.Over.faceA_cocycle_witnessComponent}
  \leanok
  \dcref{custom}

  \uses{pa-def:witnessComponent, pa-lem:tripleAwayEquiv_faceA, pa-thm:witnessBasicSection_amitsur,
    pa-def:indexwiseMaps}
  The witness components satisfy the index-wise cocycle identity:
  \(\mathrm{face}_{23}(w_{jl}) \cdot \mathrm{face}_{12}(w_{ij}) = \mathrm{face}_{13}(w_{il})\).
\end{theorem}
\begin{proof}
  \leanok
  Apply the injective identification \(\mathrm{tripleAwayEquiv}\). By
  Lemma~\ref{pa-lem:tripleAwayEquiv_faceA} it carries each index-wise face of a witness component
  to the corresponding coface pullback of the glued component
  \(\mathrm{witnessBasicSection}\); the transported identity is then the Amitsur cocycle
  identity of the glued components (Theorem~\ref{pa-thm:witnessBasicSection_amitsur}).
\end{proof}

\begin{theorem}[The assembled witness unit is an \(A\)-descent \(1\)-cocycle]
  \label{pa-thm:isDescentCocycle_assemblyUnit}
  \lean{AlgebraicGeometry.Over.isDescentCocycle_assemblyUnit}
  \leanok
  \source{kleiman-picard}

  \uses{pa-def:assemblyUnit, pa-thm:isDescentCocycle_piAssemblyUnit,
    pa-thm:diagA_eq_one_of_componentCollapse, pa-thm:componentCollapse_witnessComponent,
    pa-thm:faceA_cocycle_witnessComponent}
  \dcref{custom}
  The composite descent unit \(\mathrm{assemblyUnit}\) is an \(A\)-descent \(1\)-cocycle
  relative to \(A \to P_{\!\bullet}\).
\end{theorem}
\begin{proof}
  \leanok
  Apply the descent-cocycle criterion for a pi-assembly
  (Theorem~\ref{pa-thm:isDescentCocycle_piAssemblyUnit}). Its index-wise cocycle hypothesis is
  Theorem~\ref{pa-thm:faceA_cocycle_witnessComponent}. Its diagonal-normalization hypothesis is
  automatic: the witness components collapse onto the Zariski cover cocycle
  (Theorem~\ref{pa-thm:componentCollapse_witnessComponent}), whose diagonal is \(1\) since the
  cover cocycle is a genuine \v Cech cocycle, so
  Theorem~\ref{pa-thm:diagA_eq_one_of_componentCollapse} gives \(\mathrm{diag}(w_{ii}) = 1\). This
  is precisely the descent construction of Kleiman's separatedness argument: the assembled unit
  \(v'_{23} v'_{12} = v'_{13}\) descends the componentwise trivializations to a single datum on
  the base.
\end{proof}

\begin{theorem}[The collapse of the assembled unit is the descent unit of the cover cocycle]
  \label{pa-thm:tensorCollapse_assemblyUnit}
  \lean{AlgebraicGeometry.Over.tensorCollapse_assemblyUnit}
  \leanok
  \dcref{custom}

  \uses{pa-def:assemblyUnit, pa-thm:tensorCollapse_piAssemblyUnit,
    pa-thm:componentCollapse_witnessComponent, pa-def:module_tensorCollapse}
  The collapse of the composite descent unit along the tower \(A \to \Gamma(\Spec B, \top) \to
  P_{\!\bullet}\) is the descent unit of the Zariski cover cocycle of \(P\):
  \[
    \pi\bigl(\mathrm{assemblyUnit}\bigr) = \mathrm{cocycleUnit}(\gamma_{\bullet}) .
  \]
\end{theorem}
\begin{proof}
  \leanok
  This is the collapse identity for a pi-assembly
  (Theorem~\ref{pa-thm:tensorCollapse_piAssemblyUnit}), whose hypothesis — that the witness
  components collapse onto the cover cocycle, \(\kappa_{ij}(w_{ij}) = \gamma_{ij}\) — is
  Theorem~\ref{pa-thm:componentCollapse_witnessComponent}; the intermediate ring of the collapse
  tower is the scheme-side base ring \(\Gamma(\Spec B, \top)\).
\end{proof}

\begin{lemma}[Faithful flatness of the global sections of \(\Spec B\)]
  \label{pa-lem:faithfullyFlat_sectionsTop}
  \lean{AlgebraicGeometry.Over.sectionsTopAlgEquiv,
    AlgebraicGeometry.Over.faithfullyFlat_sectionsTop}
  \leanok
  \dcref{custom}

  The global-sections isomorphism is an isomorphism of \(A\)-algebras \(\Gamma(\Spec B, \top)
  \cong B\) for the composed \(A\)-structure. Consequently, if \(A \to B\) is faithfully flat,
  then \(A \to \Gamma(\Spec B, \top)\) is faithfully flat.
\end{lemma}
\begin{proof}
  \leanok
  The global-sections isomorphism is a ring isomorphism commuting with the two structure maps
  from \(A\) (both send \(a\) to the image of \(a\), via restriction of \(\Gamma\Spec\)-iso),
  hence an \(A\)-algebra isomorphism. Faithful flatness of a module transports along a linear
  isomorphism, so faithful flatness of \(A \to B\) gives faithful flatness of \(A \to
  \Gamma(\Spec B, \top)\).
\end{proof}

\begin{theorem}[The class equation of the composite descent unit]
  \label{pa-thm:mapAlgebra_picClass_assemblyUnit}
  \lean{AlgebraicGeometry.Over.mapAlgebra_picClass_assemblyUnit}
  \leanok
  \source{kleiman-picard}

  \uses{pa-thm:isDescentCocycle_assemblyUnit, pa-thm:tensorCollapse_assemblyUnit,
    pa-thm:module_picClass_collapse, pa-lem:faithfullyFlat_sectionsTop, pa-thm:pi_tower,
    pa-def:module_picClass, pa-thm:trivializingFamily_pic_cocycle}
  \dcref{custom}
  Assume \(A \to B\) is faithfully flat. Then the Picard class of the composite descent unit
  maps, under the base-change homomorphism \(\Pic(A) \to \Pic(\Gamma(\Spec B, \top))\) along
  \(A \to \Gamma(\Spec B, \top)\), to the Picard class of the Zariski cover cocycle of \(P\):
  \[
    \Pic(A \to \Gamma(\Spec B, \top))\bigl(\mathrm{picClass}(\mathrm{assemblyUnit})\bigr)
      = \mathrm{pic}(\gamma) ,
  \]
  which is \([N]\) when \(\gamma\) is the trivializing cocycle of an invertible module \(N\).
\end{theorem}
\begin{proof}
  \leanok
  Set \(P_0 := \Gamma(\Spec B, \top)\). The tower \(A \to P_0 \to P_{\!\bullet}\) is faithfully
  flat: \(A \to P_0\) by Lemma~\ref{pa-lem:faithfullyFlat_sectionsTop}, and \(P_0 \to
  P_{\!\bullet}\) because \(P_{\!\bullet} = \prod_i S_i\) is the product of the localizations of
  \(P_0\) at the covering family \(P.r\) (Theorem~\ref{pa-thm:pi_tower}); hence \(A \to
  P_{\!\bullet}\) is faithfully flat. By naturality of the Picard class under descent in stages
  (Theorem~\ref{pa-thm:module_picClass_collapse}) applied to \(\mathrm{assemblyUnit}\)
  (Theorem~\ref{pa-thm:isDescentCocycle_assemblyUnit}), the base change to \(P_0\) of
  \(\mathrm{picClass}(\mathrm{assemblyUnit})\) is the Picard class over \(P_0\) of the collapse
  \(\pi(\mathrm{assemblyUnit})\). By Theorem~\ref{pa-thm:tensorCollapse_assemblyUnit} that collapse
  is the descent unit of the Zariski cover cocycle, whose Picard class is \(\mathrm{pic}(\gamma)\)
  by definition (Definition~\ref{pa-def:module_picClass}); and when \(\gamma\) trivializes an
  invertible module \(N\), \(\mathrm{pic}(\gamma) = [N]\) by
  Theorem~\ref{pa-thm:trivializingFamily_pic_cocycle}. This is the invertible descended sheaf of
  Kleiman's separatedness argument, in its \v Cech incarnation.
\end{proof}

\section{Faithfully flat descent of sections along the base change of the curve product}
\label{pa-sec:sectionsDescent}

Fix a field \(k\), a curve \(C\) over \(k\) (proper, geometrically irreducible and geometrically
reduced), and a tower \(k \to A \to B\) with \(A \to B\) faithfully flat. Write \(X_A := (C
\otimes \Spec A)_{\mathrm{left}}\), \(X_B := (C \otimes \Spec B)_{\mathrm{left}}\), and
\(\mathrm{cg} := (C \lhd \Spec(A \to B))_{\mathrm{left}} \colon X_B \to X_A\) for the base-change
morphism of the curve product; the two coprojections \(u_1, u_2 \colon (C \otimes \Spec(B
\otimes_A B))_{\mathrm{left}} \to X_B\) are the whiskered coprojections of \(B \otimes_A B\).
This section proves that sections of the structure sheaf, and of its units, descend along
\(\mathrm{cg}\): the fppf sheaf property along the whiskered base inclusion.

\begin{theorem}[The base-change square of the product is a pullback]
  \label{pa-thm:isPullback_whiskerLeft_snd}
  \lean{AlgebraicGeometry.Over.isPullback_whiskerLeft_snd}
  \leanok
  \dcref{custom}

  \uses{pa-lem:snd_left_naturality}
  For \(C, T, T'\) over \(\Spec k\) and \(g \colon T' \to T\), the square with horizontal legs
  \((C \lhd g)_{\mathrm{left}}\) and \(g_{\mathrm{left}}\) and vertical legs the two second
  projections is a pullback of schemes:
  \[
    \begin{array}{ccc}
      (C \otimes T')_{\mathrm{left}} & \xrightarrow{(C \lhd g)_{\mathrm{left}}} &
        (C \otimes T)_{\mathrm{left}} \\
      \downarrow \mathrm{pr}' & & \downarrow \mathrm{pr} \\
      T'_{\mathrm{left}} & \xrightarrow{\ g_{\mathrm{left}}\ } & T_{\mathrm{left}}.
    \end{array}
  \]
\end{theorem}
\begin{proof}
  \leanok
  Paste the two cartesian squares expressing \(C \otimes T\) and \(C \otimes T'\) as products
  over \(\Spec k\).  The first-projection triangle \((C \lhd g)_{\mathrm{left}} \circ
  (\mathrm{fst}_{C,T})_{\mathrm{left}} = (\mathrm{fst}_{C,T'})_{\mathrm{left}}\) and the base
  identity \(g_{\mathrm{left}} \circ T.\mathrm{hom} = T'.\mathrm{hom}\) exhibit the outer
  rectangle over \(C \times \Spec k\) as the product square of \(C \otimes T'\); together with
  the naturality square of the second projection (Lemma~\ref{pa-lem:snd_left_naturality}), the
  right-cancellation property of pullbacks gives that the second-projection square is itself a
  pullback.
\end{proof}

\begin{theorem}[Descent of sections along the base change of the curve product]
  \label{pa-thm:existsUnique_appLE_eq}
  \lean{AlgebraicGeometry.Over.existsUnique_appLE_eq}
  \leanok
  \dcref{custom}

  \uses{pa-thm:isPullback_whiskerLeft_snd, pa-thm:descent_existsUnique_tmul_one,
    pa-def:tensorInl_tensorInr}
  Let \(U \subseteq X_A\) be an affine open. A section \(t \in \Gamma(X_B, \mathrm{cg}^{-1}U)\)
  whose two coprojection pullbacks to \(C \otimes \Spec(B \otimes_A B)\) agree,
  \(u_1^{*} t = u_2^{*} t\), descends uniquely along \(\mathrm{cg}\): there is a unique \(s \in
  \Gamma(X_A, U)\) with \(\mathrm{cg}^{*} s = t\).
\end{theorem}
\begin{proof}
  \leanok
  Over the affine open \(U\), the base-change square (Theorem~\ref{pa-thm:isPullback_whiskerLeft_snd})
  identifies the sections of \(X_B\) over \(\mathrm{cg}^{-1}U\) with the base change \(\Gamma(X_A,
  U) \otimes_A B\), and those of \(C \otimes \Spec(B \otimes_A B)\) with \(\Gamma(X_A, U)
  \otimes_A (B \otimes_A B)\), the two coprojection pullbacks becoming the two faces
  \(\id \otimes (b \mapsto b \otimes 1)\) and \(\id \otimes (b \mapsto 1 \otimes b)\)
  (Definition~\ref{pa-def:tensorInl_tensorInr}).  Under these identifications \(t\) corresponds to
  an element \(x \in \Gamma(X_A, U) \otimes_A B\) with equal faces, so by the degree-zero Amitsur
  equalizer (Theorem~\ref{pa-thm:descent_existsUnique_tmul_one}) there is a unique \(s \in \Gamma(X_A,
  U)\) with \(s \otimes 1 = x\); this \(s\) is the unique section with \(\mathrm{cg}^{*} s = t\).
\end{proof}

\begin{theorem}[Injectivity of pullback along the base change]
  \label{pa-thm:appLE_whiskerLeft_injective}
  \lean{AlgebraicGeometry.Over.appLE_whiskerLeft_injective}
  \leanok
  \dcref{custom}

  \uses{pa-thm:isPullback_whiskerLeft_snd, pa-lem:units_separation}
  For every open \(W \subseteq X_A\), pullback of sections along \(\mathrm{cg}\) is injective,
  \(\Gamma(X_A, W) \to \Gamma(X_B, \mathrm{cg}^{-1}W)\).
\end{theorem}
\begin{proof}
  \leanok
  First over an affine open \(V\): by the base-change square
  (Theorem~\ref{pa-thm:isPullback_whiskerLeft_snd}) the pullback is \(s \mapsto s \otimes 1\) into
  \(\Gamma(X_A, V) \otimes_A B\), which is injective because \(B\) is faithfully flat over
  \(A\).  For an arbitrary open \(W\), if two sections have equal pullbacks, then their
  restrictions to each affine open \(V \subseteq W\) have equal pullbacks (pullback commutes
  with restriction), hence agree by the affine case; as the affine opens contained in \(W\)
  cover \(W\), the two sections agree by separation for the structure sheaf
  (Lemma~\ref{pa-lem:units_separation}).
\end{proof}

\begin{theorem}[Descent of unit sections along the base change]
  \label{pa-thm:exists_unitsAppLE_eq}
  \lean{AlgebraicGeometry.Over.exists_unitsAppLE_eq}
  \leanok
  \dcref{custom}

  \uses{pa-thm:existsUnique_appLE_eq, pa-thm:appLE_whiskerLeft_injective}
  Let \(U \subseteq X_A\) be an affine open. A unit \(t \in \Gamma(X_B,
  \mathrm{cg}^{-1}U)^{\times}\) whose two coprojection pullbacks agree, \(u_1^{*} t = u_2^{*}
  t\), descends to a unit: there is \(s \in \Gamma(X_A, U)^{\times}\) with \(\mathrm{cg}^{*} s =
  t\).
\end{theorem}
\begin{proof}
  \leanok
  Both \(t\) and \(t^{-1}\) satisfy the equalizer hypothesis, so by
  Theorem~\ref{pa-thm:existsUnique_appLE_eq} they descend to sections \(a, b \in \Gamma(X_A, U)\)
  with \(\mathrm{cg}^{*} a = t\), \(\mathrm{cg}^{*} b = t^{-1}\).  The pullback of \(ab\) is \(t
  t^{-1} = 1 = \mathrm{cg}^{*} 1\), so \(ab = 1\) by injectivity of the pullback over the affine
  open \(U\) (Theorem~\ref{pa-thm:appLE_whiskerLeft_injective}); likewise \(ba = 1\).  Hence \(a\)
  is a unit and \(\mathrm{cg}^{*} a = t\).
\end{proof}

\section{Saturated affine pieces of the base-change cover}
\label{pa-sec:effectivityPieces}

Retain the field \(k\), the curve \(C\) over \(k\), and the tower \(k \to A \to B\) of the
previous section, with \(X_A = (C \otimes \Spec A)_{\mathrm{left}}\), \(X_B = (C \otimes \Spec
B)_{\mathrm{left}}\), the projections \(\mathrm{pr}_A \colon X_A \to \Spec A\), \(\mathrm{pr}_B
\colon X_B \to \Spec B\), and the base-change morphism \(\mathrm{cg} \colon X_B \to X_A\). This
section refines the \(\top\)-case identification \(\Gamma(X_B, \top) \cong B\)
(Lemma~\ref{pa-lem:faithfullyFlat_sectionsTop}) to an arbitrary affine open of \(X_A\): over such an
open the section ring of the cover piece is the base change of the base section ring along \(A
\to B\). Nothing here uses flatness or properness (except where explicitly assumed); these are
the geometric pieces on which a covering-and-overlaps effectivity argument runs.

\begin{theorem}[The cover inclusion is an affine morphism]
  \label{pa-thm:isAffineHom_cg}
  \lean{AlgebraicGeometry.Over.isAffineHom_cg, AlgebraicGeometry.Over.isAffineOpen_cgPreimage}
  \leanok
  \dcref{custom}

  \uses{pa-thm:isPullback_whiskerLeft_snd}
  The base-change morphism \(\mathrm{cg} \colon X_B \to X_A\) is affine. Consequently the
  preimage \(\mathrm{cg}^{-1}U\) of every affine open \(U \subseteq X_A\) is affine.
\end{theorem}
\begin{proof}
  \leanok
  By the base-change square (Theorem~\ref{pa-thm:isPullback_whiskerLeft_snd}, with \(T = \Spec A\),
  \(T' = \Spec B\) and \(g = \Spec(A \to B)\)) the morphism \(\mathrm{cg}\) is the base change of
  \(\Spec(A \to B) \colon \Spec B \to \Spec A\) along the projection \(\mathrm{pr}_A\). The
  morphism \(\Spec(A \to B)\) is affine, being the spectrum of a ring homomorphism, and affine
  morphisms are stable under base change; hence \(\mathrm{cg}\) is affine. The preimage of an
  affine open under an affine morphism is affine.
\end{proof}

\begin{theorem}[The section rings of a saturated affine piece form a pushout square]
  \label{pa-thm:isPushout_cgSectionRing}
  \lean{AlgebraicGeometry.Over.isPushout_cgSections,
    AlgebraicGeometry.Over.isPushout_cgAlgebraMap}
  \leanok
  \dcref{custom}
  \uses{pa-thm:isAffineHom_cg, pa-thm:isPullback_whiskerLeft_snd}

  Give \(\Gamma(X_A, U)\) its \(A\)-algebra structure by pullback along the projection
  \(\mathrm{pr}_A\). For every affine open \(U \subseteq X_A\) the square of commutative rings
  \[
    \begin{array}{ccc}
      A & \longrightarrow & \Gamma(X_A, U) \\
      \downarrow & & \downarrow \mathrm{cg}^{\sharp} \\
      B & \longrightarrow & \Gamma\bigl(X_B, \mathrm{cg}^{-1}U\bigr)
    \end{array}
  \]
  --- with top and left legs the two structure maps from \(A\), right leg the pullback
  \(\mathrm{cg}^{\sharp}\), and bottom leg the pullback along \(\mathrm{pr}_B\) --- is a pushout
  (cocartesian). No flatness or properness is used.
\end{theorem}
\begin{proof}
  \leanok
  The base-change square of schemes (Theorem~\ref{pa-thm:isPullback_whiskerLeft_snd}) has one leg,
  \(\Spec(A \to B)\), affine, and \(U\) is affine; for such a fibre square the induced square of
  section rings over \(U\) --- with the two \(\top\)-corners \(\Gamma(\Spec A, \top)\) and
  \(\Gamma(\Spec B, \top)\) --- is a pushout. Re-cornering along the canonical identifications
  \(\Gamma(\Spec A, \top) \cong A\) and \(\Gamma(\Spec B, \top) \cong B\) (which intertwine the
  structure maps, as in Lemma~\ref{pa-lem:faithfullyFlat_sectionsTop}) replaces those two corners by
  \(A\) and \(B\), yielding the displayed pushout.
\end{proof}

\begin{definition}[The section-ring identification]
  \label{pa-def:pieceRingEquiv}
  \lean{AlgebraicGeometry.Over.pieceTensorIso, AlgebraicGeometry.Over.pieceRingEquiv,
    AlgebraicGeometry.Over.pieceRingEquiv_tmul_one, AlgebraicGeometry.Over.pieceRingEquiv_one_tmul,
    AlgebraicGeometry.Over.pieceRingEquiv_tmul}
  \leanok
  \dcref{custom}
  \uses{pa-thm:isPushout_cgSectionRing}

  For every affine open \(U \subseteq X_A\) there is a ring isomorphism
  \[
    \Gamma(X_A, U) \otimes_A B \;\xrightarrow{\ \sim\ }\; \Gamma\bigl(X_B, \mathrm{cg}^{-1}U\bigr),
    \qquad s \otimes b \;\longmapsto\; \mathrm{cg}^{\sharp}(s)\cdot \mathrm{pr}_B^{\sharp}(b),
  \]
  the affine-open generalization of the \(\top\)-case \(\Gamma(X_B, \top) \cong B\)
  (Lemma~\ref{pa-lem:faithfullyFlat_sectionsTop}). In particular \(s \otimes 1 \mapsto
  \mathrm{cg}^{\sharp}(s)\) and \(1 \otimes b \mapsto \mathrm{pr}_B^{\sharp}(b)\).
\end{definition}
\begin{proof}
  \leanok
  The base-change tensor product \(\Gamma(X_A, U) \otimes_A B\) is the pushout of \(\Gamma(X_A, U)
  \leftarrow A \to B\) in commutative rings. Comparing it with the pushout square of
  Theorem~\ref{pa-thm:isPushout_cgSectionRing}, which has the same span, yields a unique ring
  isomorphism between the two apices; its universal property sends the image of \(\Gamma(X_A, U)\)
  to \(\mathrm{cg}^{\sharp}(s)\) and the image of \(B\) to \(\mathrm{pr}_B^{\sharp}(b)\), hence
  \(s \otimes b\) to their product.
\end{proof}

\begin{definition}[The identification is an isomorphism of base-piece algebras]
  \label{pa-def:pieceAlgEquiv}
  \lean{AlgebraicGeometry.Over.pieceCoverAlgebra, AlgebraicGeometry.Over.pieceTensorAlgEquiv,
    AlgebraicGeometry.Over.pieceEquiv}
  \leanok
  \dcref{custom}
  \uses{pa-def:pieceRingEquiv}

  Give the cover-piece ring \(\Gamma(X_B, \mathrm{cg}^{-1}U)\) its \(\Gamma(X_A,
  U)\)-algebra structure by pullback along \(\mathrm{cg}\). Then the identification of
  Definition~\ref{pa-def:pieceRingEquiv} is an isomorphism of \(\Gamma(X_A, U)\)-algebras
  \[
    \Gamma\bigl(X_B, \mathrm{cg}^{-1}U\bigr) \;\cong\; \Gamma(X_A, U) \otimes_A B .
  \]
\end{definition}
\begin{proof}
  \leanok
  The structure map of the \(\Gamma(X_A, U)\)-algebra \(\Gamma(X_B, \mathrm{cg}^{-1}U)\) is
  \(s \mapsto \mathrm{cg}^{\sharp}(s)\), which the identification carries to \(s \otimes 1\)
  (Definition~\ref{pa-def:pieceRingEquiv}); this is the structure map of the base-change algebra
  \(\Gamma(X_A, U) \otimes_A B\). So the ring isomorphism intertwines the two structure maps from
  \(\Gamma(X_A, U)\), hence is an isomorphism of \(\Gamma(X_A, U)\)-algebras.
\end{proof}

\begin{theorem}[Faithful flatness of a saturated affine piece]
  \label{pa-thm:faithfullyFlat_pieceCover}
  \lean{AlgebraicGeometry.Over.faithfullyFlat_pieceCover}
  \leanok
  \dcref{custom}
  \uses{pa-def:pieceAlgEquiv}

  Assume \(A \to B\) is faithfully flat. Then for every affine open \(U \subseteq X_A\) the
  structure map \(\Gamma(X_A, U) \to \Gamma(X_B, \mathrm{cg}^{-1}U)\) is faithfully flat.
\end{theorem}
\begin{proof}
  \leanok
  Faithful flatness is stable under base change, so \(\Gamma(X_A, U) \to \Gamma(X_A, U)
  \otimes_A B\) is faithfully flat. Transporting along the isomorphism of \(\Gamma(X_A,
  U)\)-algebras of Definition~\ref{pa-def:pieceAlgEquiv} gives faithful flatness of \(\Gamma(X_A, U)
  \to \Gamma(X_B, \mathrm{cg}^{-1}U)\).
\end{proof}

\begin{theorem}[Restriction naturality of the identification]
  \label{pa-thm:pieceRingEquiv_naturality}
  \lean{AlgebraicGeometry.Over.pieceRingEquiv_naturality}
  \leanok
  \dcref{custom}
  \uses{pa-def:pieceRingEquiv}

  Let \(W \le V\) be affine opens of \(X_A\), so \(\mathrm{cg}^{-1}W \le \mathrm{cg}^{-1}V\). The
  identification of Definition~\ref{pa-def:pieceRingEquiv} commutes with restriction: for every \(x
  \in \Gamma(X_A, V) \otimes_A B\),
  \[
    \bigl(\mathrm{cg}^{-1}W \le \mathrm{cg}^{-1}V\bigr)^{*}\bigl(\Phi_V(x)\bigr)
      = \Phi_W\bigl((\rho^{V}_{W} \otimes \id_B)(x)\bigr),
  \]
  where \(\Phi_V, \Phi_W\) are the identifications over \(V, W\) and \(\rho^{V}_{W} \colon
  \Gamma(X_A, V) \to \Gamma(X_A, W)\) is restriction (an \(A\)-algebra homomorphism).
\end{theorem}
\begin{proof}
  \leanok
  Both sides are \(A\)-algebra homomorphisms out of \(\Gamma(X_A, V) \otimes_A B\), so it
  suffices to check them on pure tensors \(s \otimes b\). There the left side restricts
  \(\mathrm{cg}^{\sharp}(s)\cdot \mathrm{pr}_B^{\sharp}(b)\) and the right side is
  \(\mathrm{cg}^{\sharp}(\rho^{V}_{W} s)\cdot \mathrm{pr}_B^{\sharp}(b)\); the two agree because
  restriction of sections commutes with the two pullbacks \(\mathrm{cg}^{\sharp}\) and
  \(\mathrm{pr}_B^{\sharp}\) (pullback of a section, then restrict, equals restrict, then
  pullback).
\end{proof}

\begin{theorem}[Saturated affine opens form a basis]
  \label{pa-thm:exists_isAffineOpen_cg_le}
  \lean{AlgebraicGeometry.Over.exists_isAffineOpen_cg_le}
  \leanok
  \dcref{custom}
  \uses{pa-thm:isAffineHom_cg}

  Every neighbourhood \(V\) of a point \(x\) of \(X_A\) contains an affine open \(W\) with \(x \in
  W\) whose preimage \(\mathrm{cg}^{-1}W\) is again affine.
\end{theorem}
\begin{proof}
  \leanok
  The affine opens of \(X_A\) form a basis of its topology, so there is an affine open \(W\) with
  \(x \in W \le V\). Its preimage \(\mathrm{cg}^{-1}W\) is affine because \(\mathrm{cg}\) is an
  affine morphism (Theorem~\ref{pa-thm:isAffineHom_cg}).
\end{proof}

\begin{theorem}[The curve product over a proper curve is separated]
  \label{pa-thm:isSeparated_left}
  \lean{AlgebraicGeometry.Over.isSeparated_left}
  \leanok
  \dcref{custom}

  Assume \(C\) is proper over \(k\). Then \(X_A = (C \otimes \Spec A)_{\mathrm{left}}\) is a
  separated scheme.
\end{theorem}
\begin{proof}
  \leanok
  Properness of \(C\) makes its structure morphism separated, and separatedness is stable under
  base change, so the projection \(\mathrm{pr}_A \colon X_A \to \Spec A\) is separated; composing
  with the affine (hence separated) morphisms \(\Spec A \to \Spec k \to \ast\) shows the structure
  morphism of \(X_A\) is separated, i.e.\ \(X_A\) is a separated scheme.
\end{proof}

\begin{theorem}[The overlaps of the affine family are affine]
  \label{pa-thm:isAffineOpen_inf}
  \lean{AlgebraicGeometry.Over.isAffineOpen_inf}
  \leanok
  \dcref{custom}
  \uses{pa-thm:isSeparated_left}

  Assume \(C\) is proper over \(k\). Then the affine opens of \(X_A\) are closed under
  intersection: if \(V, W \subseteq X_A\) are affine open, so is \(V \cap W\).
\end{theorem}
\begin{proof}
  \leanok
  A separated scheme (Theorem~\ref{pa-thm:isSeparated_left}) has affine diagonal, and the
  intersection of two affine opens of a scheme with affine diagonal is affine.
\end{proof}

\section{The descent unit of a killed \v Cech class}
\label{pa-sec:kernelDescentUnit}

Let \(\gamma = \mathrm{lam}\) be a unit \v Cech cocycle on a pointed cover \(\mathcal E\) of
\(X_A\) whose base-change class \(\mathrm{cg}^{*}[\mathcal E, \mathrm{lam}]\) is trivial. Fixing
a trivializing \(0\)-cochain of the pulled-back cocycle, this section produces a single global
unit of \(\Spec(B \otimes_A B)\) that records, on the curve product, the coprojection ratio of
that \(0\)-cochain and whose ring avatar is a descent \(1\)-cocycle — the algebraic seed of the
descended line bundle.

\begin{theorem}[The descent unit of a killed class]
  \label{pa-thm:exists_kernelDescentUnit}
  \lean{AlgebraicGeometry.Over.exists_kernelDescentUnit}
  \leanok
  \dcref{custom}

  \uses{pa-thm:exists_global_unit_of_compatible, pa-thm:appTop_units_bijective, pa-def:tensorFaces,
    pa-lem:overSpecMap_left_degeneracy, pa-def:tensorMul, pa-def:module_descent_cocycle,
    pa-lem:tensorMul_simplicial}
  Let \(\mathrm{lam}\) be a unit \v Cech cocycle on a pointed cover \(\mathcal E\) of \(X_A\),
  and let \(\beta\) be a unit \(0\)-cochain on a pointed cover \(\mathcal A\) of \(X_B\) refining
  the base-change pullback of \(\mathcal E\), trivializing the pulled-back cocycle: on overlaps,
  \(\beta_v \cdot (\mathrm{cg}^{\sharp}\mathrm{lam})_{vv'} = \beta_{v'}\). Then there is a global
  unit \(w \in \Gamma(\Spec(B \otimes_A B), \top)^{\times}\) such that
  \begin{enumerate}
    \item on the common refinement of the two coprojection pullbacks of \(\mathcal A\), the
      curve-product pullback of \(w\) is the ratio \(u_1^{*}\beta / u_2^{*}\beta\) of the two
      coprojection pullbacks of \(\beta\); and
    \item the \(\Gamma\Spec\)-avatar of \(w\) in \((B \otimes_A B)^{\times}\) is a descent
      \(1\)-cocycle relative to \(A \to B\).
  \end{enumerate}
\end{theorem}
\begin{proof}
  \leanok
  The two coprojection pullbacks \(u_1^{*}\beta, u_2^{*}\beta\) are units over the members of
  the common refinement \(u_1^{*}\mathcal A \cap u_2^{*}\mathcal A\), and their ratio is
  compatible on overlaps: the coboundary hypothesis on \(\beta\), pulled back through the
  coprojection compatibility \(u_1 \circ \mathrm{cg} = u_2 \circ \mathrm{cg}\) of the whiskered
  faces, cancels the two \(\mathrm{lam}\)-factors and leaves the ratios agreeing on double
  overlaps.  By gluing of unit sections (Theorem~\ref{pa-thm:exists_global_unit_of_compatible})
  the ratio family glues to a global section, which by surjectivity of the global-section
  pullback (Theorem~\ref{pa-thm:appTop_units_bijective}) is the curve-product pullback of a global
  unit \(w\) on \(\Spec(B \otimes_A B)\); this is (1).

  For (2), work upstream on \(\Spec(B \otimes_A (B \otimes_A B))\).  Applying the three Amitsur
  cofaces (Definition~\ref{pa-def:tensorFaces}) to the local ratio form of \(w\) and using the
  simplicial coincidences among the coprojections, the \(\beta\)-telescope cancels exactly to
  give the cocycle identity \((f_{23}^{\sharp} w)(f_{12}^{\sharp} w) = f_{13}^{\sharp} w\) of the
  avatar.  The diagonal normalization \(\mu(w) = 1\) is the whiskered degeneracy of the tensor
  diagonal (Definition~\ref{pa-def:tensorMul}, Lemma~\ref{pa-lem:tensorMul_simplicial},
  Lemma~\ref{pa-lem:overSpecMap_left_degeneracy}), under which the two coprojection pullbacks of
  \(\beta\) coincide and their ratio is \(1\).  These are exactly the descent \(1\)-cocycle
  conditions (Definition~\ref{pa-def:module_descent_cocycle}) for the avatar of \(w\).
\end{proof}

\section{The \v Cech representative of a descended Picard class}
\label{pa-sec:descentClassRep}

The descent unit \(w\) of the previous section is the descent datum of an invertible module
\(\mathrm{descended}(w)\) over \(A\).  This section presents the \v Cech class of that module on
\(\Spec A\), together with the trivializing units of its base change to \(\Spec B\), whose
coprojection ratio is exactly \(w\).  Throughout, \(\mathrm{gS} := \Spec(A \to B)_{\mathrm{left}}
\colon \Spec B \to \Spec A\) and \(q_1, q_2 \colon \Spec(B \otimes_A B) \to \Spec B\) are the two
coprojections.

\begin{definition}[The evaluation transform of a section–module tensor]
  \label{pa-def:evalT}
  \lean{AlgebraicGeometry.Over.evalT}
  \leanok
  \dcref{custom}

  Let \(N\) be an \(A\)-submodule of \(B\) and \(V \subseteq \Spec A\), \(U \subseteq \Spec B\)
  opens with \(U \subseteq \mathrm{gS}^{-1}V\).  The \emph{evaluation transform} is the
  \(A\)-linear map
  \[
    \mathrm{evalT} \colon \Gamma(\Spec A, V) \otimes_A N \longrightarrow \Gamma(\Spec B, U),
    \qquad s \otimes m \longmapsto (\mathrm{gS}^{*} s) \cdot \bigl(\Gamma\Spec^{-1} m\bigr)
    \big|_{U},
  \]
  the product of the \(\mathrm{gS}\)-pullback of the section \(s\) with the restriction to
  \(U\) of the section-avatar of \(m \in B\).  It is compatible with restriction in both opens.
\end{definition}

\begin{lemma}[The coprojection ratio of an evaluation]
  \label{pa-lem:evalT_ratio}
  \lean{AlgebraicGeometry.Over.evalT_ratio}
  \leanok
  \dcref{custom}

  \uses{pa-def:evalT, pa-def:tensorInl_tensorInr}
  Let \(w \in \Gamma(\Spec(B \otimes_A B), \top)^{\times}\) and suppose every \(m \in N\)
  satisfies the descent equation of \(w\): \(1 \otimes m = \Gamma\Spec(w) \cdot (m \otimes 1)\)
  in \(B \otimes_A B\).  Then, over any open \(O\) below both coprojection preimages of \(U\),
  the two coprojection pullbacks of an evaluation differ by the restriction of \(w\):
  \[
    q_2^{*}\bigl(\mathrm{evalT}(x)\bigr) = w|_{O} \cdot q_1^{*}\bigl(\mathrm{evalT}(x)\bigr),
    \qquad x \in \Gamma(\Spec A, V) \otimes_A N.
  \]
\end{lemma}
\begin{proof}
  \leanok
  Both sides are additive in \(x\), so it suffices to treat a pure tensor \(s \otimes m\).  The
  two coprojections send the section factor \(\mathrm{gS}^{*} s\) to the same value, the
  pullback along the common composite \(q_1 \circ \mathrm{gS} = q_2 \circ \mathrm{gS}\); and
  they send the module factor to the two coprojection-avatars \(\Gamma\Spec^{-1}(m \otimes 1)\)
  and \(\Gamma\Spec^{-1}(1 \otimes m)\) (Definition~\ref{pa-def:tensorInl_tensorInr}).  The descent
  equation of \(w\) applied to \(m\) turns the second avatar into \(w\) times the first, and the
  central unit \(w\) factors out, giving the stated ratio.
\end{proof}

\begin{definition}[The \v Cech representative of a descended class]
  \label{pa-def:DescentClassRep}
  \lean{AlgebraicGeometry.Over.DescentClassRep}
  \leanok
  \dcref{custom}

  \uses{pa-def:module_descent_cocycle, pa-def:tensorInl_tensorInr}
  Let \(w\) be a global unit of \(\Spec(B \otimes_A B)\) whose avatar is a descent
  \(1\)-cocycle.  A \emph{\v Cech representative} of the descended class consists of a pointed
  cover \(\mathcal C\) of \(\Spec A\), a unit \v Cech cocycle \(d\) on \(\mathcal C\), and a
  family of trivializing units \(\mu_a \in \Gamma(\Spec B, \mathrm{gS}^{-1}\mathcal
  C(a))^{\times}\) of the \(\mathrm{gS}\)-pullback, subject to:
  \begin{enumerate}
    \item (\emph{ratio}) the two coprojection pullbacks of each \(\mu_a\) differ by \(w\):
      \(q_2^{*}\mu_a = w \cdot q_1^{*}\mu_a\) on any sub-open of the double preimage of
      \(\mathcal C(a)\);
    \item (\emph{transition}) on overlaps the \(\mu\)'s differ by the \(\mathrm{gS}\)-pullback
      of the cocycle value: \(\mu_a = \mathrm{gS}^{*} d_{ab} \cdot \mu_b\) on \(\mathrm{gS}^{-1}
      (\mathcal C(a) \cap \mathcal C(b))\).
  \end{enumerate}
\end{definition}

\begin{theorem}[Existence of the \v Cech representative]
  \label{pa-thm:descentClassRep}
  \lean{AlgebraicGeometry.Over.descentClassRep}
  \leanok
  \dcref{custom}

  \uses{pa-def:DescentClassRep, pa-lem:trivializingFamily_nonempty, pa-def:module_descended,
    pa-lem:evalT_ratio, pa-thm:isUnit_map_rTensor_generator, pa-def:descentMulEval}
  For every global unit \(w\) of \(\Spec(B \otimes_A B)\) whose avatar is a descent
  \(1\)-cocycle, a \v Cech representative (Definition~\ref{pa-def:DescentClassRep}) of the descended
  class exists.
\end{theorem}
\begin{proof}
  \leanok
  Let \(N := \mathrm{descended}(w)\) be the invertible \(A\)-module descended by the avatar of
  \(w\) (Definition~\ref{pa-def:module_descended}), and pass to its section-ring avatar
  \(\Gamma(\Spec A, \top) \otimes_A N\).  Choose a trivializing family of that invertible module
  (Lemma~\ref{pa-lem:trivializingFamily_nonempty}): a pointed cover \(\mathcal C\) of \(\Spec A\)
  by basic opens \(D(g_a)\) with local trivializations, giving the cover \(\mathcal C\) and its
  cover cocycle \(d\).  On the base change to \(\Spec B\), define \(\mu_a\) as the evaluation of
  the collapsed generator of the \(a\)-trivialization through the surjective evaluation of the
  base-changed descent equivalence (Definition~\ref{pa-def:descentMulEval}); this value is a unit
  by the unit-production principle (Theorem~\ref{pa-thm:isUnit_map_rTensor_generator}).  The ratio
  condition is the coprojection ratio of an evaluation (Lemma~\ref{pa-lem:evalT_ratio}), whose
  hypothesis is the descent equation of \(w\) satisfied by the elements of \(N\).  The
  transition condition follows because the trivializations are compatible across the cover and
  the transition units act on the restricted generators by the cover cocycle \(d\).
\end{proof}

\section{The cobounding units on an affine chart}
\label{pa-sec:kernelCobounding}

The trivializing \(0\)-cochain \(\beta\), the cocycle \(\mathrm{lam}\), and the trivializing
units \(\mu\) of the representative are assembled into a single unit on the base-change pullback
of each chart of \(X_A\), which then descends through \(\mathrm{cg}\).  These descended units are
the cobounding \(0\)-cochain that exhibits \([\mathcal E, \mathrm{lam}]\) as pulled back from
\(\Spec A\).

\begin{theorem}[The cobounding unit at an affine chart]
  \label{pa-thm:exists_kernelCobounding}
  \lean{AlgebraicGeometry.Over.exists_kernelCobounding}
  \leanok
  \dcref{custom}

  \uses{pa-def:DescentClassRep, pa-thm:exists_kernelDescentUnit, pa-thm:exists_unitsAppLE_eq,
    pa-lem:units_gluing}
  Keep the data \((\mathcal E, \mathrm{lam}, \beta)\) of the descent unit \(w\)
  (Theorem~\ref{pa-thm:exists_kernelDescentUnit}) and a \v Cech representative \((\mathcal C, d,
  \mu)\) of the descended class (Definition~\ref{pa-def:DescentClassRep}).  For every affine open
  \(U \subseteq X_A\) contained in \(\mathcal E(s)\) and in \(\mathrm{pr}_A^{-1}\mathcal
  C(\mathrm{pr}_A(s))\), there is a unit \(c \in \Gamma(X_A, U)^{\times}\) whose
  \(\mathrm{cg}\)-pullback has, on every small enough sub-open, the glued local form
  \[
    (\mathrm{cg}^{*} c) = \beta_x^{-1} \cdot (\mathrm{cg}^{\sharp}\mathrm{lam}_{\mathrm{cg}(x),s})^{-1}
      \cdot (\mathrm{pr}_B^{\sharp}\mu_{\mathrm{pr}_A(s)})^{-1}.
  \]
\end{theorem}
\begin{proof}
  \leanok
  On each member of the base-change pullback of \(\mathcal E\) meeting \(\mathrm{cg}^{-1}U\),
  form the displayed product of the inverse trivialization \(\beta_x\), the inverse
  \(\mathrm{cg}\)-pullback of the cross-evaluation of \(\mathrm{lam}\), and the inverse
  \(\mathrm{pr}_B\)-pullback of \(\mu\).  These local pieces agree on overlaps: the coboundary
  relation of \(\beta\) against \(\mathrm{lam}\) and the cocycle identity of \(\mathrm{lam}\)
  telescope the \(\beta\)- and \(\mathrm{lam}\)-factors, while the \(\mu\)-factor is independent
  of \(x\).  By gluing of unit sections (Lemma~\ref{pa-lem:units_gluing}) they glue to a unit \(G\)
  on \(\mathrm{cg}^{-1}U\).  Its two coprojection pullbacks to \(C \otimes \Spec(B \otimes_A B)\)
  agree: the local ratio form of the descent unit \(w\)
  (Theorem~\ref{pa-thm:exists_kernelDescentUnit}) and the ratio field of the representative
  (Definition~\ref{pa-def:DescentClassRep}) exactly account for the two coprojection ratios of the
  three factors.  Since \(U\) is affine, \(G\) therefore descends to a unit \(c \in \Gamma(X_A,
  U)^{\times}\) with \(\mathrm{cg}^{*} c = G\) (Theorem~\ref{pa-thm:exists_unitsAppLE_eq}), which is
  the claim.
\end{proof}

\section{The \v Cech kernel lemma and the close of \'etale separatedness}
\label{pa-sec:cechKernelLemma}

The pieces now compose.  The kernel lemma isolates the genuinely fppf content — a \v Cech class
killed by the cover pullback is pulled back from the base — and the class-level reduction and the
plus unfold deliver, from it, the injectivity of the unit of the \'etale plus construction:
Kleiman's separatedness of the relative Picard functor.

\begin{theorem}[The \v Cech kernel lemma]
  \label{pa-thm:exists_cechPic_map_snd_of_ker_whiskerLeft}
  \lean{AlgebraicGeometry.Over.exists_cechPic_map_snd_of_ker_whiskerLeft}
  \leanok
  \source{kleiman-picard}

  \uses{pa-thm:exists_kernelDescentUnit, pa-thm:descentClassRep, pa-thm:exists_kernelCobounding,
    pa-thm:appLE_whiskerLeft_injective, pa-lem:pointedCover_affine_refinement, pa-thm:unitsRes_injective,
    pa-cor:cechPic_mk_eq_one}
  \dcref{custom}
  Let \(C\) be proper, geometrically irreducible and geometrically reduced over \(k\), and \(A
  \to B\) faithfully flat.  Then the kernel of the base-change pullback \(\mathrm{cg}^{*} \colon
  \Pic(X_A) \to \Pic(X_B)\) is contained in the image of the projection pullback
  \(\mathrm{pr}_A^{*} \colon \Pic(\Spec A) \to \Pic(X_A)\): every class \(E\) on \(X_A\) with
  \(\mathrm{cg}^{*} E = 1\) is \(\mathrm{pr}_A^{*} M_1\) for some class \(M_1\) on \(\Spec A\).
\end{theorem}
\begin{proof}
  \leanok
  Represent \(E\) by a unit cocycle \(\mathrm{lam}\) on a pointed cover \(\mathcal E\) of
  \(X_A\).  Since \(\mathrm{cg}^{*} E = 1\), the pulled-back cocycle is a coboundary: there is a
  trivializing \(0\)-cochain \(\beta\) on a refinement of \(\mathrm{cg}^{*}\mathcal E\) with
  \(\beta_v \cdot (\mathrm{cg}^{\sharp}\mathrm{lam})_{vv'} = \beta_{v'}\).  Feed \((\mathcal E,
  \mathrm{lam}, \beta)\) to the descent unit (Theorem~\ref{pa-thm:exists_kernelDescentUnit}),
  obtaining a global unit \(w\) whose avatar is a descent \(1\)-cocycle; take a \v Cech
  representative \((\mathcal C, d, \mu)\) of its descended class
  (Theorem~\ref{pa-thm:descentClassRep}).  Refine \(\mathcal E \cap \mathrm{pr}_A^{-1}\mathcal C\)
  to a cover \(\mathcal U\) by affine opens (Lemma~\ref{pa-lem:pointedCover_affine_refinement}); on
  each chart \(U = \mathcal U(s)\), the cobounding unit
  (Theorem~\ref{pa-thm:exists_kernelCobounding}) gives \(c_s \in \Gamma(X_A, U)^{\times}\).

  It remains to check that the \(c_s\) cobound \(\mathrm{lam}\), restricted to \(\mathcal U\),
  against the \(\mathrm{pr}_A\)-pullback of the representing cocycle \(d\); that is,
  \[
    c_s|_{U_s \cap U_t} \cdot \mathrm{pr}_A^{*} d_{\,\mathrm{pr}_A s,\,\mathrm{pr}_A t}
      = \mathrm{lam}_{s t}\big|_{U_s \cap U_t} \cdot c_t|_{U_s \cap U_t}.
  \]
  Both sides pull back injectively along \(\mathrm{cg}\)
  (Theorem~\ref{pa-thm:appLE_whiskerLeft_injective}); after pullback, the glued local forms of
  \(c_s, c_t\), the \(\mu\)-transition of the representative, and the cocycle identity of
  \(\mathrm{lam}\) telescope to the required identity.  Hence \(\mathrm{lam}\) and
  \(\mathrm{pr}_A^{*} d\) are cohomologous on \(\mathcal U\), so \([\mathcal U, \mathrm{lam}] =
  \mathrm{pr}_A^{*}[\mathcal C, d]\) by refinement injectivity in degree one
  (Theorem~\ref{pa-thm:unitsRes_injective}, Corollary~\ref{pa-cor:cechPic_mk_eq_one}); that is, \(E =
  \mathrm{pr}_A^{*} M_1\) with \(M_1 = [\mathcal C, d]\).  This is the descent-datum step of
  Kleiman's separatedness argument in its \v Cech incarnation: the trivialization \(\beta\)
  descends the class from the faithfully flat cover to the base.
\end{proof}

\begin{theorem}[The class-level reduction]
  \label{pa-thm:exists_cechPic_map_snd_eq_of_ker}
  \lean{AlgebraicGeometry.Over.exists_cechPic_map_snd_eq_of_ker}
  \leanok
  \source{kleiman-picard}

  \uses{pa-thm:exists_coherentCechWitness, pa-thm:isDescentCocycle_assemblyUnit,
    pa-thm:mapAlgebra_picClass_assemblyUnit, pa-thm:cechPicEquivPic, pa-thm:cechPic_toPic_injective,
    pa-thm:cechPic_toPic_map, pa-lem:snd_left_naturality, pa-def:module_picClass}
  \dcref{custom}
  Assume the kernel property of \(\mathrm{cg}^{*}\) (Theorem~\ref{pa-thm:exists_cechPic_map_snd_of_ker_whiskerLeft}).
  Let \(L\) be a \v Cech Picard class on \(X_A\) and \(N_0\) a class on \(\Spec B\) with
  \(\mathrm{pr}_B^{*} N_0 = \mathrm{cg}^{*} L\).  Then \(L\) is pulled back from the base:
  \(L = \mathrm{pr}_A^{*} M\) for some class \(M\) on \(\Spec A\).
\end{theorem}
\begin{proof}
  \leanok
  From the hypothesis \(\mathrm{pr}_B^{*} N_0 = \mathrm{cg}^{*} L\), the coherent \v Cech witness
  (Theorem~\ref{pa-thm:exists_coherentCechWitness}) provides an Amitsur-coherent trivialization of
  the two coprojection pullbacks of \(N_0\), and the assembly of the previous chapter builds from
  it a composite \(A\)-descent unit whose avatar is a descent \(1\)-cocycle
  (Theorem~\ref{pa-thm:isDescentCocycle_assemblyUnit}).  Let \(M\) be the \v Cech class on \(\Spec
  A\) corresponding, through the \v Cech–Picard dictionary
  (Theorem~\ref{pa-thm:cechPicEquivPic}), to the Picard class of the descended module
  (Definition~\ref{pa-def:module_picClass}) of that composite unit.  By the class equation of the
  assembly (Theorem~\ref{pa-thm:mapAlgebra_picClass_assemblyUnit}) and naturality of the dictionary
  in the affine scheme (Theorem~\ref{pa-thm:cechPic_toPic_map}), the base change \(\mathrm{gS}^{*}
  M\) has the same descent homomorphism image as \(N_0\); since that homomorphism is injective
  (Theorem~\ref{pa-thm:cechPic_toPic_injective}), \(\mathrm{gS}^{*} M = N_0\) on the nose.  Then,
  using the projection base square (Lemma~\ref{pa-lem:snd_left_naturality}),
  \[
    \mathrm{cg}^{*}\bigl(\mathrm{pr}_A^{*} M\bigr) = \mathrm{pr}_B^{*}\bigl(\mathrm{gS}^{*} M\bigr)
      = \mathrm{pr}_B^{*} N_0 = \mathrm{cg}^{*} L ,
  \]
  so \(L \cdot (\mathrm{pr}_A^{*} M)^{-1}\) is killed by \(\mathrm{cg}^{*}\).  The kernel
  property yields \(M_1\) with \(\mathrm{pr}_A^{*} M_1 = L \cdot (\mathrm{pr}_A^{*} M)^{-1}\),
  whence \(L = \mathrm{pr}_A^{*}(M \cdot M_1)\).
\end{proof}

\begin{theorem}[The plus unfold of \'etale separatedness]
  \label{pa-thm:PicEtAff_unit_injective_of_ker}
  \lean{AlgebraicGeometry.PicEtAff.unit_injective_of_ker}
  \leanok
  \dcref{custom}

  \uses{pa-def:PicEtAff_unit, pa-thm:PicEtAff_mk_eq_mk_iff, pa-def:PicEtAff_one,
    pa-thm:exists_cechPic_map_snd_eq_of_ker, pa-def:relPicAlgMap, pa-def:descentClasses}
  Fix \(A\).  Suppose the kernel property of the base-change pullback holds for every tower \(A
  \to B\) with \(A \to B\) faithfully flat.  Then the unit \(\eta_A \colon \mathrm{relPic}(C,A)
  \to \mathrm{PicEtAff}(C,A)\) of the plus construction (Definition~\ref{pa-def:PicEtAff_unit}) is
  injective.
\end{theorem}
\begin{proof}
  \leanok
  As \(\eta_A\) is a group homomorphism it suffices to show its kernel is trivial.  Let \(x =
  [\Spec A, L] \in \mathrm{relPic}(C,A)\) with \(\eta_A(x) = 1\).  Unwinding the plus-class
  equality \([\mathrm{self}(A), \iota^{*}x] = 1\) through the presentation criterion
  (Theorem~\ref{pa-thm:PicEtAff_mk_eq_mk_iff}) and the essential uniqueness of \(A\)-algebra maps
  out of \(A\), the datum reduces to a single presented \'etale cover \(A \to B_G\): the base
  change \((\mathrm{ofId}\ A\ B_G)^{*}L\) lies in the descent-trivial subgroup of
  \(\mathrm{relPic}(C,B_G)\) (Definition~\ref{pa-def:relPicAlgMap},
  Definition~\ref{pa-def:descentClasses}).  Concretely there is a class \(N_0\) on \(\Spec B_G\)
  with \(\mathrm{pr}_{B_G}^{*} N_0 = \mathrm{cg}^{*} L\) — the hypothesis of the reduction.  Since
  a presented \'etale cover is in particular faithfully flat, the reduction
  (Theorem~\ref{pa-thm:exists_cechPic_map_snd_eq_of_ker}) pulls \(L\) back from the base, \(L =
  \mathrm{pr}_A^{*} M\); that is exactly \(x = 1\) in \(\mathrm{relPic}(C,A)\), using the
  normalization of the unit at the identity (Definition~\ref{pa-def:PicEtAff_one}).
\end{proof}

\begin{theorem}[\'Etale separatedness of the relative Picard functor]
  \label{pa-thm:PicEtAff_unit_injective}
  \lean{AlgebraicGeometry.PicEtAff.unit_injective}
  \leanok
  \source{kleiman-picard}

  \uses{pa-thm:PicEtAff_unit_injective_of_ker, pa-thm:exists_cechPic_map_snd_of_ker_whiskerLeft}
  \dcref{2.5,custom}
  Let \(C\) be a curve over \(k\) that is proper, geometrically irreducible and geometrically
  reduced, and let \(A\) be any \(k\)-algebra.  Then the unit \(\eta_A \colon
  \mathrm{relPic}(C,A) \to \mathrm{PicEtAff}(C,A)\) of the one-step \'etale plus construction is
  injective.  Equivalently, the relative Picard functor of \(C\) is separated for the \'etale
  topology — the separatedness half of Kleiman, \emph{The Picard Scheme}, Theorem 2.5(1).
\end{theorem}
\begin{proof}
  \leanok
  The kernel property of the base-change pullback holds for every faithfully flat tower \(A \to
  B\) by the \v Cech kernel lemma
  (Theorem~\ref{pa-thm:exists_cechPic_map_snd_of_ker_whiskerLeft}).  Feeding it to the plus unfold
  (Theorem~\ref{pa-thm:PicEtAff_unit_injective_of_ker}) gives the injectivity of \(\eta_A\).  This
  is exactly Part 1 of Kleiman's comparison theorem in the \'etale case: the natural map
  \(\Pic_{X/S} \hookrightarrow \Pic_{(X/S)\mathrm{\acute et}}\) is injective, so a single
  \'etale plus computes the sheafification and \(\mathrm{PicEtAff}\) is separated on the \'etale
  site of affines.
\end{proof}


\section{Unit sections across an open immersion}

Throughout this section \(w \colon Z \to Y\) is an open immersion of schemes. Over an open
\(V \subseteq Y\) contained in the image of \(w\), pullback of sections along \(w\) is
bijective; this transports \emph{unit} sections between \(Y\) and \(Z\) and is the substrate
for decomposing \v Cech Picard classes along clopen partitions.

\begin{lemma}[Units over an empty open form a trivial group]
  \label{pa-lem:subsingleton_units_empty}
  \lean{AlgebraicGeometry.Scheme.subsingleton_units_of_le_bot}
  \leanok
  \dcref{custom}

  For a scheme \(X\) and an open \(O \subseteq X\) with \(O \subseteq \emptyset\) (that is,
  \(O = \emptyset\)), the unit group \(\Gamma(X, O)^{\times}\) is trivial (a subsingleton).
\end{lemma}
\begin{proof}
  \leanok
  Since \(O = \emptyset\), the ring \(\Gamma(X, O)\) is the zero ring, in which \(0 = 1\) is the
  only element; hence its group of units has a single element.
\end{proof}

\begin{lemma}[Sections pull back isomorphically over an open inside the image]
  \label{pa-lem:isIso_appLE_of_le_opensRange}
  \lean{AlgebraicGeometry.Scheme.Hom.isIso_appLE_of_le_opensRange}
  \leanok
  \dcref{custom}

  Let \(V \subseteq Y\) be an open with \(V \subseteq \operatorname{range}(w)\). Then pullback of
  sections along \(w\) is an isomorphism of rings
  \(\Gamma(Y, V) \xrightarrow{\ \sim\ } \Gamma(Z, w^{-1}V)\).
\end{lemma}
\begin{proof}
  \leanok
  An open immersion \(w\) identifies \(Z\) with the open subscheme \(\operatorname{range}(w)\) of
  \(Y\); for an open \(V\) contained in that image the comparison map on sections is the
  restriction isomorphism of the structure sheaf between \(V\) and its homeomorphic preimage
  \(w^{-1}V\). Hence pullback of sections over \(V\) is an isomorphism of rings.
\end{proof}

\begin{definition}[Transport of unit sections across an open immersion]
  \label{pa-def:unitsPreimageEquiv}
  \lean{AlgebraicGeometry.Scheme.Hom.unitsPreimageEquiv}
  \leanok
  \dcref{custom}

  \uses{pa-lem:isIso_appLE_of_le_opensRange}
  For an open \(V \subseteq Y\) with \(V \subseteq \operatorname{range}(w)\), the isomorphism of
  unit groups
  \[
    \Gamma(Y, V)^{\times} \xrightarrow{\ \sim\ } \Gamma(Z, w^{-1}V)^{\times}
  \]
  induced by pullback of sections along \(w\).
\end{definition}
\begin{proof}
  \leanok
  Pullback of sections along \(w\) is an isomorphism of rings over \(V\)
  (Lemma~\ref{pa-lem:isIso_appLE_of_le_opensRange}); an isomorphism of rings restricts to an
  isomorphism of their groups of units.
\end{proof}

\section{Decomposition of the \v Cech Picard group along a clopen partition}

A \emph{clopen partition} of a scheme \(Y\) is a family of open immersions
\(w_i \colon Z_i \to Y\), \(i \in \iota\), whose ranges are pairwise disjoint
(\(\operatorname{range}(w_i) \cap \operatorname{range}(w_j) = \emptyset\) for \(i \ne j\)) and
cover \(Y\); each range is then open with open complement. This section proves that the \v Cech
Picard group of \(Y\) decomposes along such a partition: separation along an arbitrary clopen
partition, and gluing along a finite one. Throughout, \(\Pic(Y)\) is the definitional \v Cech
Picard group recalled in Section~\ref{pa-sec:etaleSeparatedness}, with pullback
\(w^{*} \colon \Pic(Y) \to \Pic(Z)\) along a morphism \(w \colon Z \to Y\).

\begin{theorem}[Separation of \v Cech Picard classes along a clopen partition]
  \label{pa-thm:cechPic_clopen_separation}
  \lean{AlgebraicGeometry.Scheme.CechPic.eq_of_map_eq_of_clopen}
  \leanok
  \dcref{custom}

  \uses{pa-def:unitsPreimageEquiv, pa-lem:subsingleton_units_empty, pa-cor:cechPic_mk_injective}
  Let \((w_i)_{i \in \iota}\) be a clopen partition of \(Y\). Two \v Cech Picard classes
  \(L, L' \in \Pic(Y)\) with \(w_i^{*} L = w_i^{*} L'\) for every \(i\) are equal.
\end{theorem}
\begin{proof}
  \leanok
  \uses{pa-lem:units_separation, pa-lem:units_gluing}
  Choose for each point \(y \in Y\) an index \(p(y)\) with \(y \in \operatorname{range}(w_{p(y)})\).
  Represent \(L = [\mathcal U_1, a_1]\) and \(L' = [\mathcal U_2, a_2]\) and pass to the common
  adapted pointed cover \(\mathcal V(y) := \mathcal U_1(y) \cap \mathcal U_2(y) \cap
  \operatorname{range}(w_{p(y)})\), so that \(\mathcal V(y) \subseteq \operatorname{range}(w_{p(y)})\)
  for every \(y\); it suffices to show the ratio cocycle \(\gamma\) representing \(L/L'\) on
  \(\mathcal V\) is a coboundary.

  \emph{Triviality on each piece.} Fix \(i\). Since \(w_i^{*}L = w_i^{*}L'\), the pullback classes
  \(w_i^{*}[\mathcal V, a_1|_{\mathcal V}]\) and \(w_i^{*}[\mathcal V, a_2|_{\mathcal V}]\) are
  represented on the pulled-back pointed cover \(\mathcal V^{(i)} := w_i^{-1}\mathcal V\) by the
  pulled-back cocycles, and they agree in \(\Pic(Z_i)\); by refinement injectivity
  (Corollary~\ref{pa-cor:cechPic_mk_injective}) already the pulled-back ratio cocycle
  \(w_i^{\sharp}\gamma\) is a coboundary on \(\mathcal V^{(i)}\): there are units
  \(\alpha^{(i)}_z \in \struct{Z_i}(\mathcal V^{(i)}(z))^{\times}\) with
  \(\alpha^{(i)}_z \cdot (w_i^{\sharp}\gamma)_{zz'} = \alpha^{(i)}_{z'}\) on the overlaps.

  \emph{Gluing the piece coboundaries.} For each \(y \in Y\), let \(z_0\) be the unique point of
  \(Z_{p(y)}\) over \(y\) (an open immersion is injective on points), and set
  \(\beta_y := \Phi^{-1}\bigl(\alpha^{(p(y))}_{z_0}|_{w_{p(y)}^{-1}\mathcal V(y)}\bigr) \in
  \Gamma(Y, \mathcal V(y))^{\times}\), where \(\Phi\) is the unit transport of
  Definition~\ref{pa-def:unitsPreimageEquiv} across \(w_{p(y)}\) (applicable since
  \(\mathcal V(y) \subseteq \operatorname{range}(w_{p(y)})\)). The choice of index is forced: if
  \(y \in \operatorname{range}(w_i)\) too, then \(y \in \operatorname{range}(w_i) \cap
  \operatorname{range}(w_{p(y)})\), which is empty unless \(i = p(y)\); and the choice of preimage
  point is forced by injectivity.

  \emph{\(\beta\) cobounds \(\gamma\).} Fix \(y, y'\). If \(p(y') = p(y)\), both members lie in the
  single piece \(Z_{p(y)}\); the coboundary relation for \(\alpha^{(p(y))}\) pulls back injectively
  through \(w_{p(y)}\) on the overlap \(\mathcal V(y) \cap \mathcal V(y') \subseteq
  \operatorname{range}(w_{p(y)})\) (pullback of unit sections is injective on opens inside the
  range, Definition~\ref{pa-def:unitsPreimageEquiv}), yielding
  \(\beta_y \cdot \gamma_{yy'} = \beta_{y'}\) there. If \(p(y') \ne p(y)\), then
  \(\mathcal V(y) \cap \mathcal V(y') \subseteq \operatorname{range}(w_{p(y)}) \cap
  \operatorname{range}(w_{p(y')}) = \emptyset\), so the units over the overlap form a trivial group
  (Lemma~\ref{pa-lem:subsingleton_units_empty}) and the relation holds vacuously. Assembling over the
  cover of \(\mathcal V(y) \cap \mathcal V(y')\) by these two cases and separating
  (Lemma~\ref{pa-lem:units_separation}; gluing across pieces uses Lemma~\ref{pa-lem:units_gluing} in the
  construction of \(\beta\)), \(\gamma = \delta\beta\) is a coboundary, so \(L/L' = 1\), i.e.
  \(L = L'\).
\end{proof}

\begin{theorem}[The one-piece extension across a clopen decomposition]
  \label{pa-thm:cechPic_extendClass}
  \lean{AlgebraicGeometry.Scheme.CechPic.extendClass,
    AlgebraicGeometry.Scheme.CechPic.map_extendClass,
    AlgebraicGeometry.Scheme.CechPic.map_extendClass_of_le_compl}
  \leanok
  \dcref{custom}

  \uses{pa-def:unitsPreimageEquiv, pa-lem:subsingleton_units_empty}
  Let \(w \colon Z \to Y\) be an open immersion whose range is complemented by an open
  \(\Omega' \subseteq Y\): \(\operatorname{range}(w) \cap \Omega' = \emptyset\) and
  \(\operatorname{range}(w) \cup \Omega' = Y\). For every \(L_0 \in \Pic(Z)\) there is a class
  \(\widetilde L \in \Pic(Y)\) with
  \[
    w^{*}\widetilde L = L_0, \qquad v^{*}\widetilde L = 1 \text{ for every } v \colon W \to Y
    \text{ with image contained in } \Omega'.
  \]
\end{theorem}
\begin{proof}
  \leanok
  Represent \(L_0 = [\mathcal U_0, \gamma_0]\) by a unit cocycle on a pointed cover of \(Z\). Build
  the \emph{extension cover} \(\mathcal E\) of \(Y\): at a point \(y \in \operatorname{range}(w)\)
  take the \(w\)-image of the member of \(\mathcal U_0\) at the unique preimage of \(y\); at
  \(y \notin \operatorname{range}(w)\) take \(\Omega'\) (a legitimate pointed cover, since
  \(\operatorname{range}(w) \cup \Omega' = Y\)). Define its pair values: on range--range overlaps
  the transported cocycle value \(\Phi^{-1}\bigl((\gamma_0)_{\cdot\cdot}\bigr)\) of
  Definition~\ref{pa-def:unitsPreimageEquiv} across \(w\), and \(1\) whenever either argument lies off
  the range. These satisfy the cocycle identity: on triple overlaps of range members it is the
  transported cocycle identity of \(\gamma_0\); every other triple overlap meets a range member and
  a non-range member, hence lies in \(\operatorname{range}(w) \cap \Omega' = \emptyset\) and carries
  only the trivial unit (Lemma~\ref{pa-lem:subsingleton_units_empty}), so the identity degenerates.
  This produces \(\widetilde L := [\mathcal E, \text{extension cocycle}] \in \Pic(Y)\).

  \emph{Restriction on the piece.} The given cover \(\mathcal U_0\) refines the pullback
  \(w^{-1}\mathcal E\) (the preimage of the image of a member is that member), and along this
  refinement the pulled-back extension cocycle restricts to \(\gamma_0\) (its range--range values
  are the transported \(\gamma_0\)). Hence \(w^{*}\widetilde L = [\mathcal U_0, \gamma_0] = L_0\).

  \emph{Restriction off the piece.} For \(v \colon W \to Y\) with image in \(\Omega'\), every point
  \(v(x)\) lies off \(\operatorname{range}(w)\) (else it would lie in
  \(\operatorname{range}(w) \cap \Omega' = \emptyset\)); so every pulled-back pair value is \(1\),
  the pulled-back cocycle is trivial, and \(v^{*}\widetilde L = 1\).
\end{proof}

\begin{theorem}[Gluing of \v Cech Picard classes along a finite clopen partition]
  \label{pa-thm:cechPic_clopen_gluing}
  \lean{AlgebraicGeometry.Scheme.CechPic.exists_map_eq_of_clopen}
  \leanok
  \dcref{custom}

  \uses{pa-thm:cechPic_extendClass}
  Let \((w_i)_{i \in \iota}\) be a \emph{finite} clopen partition of \(Y\). Every family
  \((L_i)_{i \in \iota}\) with \(L_i \in \Pic(Z_i)\) is the family of restrictions of a single
  class \(M \in \Pic(Y)\): \(w_i^{*} M = L_i\) for all \(i\).
\end{theorem}
\begin{proof}
  \leanok
  For each \(i\), set \(\Omega'_i := \bigcup_{j \ne i} \operatorname{range}(w_j)\). By pairwise
  disjointness \(\operatorname{range}(w_i) \cap \Omega'_i = \emptyset\), and since the ranges cover
  \(Y\), \(\operatorname{range}(w_i) \cup \Omega'_i = Y\). Let \(M_i \in \Pic(Y)\) be the one-piece
  extension of \(L_i\) across \((w_i, \Omega'_i)\) (Theorem~\ref{pa-thm:cechPic_extendClass}). Put
  \(M := \prod_{i} M_i\) (finite product). For each \(i\), pullback along \(w_i\) is a homomorphism,
  so \(w_i^{*} M = \prod_j w_i^{*} M_j\). The factor \(j = i\) is \(w_i^{*} M_i = L_i\); every factor
  \(j \ne i\) is \(1\), because \(\operatorname{range}(w_i) \subseteq \Omega'_j\) means \(w_i\) has
  image in the complement of the \(j\)-th piece, so \(w_i^{*}M_j = 1\)
  (Theorem~\ref{pa-thm:cechPic_extendClass}). Hence \(w_i^{*} M = L_i\).
\end{proof}

\section{The relative Picard group of a finite product of test algebras}

Fix a field \(k\) and a curve \(C\) over \(k\). For a finite family \((B_i)_{i \in \iota}\) of
commutative \(k\)-algebras with product \(P := \prod_i B_i\), the spectrum of \(P\) is the disjoint
union of the spectra of the factors, clopen-partitioned by the basic opens of the coordinate
idempotents; base-changing along the curve, the whiskered coordinate projections clopen-partition
the curve-product. Through the clopen decomposition of the previous section this identifies the
relative Picard group of the product test with the product of the relative Picard groups of the
factors. Write \(e_i \in P\) for the coordinate idempotent (\(1\) in slot \(i\), \(0\) elsewhere)
and \(\pi_i \colon P \to B_i\) for the projection.

\begin{lemma}[Range of a whiskered open immersion]
  \label{pa-lem:range_whiskerLeft}
  \lean{AlgebraicGeometry.Over.range_whiskerLeft}
  \leanok
  \dcref{custom}

  \uses{pa-thm:isPullback_whiskerLeft_snd}
  Let \(g \colon T' \to T\) be a morphism of test objects with \(g_{\mathrm{left}}\) an open
  immersion. Then the whiskered morphism \((C \lhd g)_{\mathrm{left}}\) is an open immersion, and
  its range is the preimage under the second projection \(\mathrm{pr} \colon (C \otimes
  T)_{\mathrm{left}} \to T_{\mathrm{left}}\) of the range of \(g_{\mathrm{left}}\):
  \[
    \operatorname{range}\bigl((C \lhd g)_{\mathrm{left}}\bigr) = \mathrm{pr}^{-1}\bigl(
    \operatorname{range}(g_{\mathrm{left}})\bigr).
  \]
\end{lemma}
\begin{proof}
  \leanok
  The whiskering square with horizontal legs \((C \lhd g)_{\mathrm{left}}\), \(g_{\mathrm{left}}\)
  and vertical legs the two second projections is a pullback of schemes
  (Theorem~\ref{pa-thm:isPullback_whiskerLeft_snd}). Open immersions are stable under base change, so
  \((C \lhd g)_{\mathrm{left}}\), being the base change of \(g_{\mathrm{left}}\) along
  \(\mathrm{pr}\), is an open immersion; and the range of the base change of an open immersion along
  \(\mathrm{pr}\) is the \(\mathrm{pr}\)-preimage of its range.
\end{proof}

\begin{lemma}[A factor is the localization away from the coordinate idempotent]
  \label{pa-lem:isLocalization_away_piFactor}
  \lean{AlgebraicGeometry.isLocalization_away_piFactor}
  \leanok
  \dcref{custom}

  For a finite product \(P = \prod_j B_j\) and an index \(i\), the projection \(\pi_i \colon P \to
  B_i\) exhibits \(B_i\) as the localization of \(P\) away from the coordinate idempotent \(e_i\).
\end{lemma}
\begin{proof}
  \leanok
  The element \(e_i\) is idempotent, \(e_i^2 = e_i\). Two elements \(x, y \in P\) have the same
  image in \(B_i\) iff they agree in slot \(i\), iff \(e_i x = e_i y\) (multiplication by \(e_i\)
  zeroes every other slot). And every \(b \in B_i\) is the image of \(\mathrm{single}_i(b) \in P\).
  By the criterion for a localization away from an idempotent, \(B_i = P[e_i^{-1}]\).
\end{proof}

\begin{lemma}[The coordinate projections clopen-partition the spectrum of a product]
  \label{pa-lem:opensRange_eval_clopen}
  \lean{AlgebraicGeometry.isOpenImmersion_overSpecMap_evalAlgHom,
    AlgebraicGeometry.opensRange_overSpecMap_evalAlgHom_disjoint,
    AlgebraicGeometry.exists_mem_opensRange_overSpecMap_evalAlgHom}
  \leanok
  \dcref{custom}

  \uses{pa-lem:isLocalization_away_piFactor, pa-def:overSpecMap}
  The morphisms \(\Spec \pi_i \colon \Spec B_i \to \Spec P\) (for the finite family
  \((B_i)_{i \in \iota}\)) are open immersions with pairwise disjoint ranges covering \(\Spec P\):
  the range of \(\Spec \pi_i\) is the basic open \(D(e_i)\), and \(D(e_i) \cap D(e_j) = \emptyset\)
  for \(i \ne j\) while \(\bigcup_i D(e_i) = \Spec P\).
\end{lemma}
\begin{proof}
  \leanok
  By Lemma~\ref{pa-lem:isLocalization_away_piFactor}, \(\pi_i\) realizes \(B_i\) as a localization of
  \(P\) away from \(e_i\), so \(\Spec \pi_i\) is an open immersion with range the basic open
  \(D(e_i)\) (the spectrum of a localization away from an element is that element's basic open). For
  \(i \ne j\), \(e_i e_j = 0\), whence \(D(e_i) \cap D(e_j) = D(e_i e_j) = D(0) = \emptyset\).
  Finally \(\sum_i e_i = 1\); since \(\iota\) is finite, no prime of \(P\) contains all the \(e_i\)
  (it would then contain \(1\)), so every prime lies in some \(D(e_i)\), i.e. the ranges cover.
\end{proof}

\begin{theorem}[Separation of relative Picard classes over a finite product]
  \label{pa-thm:relPic_pi_separation}
  \lean{AlgebraicGeometry.relPic.eq_of_pi_proj_eq}
  \leanok
  \dcref{custom}

  \uses{pa-thm:cechPic_clopen_separation, pa-thm:cechPic_clopen_gluing, pa-lem:range_whiskerLeft,
    pa-lem:opensRange_eval_clopen, pa-def:relPicAlgMap, pa-lem:snd_left_naturality}
  For a finite family \((B_i)_{i \in \iota}\) with product \(P\), two relative Picard classes
  \(\zeta, \zeta' \in \mathrm{relPic}(C, P)\) with \(\pi_i^{*}\zeta = \pi_i^{*}\zeta'\) in
  \(\mathrm{relPic}(C, B_i)\) for every \(i\) are equal.
\end{theorem}
\begin{proof}
  \leanok
  Represent \(\zeta, \zeta'\) by \v Cech Picard classes \(L, L'\) on the curve-product
  \((C \otimes \Spec P)_{\mathrm{left}}\) (a relative Picard class is a class on the curve-product
  modulo those pulled back from the base along \(\mathrm{pr}\)); restriction along \(\pi_i\)
  corresponds to \v Cech pullback of the underlying class along the whiskered projection
  \((C \lhd \Spec \pi_i)_{\mathrm{left}}\). The hypotheses give, for each \(i\), a class \(N_i\) on
  \(\Spec B_i\) with \(\mathrm{pr}^{*}N_i = (C \lhd \Spec\pi_i)_{\mathrm{left}}^{*}(L/L')\) (the
  restriction of \(L/L'\) to the \(i\)-th piece is pulled back from the base). Glue the \(N_i\) to a
  single class \(M\) on \(\Spec P\) by the clopen gluing (Theorem~\ref{pa-thm:cechPic_clopen_gluing})
  for the coordinate partition of \(\Spec P\) (Lemma~\ref{pa-lem:opensRange_eval_clopen}), so that
  \(\Spec\pi_i^{*} M = N_i\).

  It remains to show \(\mathrm{pr}^{*} M = L/L'\) on the curve-product; then \(\zeta = \zeta'\). The
  whiskered coordinate projections \((C \lhd \Spec\pi_i)_{\mathrm{left}}\) form a clopen partition
  of \((C \otimes \Spec P)_{\mathrm{left}}\): their ranges are the \(\mathrm{pr}\)-preimages of the
  ranges of the coordinate projections (Lemma~\ref{pa-lem:range_whiskerLeft}), which partition
  \(\Spec P\) (Lemma~\ref{pa-lem:opensRange_eval_clopen}). By separation along this partition
  (Theorem~\ref{pa-thm:cechPic_clopen_separation}) it suffices to check equal pullback along each
  whiskered projection. Using the base square of the projection
  (Lemma~\ref{pa-lem:snd_left_naturality}) and functoriality of \v Cech pullback,
  \[
    (C \lhd \Spec\pi_i)_{\mathrm{left}}^{*}(\mathrm{pr}^{*} M)
      = \mathrm{pr}^{*}\bigl(\Spec\pi_i^{*} M\bigr)
      = \mathrm{pr}^{*} N_i
      = (C \lhd \Spec\pi_i)_{\mathrm{left}}^{*}(L/L'),
  \]
  as required.
\end{proof}

\begin{theorem}[Gluing of relative Picard classes over a finite product]
  \label{pa-thm:relPic_pi_gluing}
  \lean{AlgebraicGeometry.relPic.exists_pi_lift}
  \leanok
  \dcref{custom}

  \uses{pa-thm:cechPic_clopen_gluing, pa-lem:range_whiskerLeft, pa-lem:opensRange_eval_clopen,
    pa-def:relPicAlgMap}
  For a finite family \((B_i)_{i \in \iota}\) with product \(P\), every family
  \((\xi_i)_{i \in \iota}\) with \(\xi_i \in \mathrm{relPic}(C, B_i)\) is the family of restrictions
  of a class \(\zeta \in \mathrm{relPic}(C, P)\): \(\pi_i^{*}\zeta = \xi_i\) for all \(i\).
\end{theorem}
\begin{proof}
  \leanok
  Represent each \(\xi_i\) by a \v Cech Picard class \(L_i\) on the curve-product
  \((C \otimes \Spec B_i)_{\mathrm{left}}\). The whiskered coordinate projections
  \((C \lhd \Spec\pi_i)_{\mathrm{left}}\) form a clopen partition of
  \((C \otimes \Spec P)_{\mathrm{left}}\), their ranges being the \(\mathrm{pr}\)-preimages
  (Lemma~\ref{pa-lem:range_whiskerLeft}) of the coordinate ranges that partition \(\Spec P\)
  (Lemma~\ref{pa-lem:opensRange_eval_clopen}). By the clopen gluing
  (Theorem~\ref{pa-thm:cechPic_clopen_gluing}) there is a class \(M\) on the curve-product with
  \((C \lhd \Spec\pi_i)_{\mathrm{left}}^{*} M = L_i\) for all \(i\). Let \(\zeta\) be the relative
  Picard class of \(M\). Then \(\pi_i^{*}\zeta\) is the relative Picard class of
  \((C \lhd \Spec\pi_i)_{\mathrm{left}}^{*} M = L_i\), which is \(\xi_i\).
\end{proof}

\section{Zariski separation of the \'etale plus construction}

Throughout, \(k\) is a field and \(C\) is a curve over \(k\) (proper, geometrically irreducible,
geometrically reduced). This section proves that the affine \'etale plus \(\mathrm{PicEtAff}(C, -)\)
is a separated presheaf for the Zariski topology of affine tests: two plus classes over \(A\)
agreeing on a finite covering family of localizations agree. The argument is instance-free of
\'etale separatedness; it consumes only the product decomposition of the relative Picard group.

\begin{theorem}[Unfolding an equality of restricted plus classes]
  \label{pa-thm:PicEtAff_exists_relPicAlgMap_eq}
  \lean{AlgebraicGeometry.PicEtAff.exists_relPicAlgMap_eq_of_mapAlg_eq}
  \leanok
  \dcref{custom}

  \uses{pa-thm:PicEtAff_mk_eq_mk_iff, pa-def:PicEtAff_map, pa-def:etale_cover_baseChange,
    pa-def:etale_cover_baseChangeInclude, pa-def:relPicAlgMap, pa-lem:relPicAlgMap_functor}
  Let \(\varphi_1 \colon A_1 \to R\) and \(\varphi_2 \colon A_2 \to R\) be homomorphisms of
  \(k\)-algebras, \(E_1\) a presented \'etale cover of \(A_1\) and \(E_2\) of \(A_2\), and
  \(\xi_1 \in \mathrm{Desc}(C, E_1)\), \(\xi_2 \in \mathrm{Desc}(C, E_2)\). If \(\varphi_1^{*}
  [E_1, \xi_1] = \varphi_2^{*}[E_2, \xi_2]\) in \(\mathrm{PicEtAff}(C, R)\), then there are a
  presented \'etale cover \(H\) of \(R\) and homomorphisms of \(k\)-algebras \(m_1 \colon B_{E_1}
  \to B_H\), \(m_2 \colon B_{E_2} \to B_H\), semilinear over \(\varphi_1, \varphi_2\)
  (\(m_1 \circ \iota_{E_1} = \iota_H \circ \varphi_1\) and \(m_2 \circ \iota_{E_2} = \iota_H \circ
  \varphi_2\)), with \(m_1^{*}\xi_1 = m_2^{*}\xi_2\) in \(\mathrm{relPic}(C, B_H)\).
\end{theorem}
\begin{proof}
  \leanok
  Regard \(R\) as an \(A_1\)- and \(A_2\)-algebra through \(\varphi_1, \varphi_2\). Then
  \(\varphi_i^{*}[E_i, \xi_i] = [(E_i)_R, \kappa_{E_i}^{*}\xi_i]\) is the restriction along the base
  extension \(A_i \to R\) of the plus class (Definition~\ref{pa-def:PicEtAff_map}), with \((E_i)_R\)
  the base change of \(E_i\) (Definition~\ref{pa-def:etale_cover_baseChange}) and \(\kappa_{E_i}\) its
  canonical inclusion (Definition~\ref{pa-def:etale_cover_baseChangeInclude}). The hypothesis
  \([(E_1)_R, \kappa^{*}\xi_1] = [(E_2)_R, \kappa^{*}\xi_2]\) unfolds by the presentation criterion
  (Theorem~\ref{pa-thm:PicEtAff_mk_eq_mk_iff}): there are a presented \'etale cover \(H\) of \(R\) and
  homomorphisms of \(R\)-algebras \(r_i \colon B_{(E_i)_R} \to B_H\) with
  \((r_1)_{*}(\kappa^{*}\xi_1) = (r_2)_{*}(\kappa^{*}\xi_2)\). Set \(m_i := r_i \circ \kappa_{E_i}
  \colon B_{E_i} \to B_H\). Since \(\kappa_{E_i}\) and \(r_i\) both cover the base inclusions,
  \(m_i \circ \iota_{E_i} = \iota_H \circ \varphi_i\); and by functoriality of restriction
  (Lemma~\ref{pa-lem:relPicAlgMap_functor}), \(m_i^{*}\xi_i = r_i^{*}(\kappa_{E_i}^{*}\xi_i)\), so the
  refinement equality gives \(m_1^{*}\xi_1 = m_2^{*}\xi_2\).
\end{proof}

\begin{theorem}[A finite covering family of localizations assembles into an \'etale cover]
  \label{pa-thm:PicEtAff_exists_etaleCover_pi}
  \lean{AlgebraicGeometry.PicEtAff.exists_etaleCover_pi}
  \leanok
  \dcref{custom}

  \uses{pa-def:etale_cover_of, pa-lem:etale_cover_ofEquiv, pa-def:etale_cover}
  Let \(A\) be a \(k\)-algebra and \((g_i)_{i \in \iota}\) a finite family with
  \(\sum_i (g_i) = (1)\), and let \(S_i\) be the localization of \(A\) away from \(g_i\). Given
  presented \'etale covers \(K_i\) of \(S_i\), the product \(\prod_i B_{K_i}\) is the carrier of a
  presented \'etale cover \(W\) of \(A\), with an isomorphism of \(A\)-algebras \(B_W \cong
  \prod_i B_{K_i}\).
\end{theorem}
\begin{proof}
  \leanok
  Each \(S_i\), a localization of \(A\) away from \(g_i\), is \'etale over \(A\); each \(B_{K_i}\)
  is \'etale over \(S_i\), hence over \(A\); so the finite product \(\prod_i B_{K_i}\) is \'etale
  over \(A\). It remains that \(\Spec(\prod_i B_{K_i}) \to \Spec A\) is surjective; then
  \(\prod_i B_{K_i}\) is a presented \'etale cover of \(A\) (Definition~\ref{pa-def:etale_cover_of},
  Lemma~\ref{pa-lem:etale_cover_ofEquiv}). Given a prime \(p\) of \(A\): since
  \(\sum_i (g_i) = (1)\), some \(g_i \notin p\), so \(p\) lifts to \(\Spec S_i\) (the spectrum of
  the localization away from \(g_i\) is the basic open \(D(g_i) \ni p\)); as \(K_i\) is a cover of
  \(S_i\) it is surjective on spectra (Definition~\ref{pa-def:etale_cover}), so \(p\) lifts further to
  \(\Spec B_{K_i}\), and thence to \(\Spec(\prod_j B_{K_j})\) through the \(i\)-th coordinate. Hence
  the structure map is surjective.
\end{proof}

\begin{theorem}[Zariski separation of the \'etale plus construction]
  \label{pa-thm:PicEtAff_eq_of_away_eq}
  \lean{AlgebraicGeometry.PicEtAff.eq_of_away_eq}
  \leanok
  \dcref{custom}

  \uses{pa-thm:relPic_pi_separation, pa-thm:PicEtAff_exists_relPicAlgMap_eq,
    pa-thm:PicEtAff_exists_etaleCover_pi, pa-lem:PicEtAff_ind, pa-thm:PicEtAff_mk_one, pa-def:PicEtAff_map}
  Let \(A\) be a \(k\)-algebra, \((g_i)_{i \in \iota}\) a finite family with \(\sum_i (g_i) = (1)\),
  and \(S_i\) the localization of \(A\) away from \(g_i\). Two plus classes \(x, y \in
  \mathrm{PicEtAff}(C, A)\) with equal restrictions along every \(A \to S_i\) are equal.
\end{theorem}
\begin{proof}
  \leanok
  Restriction is a homomorphism, so it suffices to show a class \(z\) with trivial restriction to
  every \(S_i\) is trivial. Represent \(z = [G, \zeta]\) (Lemma~\ref{pa-lem:PicEtAff_ind}). For each
  \(i\), comparing the restriction of \([G, \zeta]\) with that of the trivial class
  \([\mathrm{self}(A), 1]\) over \(S_i\) and unfolding (Theorem~\ref{pa-thm:PicEtAff_exists_relPicAlgMap_eq}),
  there are a presented \'etale cover \(K_i\) of \(S_i\) and a homomorphism \(m_i \colon B_G \to
  B_{K_i}\), semilinear over \(A \to S_i\), with \(m_i^{*}\zeta = 1\) in \(\mathrm{relPic}(C,
  B_{K_i})\). Assemble the product cover \(W\) of \(A\) from the \(K_i\)
  (Theorem~\ref{pa-thm:PicEtAff_exists_etaleCover_pi}), \(B_W \cong \prod_i B_{K_i}\) over \(A\). The
  joint map \(\langle m_i \rangle \colon B_G \to \prod_i B_{K_i}\) pushes \(\zeta\) to a class whose
  \(i\)-th factor is \(\pi_i^{*}(\langle m_i\rangle^{*}\zeta) = m_i^{*}\zeta = 1\); by the product
  separation (Theorem~\ref{pa-thm:relPic_pi_separation}) the pushed class is trivial in
  \(\mathrm{relPic}(C, \prod_i B_{K_i})\), hence in \(\mathrm{relPic}(C, B_W)\) after transport
  through the carrier isomorphism. As the joint map is an \(A\)-algebra refinement \(B_G \to B_W\),
  \([G, \zeta] = [W, 1] = 1\) (Theorem~\ref{pa-thm:PicEtAff_mk_one}). So \(z = 1\).
\end{proof}

\section{Restriction of relative Picard classes to an \'etale cover is injective}

This section extracts the usable corollary of the (C1) \'etale-separatedness theorem for the
Layer-2 gluing arguments. Throughout \(C\) is a curve over \(k\) (proper, geometrically irreducible,
geometrically reduced) and \(R\) is a commutative \(k\)-algebra.

\begin{theorem}[The unit represented on an arbitrary cover]
  \label{pa-thm:PicEtAff_unit_eq_mk}
  \lean{AlgebraicGeometry.PicEtAff.unit_eq_mk}
  \leanok
  \dcref{custom}

  \uses{pa-def:PicEtAff_unit, pa-lem:PicEtAff_tautological_mem_descentClasses, pa-def:PicEtAff_mk,
    pa-def:descentMap, pa-lem:PicEtAff_mk_descentMap, pa-def:etale_cover_fromSelf, pa-def:relPicAlgMap}
  For a presented \'etale cover \(E\) of \(R\) and \(z \in \mathrm{relPic}(C, R)\), the unit
  \(\eta_R(z)\) is represented on \(E\) by the tautological descent datum of \(z\):
  \[
    \eta_R(z) = [E, \iota_E^{*}z], \qquad \iota_E^{*}z \in \mathrm{Desc}(C, E),
  \]
  where \(\iota_E \colon R \to B_E\) is the structure map.
\end{theorem}
\begin{proof}
  \leanok
  By definition (Definition~\ref{pa-def:PicEtAff_unit}) \(\eta_R(z) = [\mathrm{self}(R), \iota^{*}z]\)
  with \(\iota^{*}z\) the tautological datum on the trivial cover
  (Lemma~\ref{pa-lem:PicEtAff_tautological_mem_descentClasses}). Transport along the canonical map
  \(\mathrm{fromSelf} \colon B_{\mathrm{self}(R)} \to B_E\) (Definition~\ref{pa-def:etale_cover_fromSelf})
  carries the tautological datum on \(\mathrm{self}(R)\) to the tautological datum
  \(\iota_E^{*}z\) on \(E\): both are the restriction of \(z\) along the structure map, by
  functoriality of restriction, and the composite structure map \(R \to B_{\mathrm{self}(R)} \to
  B_E\) is \(\iota_E\) (Definition~\ref{pa-def:relPicAlgMap}). By
  Lemma~\ref{pa-lem:PicEtAff_mk_descentMap}, transporting a representative along a refinement map does
  not change the plus class, so \([E, \iota_E^{*}z] = [E, \mathrm{fromSelf}_{*}(\iota^{*}z)] =
  [\mathrm{self}(R), \iota^{*}z] = \eta_R(z)\).
\end{proof}

\begin{theorem}[Restriction to an \'etale cover is injective --- the (C1) consumable]
  \label{pa-thm:relPicAlgMap_injective_of_etaleCover}
  \lean{AlgebraicGeometry.relPicAlgMap_injective_of_etaleCover}
  \leanok
  \dcref{custom}

  \uses{pa-thm:PicEtAff_unit_injective, pa-thm:PicEtAff_unit_eq_mk, pa-def:relPicAlgMap, pa-def:PicEtAff_mk}
  For a presented \'etale cover \(E\) of \(R\), restriction of relative Picard classes along the
  structure map \(\iota_E \colon R \to B_E\) is injective: if \(x, y \in \mathrm{relPic}(C, R)\)
  satisfy \(\iota_E^{*}x = \iota_E^{*}y\) in \(\mathrm{relPic}(C, B_E)\), then \(x = y\). This is
  what licenses every gluing step of the Layer-2 functor.
\end{theorem}
\begin{proof}
  \leanok
  By the \'etale separatedness of the relative Picard functor
  (Theorem~\ref{pa-thm:PicEtAff_unit_injective}), the unit \(\eta_R\) is injective, so it suffices to
  show \(\eta_R(x) = \eta_R(y)\). Represent both on \(E\) (Theorem~\ref{pa-thm:PicEtAff_unit_eq_mk}):
  \(\eta_R(x) = [E, \iota_E^{*}x]\) and \(\eta_R(y) = [E, \iota_E^{*}y]\). Since \(\iota_E^{*}x =
  \iota_E^{*}y\) by hypothesis, these classes coincide, so \(\eta_R(x) = \eta_R(y)\) and hence
  \(x = y\).
\end{proof}

\section{Zariski gluing of the \'etale plus construction}

Throughout, \(C\) is a curve over \(k\) (proper, geometrically irreducible, geometrically reduced).
This section proves the gluing half of the Zariski sheaf property of \(\mathrm{PicEtAff}(C, -)\) on
affine tests. Its engine, the pair condition, is the single point where the Layer-2 functoriality
consumes the (C1) \'etale separatedness.

\begin{theorem}[The pair condition]
  \label{pa-thm:PicEtAff_pair_condition}
  \lean{AlgebraicGeometry.PicEtAff.relPicAlgMap_tensor_eq_of_compat}
  \leanok
  \dcref{custom}

  \uses{pa-thm:relPicAlgMap_injective_of_etaleCover, pa-thm:PicEtAff_exists_relPicAlgMap_eq,
    pa-thm:relPicAlgMap_congr, pa-def:etale_cover_baseChange, pa-def:relPicAlgMap}
  Let \(S_1, S_2\) be localizations of a \(k\)-algebra \(A\) away from \(g_1, g_2\), and \(T_v\) a
  localization away from \(g_1 g_2\), with \(A\)-algebra homomorphisms \(\varphi_1 \colon S_1 \to
  T_v\), \(\varphi_2 \colon S_2 \to T_v\). Let \(E_1\) be a presented \'etale cover of \(S_1\),
  \(E_2\) of \(S_2\), and \(\xi_1 \in \mathrm{Desc}(C, E_1)\), \(\xi_2 \in \mathrm{Desc}(C, E_2)\).
  If \(\varphi_1^{*}[E_1, \xi_1] = \varphi_2^{*}[E_2, \xi_2]\) in \(\mathrm{PicEtAff}(C, T_v)\), then
  the two pullbacks of the underlying classes into the \(A\)-tensor product agree:
  \[
    \mathrm{inl}^{*}\xi_1 = \mathrm{inr}^{*}\xi_2 \quad\text{in } \mathrm{relPic}(C, B_{E_1}
    \otimes_A B_{E_2}),
  \]
  where \(\mathrm{inl} \colon B_{E_1} \to B_{E_1} \otimes_A B_{E_2}\) and \(\mathrm{inr} \colon
  B_{E_2} \to B_{E_1} \otimes_A B_{E_2}\) are the two coprojections.
\end{theorem}
\begin{proof}
  \leanok
  Unfold the overlap compatibility (Theorem~\ref{pa-thm:PicEtAff_exists_relPicAlgMap_eq}): there are a
  presented \'etale cover \(H\) of \(T_v\) and homomorphisms \(m_1 \colon B_{E_1} \to B_H\),
  \(m_2 \colon B_{E_2} \to B_H\), semilinear over \(\varphi_1, \varphi_2\), with \(m_1^{*}\xi_1 =
  m_2^{*}\xi_2\) in \(\mathrm{relPic}(C, B_H)\).

  In \(B_{E_1} \otimes_A B_{E_2}\), the element \(g_1 g_2\) is a unit: \(g_1\) is a unit in
  \(B_{E_1}\) and \(g_2\) in \(B_{E_2}\) (both localized-away), and their images multiply to
  \(g_1 g_2\). Since \(T_v\) is the localization of \(A\) away from \(g_1 g_2\), there is a unique
  homomorphism of \(A\)-algebras \(\rho \colon T_v \to B_{E_1} \otimes_A B_{E_2}\). Base-change \(H\)
  along \(\rho\) to a presented \'etale cover \(W := H_{B_{E_1} \otimes_A B_{E_2}}\)
  (Definition~\ref{pa-def:etale_cover_baseChange}), with canonical inclusion \(\theta \colon B_H \to
  B_W\). By the (C1) consumable (Theorem~\ref{pa-thm:relPicAlgMap_injective_of_etaleCover}), restriction
  along the structure map \(B_{E_1} \otimes_A B_{E_2} \to B_W\) is injective; it therefore suffices
  to show that the two sides become equal after transport into \(B_W\).

  Consider the two routes \(B_{E_1} \to B_W\). The first is \(\mathrm{inl}\) followed by the
  structure map into \(B_W\); the second is \(\theta \circ m_1\). Both are homomorphisms of
  \(S_1\)-algebras: for the second, \(m_1\) is \(\varphi_1\)-semilinear and \(\rho \circ \varphi_1\)
  equals the \(S_1\)-structure map \(S_1 \to B_{E_1} \otimes_A B_{E_2}\) (two homomorphisms of
  \(A\)-algebras out of the localization \(S_1\), hence equal). By refinement independence
  (Theorem~\ref{pa-thm:relPicAlgMap_congr}, applied to the two \(S_1\)-algebra maps out of the cover
  \(E_1\) of \(S_1\) and the descent class \(\xi_1\)), the two routes transport \(\xi_1\) to the
  same class in \(\mathrm{relPic}(C, B_W)\). Symmetrically the two routes \(B_{E_2} \to B_W\)
  (namely \(\mathrm{inr}\) into \(B_W\) and \(\theta \circ m_2\)) transport \(\xi_2\) equally. Now
  \[
    (\theta \circ m_1)^{*}\xi_1 = \theta^{*}(m_1^{*}\xi_1) = \theta^{*}(m_2^{*}\xi_2)
      = (\theta \circ m_2)^{*}\xi_2 ,
  \]
  so, chaining the two refinement independences, the transports into \(B_W\) of
  \(\mathrm{inl}^{*}\xi_1\) and \(\mathrm{inr}^{*}\xi_2\) agree; by injectivity,
  \(\mathrm{inl}^{*}\xi_1 = \mathrm{inr}^{*}\xi_2\).
\end{proof}

\begin{lemma}[Restriction along an algebra isomorphism is injective]
  \label{pa-lem:relPicAlgMap_algEquiv_injective}
  \lean{AlgebraicGeometry.relPicAlgMap_algEquiv_injective}
  \leanok
  \dcref{custom}

  \uses{pa-def:relPicAlgMap, pa-lem:relPicAlgMap_functor}
  For an isomorphism of \(k\)-algebras \(e \colon A \xrightarrow{\ \sim\ } B\), restriction of
  relative Picard classes along \(e\), \(e^{*} \colon \mathrm{relPic}(C, A) \to \mathrm{relPic}(C,
  B)\), is injective.
\end{lemma}
\begin{proof}
  \leanok
  Restriction along \(e^{-1}\) is a left inverse: \((e^{-1})^{*} \circ e^{*} = (e^{-1} \circ e)^{*}
  = (\id_A)^{*} = \id\), by functoriality (Lemma~\ref{pa-lem:relPicAlgMap_functor}). A map with a left
  inverse is injective.
\end{proof}

\begin{lemma}[Computing a restricted plus class on a given cover]
  \label{pa-lem:PicEtAff_mapAlg_mk_eq_mk}
  \lean{AlgebraicGeometry.PicEtAff.mapAlg_mk_eq_mk}
  \leanok
  \dcref{custom}

  \uses{pa-thm:PicEtAff_mk_eq_mk_iff, pa-def:PicEtAff_map, pa-def:PicEtAff_mapAlg,
    pa-def:etale_cover_baseChange, pa-def:relPicAlgMap}
  Let \(\varphi \colon A_1 \to R\) be a homomorphism of \(k\)-algebras, \(E_1\) a presented \'etale
  cover of \(A_1\), \(F\) a presented \'etale cover of \(R\), \(\xi \in \mathrm{Desc}(C, E_1)\) and
  \(\eta \in \mathrm{Desc}(C, F)\). Suppose there is a homomorphism of \(k\)-algebras \(n \colon
  B_{E_1} \to B_F\), semilinear over \(\varphi\) (\(n \circ \iota_{E_1} = \iota_F \circ \varphi\)),
  with \(n^{*}\xi = \eta\) in \(\mathrm{relPic}(C, B_F)\). Then \(\varphi^{*}[E_1, \xi] = [F, \eta]\).
\end{lemma}
\begin{proof}
  \leanok
  Regard \(R\) as an \(A_1\)-algebra through \(\varphi\), so that \(\varphi^{*}[E_1, \xi] =
  [(E_1)_R, \kappa_{E_1}^{*}\xi]\) (Definition~\ref{pa-def:PicEtAff_map}, with \((E_1)_R\) the base
  change, Definition~\ref{pa-def:etale_cover_baseChange}). The semilinear map \(n\) extends uniquely to
  an \(R\)-algebra refinement \(n_R \colon B_{(E_1)_R} \to B_F\) with \(n_R \circ \kappa_{E_1} = n\).
  By the presentation criterion (Theorem~\ref{pa-thm:PicEtAff_mk_eq_mk_iff}, with the common cover
  \(F\), the map \(n_R\) out of \((E_1)_R\) and the identity out of \(F\)), \([(E_1)_R,
  \kappa_{E_1}^{*}\xi] = [F, (n_R)_{*}(\kappa_{E_1}^{*}\xi)]\); and \((n_R)_{*}(\kappa_{E_1}^{*}\xi)
  = n^{*}\xi = \eta\) by functoriality (Definition~\ref{pa-def:relPicAlgMap}). Hence
  \(\varphi^{*}[E_1, \xi] = [F, \eta]\).
\end{proof}

\begin{theorem}[Zariski gluing of the \'etale plus construction]
  \label{pa-thm:PicEtAff_exists_mapAlg_eq_of_compat}
  \lean{AlgebraicGeometry.PicEtAff.exists_mapAlg_eq_of_compat}
  \leanok
  \dcref{custom}

  \uses{pa-thm:PicEtAff_pair_condition, pa-thm:PicEtAff_exists_etaleCover_pi, pa-thm:relPic_pi_gluing,
    pa-thm:relPic_pi_separation, pa-lem:PicEtAff_mapAlg_mk_eq_mk, pa-lem:relPicAlgMap_algEquiv_injective,
    pa-lem:PicEtAff_ind}
  Let \(A\) be a \(k\)-algebra, \((g_i)_{i \in \iota}\) a finite family with \(\sum_i (g_i) = (1)\),
  \(S_i\) the localization of \(A\) away from \(g_i\), and \(T_{ij}\) the localization away from
  \(g_i g_j\). A family \((x_i)_{i \in \iota}\) with \(x_i \in \mathrm{PicEtAff}(C, S_i)\),
  compatible on the overlaps (the restrictions of \(x_i\) and \(x_j\) to \(T_{ij}\) agree for all
  \(i, j\)), is the family of restrictions of a plus class \(z \in \mathrm{PicEtAff}(C, A)\):
  restriction of \(z\) to \(S_i\) is \(x_i\). Together with Theorem~\ref{pa-thm:PicEtAff_eq_of_away_eq}
  this is the Zariski sheaf property of \(\mathrm{PicEtAff}(C, -)\) on affine tests.
\end{theorem}
\begin{proof}
  \leanok
  Represent \(x_i = [E_i, \xi_i]\) (Lemma~\ref{pa-lem:PicEtAff_ind}). For each \(i, j\), the overlap
  compatibility over \(T_{ij}\) yields, by the pair condition
  (Theorem~\ref{pa-thm:PicEtAff_pair_condition}), \(\mathrm{inl}^{*}\xi_i = \mathrm{inr}^{*}\xi_j\) in
  \(\mathrm{relPic}(C, B_{E_i} \otimes_A B_{E_j})\). Lift the underlying classes to the product: by
  the product gluing (Theorem~\ref{pa-thm:relPic_pi_gluing}) there is \(\zeta \in \mathrm{relPic}(C,
  \prod_i B_{E_i})\) with \(\pi_i^{*}\zeta = \xi_i\). Assemble the product cover \(W\) of \(A\) from
  the \(E_i\) (Theorem~\ref{pa-thm:PicEtAff_exists_etaleCover_pi}), \(B_W \cong \prod_i B_{E_i}\) over
  \(A\); transport \(\zeta\) to \(\zeta_W \in \mathrm{relPic}(C, B_W)\).

  \emph{\(\zeta_W\) is a descent class on \(W\).} The descent condition \(\mathrm{inl}^{*}\zeta_W =
  \mathrm{inr}^{*}\zeta_W\) in \(\mathrm{relPic}(C, B_W \otimes_A B_W)\) is checked through the
  \(\pi\)-decomposition of the tensor square, \(B_W \otimes_A B_W \cong \prod_{(i, j)} B_{E_i}
  \otimes_A B_{E_j}\): by injectivity of transport along this isomorphism
  (Lemma~\ref{pa-lem:relPicAlgMap_algEquiv_injective}) and product separation over the index set
  \(\iota \times \iota\) (Theorem~\ref{pa-thm:relPic_pi_separation}), it suffices to compare
  components. The \((i, j)\)-component of \(\mathrm{inl}^{*}\zeta_W\) is \(\mathrm{inl}^{*}
  (\pi_i^{*}\zeta) = \mathrm{inl}^{*}\xi_i\) and of \(\mathrm{inr}^{*}\zeta_W\) is
  \(\mathrm{inr}^{*}(\pi_j^{*}\zeta) = \mathrm{inr}^{*}\xi_j\); these agree by the pair conditions.

  \emph{The glued class.} Set \(z := [W, \zeta_W]\). Its restriction to \(S_i\) is \(x_i\): the
  \(i\)-th projection \(B_W \cong \prod_j B_{E_j} \to B_{E_i}\) is a homomorphism semilinear over
  \(A \to S_i\) carrying \(\zeta_W\) onto \(\pi_i^{*}\zeta = \xi_i\), so by
  Lemma~\ref{pa-lem:PicEtAff_mapAlg_mk_eq_mk}, \((A \to S_i)^{*}[W, \zeta_W] = [E_i, \xi_i] = x_i\).
\end{proof}

\section{The \'etale-sheafified relative Picard functor: the affine-opens limit}

The frozen group-object interface forces the relative Picard functor to exist on \emph{all} of
\(\Over(\Spec k)\), not only on affine tests. Following the resolution of the Layer-2 vehicle
question, the extension is the \emph{bespoke affine-opens limit}: a section over a test object
\(T\) is a compatible family of affine-level plus classes over the affine opens of
\(T_{\mathrm{left}}\). The index poset is \(\mathrm{Type}\)-small, which is the point of the
vehicle. Throughout, \(k\) is a field and \(C\) is a curve over \(k\).

\begin{definition}[Transport of plus classes along an algebra isomorphism]
  \label{pa-def:PicEtAff_congr}
  \lean{AlgebraicGeometry.PicEtAff.congr}
  \leanok
  \dcref{custom}

  \uses{pa-def:PicEtAff_map, pa-thm:PicEtAff_map_map, pa-thm:PicEtAff_map_id}
  For an isomorphism of \(k\)-algebras \(e \colon A \xrightarrow{\ \sim\ } B\), restriction of plus
  classes along \(e\) is a group isomorphism \(\mathrm{PicEtAff}(C, A) \xrightarrow{\ \sim\ }
  \mathrm{PicEtAff}(C, B)\), with inverse restriction along \(e^{-1}\).
\end{definition}
\begin{proof}
  \leanok
  Restriction along \(e\) and along \(e^{-1}\) are homomorphisms of groups
  (Definition~\ref{pa-def:PicEtAff_map}); their two composites are restriction along
  \(e^{-1} \circ e = \id_A\) and along \(e \circ e^{-1} = \id_B\), each the identity by transitivity
  of restriction (Theorem~\ref{pa-thm:PicEtAff_map_map}) and its normalization at the identity
  (Theorem~\ref{pa-thm:PicEtAff_map_id}). Hence the two are mutually inverse.
\end{proof}

\begin{definition}[The affine identification of global sections]
  \label{pa-def:overSpecGammaTop}
  \lean{AlgebraicGeometry.Over.overSpecΓTopAlgEquiv}
  \leanok
  \dcref{custom}

  For a \(k\)-algebra \(A\), the global sections of the affine test \(\Spec A\) over its top open
  are identified with \(A\) as a \(k\)-algebra: \(\Gamma((\Spec A)_{\mathrm{left}}, \top)
  \xrightarrow{\ \sim\ } A\), through the canonical isomorphism \(\Gamma(\Spec A, \top) \cong A\) of
  the affine scheme.
\end{definition}
\begin{proof}
  \leanok
  The canonical isomorphism of rings \(\Gamma(\Spec A, \top) \cong A\) intertwines the two
  \(k\)-algebra structures: the structure map of the sections and that of \(A\) both come from the
  \(k\)-algebra structure of \(A\) through this same isomorphism. Hence it is an isomorphism of
  \(k\)-algebras.
\end{proof}

\begin{definition}[The \'etale-sheafified relative Picard functor at a test object]
  \label{pa-def:picEt}
  \lean{AlgebraicGeometry.picEtSubgroup, AlgebraicGeometry.picEt}
  \leanok
  \source{kleiman-picard}

  \uses{pa-def:PicEtAff, pa-def:PicEtAff_map}
  \dcref{custom}
  For a test object \(T\) of \(\Over(\Spec k)\), \(\mathrm{picEt}(C, T)\) is the group of
  \emph{compatible families}: an element assigns to each affine open \(V \subseteq T_{\mathrm{left}}\)
  a plus class \(s_V \in \mathrm{PicEtAff}(C, \Gamma(T, V))\) such that for affine opens
  \(U \subseteq V\), the value \(s_U\) is the restriction of \(s_V\) along the section-restriction
  homomorphism \(\Gamma(T, V) \to \Gamma(T, U)\) (Definition~\ref{pa-def:PicEtAff_map}). Componentwise
  product makes \(\mathrm{picEt}(C, T)\) a commutative group, a subgroup of the product
  \(\prod_{V} \mathrm{PicEtAff}(C, \Gamma(T, V))\); evaluation at an affine open \(V\) is the
  homomorphism \(s \mapsto s_V\). On affine tests this is the plus itself (the affine comparison,
  Definition~\ref{pa-def:picEtAffineEquiv}); on general tests it is the canonical Zariski-continuous
  extension --- the realization on all test objects of the \'etale sheafification
  \(\Pic_{(X/S)\mathrm{\acute et}}\) of the relative Picard functor.
\end{definition}
\begin{proof}
  \leanok
  The compatibility condition is closed under the componentwise group operations: the constant
  family \(1\) is compatible (restriction preserves the identity); and if \(s, t\) are compatible,
  so are \(s \cdot t\) and \(s^{-1}\), because restriction of plus classes is a homomorphism of
  groups (Definition~\ref{pa-def:PicEtAff_map}). Hence the compatible families form a subgroup of the
  product group, which is therefore a commutative group.
\end{proof}

\begin{definition}[The affine comparison isomorphism]
  \label{pa-def:picEtAffineEquiv}
  \lean{AlgebraicGeometry.picEtAffineEquiv}
  \leanok
  \dcref{custom}

  \uses{pa-def:picEt, pa-def:PicEtAff_congr, pa-def:overSpecGammaTop, pa-def:PicEtAff_map}
  On an affine test \(\Spec A\), evaluation at the top affine open \(\top\) is a group isomorphism
  \[
    \mathrm{picEt}(C, \Spec A) \xrightarrow{\ \sim\ } \mathrm{PicEtAff}(C, A),
  \]
  transported through the identification \(\Gamma(\top) \cong A\) of
  Definition~\ref{pa-def:overSpecGammaTop}.
\end{definition}
\begin{proof}
  \leanok
  The affine-opens poset of \(\Spec A\) has a terminal element, the top open \(\top\) (all of
  \(\Spec A\), an affine open containing every affine open), so a compatible family is determined by
  its value at \(\top\). The forward map sends \(s\) to \(s_\top\) carried through \(\Gamma(\top)
  \cong A\) (Definition~\ref{pa-def:overSpecGammaTop}, Definition~\ref{pa-def:PicEtAff_congr}). The
  backward map sends \(x \in \mathrm{PicEtAff}(C, A)\) to the family whose value at \(V\) is \(x\)
  carried to \(\Gamma(\top)\) and restricted to \(\Gamma(V)\) (Definition~\ref{pa-def:PicEtAff_map});
  this family is compatible because restriction is functorial (restricting from \(\top\) to \(V\)
  and then to \(U\) is restricting to \(U\)). The two are mutually inverse: forward-then-backward
  recovers a family \(s\) from its \(\top\)-value, since \(s_V\) is the restriction of \(s_\top\) by
  the compatibility of \(s\) at \(V \subseteq \top\); backward-then-forward returns \(x\), since
  restriction from \(\top\) to \(\top\) is the identity and \(\Gamma(\top) \cong A\) cancels. Both
  maps are homomorphisms of groups, so this is a group isomorphism.
\end{proof}

\section{Functoriality by the unique-pullback-value characterization}

The Layer-2 functor \(\mathrm{picEt}\) is functorial along an \emph{arbitrary} morphism
\(f \colon T' \to T\) of test objects. The value of the restricted family at an affine open
\(W \subseteq T'_{\mathrm{left}}\) --- whose \(f\)-image need not lie in any affine open of
\(T_{\mathrm{left}}\) --- is glued over a finite basic-open refinement of \(W\) subordinate to the
preimages of affine opens, using the Zariski sheaf property of the affine \'etale Picard functor.
Throughout, \(C\) is a curve over \(k\).

\begin{theorem}[Finite basic subcovers subordinate to a constraint]
  \label{pa-thm:exists_basic_subcover}
  \lean{AlgebraicGeometry.Scheme.exists_basic_subcover}
  \leanok
  \dcref{custom}

  Let \(W\) be an affine open of a scheme \(X\) and \(P\) a property of opens such that every point
  of \(W\) has an open neighbourhood inside \(W\) satisfying \(P\). Then there is a finite family of
  sections \(r \in \Gamma(X, W)\) spanning the unit ideal whose basic opens \(D(r)\) each lie inside
  some open satisfying \(P\).
\end{theorem}
\begin{proof}
  \leanok
  For each \(w \in W\), choose an open \(U_0 \ni w\) with \(P(U_0)\); since the basic opens of the
  affine \(W\) form a basis of its topology, choose \(r_w \in \Gamma(X, W)\) with \(w \in D(r_w)
  \subseteq U_0\). The \(D(r_w)\) cover \(W\); as \(W\) is quasi-compact (being affine), finitely
  many \(D(r_{w_1}), \dots, D(r_{w_n})\) already cover \(W\), which means the \(r_{w_j}\) span the
  unit ideal of \(\Gamma(X, W)\). Each \(D(r_{w_j}) \subseteq U_0\) lies inside an open with \(P\).
\end{proof}

\begin{theorem}[Zariski separation of the plus over the section rings of a basic cover]
  \label{pa-thm:picEt_eq_of_basic_eq}
  \lean{AlgebraicGeometry.picEt.eq_of_basic_eq}
  \leanok
  \dcref{custom}

  \uses{pa-thm:PicEtAff_eq_of_away_eq}
  Let \(W\) be an affine open of a test object \(T\) and \((r_i)_{i \in s}\) a finite family in
  \(\Gamma(T, W)\) spanning the unit ideal. Two plus classes \(x, y \in \mathrm{PicEtAff}(C,
  \Gamma(T, W))\) whose restrictions to \(\Gamma(T, D(r_i))\) agree for every \(i\) are equal.
\end{theorem}
\begin{proof}
  \leanok
  The section ring \(\Gamma(T, D(r_i))\) of a basic open of the affine open \(W\) is the
  localization of \(\Gamma(T, W)\) away from \(r_i\), and the \(r_i\) span the unit ideal. Apply the
  Zariski separation of the plus construction (Theorem~\ref{pa-thm:PicEtAff_eq_of_away_eq}) to this
  finite covering family of localizations of \(\Gamma(T, W)\).
\end{proof}

\begin{definition}[The pullback-value property]
  \label{pa-def:picEt_IsPullbackValue}
  \lean{AlgebraicGeometry.picEt.IsPullbackValue}
  \leanok
  \dcref{custom}

  \uses{pa-def:picEt, pa-def:PicEtAff_map}
  Let \(f \colon T' \to T\) be a morphism of test objects, \(s \in \mathrm{picEt}(C, T)\), and
  \(W\) an affine open of \(T'_{\mathrm{left}}\). A plus class \(v \in \mathrm{PicEtAff}(C,
  \Gamma(T', W))\) has the \emph{pullback-value property} for \((f, s, W)\) if, for every affine
  open \(W_0 \subseteq W\) contained in the preimage \(f^{-1}V\) of some affine open \(V\) of
  \(T_{\mathrm{left}}\), the restriction of \(v\) to \(\Gamma(T', W_0)\) equals the pullback of
  \(s_V\) along the section-pullback homomorphism \(\Gamma(T, V) \to \Gamma(T', W_0)\) (pullback of
  sections along \(f\), a \(k\)-algebra map; Definition~\ref{pa-def:PicEtAff_map}).
\end{definition}

\begin{theorem}[Choice independence]
  \label{pa-thm:picEt_mapAlg_appLE_eq}
  \lean{AlgebraicGeometry.picEt.mapAlg_appLE_eq}
  \leanok
  \dcref{custom}

  \uses{pa-thm:picEt_eq_of_basic_eq, pa-thm:exists_basic_subcover, pa-def:picEt}
  Let \(f \colon T' \to T\), \(s \in \mathrm{picEt}(C, T)\), \(W_0\) an affine open of
  \(T'_{\mathrm{left}}\), and \(V, V'\) affine opens of \(T_{\mathrm{left}}\) with \(W_0 \subseteq
  f^{-1}V\) and \(W_0 \subseteq f^{-1}V'\). Then the pullback of \(s_V\) and the pullback of
  \(s_{V'}\) to \(\Gamma(T', W_0)\) (along the two section-pullback maps) agree.
\end{theorem}
\begin{proof}
  \leanok
  Refine \(W_0\) by basic opens: for each point \(w \in W_0\), the image \(f(w) \in V \cap V'\) lies
  in an affine open \(V'' \subseteq V \cap V'\) of \(T_{\mathrm{left}}\); choose a basic open
  \(D(r) \ni w\) of \(W_0\) with \(f(D(r)) \subseteq V''\), i.e. \(D(r) \subseteq f^{-1}V''\). By
  Theorem~\ref{pa-thm:exists_basic_subcover} a finite family of such \(r\) spans the unit ideal. On
  each \(D(r)\), the pullback of \(s_V\) restricted to \(\Gamma(T', D(r))\) factors as the pullback
  of \(s_{V''}\): using the compatibility \(s_{V''} = s_V|_{V''}\) (as \(V'' \subseteq V\)) and the
  functoriality of section pullback and restriction, it equals the pullback of \(s_{V''}\) along
  \(f\) from \(V''\) to \(D(r)\); likewise for \(s_{V'}\). Hence the two pullbacks agree on every
  \(D(r)\), so by Theorem~\ref{pa-thm:picEt_eq_of_basic_eq} they agree on \(\Gamma(T', W_0)\).
\end{proof}

\begin{theorem}[Uniqueness of the pullback value]
  \label{pa-thm:picEt_pullbackValue_unique}
  \lean{AlgebraicGeometry.picEt.pullbackValue_unique}
  \leanok
  \dcref{custom}

  \uses{pa-def:picEt_IsPullbackValue, pa-thm:picEt_eq_of_basic_eq, pa-thm:exists_basic_subcover}
  For fixed \((f, s, W)\), a plus class with the pullback-value property is unique: if \(v, v'\)
  both have the property, then \(v = v'\).
\end{theorem}
\begin{proof}
  \leanok
  Refine \(W\) by basic opens \(D(r)\) each contained in the preimage of an affine open of
  \(T_{\mathrm{left}}\): every point of \(W\) maps into some affine open \(V\) of
  \(T_{\mathrm{left}}\), so has a basic-open neighbourhood in \(W\) inside \(f^{-1}V\; \)by
  Theorem~\ref{pa-thm:exists_basic_subcover} a finite such family spans the unit ideal. On each
  \(D(r) \subseteq f^{-1}V\), both \(v\) and \(v'\) restrict to the pullback of \(s_V\) (the
  defining property, Definition~\ref{pa-def:picEt_IsPullbackValue}), hence agree. By
  Theorem~\ref{pa-thm:picEt_eq_of_basic_eq}, \(v = v'\).
\end{proof}

\begin{theorem}[Existence of the pullback value]
  \label{pa-thm:picEt_exists_isPullbackValue}
  \lean{AlgebraicGeometry.picEt.exists_isPullbackValue,
    AlgebraicGeometry.picEt.existsUnique_isPullbackValue}
  \leanok
  \dcref{custom}

  \uses{pa-def:picEt_IsPullbackValue, pa-thm:PicEtAff_exists_mapAlg_eq_of_compat,
    pa-thm:picEt_eq_of_basic_eq, pa-thm:exists_basic_subcover, pa-thm:picEt_mapAlg_appLE_eq,
    pa-thm:picEt_pullbackValue_unique}
  For every \(f \colon T' \to T\), \(s \in \mathrm{picEt}(C, T)\) and affine open \(W\) of
  \(T'_{\mathrm{left}}\), there exists a unique plus class \(v \in \mathrm{PicEtAff}(C, \Gamma(T',
  W))\) with the pullback-value property for \((f, s, W)\).
\end{theorem}
\begin{proof}
  \leanok
  Choose (Theorem~\ref{pa-thm:exists_basic_subcover}) a finite family \((r)\) in \(\Gamma(T', W)\)
  spanning the unit ideal such that each \(D(r)\) lies in the preimage \(f^{-1}V_r\) of an affine
  open \(V_r\) of \(T_{\mathrm{left}}\) (every point of \(W\) maps into some affine open). On
  \(D(r)\) take the local value to be the pullback of \(s_{V_r}\) along \(f\). On an overlap
  \(D(r) \cap D(r') = D(r r')\), the two local values agree: both restrict, through the affine
  opens \(V_r, V_{r'}\), to the same plus class by choice independence
  (Theorem~\ref{pa-thm:picEt_mapAlg_appLE_eq}). The section rings \(\Gamma(T', D(r))\) are localizations
  of \(\Gamma(T', W)\) away from the \(r\), which span the unit ideal, and the overlaps are
  localizations away from \(r r'\); so by the Zariski gluing of the plus construction
  (Theorem~\ref{pa-thm:PicEtAff_exists_mapAlg_eq_of_compat}) the local values glue to a plus class
  \(v \in \mathrm{PicEtAff}(C, \Gamma(T', W))\).

  This \(v\) has the pullback-value property: for an affine \(W_0 \subseteq W\) inside \(f^{-1}V\),
  refine \(W_0\) by basic opens \(D(q)\) each inside some \(D(r)\)
  (Theorem~\ref{pa-thm:exists_basic_subcover}). On \(D(q)\), \(v\) restricts through the glued datum on
  \(D(r)\) to the pullback of \(s_{V_r}\), which by choice independence
  (Theorem~\ref{pa-thm:picEt_mapAlg_appLE_eq}) equals the pullback of \(s_V\); by
  Theorem~\ref{pa-thm:picEt_eq_of_basic_eq}, the restriction of \(v\) to \(\Gamma(T', W_0)\) equals the
  pullback of \(s_V\). Uniqueness is Theorem~\ref{pa-thm:picEt_pullbackValue_unique}.
\end{proof}

\begin{definition}[Restriction of \(\mathrm{picEt}\) along an arbitrary morphism]
  \label{pa-def:picEtMap}
  \lean{AlgebraicGeometry.picEtMapVal, AlgebraicGeometry.picEtMap}
  \leanok
  \dcref{custom}

  \uses{pa-thm:picEt_exists_isPullbackValue, pa-def:picEt}
  For a morphism \(f \colon T' \to T\) of test objects, the homomorphism of groups \(f^{*} \colon
  \mathrm{picEt}(C, T) \to \mathrm{picEt}(C, T')\) whose value at an affine open \(W\) of
  \(T'_{\mathrm{left}}\) is the unique plus class with the pullback-value property for \((f, s, W)\)
  (Theorem~\ref{pa-thm:picEt_exists_isPullbackValue}).
\end{definition}
\begin{proof}
  \leanok
  The values form a compatible family: for affine opens \(U \subseteq W\), the restriction of the
  value at \(W\) still has the pullback-value property at \(U\) (the property is stable under
  further restriction, by functoriality of section pullback and restriction), so by uniqueness it
  equals the value at \(U\); thus \((W \mapsto v_W)\) is a section of \(\mathrm{picEt}(C, T')\). It
  is a homomorphism: the pullback-value property is preserved by the componentwise group operations
  (pullback of \(s\) is a homomorphism), so by uniqueness the value at \(s \cdot t\) is the product
  of the values and the value at \(1\) is \(1\).
\end{proof}

\begin{theorem}[The functor laws]
  \label{pa-thm:picEtMap_functor_laws}
  \lean{AlgebraicGeometry.picEtMap_id, AlgebraicGeometry.picEtMap_comp}
  \leanok
  \dcref{custom}

  \uses{pa-def:picEtMap, pa-thm:picEt_exists_isPullbackValue, pa-thm:exists_basic_subcover,
    pa-thm:picEt_eq_of_basic_eq, pa-def:picEt}
  Restriction along the identity is the identity, \((\id_T)^{*} = \id\) on \(\mathrm{picEt}(C, T)\);
  and for \(f \colon T' \to T\) and \(g \colon T'' \to T'\), restriction along the composite is the
  composite of the restrictions, \((f \circ g)^{*} = g^{*} \circ f^{*}\).
\end{theorem}
\begin{proof}
  \leanok
  \emph{Identity.} The value of \((\id_T)^{*}s\) at an affine open \(W\) is the unique class with
  the pullback-value property; \(s_W\) itself has it, since for \(W_0 \subseteq W\) inside
  \(\id^{-1}V = V\) the restriction of \(s_W\) is \(s_{W_0}\), the pullback of \(s_V\) along the
  identity (by compatibility of \(s\)). By uniqueness (Theorem~\ref{pa-thm:picEt_exists_isPullbackValue})
  the value is \(s_W\), so \((\id_T)^{*}s = s\).

  \emph{Composite.} It suffices to show the value \((g^{*}(f^{*}s))_W\) has the pullback-value
  property for \((f \circ g, s, W)\); uniqueness then identifies it with \(((f \circ g)^{*}s)_W\).
  Fix an affine \(W \subseteq T''_{\mathrm{left}}\) and \(W_0 \subseteq W\) inside
  \((f \circ g)^{-1}V\). Refine \(W_0\) by basic opens \(D(q)\) whose \(g\)-image lies in an affine
  open \(V'' \subseteq f^{-1}V\) of \(T'_{\mathrm{left}}\)
  (Theorem~\ref{pa-thm:exists_basic_subcover}). On \(D(q)\): the value of \(f^{*}s\) at \(V''\) is the
  pullback of \(s_V\) (the pullback-value property of \(f^{*}s\), as \(V'' \subseteq f^{-1}V\)), and
  the value of \(g^{*}(f^{*}s)\) at \(D(q)\) is the pullback of \((f^{*}s)_{V''}\) along \(g\), which
  by functoriality of section pullback equals the pullback of \(s_V\) along \(f \circ g\). By
  Theorem~\ref{pa-thm:picEt_eq_of_basic_eq}, the value of \(g^{*}(f^{*}s)\) at \(W\) has the
  pullback-value property for \((f \circ g, s, W)\), as required.
\end{proof}

\begin{definition}[The \'etale-sheafified relative Picard functor]
  \label{pa-def:picEtFunctor}
  \lean{AlgebraicGeometry.picEtFunctor}
  \leanok
  \dcref{custom}

  \uses{pa-def:picEt, pa-def:picEtMap, pa-thm:picEtMap_functor_laws}
  For a curve \(C\) over \(k\) (proper, geometrically irreducible, geometrically reduced), the
  contravariant functor
  \[
    \mathrm{picEt}(C, -) \colon \bigl(\Over(\Spec k)\bigr)^{\mathrm{op}} \to \mathbf{CommGrp}
  \]
  sending a test object \(T\) to the commutative group \(\mathrm{picEt}(C, T)\)
  (Definition~\ref{pa-def:picEt}) and a morphism \(f\) to the restriction homomorphism \(f^{*}\)
  (Definition~\ref{pa-def:picEtMap}).
\end{definition}
\begin{proof}
  \leanok
  The object and morphism assignments are given by Definition~\ref{pa-def:picEt} and
  Definition~\ref{pa-def:picEtMap}; the identity and composition laws of a functor are exactly the two
  functor laws of restriction (Theorem~\ref{pa-thm:picEtMap_functor_laws}).
\end{proof}

\begin{theorem}[Naturality of the affine comparison]
  \label{pa-thm:picEtAffineEquiv_naturality}
  \lean{AlgebraicGeometry.picEtAffineEquiv_naturality}
  \leanok
  \dcref{custom}

  \uses{pa-def:picEtAffineEquiv, pa-def:picEtMap, pa-def:PicEtAff_map, pa-def:overSpecGammaTop,
    pa-thm:picEt_exists_isPullbackValue}
  The affine comparison is natural in the test algebra: for a homomorphism of \(k\)-algebras
  \(\varphi \colon A \to B\) and \(s \in \mathrm{picEt}(C, \Spec A)\), the affine comparison of
  \((\Spec\varphi)^{*}s\) equals the restriction along \(\varphi\) of the affine comparison of
  \(s\). Under the comparison isomorphisms (Definition~\ref{pa-def:picEtAffineEquiv}), restriction along
  \(\Spec\varphi\) corresponds to restriction of plus classes along \(\varphi\).
\end{theorem}
\begin{proof}
  \leanok
  Both sides are computed at the top affine opens. Since \(\top_B \subseteq (\Spec\varphi)^{-1}
  \top_A\), the value of \((\Spec\varphi)^{*}s\) at \(\top_B\) is, by the pullback-value property
  (Theorem~\ref{pa-thm:picEt_exists_isPullbackValue}), the pullback of \(s_{\top_A}\) along the
  section-pullback map \(\Gamma(\Spec A, \top) \to \Gamma(\Spec B, \top)\). Under the identifications
  \(\Gamma(\top) \cong A\) and \(\Gamma(\top) \cong B\) (Definition~\ref{pa-def:overSpecGammaTop}), that
  section-pullback map corresponds to \(\varphi\). Transporting through the comparison isomorphisms
  (Definition~\ref{pa-def:picEtAffineEquiv}), the value corresponds to the restriction along
  \(\varphi\) (Definition~\ref{pa-def:PicEtAff_map}) of the affine comparison of \(s\).
\end{proof}

\begin{theorem}[Uniqueness of gluings for \(\mathrm{picEt}\)]
  \label{pa-thm:picEt_ext_of_le_cover}
  \lean{AlgebraicGeometry.picEt.ext_of_le_cover}
  \leanok
  \dcref{custom}

  \uses{pa-def:picEt, pa-thm:picEt_eq_of_basic_eq, pa-thm:exists_basic_subcover}
  Two sections \(s, t \in \mathrm{picEt}(C, T)\) agreeing on every affine open subordinate to an
  open cover of \(T_{\mathrm{left}}\) are equal: if \((O_i)_{i}\) is a family of opens covering
  \(T_{\mathrm{left}}\) and \(s_U = t_U\) for every affine open \(U\) contained in some \(O_i\),
  then \(s = t\). This is the separation half of the sheaf property of the Layer-2 functor.
\end{theorem}
\begin{proof}
  \leanok
  Fix an affine open \(U\) of \(T_{\mathrm{left}}\). Refine \(U\) by basic opens \(D(r)\) each
  contained in some \(O_i\): every point of \(U\) lies in some \(O_i\), so has a basic-open
  neighbourhood in \(U\) inside \(O_i\); by Theorem~\ref{pa-thm:exists_basic_subcover} a finite such
  family spans the unit ideal. On \(D(r) \subseteq O_i\), the restriction of \(s_U\) equals
  \(s_{D(r)} = t_{D(r)}\) (the hypothesis, as \(D(r)\) is an affine open inside \(O_i\)), which
  equals the restriction of \(t_U\) (compatibility of \(s, t\) at \(D(r) \subseteq U\)). By
  Theorem~\ref{pa-thm:picEt_eq_of_basic_eq}, \(s_U = t_U\). As \(U\) was an arbitrary affine open,
  \(s = t\).
\end{proof}

\section{Rigidification along a section of the projection}
\label{pa-sec:rigidification}

Let \(C\) be an object of \(\Over(\Spec k)\) (the curve). For a test object \(T\) in
\(\Over(\Spec k)\), a \(T\)-point \(\sigma \colon T \to C\) determines a section of the
projection \(\mathrm{pr} = \mathrm{snd} \colon C \otimes T \to T\); rigidifying a \v Cech Picard
class along that section is the cocycle recast of Kleiman's \(g\)-rigidification. These are the
inputs to the \((C2)\) comparison of the next section.

\begin{definition}[The section of the projection attached to a point]
  \label{pa-def:sectionOfPoint}
  \lean{AlgebraicGeometry.Over.sectionOfPoint, AlgebraicGeometry.Over.sectionOfPoint_fst,
    AlgebraicGeometry.Over.sectionOfPoint_snd, AlgebraicGeometry.Over.sectionOfPoint_naturality}
  \leanok
  \source{kleiman-picard}

  \dcref{2.5,custom}
  For a \(T\)-point \(\sigma \colon T \to C\), the \emph{section of the projection}
  \(g_\sigma \colon T \to C \otimes T\) is the map \(\langle \sigma, \id_T \rangle\) into the
  product, characterised by
  \[
    \mathrm{fst} \circ g_\sigma = \sigma, \qquad \mathrm{snd} \circ g_\sigma = \id_T .
  \]
  It is natural in the test object: for \(h \colon T' \to T\), \(h \circ g_\sigma = g_{h \circ
  \sigma} \circ (C \lhd h)\). This is Kleiman's \(g_T\) (\emph{The Picard Scheme}, proof of
  Theorem 2.5).
\end{definition}
\begin{proof}
  \leanok
  By the universal property of the cartesian product, \(T\)-points of \(C\) and sections of the
  projection \(\mathrm{snd}\) are the same datum: \(\langle \sigma, \id_T\rangle\) is the unique
  map with the two stated composites. Naturality is checked on the two projections, using
  \(\mathrm{fst} \circ (C \lhd h) = h \circ \mathrm{fst}\) and \(\mathrm{snd} \circ (C \lhd h) =
  \mathrm{snd}\).
\end{proof}

\begin{definition}[Rigidified \v Cech Picard class]
  \label{pa-def:isRigidified}
  \lean{AlgebraicGeometry.Over.IsRigidified}
  \leanok
  \source{kleiman-picard}

  \uses{pa-def:sectionOfPoint}
  \dcref{custom}
  A \v Cech Picard class \(L\) on the product \(C \otimes T\) is \emph{rigidified} along a point
  \(\sigma \colon T \to C\) when its pullback along the section of the projection is trivial,
  \[
    g_\sigma^{*} L = 1 .
  \]
  This is the cocycle form of Kleiman's \(g\)-rigidification (\(\mathrm{df:rgd}\)): on the
  cocycle model an isomorphism class carries no datum, so ``a trivialisation \(u \colon
  \mathcal O_T \xrightarrow{\ \sim\ } g_\sigma^{*} L\) exists'' is exactly \(g_\sigma^{*} L = 1\).
\end{definition}

\begin{theorem}[Rigidity of automorphisms of a rigidified pair]
  \label{pa-thm:unitsAppTop_sectionOfPoint_bijective}
  \lean{AlgebraicGeometry.Over.unitsAppTop_sectionOfPoint_bijective,
    AlgebraicGeometry.Over.isIso_appTop_sectionOfPoint}
  \leanok
  \source{kleiman-picard}

  \uses{pa-def:sectionOfPoint, pa-thm:isIso_appLE_snd}
  \dcref{custom}
  Let the curve \(C\) be proper, geometrically irreducible and geometrically reduced. For every
  point \(\sigma \colon \Spec A \to C\), pullback of global unit sections along the section of the
  projection,
  \[
    g_\sigma^{\sharp} \colon \Gamma(\Spec A, \top)^{\times}
      \;\longrightarrow\; \Gamma(X_A, \top)^{\times},
  \]
  is a bijection. In particular, a global unit of the curve \(X_A\) whose restriction along the
  section is \(1\) is itself \(1\). This is the cocycle form of Kleiman's automorphism rigidity: an
  automorphism of a rigidified pair is trivial.
\end{theorem}
\begin{proof}
  \leanok
  Because \(C\) is proper, geometrically irreducible and geometrically reduced, the degree-zero
  cohomology of the relative curve is the test ring itself: pullback of global sections along the
  projection \(\mathrm{pr} = \mathrm{snd} \colon X_A \to \Spec A\) is a ring isomorphism
  \(\Gamma(\Spec A, \top) \xrightarrow{\sim} \Gamma(X_A, \top)\)
  (Theorem~\ref{pa-thm:isIso_appLE_snd}, the case \(H^0(X_A, \struct{}) = A\); this is Kleiman's
  hypothesis \(\struct S \xrightarrow{\sim} f_{*}\struct X\) holding universally in the test ring).
  The section retracts the projection, \(\mathrm{pr} \circ g_\sigma = \id\)
  (Definition~\ref{pa-def:sectionOfPoint}), and pullback of sections is contravariant, so on global
  sections \(g_\sigma^{\sharp} \circ \mathrm{pr}^{\sharp} = (\mathrm{pr} \circ g_\sigma)^{\sharp} =
  \id\). As \(\mathrm{pr}^{\sharp}\) is an isomorphism, \(g_\sigma^{\sharp}\) is its inverse, hence
  an isomorphism, hence bijective on the unit groups. Kleiman's argument is exactly this: an
  automorphism \(v\) of a pair \((L, u)\) rigidified along \(\sigma\) satisfies
  \(g_\sigma^{\sharp} v = 1\), while \(v\) lies in \(H^0(\struct{X_A}) = H^0(\struct{\Spec A})\); the
  displayed isomorphism forces \(v = 1\).
\end{proof}

\begin{lemma}[The retraction identity]
  \label{pa-lem:cechPicMap_sectionOfPoint_snd}
  \lean{AlgebraicGeometry.Over.cechPicMap_sectionOfPoint_snd}
  \leanok
  \source{kleiman-picard}

  \uses{pa-def:sectionOfPoint}
  \dcref{custom}
  For a point \(\sigma \colon T \to C\), pullback of \v Cech Picard classes along the section
  of the projection retracts pullback along the projection: for every \(N \in \Pic(T)\),
  \[
    g_\sigma^{*}\,\mathrm{snd}^{*} N = N .
  \]
  This is \(g_T^{*} f_T^{*} = 1\) in Kleiman's notation.
\end{lemma}
\begin{proof}
  \leanok
  Pullback of \v Cech classes is contravariantly functorial, and \(\mathrm{snd} \circ
  g_\sigma = \id_T\) (Definition~\ref{pa-def:sectionOfPoint}), so the composite pullback is the
  pullback along the identity, which is the identity.
\end{proof}

\begin{lemma}[Stability of rigidification under restriction]
  \label{pa-lem:isRigidified_cechPicMap}
  \lean{AlgebraicGeometry.Over.isRigidified_one, AlgebraicGeometry.Over.IsRigidified.cechPicMap}
  \leanok
  \dcref{custom}

  \uses{pa-def:isRigidified, pa-def:sectionOfPoint}
  The trivial class is rigidified along every point. Moreover, for \(h \colon T' \to T\) and a
  class \(L\) on \(C \otimes T\) rigidified along \(\sigma \colon T \to C\), the pullback
  \((C \lhd h)^{*} L\) is rigidified along the restricted point \(h \circ \sigma\).
\end{lemma}
\begin{proof}
  \leanok
  The first claim is \(g_\sigma^{*} 1 = 1\). For the second, the naturality of the section in
  the test object (Definition~\ref{pa-def:sectionOfPoint}) gives \(h \circ g_\sigma =
  g_{h \circ \sigma} \circ (C \lhd h)\), so
  \[
    g_{h \circ \sigma}^{*} (C \lhd h)^{*} L
      = \bigl((C \lhd h) \circ g_{h \circ \sigma}\bigr)^{*} L
      = (g_\sigma \circ h)^{*} L
      = h^{*} \bigl(g_\sigma^{*} L\bigr) = h^{*} 1 = 1 . \qedhere
  \]
\end{proof}

\begin{theorem}[Kleiman's identification, injectivity: rigidified representatives are unique]
  \label{pa-thm:isRigidified_eq_of_relPicMk_eq}
  \lean{AlgebraicGeometry.Over.IsRigidified.eq_of_relPicMk_eq,
    AlgebraicGeometry.Over.IsRigidified.eq_one_of_relPicMk_eq_one}
  \leanok
  \source{kleiman-picard}

  \uses{pa-def:isRigidified, pa-lem:cechPicMap_sectionOfPoint_snd, pa-def:relPic}
  \dcref{custom}
  Two \v Cech Picard classes on \(C \otimes T\), both rigidified along a point
  \(\sigma \colon T \to C\), which represent the same relative Picard class, are equal on the
  nose. In particular a rigidified representative of the identity relative class is \(1\).
  This is the injectivity half of Kleiman's Lemma \(\mathrm{lm{:}idn}\) (rigidified pairs are
  carried isomorphically onto \(\Pic_{X/S}(T)\)), in cocycle form.
\end{theorem}
\begin{proof}
  \leanok
  Let \(L, L'\) be the two classes. By the exact coset relation
  (Definition~\ref{pa-def:relPic}) the ratio is pulled back from the base: \(L/L' =
  \mathrm{snd}^{*} N\) for some \(N \in \Pic(T)\). Apply \(g_\sigma^{*}\): by the retraction
  identity (Lemma~\ref{pa-lem:cechPicMap_sectionOfPoint_snd}),
  \[
    N = g_\sigma^{*}\,\mathrm{snd}^{*} N = g_\sigma^{*}(L/L')
      = g_\sigma^{*}L \,/\, g_\sigma^{*}L' = 1/1 = 1 ,
  \]
  using both rigidifications. Hence \(L/L' = \mathrm{snd}^{*}1 = 1\), i.e.\ \(L = L'\). The
  special case takes \(L' = 1\), which is rigidified.
\end{proof}

\begin{theorem}[Kleiman's identification, surjectivity: existence of a rigidified representative]
  \label{pa-thm:relPic_exists_isRigidified_rep}
  \lean{AlgebraicGeometry.relPic.exists_isRigidified_rep}
  \leanok
  \source{kleiman-picard}

  \uses{pa-def:isRigidified, pa-def:relPic, pa-lem:cechPicMap_sectionOfPoint_snd}
  \dcref{custom}
  Every relative Picard class on \(C \otimes T\) has a representative rigidified along any
  given point \(\sigma \colon T \to C\). This is the surjectivity half of Kleiman's Lemma
  \(\mathrm{lm{:}idn}\), in cocycle form.
\end{theorem}
\begin{proof}
  \leanok
  Let \(M\) be any \v Cech representative of the class (the projection to the quotient is
  surjective, Definition~\ref{pa-def:relPic}). Kleiman's identity-rigidification twist
  \[
    L := M \cdot \bigl(\mathrm{snd}^{*} g_\sigma^{*} M\bigr)^{-1}
  \]
  (\(\mathcal L := \mathcal M \otimes (f_T^{*} g_T^{*} \mathcal M)^{-1}\)) works. The twist factor is pulled
  back from the base, so \(L\) represents the same relative class. And by the retraction
  identity (Lemma~\ref{pa-lem:cechPicMap_sectionOfPoint_snd}),
  \(g_\sigma^{*} L = g_\sigma^{*} M \cdot (g_\sigma^{*} M)^{-1} = 1\), so \(L\) is
  rigidified along \(\sigma\).
\end{proof}

\begin{theorem}[The rigidified transport keystone]
  \label{pa-thm:isRigidified_cechPicMap_congr}
  \lean{AlgebraicGeometry.Over.IsRigidified.cechPicMap_congr}
  \leanok
  \source{kleiman-picard}

  \uses{pa-def:isRigidified, pa-thm:relPicAlgMap_congr, pa-thm:isRigidified_eq_of_relPicMk_eq,
    pa-lem:isRigidified_cechPicMap, pa-def:descentClasses, pa-def:relPicMap, pa-def:overSpecMap}
  \dcref{custom}
  Let \(A\) be a \(k\)-algebra, \(E\) a presented \'etale cover of \(A\), \(R\) an
  \(A\)-algebra, \(\sigma \colon \Spec A \to C\) a point, and \(L\) a \v Cech Picard class on
  \(C \otimes \Spec B_E\) such that the relative class of \(L\) is a descent class on \(E\)
  and \(L\) is rigidified along the point of \(\Spec B_E\) induced by \(\sigma\). Then for
  \emph{any} two \(A\)-algebra maps \(j_1, j_2 \colon B_E \to R\), the two pullbacks of \(L\)
  agree \emph{on the nose}:
  \[
    (C \lhd \Spec j_1)^{*} L = (C \lhd \Spec j_2)^{*} L \quad\text{in } \Pic(C \otimes \Spec R).
  \]
\end{theorem}
\begin{proof}
  \leanok
  The plus-construction keystone \(\mathrm{relPicAlgMap\_congr}\)
  (Theorem~\ref{pa-thm:relPicAlgMap_congr}) gives the equality of the two pullbacks as
  \emph{relative} Picard classes, since the class of \(L\) is a descent class and the
  \(j_i\) are \(A\)-algebra maps. Both pullbacks are rigidified along the point of
  \(\Spec R\) induced by \(\sigma\): the \(j_i\) are \(A\)-algebra maps, so the induced point
  of \(\Spec R\) restricts, along \(\Spec j_i\), the induced point of \(\Spec B_E\), and
  rigidification is stable under restriction (Lemma~\ref{pa-lem:isRigidified_cechPicMap}).
  Kleiman's injectivity (Theorem~\ref{pa-thm:isRigidified_eq_of_relPicMk_eq}) upgrades the
  coset equality of two rigidified classes to an equality of \v Cech classes.
\end{proof}

\begin{theorem}[The descent equation of a rigidified representative holds on the nose]
  \label{pa-thm:isRigidified_cechPicMap_doubleInl_eq_doubleInr}
  \lean{AlgebraicGeometry.Over.IsRigidified.cechPicMap_doubleInl_eq_doubleInr}
  \leanok
  \source{kleiman-picard}

  \uses{pa-thm:isRigidified_cechPicMap_congr, pa-def:doubleInl_doubleInr}
  \dcref{custom}
  In the situation of Theorem~\ref{pa-thm:isRigidified_cechPicMap_congr}, the two pullbacks of
  the rigidified representative \(L\) to the double cover \(\Spec(B_E \otimes_A B_E)\), along
  the two coprojections, are \emph{equal \v Cech Picard classes} — not merely equal relative
  classes. This is Kleiman's rigidified pair \((\mathcal L', u')\) having a canonical comparison on
  the double cover, the input to the descent-datum construction in the proof of his
  comparison theorem.
\end{theorem}
\begin{proof}
  \leanok
  The two coprojections \(B_E \rightrightarrows B_E \otimes_A B_E\)
  (Definition~\ref{pa-def:doubleInl_doubleInr}) are \(A\)-algebra maps out of the cover carrier
  into the common \(A\)-algebra \(R := B_E \otimes_A B_E\); apply the transport keystone
  (Theorem~\ref{pa-thm:isRigidified_cechPicMap_congr}) to them.
\end{proof}

\section{The \texorpdfstring{\(\sigma\)}{sigma}-normalized comparison cochain}
\label{pa-sec:normalizedComparison}

We now fix a \(k\)-algebra \(A\) with a further \(A\)-algebra \(B\) (the carrier of an \'etale
cover in the intended application) and a curve point \(\sigma \colon \Spec A \to C\).
Write \(X_R := (C \otimes \Spec R)_{\mathrm{left}}\) for the curve over \(\Spec R\). Let
\(q_1, q_2 \colon \Spec(B \otimes_A B) \to \Spec B\) be the coprojections
(Definition~\ref{pa-def:tensorInl_tensorInr}), \(u_1, u_2 \colon X_{B \otimes_A B} \to X_B\) their
whiskerings \(C \lhd q_i\), and \(f_{12}, f_{13}, f_{23}\), \(w_{ij} = C \lhd f_{ij}\) the cofaces
of Definition~\ref{pa-def:tensorFaces}. Write \(s_B, s_q\) for the sections of the projection
(Definition~\ref{pa-def:sectionOfPoint}) attached to the base changes of \(\sigma\) to \(B\) and
\(B \otimes_A B\), and \(p_2 \colon X_{B \otimes_A B} \to \Spec(B \otimes_A B)\) for the
projection. Throughout, \(s^{\sharp}\) denotes pullback of unit cochains as in
Section~\ref{pa-sec:cochainPullback}.

This section packages, and constructs with \emph{no geometric hypotheses}, a cochain-level
realisation of the descent equation \(u_1^{*} L = u_2^{*} L\), normalized along \(\sigma\).

\begin{lemma}[The section squares]
  \label{pa-lem:sectionSquares}
  \lean{AlgebraicGeometry.Over.overSpecMap_comp_point,
    AlgebraicGeometry.Over.sectionOfPoint_left_comp_whiskerLeft,
    AlgebraicGeometry.Over.sectionOfPoint_left_comp_snd,
    AlgebraicGeometry.Over.sectionSq_comp_inl, AlgebraicGeometry.Over.sectionSq_comp_inr,
    AlgebraicGeometry.Over.sectionCb_comp_face₁₂, AlgebraicGeometry.Over.sectionCb_comp_face₁₃,
    AlgebraicGeometry.Over.sectionCb_comp_face₂₃}
  \leanok
  \dcref{custom}

  \uses{pa-def:sectionOfPoint, pa-def:overSpecMap, pa-def:tensorInl_tensorInr, pa-def:tensorFaces}
  The section of the projection is natural across the base-change coprojections and cofaces: as
  maps of curve products,
  \[
    s_q \circ u_1 = q_1 \circ s_B, \quad s_q \circ u_2 = q_2 \circ s_B, \quad
    s_{\mathrm{cb}} \circ w_{ij} = f_{ij} \circ s_q \ (ij \in \{12,13,23\}),
  \]
  and \(s_q\) retracts the projection, \(s_q \circ p_2 = \id\).
\end{lemma}
\begin{proof}
  \leanok
  Each square is the naturality of the section (Definition~\ref{pa-def:sectionOfPoint}) applied to
  the relevant \(A\)-algebra map: the coprojections \(q_i\) and cofaces \(f_{ij}\) are
  \(\Spec\) of \(A\)-algebra homomorphisms out of \(B\) respectively \(B \otimes_A B\)
  (Definitions~\ref{pa-def:tensorInl_tensorInr}, \ref{pa-def:tensorFaces}), and restricting the
  base-changed point along such a map lands at the base-changed point over the target algebra;
  whiskering with \(C\) and passing to the section turns this point-compatibility into the stated
  square. The retraction is \(\mathrm{snd} \circ g_\sigma = \id\) of
  Definition~\ref{pa-def:sectionOfPoint}, whiskered to the curve product.
\end{proof}

\begin{lemma}[Pullback of a one-sided trivializing relation]
  \label{pa-lem:unitsAppLE_trivializing_rel}
  \lean{AlgebraicGeometry.Scheme.Hom.unitsAppLE_trivializing_rel}
  \leanok
  \dcref{custom}

  Let \(\varphi \colon X \to Y\) and \(r \colon Y \to Z\). Suppose a unit \(0\)-cochain \(a\) on a
  pointed cover of \(Y\) trivializes the pullback along \(r\) of a pair-indexed unit family \(E\):
  for all \(y, y'\), on the overlap,
  \[
    a_y \cdot \bigl(r^{\sharp} E\bigr)_{y y'} = a_{y'} .
  \]
  Then for every \(x, x' \in X\) and every open \(O \subseteq \varphi^{-1}(\cdots)\) inside the
  relevant overlap pullback, the pulled-back cochain \(\varphi^{\sharp} a\) trivializes the family
  pulled back along the composite:
  \[
    \varphi^{\sharp}_O\, a_{\varphi x} \cdot \bigl((r\varphi)^{\sharp}_O E\bigr)_{x x'}
      = \varphi^{\sharp}_O\, a_{\varphi x'} .
  \]
  This is the one-sided companion of Lemma~\ref{pa-lem:unitsAppLE_coboundary_rel}.
\end{lemma}
\begin{proof}
  \leanok
  Apply the pullback homomorphism \(\varphi^{\sharp}_O\) to the given relation at \(\varphi x,
  \varphi x'\). Being a group homomorphism on units it distributes over the product; by
  functoriality \(\varphi^{\sharp}_O(r^{\sharp} E) = (r\varphi)^{\sharp}_O E\) merges the pullback,
  while \(\varphi^{\sharp}_O a\) is left unchanged.
\end{proof}

\begin{lemma}[Compatibility of a ratio cochain]
  \label{pa-lem:unitsAppLE_trivializing_ratio_compat}
  \lean{AlgebraicGeometry.Scheme.Hom.unitsAppLE_trivializing_ratio_compat}
  \leanok
  \dcref{custom}

  \uses{pa-lem:unitsAppLE_coboundary_rel, pa-lem:unitsAppLE_trivializing_rel, pa-lem:unitsAppLE_congr_hom}
  Let \(\varphi_t \colon X \to Y_c\), \(\varphi_1, \varphi_2 \colon X \to Y_a\), \(r_1, r_2
  \colon Y_c \to Z\), \(s_1 \colon Y_a \to Z\). Suppose a \(0\)-cochain \(c\) on \(Y_c\) cobounds
  the comparison of a pair-indexed family \(E\) pulled back along \(r_1, r_2\)
  (\ref{pa-lem:unitsAppLE_coboundary_rel}), and a \(0\)-cochain \(a\) on \(Y_a\) trivializes \(E\)
  pulled back along \(s_1\) (\ref{pa-lem:unitsAppLE_trivializing_rel}). If the pulled-back composites
  agree, \(\varphi_t \circ r_1 = \varphi_1 \circ s_1\) and \(\varphi_t \circ r_2 = \varphi_2
  \circ s_1\), then the ratio cochain
  \[
    \varphi_t^{\sharp} c \cdot \frac{\varphi_2^{\sharp} a}{\varphi_1^{\sharp} a}
  \]
  on \(X\) is \v Cech-compatible: its values at \(x\) and \(x'\) agree on overlaps.
\end{lemma}
\begin{proof}
  \leanok
  Pull the coboundary relation for \(c\) back along \(\varphi_t\)
  (\ref{pa-lem:unitsAppLE_coboundary_rel}) and the trivializing relation for \(a\) back along
  \(\varphi_1\) and \(\varphi_2\) (\ref{pa-lem:unitsAppLE_trivializing_rel}). The hypotheses
  \(\varphi_t r_i = \varphi_i s_1\) identify the two families through
  Lemma~\ref{pa-lem:unitsAppLE_congr_hom}; combining the pulled-back relations, the comparison of
  the ratio cochain at \(x\) and \(x'\) reduces to an identity of units, giving compatibility.
\end{proof}

\begin{definition}[The normalization cover]
  \label{pa-def:normalizationCover}
  \lean{AlgebraicGeometry.Over.normalizationCover,
    AlgebraicGeometry.Over.normalizationCover_le_section,
    AlgebraicGeometry.Over.normalizationCover_le_inl,
    AlgebraicGeometry.Over.normalizationCover_le_inr}
  \leanok
  \dcref{custom}

  \uses{pa-def:sectionOfPoint}
  For a pointed cover \(\mathcal N\) of the curve \(X_B\) and a pointed cover \(\mathcal W\) of
  the curve \(X_{B \otimes_A B}\), the \emph{normalization cover} of \(\Spec(B \otimes_A B)\) is
  \[
    \mathcal W \text{ pulled back along } s_q, \ \text{ intersected with the two coprojection
    pullbacks } q_1^{*}, q_2^{*} \text{ of } \mathcal N \text{ pulled back along } s_B .
  \]
  It is the cover on which the \(\sigma\)-normalization below is stated; it lies below \(s_q^{-1}
  \mathcal W\) and below each \(q_i^{-1}(s_B^{-1}\mathcal N)\).
\end{definition}

\begin{definition}[The \(\sigma\)-normalized comparison datum]
  \label{pa-def:normalizedCechComparison}
  \lean{AlgebraicGeometry.NormalizedCechComparison}
  \leanok
  \dcref{custom}

  \uses{pa-def:sectionOfPoint, pa-def:isRigidified, pa-def:normalizationCover}
  Fix a unit \v Cech cocycle \(\gamma\) on a pointed cover \(\mathcal N\) of \(X_B\), representing
  a class \(L = [\gamma]\). A \emph{\(\sigma\)-normalized comparison} of \(\gamma\) packages:
  \begin{itemize}
    \item a trivializing \(0\)-cochain \(\rho\) of the section-pullback \(s_B^{\sharp}\gamma\) on
      the cover \(s_B^{-1}\mathcal N\) of \(\Spec B\): \(\rho_b \cdot (s_B^{\sharp}\gamma)_{b b'} =
      \rho_{b'}\) on overlaps — the cochain shadow of rigidification (\ref{pa-def:isRigidified})
      along the section over \(B\);
    \item a pointed cover \(\mathcal W\) of \(X_{B \otimes_A B}\) refining both coprojection
      pullbacks of \(\mathcal N\), and a comparison \(0\)-cochain \(\theta\) on it, with
    \item \textbf{(N1)} \(\theta_x \cdot (u_1^{\sharp}\gamma)_{x y} = (u_2^{\sharp}\gamma)_{x y}
      \cdot \theta_y\) on overlaps — \(\theta\) cobounds the two coprojection pullbacks of
      \(\gamma\), realising \(u_1^{*} L = u_2^{*} L\) at the cochain level; and
    \item \textbf{(N2)} \((s_q^{\sharp}\theta)_y \cdot (q_2^{\sharp}\rho)_y = (q_1^{\sharp}\rho)_y\)
      on the normalization cover (\ref{pa-def:normalizationCover}) — the \(\sigma\)-restriction of
      \(\theta\) is the \(\rho\)-comparison \(q_1^{\sharp}\rho / q_2^{\sharp}\rho\).
  \end{itemize}
  The normalization is \(\rho\)-relative: rigidification makes \(s_B^{\sharp}\gamma\) trivial only
  as a class, giving the trivializing cochain \(\rho\), so ``\(\sigma\)-restriction \(=1\)'' is
  replaced by (N2). Every step is stated pointwise via unit pullback and is blind to
  disconnectedness of \(\Spec(B \otimes_A B)\).
\end{definition}

\begin{theorem}[The \(\sigma\)-defect glues to a global unit]
  \label{pa-thm:exists_glued_sigmaDefect}
  \lean{AlgebraicGeometry.Over.sigmaDefectCochain,
    AlgebraicGeometry.Over.exists_glued_sigmaDefect}
  \leanok
  \dcref{custom}

  \uses{pa-def:normalizationCover, pa-lem:sectionSquares,
    pa-lem:unitsAppLE_trivializing_ratio_compat, pa-thm:exists_global_unit_of_compatible}
  Let \(\gamma\) be a unit cocycle on \(\mathcal N\), let \(\theta_0\) be a comparison cochain on
  a refinement \(\mathcal W\) satisfying (N1), and let \(\rho\) be a trivializing cochain of
  \(s_B^{\sharp}\gamma\). The \emph{\(\sigma\)-defect cochain}
  \[
    s_q^{\sharp}\theta_0 \cdot \frac{q_2^{\sharp}\rho}{q_1^{\sharp}\rho}
  \]
  on the normalization cover is \v Cech-compatible, hence glues to a global unit \(\delta \in
  \Gamma(\Spec(B \otimes_A B), \top)^{\times}\).
\end{theorem}
\begin{proof}
  \leanok
  Compatibility is the ratio-compatibility Lemma~\ref{pa-lem:unitsAppLE_trivializing_ratio_compat}
  applied with \(\varphi_t = s_q\), \(\varphi_i = q_i\), \(r_i = u_i\), \(s_1 = s_B\): (N1) is
  the coboundary relation for \(\theta_0\), the relation for \(\rho\) is the one-sided
  trivialization, and the composite coincidences \(s_q \circ u_i = q_i \circ s_B\) are the
  section squares (\ref{pa-lem:sectionSquares}). A compatible unit family on a pointed cover glues to
  a global unit by the sheaf axiom (\ref{pa-thm:exists_global_unit_of_compatible}).
\end{proof}

\begin{theorem}[Existence of the \(\sigma\)-normalized comparison]
  \label{pa-thm:exists_normalizedCechComparison}
  \lean{AlgebraicGeometry.Over.correctedComparison,
    AlgebraicGeometry.Over.exists_normalizedCechComparison}
  \leanok
  \dcref{custom}

  \uses{pa-def:isRigidified, pa-def:normalizedCechComparison, pa-thm:exists_glued_sigmaDefect,
    pa-lem:sectionSquares, pa-cor:cechPic_mk_eq_one,
    pa-thm:isRigidified_cechPicMap_doubleInl_eq_doubleInr}
  Let \(\gamma\) be a unit cocycle on a pointed cover \(\mathcal N\) of \(X_B\) whose class \(L =
  [\gamma]\) is rigidified along the section over \(B\) (\ref{pa-def:isRigidified}) and satisfies the
  descent equation \(u_1^{*} L = u_2^{*} L\) — as holds for any rigidified representative of a
  descent class, by Theorem~\ref{pa-thm:isRigidified_cechPicMap_doubleInl_eq_doubleInr}. Then a
  \(\sigma\)-normalized comparison of \(\gamma\) exists. No properness, geometric
  irreducibility or reducedness, or faithful flatness is required.
\end{theorem}
\begin{proof}
  \leanok
  Rigidification gives \(g_{\sigma_B}^{*} L = 1\); on cocycles (\ref{pa-cor:cechPic_mk_eq_one}) this
  means the section-pullback \(s_B^{\sharp}\gamma\) is a coboundary, so a trivializing cochain
  \(\rho\) exists on \(s_B^{-1}\mathcal N\) itself (no refinement). The descent equation
  \(u_1^{*} L = u_2^{*} L\) means the two coprojection pullbacks of \(\gamma\) cobound on a common
  refinement \(\mathcal W\), yielding a comparison cochain \(\theta_0\) with (N1). By
  Theorem~\ref{pa-thm:exists_glued_sigmaDefect} the \(\sigma\)-defect of \(\theta_0\) glues to a
  global unit \(\delta\) of \(\Spec(B \otimes_A B)\). Correct \(\theta_0\) by dividing out the
  projection pullback \(p_2^{\sharp}\delta\): this preserves (N1) (a global factor cancels from
  both sides of the coboundary relation), and, because \(s_q\) retracts \(p_2\)
  (\ref{pa-lem:sectionSquares}), it makes the \(\sigma\)-restriction of the corrected cochain equal
  to \(q_1^{\sharp}\rho / q_2^{\sharp}\rho\) on the nose — that is (N2). The corrected data is the
  required datum.
\end{proof}

\section{Coherence and uniqueness of the
  \texorpdfstring{\(\sigma\)}{sigma}-normalized comparison}
\label{pa-sec:comparisonCoherence}

Fix a unit \v Cech cocycle \(\gamma\) on a pointed cover \(\mathcal N\) of \(X_B\) and a
\(\sigma\)-normalized comparison of \(\gamma\)
(Definition~\ref{pa-def:normalizedCechComparison}), with cover \(\mathcal W\) of \(X_{B \otimes_A B}\)
and comparison cochain \(\theta\). We show that \(\theta\) satisfies the Amitsur cocycle identity on
the triple base \emph{on the nose}, is unique, and is trivial on the diagonal --- the three
inputs the per-piece descent of the next section runs on. All three are the cochain form of
Kleiman's automorphism-rigidity argument (Lemma~\ref{pa-thm:unitsAppTop_sectionOfPoint_bijective}):
glue a defect to a global unit, kill its \(\sigma\)-restriction by the normalization, and apply
rigidity at the relevant affine test.

Write \(X_{B \otimes_A (B \otimes_A B)}\) for the curve over the triple base. Its three
tensor-slot \emph{insertions} \(v_1, v_2, v_3 \colon X_{B \otimes_A (B \otimes_A B)} \to X_B\) are
the whiskerings \(C \lhd (\Spec\ \text{of})\) of the \(A\)-algebra maps \(x \mapsto x \otimes 1
\otimes 1\), \(x \mapsto 1 \otimes x \otimes 1\), \(x \mapsto 1 \otimes 1 \otimes x\).

\begin{definition}[The Amitsur product cover and its insertions]
  \label{pa-def:amitsurProductCover}
  \lean{AlgebraicGeometry.Over.amitsurProductCover, AlgebraicGeometry.Over.amitsurInsertion₁,
    AlgebraicGeometry.Over.amitsurInsertion₂, AlgebraicGeometry.Over.amitsurInsertion₃}
  \leanok
  \dcref{custom}

  \uses{pa-def:amitsurCover, pa-def:tensorInl_tensorInr, pa-def:tensorFaces}
  For a pointed cover \(\mathcal W\) of the curve \(X_{B \otimes_A B}\), the \emph{Amitsur product
  cover} of \(X_{B \otimes_A (B \otimes_A B)}\) is the common refinement of the three coface
  pullbacks,
  \[
    w_{23}^{-1}\mathcal W \ \sqcap\ w_{12}^{-1}\mathcal W \ \sqcap\ w_{13}^{-1}\mathcal W ,
  \]
  the curve-product mirror of the Amitsur cover (Definition~\ref{pa-def:amitsurCover}). The three
  insertions \(v_1, v_2, v_3\) are the pairwise-coincident values of the six coface--coprojection
  composites \(w_{ij} \circ u_k \colon X_{B \otimes_A (B \otimes_A B)} \to X_B\):
  \[
    v_1 = w_{12} \circ u_1 = w_{13} \circ u_1, \quad
    v_2 = w_{23} \circ u_1 = w_{12} \circ u_2, \quad
    v_3 = w_{23} \circ u_2 = w_{13} \circ u_2 .
  \]
\end{definition}
\begin{proof}
  \leanok
  The cover is a finite intersection of pullbacks of the pointed cover \(\mathcal W\), hence a
  pointed cover. Each of the three displayed pairs of composites coincides by the simplicial
  identities of the tensor tower relating the cofaces \(f_{ij}\) (Definition~\ref{pa-def:tensorFaces})
  to the coprojections \(x \mapsto x \otimes 1\), \(x \mapsto 1 \otimes x\)
  (Definition~\ref{pa-def:tensorInl_tensorInr}): whiskering to the curve preserves those equalities,
  and each \(v_i\) is the whiskered \(\Spec\) of one tensor-slot insertion.
\end{proof}

\begin{theorem}[The Amitsur defect glues to a global unit]
  \label{pa-thm:exists_glued_productDefect}
  \lean{AlgebraicGeometry.Over.productDefectCochain,
    AlgebraicGeometry.Over.exists_glued_productDefect}
  \leanok
  \dcref{custom}

  \uses{pa-def:amitsurProductCover, pa-def:normalizedCechComparison, pa-thm:exists_global_unit_of_compatible}
  Let \(\gamma\) be a unit cocycle on \(\mathcal N\) and \(\theta_0\) a comparison cochain on a
  refinement \(\mathcal W\) satisfying \textbf{(N1)}, i.e.\ \(\theta_0\) cobounds the two
  coprojection pullbacks of \(\gamma\). The \emph{Amitsur defect}
  \[
    \varepsilon \;=\; \frac{(w_{23}^{\sharp}\theta_0)\,(w_{12}^{\sharp}\theta_0)}
      {w_{13}^{\sharp}\theta_0}
  \]
  --- a unit \(0\)-cochain on the Amitsur product cover --- is \v Cech-compatible, hence glues to a
  global unit \(\bar\varepsilon \in \Gamma(X_{B \otimes_A (B \otimes_A B)}, \top)^{\times}\). No
  geometric hypotheses are used.
\end{theorem}
\begin{proof}
  \leanok
  Relation \textbf{(N1)} reads \(\theta_{0,x} \cdot (u_1^{\sharp}\gamma)_{xy} =
  (u_2^{\sharp}\gamma)_{xy} \cdot \theta_{0,y}\) on overlaps. Pull it back along each of the three
  cofaces \(w_{23}, w_{12}, w_{13}\); by Definition~\ref{pa-def:amitsurProductCover} the
  coface--coprojection composites appearing are exactly the three insertions, so the pulled-back
  relations express the overlap discrepancy of each factor \(w_{ij}^{\sharp}\theta_0\) through the
  cocycle \(\gamma\) evaluated at the insertions \(v_1, v_2, v_3\). The three insertion values are
  common to the numerator and denominator of \(\varepsilon\); multiplying the three relations
  telescopically, the \(\gamma\)-terms cancel and the value of \(\varepsilon\) at \(z\) agrees with
  its value at \(z'\) on the overlap. A compatible unit family on a pointed cover glues to a global
  unit by the sheaf axiom (Theorem~\ref{pa-thm:exists_global_unit_of_compatible}).
\end{proof}

\begin{theorem}[Coherence of the \(\sigma\)-normalized comparison]
  \label{pa-thm:normalizedComparison_coherent}
  \lean{AlgebraicGeometry.NormalizedCechComparison.coherent}
  \leanok
  \source{kleiman-picard}

  \uses{pa-def:normalizedCechComparison, pa-thm:exists_glued_productDefect, pa-lem:sectionSquares,
    pa-thm:unitsAppTop_sectionOfPoint_bijective, pa-def:module_descent_cocycle}
  \dcref{custom}
  Let the curve \(C\) be proper, geometrically irreducible and geometrically reduced. Then the
  comparison cochain \(\theta\) of a \(\sigma\)-normalized comparison of \(\gamma\) satisfies the
  Amitsur cocycle identity on the nose:
  \[
    (w_{23}^{\sharp}\theta)\,(w_{12}^{\sharp}\theta) \;=\; w_{13}^{\sharp}\theta ,
  \]
  in the face convention of the descent \(1\)-cocycle (Definition~\ref{pa-def:module_descent_cocycle}).
  No faithful flatness is required.
\end{theorem}
\begin{proof}
  \leanok
  By \textbf{(N1)} the Amitsur defect \(\varepsilon\) of \(\theta\) glues to a global unit
  \(\bar\varepsilon\) of \(X_{B \otimes_A (B \otimes_A B)}\)
  (Theorem~\ref{pa-thm:exists_glued_productDefect}). We claim its \(\sigma\)-restriction is \(1\).
  Pull the glued identity back along the section \(s_{\mathrm{cb}}\) over the triple base: by the
  section squares \(s_{\mathrm{cb}} \circ w_{ij} = f_{ij} \circ s_q\)
  (Lemma~\ref{pa-lem:sectionSquares}) the three coface pullbacks of \(\theta\) become the coface
  pullbacks \(f_{ij}^{\sharp}\) of the \(\sigma\)-restriction of \(\theta\). Now \textbf{(N2)}
  rewrites the \(\sigma\)-restriction of \(\theta\) as the \(\rho\)-comparison \(q_1^{\sharp}\rho /
  q_2^{\sharp}\rho\); substituting, the three coface pullbacks become three ratios of
  \(\rho\)-terms whose endpoints match up through the simplicial coincidences \(f_{12} \circ q_2 =
  f_{23} \circ q_1\), \(f_{13} \circ q_1 = f_{12} \circ q_1\), \(f_{13} \circ q_2 = f_{23} \circ
  q_2\). These telescope to \(1\); hence the \(\sigma\)-restriction of \(\bar\varepsilon\) is \(1\).
  Finally, at the affine test \(B \otimes_A (B \otimes_A B)\), rigidity
  (Theorem~\ref{pa-thm:unitsAppTop_sectionOfPoint_bijective}) forces a global unit with trivial
  \(\sigma\)-restriction to be \(1\), so \(\bar\varepsilon = 1\), i.e.\ \(\varepsilon = 1\), which is
  the claimed identity. This is Kleiman's descent-cocycle step: the automorphism
  \(v_{13}^{-1}v_{23}v_{12}\) of the first-projection pullback of the rigidified pair is trivial.
\end{proof}

\begin{theorem}[Uniqueness of the \(\sigma\)-normalized comparison]
  \label{pa-thm:normalizedComparison_unique}
  \lean{AlgebraicGeometry.Over.normalizedComparison_unique}
  \leanok
  \source{kleiman-picard}

  \uses{pa-def:normalizedCechComparison, pa-thm:exists_global_unit_of_compatible,
    pa-thm:unitsAppTop_sectionOfPoint_bijective}
  \dcref{custom}
  Let the curve \(C\) be proper, geometrically irreducible and geometrically reduced. Two comparison
  cochains \(\theta, \theta'\) on the same cover \(\mathcal W\), each satisfying \textbf{(N1)}
  against the same cocycle \(\gamma\) and \textbf{(N2)} against the same trivializing cochain
  \(\rho\), are equal.
\end{theorem}
\begin{proof}
  \leanok
  The ratio \(\theta/\theta'\) has trivial coboundary: since \(\theta\) and \(\theta'\) cobound the
  same comparison of \(\gamma\) (both satisfy \textbf{(N1)}), the overlap discrepancies of
  \(\theta\) and \(\theta'\) coincide and cancel in the ratio. Hence \(\theta/\theta'\) is \v
  Cech-compatible and glues to a global unit \(\chi\) of \(X_{B \otimes_A B}\)
  (Theorem~\ref{pa-thm:exists_global_unit_of_compatible}). Its \(\sigma\)-restriction is \(1\): by
  \textbf{(N2)} both \(s_q^{\sharp}\theta\) and \(s_q^{\sharp}\theta'\) equal the same
  \(\rho\)-comparison \(q_1^{\sharp}\rho / q_2^{\sharp}\rho\) on the normalization cover, so their
  ratio restricts to \(1\). Rigidity (Theorem~\ref{pa-thm:unitsAppTop_sectionOfPoint_bijective}) at the
  affine test \(B \otimes_A B\) forces \(\chi = 1\); reading this pointwise gives \(\theta =
  \theta'\).
\end{proof}

\begin{theorem}[The diagonal pullback of a witness glues to a global unit]
  \label{pa-thm:exists_glued_diagonalSeed}
  \lean{AlgebraicGeometry.Over.exists_glued_diagonalSeed,
    AlgebraicGeometry.tensorLmul'_comp_tensorInl, AlgebraicGeometry.tensorLmul'_comp_tensorInr,
    AlgebraicGeometry.Over.whiskerLeft_lmul'_inl, AlgebraicGeometry.Over.whiskerLeft_lmul'_inr}
  \leanok
  \dcref{custom}

  \uses{pa-def:normalizedCechComparison, pa-def:tensorInl_tensorInr, pa-thm:exists_global_unit_of_compatible}
  The multiplication \(\nabla \colon B \otimes_A B \to B\) retracts both coprojections, \(\nabla
  \circ q_1 = \id = \nabla \circ q_2\); whiskered to the curve this is \(u_1 \circ \Delta = \id =
  u_2 \circ \Delta\), where \(\Delta \colon X_B \to X_{B \otimes_A B}\) is the diagonal (the
  whiskered \(\Spec\) of \(\nabla\)). Consequently, for a comparison cochain \(\theta_0\) satisfying
  \textbf{(N1)}, the diagonal pullback \(\Delta^{\sharp}\theta_0\) is \v Cech-compatible and glues to
  a global unit \(\bar d \in \Gamma(X_B, \top)^{\times}\). No geometric hypotheses.
\end{theorem}
\begin{proof}
  \leanok
  The retractions are the identities \(\nabla(b \otimes 1) = b\) and \(\nabla(1 \otimes b) = b\);
  whiskering the resulting equalities of \(A\)-algebra maps to the curve gives \(u_1 \circ \Delta =
  \id = u_2 \circ \Delta\). Pull \textbf{(N1)} back along \(\Delta\): both insertion terms of the
  pulled relation are \(\gamma\) evaluated along \(u_1 \circ \Delta = u_2 \circ \Delta\), hence
  coincide and cancel, leaving \(\Delta^{\sharp}\theta_0\) with matching overlap values. Gluing by
  the sheaf axiom (Theorem~\ref{pa-thm:exists_global_unit_of_compatible}) yields \(\bar d\).
\end{proof}

\begin{theorem}[Triviality of the comparison on the diagonal]
  \label{pa-thm:normalizedComparison_diagonal_eq_one}
  \lean{AlgebraicGeometry.Over.normalizedComparison_diagonal_eq_one}
  \leanok
  \source{kleiman-picard}

  \uses{pa-def:normalizedCechComparison, pa-thm:exists_glued_diagonalSeed, pa-lem:sectionSquares,
    pa-thm:unitsAppTop_sectionOfPoint_bijective}
  \dcref{custom}
  Let the curve \(C\) be proper, geometrically irreducible and geometrically reduced. The pullback
  of the comparison cochain \(\theta\) of a \(\sigma\)-normalized comparison along the diagonal
  \(\Delta\) of \(X_{B \otimes_A B}\) is trivial: \(\Delta^{\sharp}\theta = 1\). This is the
  cochain source of the diagonal normalization \(\mu(u) = 1\) of the descent \(1\)-cocycle.
\end{theorem}
\begin{proof}
  \leanok
  Glue the diagonal seed of \(\theta\) to a global unit \(\bar d\) of \(X_B\)
  (Theorem~\ref{pa-thm:exists_glued_diagonalSeed}). Its \(\sigma\)-restriction is \(1\): the section
  square relating \(s_B\), \(\Delta\), the \(\Spec\) of \(\nabla\) and \(s_q\)
  (Lemma~\ref{pa-lem:sectionSquares}) turns it into the \(\nabla\)-pullback of the
  \(\sigma\)-restriction of \(\theta\), which \textbf{(N2)} makes the \(\nabla\)-pullback of the
  \(\rho\)-comparison \(q_1^{\sharp}\rho / q_2^{\sharp}\rho\); the two legs coincide on the diagonal
  because the multiplication retracts both coprojections
  (Theorem~\ref{pa-thm:exists_glued_diagonalSeed}), so the comparison collapses to \(1\). Rigidity
  (Theorem~\ref{pa-thm:unitsAppTop_sectionOfPoint_bijective}) at the affine test \(B\) forces \(\bar d
  = 1\); pointwise this is \(\Delta^{\sharp}\theta = 1\).
\end{proof}

\section{The per-piece descent datum}
\label{pa-sec:effectivityDescentDatum}

The \(\sigma\)-normalized comparison is descended piece by piece over a \(\mathrm{cg}\)-saturated
affine cover of the base curve \(X_A\). On a piece \(V\) the cover-piece ring is \(M := S \otimes_A
B\) with \(S := \Gamma(X_A, U)\) the base-piece ring (Definition~\ref{pa-def:pieceAlgEquiv}), and the
descent along \(S \to M\) is governed by the piece restriction of the comparison. This section
records the pure-algebra bridge: the identification of the descent-cocycle ring with the geometric
double-base ring, and the extraction of a descent \(1\)-cocycle from a comparison unit in geometric
coordinates. Throughout, \(A\) is a commutative ring, \(S\) and \(B\) are \(A\)-algebras, and \(M =
S \otimes_A B\).

\begin{definition}[The descent-cocycle ring identification]
  \label{pa-def:pieceDescentEquiv}
  \lean{Algebra.TensorProduct.pieceDescentEquiv, Algebra.TensorProduct.pieceDescentTripleEquiv}
  \leanok
  \dcref{custom}

  \uses{pa-def:module_descent_cocycle}
  The ring \(M \otimes_S M\) in which a descent \(1\)-cocycle relative to \(S \to M\)
  (Definition~\ref{pa-def:module_descent_cocycle}) lives is the base change to \(S\) of the whole-cover
  descent-cocycle ring \(B \otimes_A B\): there is an \(S\)-algebra isomorphism
  \[
    M \otimes_S M \;\xrightarrow{\ \sim\ }\; S \otimes_A (B \otimes_A B), \qquad
    (s_1 \otimes b_1) \otimes (s_2 \otimes b_2) \;\longmapsto\; (s_2 s_1) \otimes (b_1 \otimes b_2),
  \]
  together with its triple analogue \(M \otimes_S (M \otimes_S M) \xrightarrow{\sim} S \otimes_A (B
  \otimes_A (B \otimes_A B))\).
\end{definition}
\begin{proof}
  \leanok
  The map is the composite of the cancellation of the middle base ring, \((S \otimes_A B)
  \otimes_S (S \otimes_A B) \cong S \otimes_A (S \otimes_A B) \cong (S \otimes_A S) \otimes_A B\)
  contracted along \(S \otimes_A S \to S\), with the associativity isomorphism of the iterated
  tensor product; the displayed value on pure tensors is read off from these two standard
  isomorphisms. The triple form is obtained by applying the same identification in the inner two
  factors and then cancelling and reassociating.
\end{proof}

\begin{lemma}[Compatibility with the multiplication and the cofaces]
  \label{pa-lem:pieceDescentEquiv_faces}
  \lean{Algebra.TensorProduct.pieceDescentEquiv_lmul',
    Algebra.TensorProduct.pieceDescentTripleEquiv_descentFace₂₃,
    Algebra.TensorProduct.pieceDescentTripleEquiv_descentFace₁₂,
    Algebra.TensorProduct.pieceDescentTripleEquiv_descentFace₁₃}
  \leanok
  \dcref{custom}

  \uses{pa-def:pieceDescentEquiv, pa-def:module_descent_cocycle}
  Under the identification of Definition~\ref{pa-def:pieceDescentEquiv}, the multiplication
  \(\mu \colon M \otimes_S M \to M\) corresponds to \(\id_S \otimes \mu\) on \(S \otimes_A (B
  \otimes_A B)\), and the three cofaces \(\partial_{23}, \partial_{12}, \partial_{13}\) of \(M
  \otimes_S M\) correspond to \(\id_S \otimes \partial_{23}, \id_S \otimes \partial_{12}, \id_S
  \otimes \partial_{13}\) on \(S \otimes_A (B \otimes_A (B \otimes_A B))\).
\end{lemma}
\begin{proof}
  \leanok
  Each is an identity between two \(S\)-algebra maps out of \(M \otimes_S M\); it therefore suffices
  to check it on pure tensors, where both sides evaluate to the same element by the computation of
  Definition~\ref{pa-def:pieceDescentEquiv}, and to extend by bilinearity.
\end{proof}

\begin{theorem}[A coherent comparison unit is a descent \(1\)-cocycle]
  \label{pa-thm:isDescentCocycle_comparisonDescentUnit}
  \lean{Module.comparisonDescentUnit, Module.isDescentCocycle_comparisonDescentUnit}
  \leanok
  \dcref{custom}

  \uses{pa-def:pieceDescentEquiv, pa-lem:pieceDescentEquiv_faces, pa-def:module_descent_cocycle}
  Let \(v \in (S \otimes_A (B \otimes_A B))^{\times}\) be a \emph{comparison unit} that is
  diagonal-normalized, \((\id_S \otimes \mu)(v) = 1\), and satisfies the Amitsur identity through
  the three cofaces,
  \[
    (\id_S \otimes \partial_{23})(v) \cdot (\id_S \otimes \partial_{12})(v) = (\id_S \otimes
    \partial_{13})(v) .
  \]
  Then the unit of \(M \otimes_S M\) that corresponds to \(v\) under
  Definition~\ref{pa-def:pieceDescentEquiv} --- the \emph{comparison descent unit} --- is a descent
  \(1\)-cocycle relative to \(S \to M\) (Definition~\ref{pa-def:module_descent_cocycle}).
\end{theorem}
\begin{proof}
  \leanok
  Transport the two hypotheses through the identification. By
  Lemma~\ref{pa-lem:pieceDescentEquiv_faces} the normalization \(\mu\) of the comparison descent unit
  is carried to \((\id_S \otimes \mu)(v) = 1\), which is the normalization condition; and the
  cocycle product \(\partial_{23} \cdot \partial_{12}\) of the comparison descent unit is carried,
  through the triple identification, to \((\id_S \otimes \partial_{23})(v) \cdot (\id_S \otimes
  \partial_{12})(v) = (\id_S \otimes \partial_{13})(v)\), the image of \(\partial_{13}\). As the
  triple identification is injective, the cocycle identity holds.
\end{proof}

\begin{definition}[The descended per-piece class and its cover pullback]
  \label{pa-def:comparisonDescentClass}
  \lean{Module.comparisonDescentClass, Module.mapAlgebra_comparisonDescentClass_eq_one,
    Module.IsDescentCocycle.mapAlgebra_picClass_eq_one}
  \leanok
  \dcref{custom}

  \uses{pa-thm:isDescentCocycle_comparisonDescentUnit, pa-def:module_picClass, pa-def:module_descended}
  Assume \(A \to B\) is faithfully flat, and let \(v\) be a coherent comparison unit as in
  Theorem~\ref{pa-thm:isDescentCocycle_comparisonDescentUnit}. The \emph{descended per-piece class}
  \(M_V \in \operatorname{Pic}(S)\) is the Picard class (Definition~\ref{pa-def:module_picClass}) of the
  descent \(1\)-cocycle of the comparison descent unit. Its image under \(\operatorname{Pic}(S) \to
  \operatorname{Pic}(M)\) is trivial.
\end{definition}
\begin{proof}
  \leanok
  Faithful flatness passes from \(A \to B\) to \(S \to M = S \otimes_A B\) by base change, so the
  descent \(1\)-cocycle of the comparison descent unit has a well-defined Picard class
  (Definition~\ref{pa-def:module_picClass}). Base change to \(M\) sends the descended module to \(M\)
  itself: the descended module of a cocycle is, by construction, the submodule of \(M\) on which the
  two coactions agree (Definition~\ref{pa-def:module_descended}), and tensoring up to \(M\) recovers
  \(M\) as a free rank-one module through the descent isomorphism. Hence the class of \(M_V\) in
  \(\operatorname{Pic}(M)\) is that of a free module, namely \(1\).
\end{proof}

\begin{definition}[The geometric descent-cocycle ring identification]
  \label{pa-def:pieceDescentGeom}
  \lean{AlgebraicGeometry.Over.pieceDescentTensorEquiv,
    AlgebraicGeometry.Over.pieceDescentCoverEquiv}
  \leanok
  \dcref{custom}

  \uses{pa-def:pieceAlgEquiv, pa-def:pieceDescentEquiv}
  For a \(\mathrm{cg}\)-saturated affine piece \(U \subseteq X_A\), with base-piece ring \(\Gamma(X_A,
  U)\) and cover-piece ring \(\Gamma(X_B, \mathrm{cg}^{-1}U)\), the descent-cocycle ring
  \(\Gamma(X_B, \mathrm{cg}^{-1}U) \otimes_{\Gamma(X_A, U)} \Gamma(X_B, \mathrm{cg}^{-1}U)\) is
  identified, as a \(\Gamma(X_A, U)\)-algebra, with \(\Gamma(X_A, U) \otimes_A (B \otimes_A B)\), and
  hence with the section ring \(\Gamma(X_{B \otimes_A B}, \mathrm{cgq}^{-1}U)\) of the curve over the
  descent-cocycle base --- the ring on whose spectrum the piece restriction of the comparison lives.
\end{definition}
\begin{proof}
  \leanok
  Apply the piece identification \(\Gamma(X_B, \mathrm{cg}^{-1}U) \cong \Gamma(X_A, U) \otimes_A B\)
  (Definition~\ref{pa-def:pieceAlgEquiv}) to each tensor factor, then the algebraic descent-cocycle
  ring identification (Definition~\ref{pa-def:pieceDescentEquiv}) with \(S = \Gamma(X_A, U)\); this
  gives the first identification. For the second, apply the piece identification at the test algebra
  \(B \otimes_A B\), giving \(\Gamma(X_A, U) \otimes_A (B \otimes_A B) \cong \Gamma(X_{B \otimes_A
  B}, \mathrm{cgq}^{-1}U)\).
\end{proof}

\section{The per-piece comparison unit and its descended class}
\label{pa-sec:effectivityComparisonUnit}

We continue with the \((C2)\) data: a commutative ring \(A\), an \(A\)-algebra \(B\), the curve
\(C\) over \(\Spec k\), and the cover inclusion \(\mathrm{cg} \colon X_B \to X_A\) (the whiskering
\(C \lhd \Spec(A \to B)\)), with \(\mathrm{cgq} \colon X_{B \otimes_A B} \to X_A\) its double, and
\(u_1, u_2 \colon X_{B \otimes_A B} \to X_B\) the coprojections. A rigidified class \(L =
[\gamma]\) on \(X_B\) is descended \emph{piece by piece} over \(\mathrm{cg}\)-saturated affine opens
\(V\) of \(X_A\). This section manufactures, on each such piece equipped with a trivialization of
\(L\), the comparison unit that the per-piece descent extraction (brick E2) consumes, and its
descended class.

\begin{definition}[A piece trivialization]
  \label{pa-def:pieceTrivialization}
  \lean{AlgebraicGeometry.Over.PieceTrivialization, AlgebraicGeometry.Over.whiskerLeft_inl_comp_cover,
    AlgebraicGeometry.Over.whiskerLeft_inr_comp_cover, AlgebraicGeometry.Over.cgqPreimage_le_inl,
    AlgebraicGeometry.Over.cgqPreimage_le_inr}
  \leanok
  \dcref{custom}

  \uses{pa-def:tensorInl_tensorInr, pa-cor:cechPic_mk_eq_one}
  A \emph{piece trivialization} of a representing cocycle \(\gamma\) on a pointed cover \(\mathcal N\)
  of \(X_B\), over a piece \(V \subseteq X_A\), is a trivializing \(0\)-cochain \(t\) of \(\gamma\)
  on the \(\mathrm{cg}^{-1}V\)-trimmed opens: units \(t_b \in \Gamma(X_B, \mathcal N(b) \cap
  \mathrm{cg}^{-1}V)^{\times}\) with \(t_b \cdot \gamma_{b b'} = t_{b'}\) on the trimmed pairwise
  overlaps. It is the cochain witness that the class \(L = [\gamma]\) is trivial on the cover piece
  \(\mathrm{cg}^{-1}V\) (\ref{pa-cor:cechPic_mk_eq_one}); its existence over a covering family of pieces
  is the piece-selection input consumed by the final splice. The two coprojections both lie over the
  cover inclusions, \(u_1 \circ \mathrm{cg} = \mathrm{cgq} = u_2 \circ \mathrm{cg}\) as maps to
  \(X_A\), so the double piece is bounded by each coprojection preimage of the piece,
  \(\mathrm{cgq}^{-1}V \le u_i^{-1}(\mathrm{cg}^{-1}V)\).
\end{definition}
\begin{proof}
  \leanok
  The two coprojection identities are the whiskerings of \(q_i \circ (A \to B) = (A \to B \otimes_A
  B)\) as \(A\)-algebra maps (Definition~\ref{pa-def:tensorInl_tensorInr}); whiskering to the curve and
  passing to preimages gives \(\mathrm{cgq}^{-1}V \le u_i^{-1}(\mathrm{cg}^{-1}V)\).
\end{proof}

\begin{theorem}[The trivialization-twisted witness glues to a unit]
  \label{pa-thm:exists_glued_unitsTrivTwist}
  \lean{AlgebraicGeometry.unitsTrivTwistCochain, AlgebraicGeometry.unitsTrivTwistCochain_compat,
    AlgebraicGeometry.exists_glued_unitsTrivTwist,
    AlgebraicGeometry.Scheme.Hom.unitsAppLE_glued_trivTwist}
  \leanok
  \dcref{custom}

  \uses{pa-thm:exists_global_unit_of_compatible}
  Over abstract schemes, let \(r_1, r_2 \colon Y \to Z\), let \(\mathcal C\) be a pointed cover of
  \(Y\) and \(\mathcal N\) one of \(Z\), let \(c\) be a comparison \(0\)-cochain on \(\mathcal C\)
  and \(E\) a pair-indexed unit family on \(\mathcal N\), and let \(t\) trivialize \(E\) on the
  \(DZ\)-trimmed opens of \(\mathcal N\), for opens \(DZ \subseteq Z\), \(DY \subseteq Y\) with
  \(DY \le r_i^{-1}DZ\). Define the \emph{twisted value} at \(y\), on \(\mathcal C(y) \cap DY\),
  \[
    (c_y) \cdot \bigl(r_2^{\sharp}t\bigr)_y \cdot \bigl(r_1^{\sharp}t\bigr)_y^{-1} .
  \]
  If \(c\) cobounds the two pullbacks of \(E\) (the relation \(c_x \cdot (r_1^{\sharp}E)_{xy} =
  (r_2^{\sharp}E)_{xy} \cdot c_y\)) and \(t\) satisfies the trivializing relation, then the twisted
  values are \v Cech-compatible on the trimmed cover \(\{\mathcal C(y) \cap DY\}\) of \(DY\), hence
  glue to a global unit \(v \in \Gamma(Y, DY)^{\times}\); and for a further map \(\varphi\) into
  \(Y\) with \(\varphi \circ r_i = m_i\), the \(\varphi\)-pullback of \(v\) expands into the pulled
  witness value \(\varphi^{\sharp}c\) times the two \(m_i\)-pullbacks of \(t\).
\end{theorem}
\begin{proof}
  \leanok
  Compatibility is a telescoping identity: on a pairwise overlap the coboundary relation for \(c\)
  and the two pullbacks of the trivializing relation for \(t\) (along \(r_1\) and \(r_2\)) combine so
  that the \(E\)-terms cancel and the twisted value at \(y\) matches that at \(y'\). A compatible
  unit family on a pointed cover glues to a global unit by the sheaf axiom
  (Theorem~\ref{pa-thm:exists_global_unit_of_compatible}). The pullback expansion is the defining
  gluing relation pulled back along \(\varphi\), with the composites \(\varphi \circ r_i = m_i\)
  merging the two pullbacks.
\end{proof}

\begin{theorem}[The per-piece comparison unit]
  \label{pa-thm:pieceComparisonUnit}
  \lean{AlgebraicGeometry.Over.pieceComparisonUnit, AlgebraicGeometry.Over.pieceComparisonUnit_spec,
    AlgebraicGeometry.Over.pieceComparisonUnit_unique}
  \leanok
  \dcref{custom}

  \uses{pa-thm:exists_glued_unitsTrivTwist, pa-def:pieceTrivialization, pa-def:normalizedCechComparison}
  Let \(W\) be a \(\sigma\)-normalized comparison of \(\gamma\)
  (Definition~\ref{pa-def:normalizedCechComparison}) with comparison cochain \(\theta\), and \(T\) a
  piece trivialization of \(\gamma\) on \(V\). The \emph{per-piece comparison unit}
  \[
    v_V \;\in\; \Gamma\bigl(X_{B \otimes_A B},\, \mathrm{cgq}^{-1}V\bigr)^{\times}
  \]
  is the global unit glued from the values \(\theta_x \cdot (u_2^{\sharp}T)(x) /
  (u_1^{\sharp}T)(x)\) on the trimmed comparison cover: on each member it restricts to that twisted
  witness value. It is the unique unit of the double piece with this restriction property.
\end{theorem}
\begin{proof}
  \leanok
  Apply Theorem~\ref{pa-thm:exists_glued_unitsTrivTwist} with \(r_1, r_2 = u_1, u_2\), comparison
  cover \(\mathcal C = W.\mathrm{cover}\), representing cover \(\mathcal N\), comparison cochain
  \(c = \theta\), pair family \(E = \gamma\), trimming opens \(DZ = \mathrm{cg}^{-1}V\), \(DY =
  \mathrm{cgq}^{-1}V\), and trivialization \(t = T\): the coboundary relation for \(\theta\) is the
  \((N1)\) witness relation of the normalized comparison, and the trivializing relation is that of
  the piece trivialization (Definition~\ref{pa-def:pieceTrivialization}). Gluing yields \(v_V\) and its
  defining restriction property. Uniqueness is the separation half of the sheaf axiom: two units of
  \(\mathrm{cgq}^{-1}V\) with the same restrictions to a cover agree.
\end{proof}

\begin{theorem}[Naturality of the section-ring identification in the algebra]
  \label{pa-thm:pieceRingEquiv_whiskerLeft_naturality}
  \lean{AlgebraicGeometry.Over.pieceRingEquiv_whiskerLeft_naturality,
    AlgebraicGeometry.Over.map_pieceEquiv_eq_pieceEquiv_whiskerLeft,
    AlgebraicGeometry.Over.whiskerLeft_overSpecMap_comp_cover,
    AlgebraicGeometry.Over.coverPreimage_le_whiskerLeft,
    AlgebraicGeometry.Over.restrictScalars_comp_ofId,
    AlgebraicGeometry.Over.pieceRingEquiv_pieceEquiv}
  \leanok
  \dcref{custom}

  \uses{pa-def:pieceRingEquiv, pa-def:pieceAlgEquiv, pa-thm:pieceRingEquiv_naturality}
  The section-ring identification \(\Gamma(X_R, \mathrm{cg}_R^{-1}V) \cong \Gamma(X_A, V)
  \otimes_A R\) (Definition~\ref{pa-def:pieceRingEquiv}, \(R\)-uniform) is natural in the base algebra:
  for an \(A\)-algebra map \(j \colon R \to R'\), the whiskered map \(W_j = (C \lhd \Spec\, j) \colon
  X_{R'} \to X_R\) lies over the cover structure maps (\(W_j \circ \mathrm{cg}_R = \mathrm{cg}_{R'}\)),
  and pulling identified sections back along \(W_j\) is \(\id_{\Gamma(X_A, V)} \otimes j\) through the
  identifications of the two pieces:
  \[
    W_j^{\sharp}\bigl(\Phi^R_V(x)\bigr) \;=\; \Phi^{R'}_V\bigl((\id \otimes j)(x)\bigr) .
  \]
\end{theorem}
\begin{proof}
  \leanok
  Both sides are \(A\)-algebra homomorphisms out of \(\Gamma(X_A, V) \otimes_A R\), so it suffices to
  check them on a pure tensor \(s \otimes r\), then extend by bilinearity. There the identification
  sends \(s \otimes r\) to \(\mathrm{cg}_R^{\sharp}(s) \cdot \mathrm{pr}_R^{\sharp}(r)\)
  (Definition~\ref{pa-def:pieceRingEquiv}). The \(\mathrm{cg}\)-factor pulls back along \(W_j\) by the
  composite \(W_j \circ \mathrm{cg}_R = \mathrm{cg}_{R'}\) (as \(A\)-algebra maps, \(j \circ (A \to R)
  = (A \to R')\)), giving \(\mathrm{cg}_{R'}^{\sharp}(s)\). The projection factor pulls back along the
  square \(W_j \circ \mathrm{pr}_R = \mathrm{pr}_{R'} \circ \Spec\, j\), and the \(\Spec\)-side factor
  is the naturality of \(\Gamma\)-of-\(\Spec\) applied to \(j\), giving \(\mathrm{pr}_{R'}^{\sharp}(j
  r)\). Their product is \(\Phi^{R'}_V(s \otimes j r) = \Phi^{R'}_V((\id \otimes j)(s \otimes r))\).
  This is the algebra-variable mirror of the open-variable restriction seam
  (Theorem~\ref{pa-thm:pieceRingEquiv_naturality}).
\end{proof}

\begin{theorem}[Diagonal triviality of the comparison unit]
  \label{pa-thm:pieceComparisonUnit_diagonal_eq_one}
  \lean{AlgebraicGeometry.Over.pieceComparisonUnit_diagonal_eq_one}
  \leanok
  \source{kleiman-picard}

  \uses{pa-thm:pieceComparisonUnit, pa-thm:normalizedComparison_diagonal_eq_one,
    pa-thm:exists_glued_diagonalSeed}
  \dcref{custom}
  Let \(C\) be proper, geometrically irreducible and geometrically reduced. The pullback of the
  per-piece comparison unit \(v_V\) along the diagonal \(\Delta = (C \lhd \Spec\, \nabla) \colon X_B
  \to X_{B \otimes_A B}\) (\(\nabla\) the multiplication of \(B \otimes_A B\)) is \(1\).
\end{theorem}
\begin{proof}
  \leanok
  Work locally on the diagonal pullback of the comparison cover. The multiplication retracts both
  coprojections, \(\Delta \circ u_1 = \id = \Delta \circ u_2\) whiskered
  (Theorem~\ref{pa-thm:exists_glued_diagonalSeed}), so the pullback expansion of \(v_V\)
  (Theorem~\ref{pa-thm:exists_glued_unitsTrivTwist}, with \(\varphi = \Delta\) and \(m_1 = m_2 = \id\))
  gives the diagonal witness value \(\Delta^{\sharp}\theta\) times two trivialization factors along
  the \emph{same} insertion, which cancel. The remaining witness factor \(\Delta^{\sharp}\theta = 1\)
  is the landed diagonal triviality of the \(\sigma\)-normalized comparison
  (Theorem~\ref{pa-thm:normalizedComparison_diagonal_eq_one}). So the local value is \(1\); by sheaf
  separation the diagonal pullback of \(v_V\) is \(1\). This is the cochain form of the diagonal
  normalization in Kleiman's descent-cocycle step.
\end{proof}

\begin{theorem}[Amitsur coherence of the comparison unit]
  \label{pa-thm:pieceComparisonUnit_coherent}
  \lean{AlgebraicGeometry.Over.pieceComparisonUnit_coherent,
    AlgebraicGeometry.Over.pieceAmitsurOpen}
  \leanok
  \source{kleiman-picard}

  \uses{pa-thm:pieceComparisonUnit, pa-thm:normalizedComparison_coherent, pa-def:amitsurProductCover}
  \dcref{custom}
  Let \(C\) be proper, geometrically irreducible and geometrically reduced. The three coface
  pullbacks of \(v_V\) to the triple piece \(\Gamma(X_{B \otimes_A (B \otimes_A B)},
  \mathrm{cgcb}^{-1}V)^{\times}\) satisfy the Amitsur cocycle identity on the nose:
  \[
    \bigl(w_{23}^{\sharp}v_V\bigr)\bigl(w_{12}^{\sharp}v_V\bigr) \;=\; w_{13}^{\sharp}v_V .
  \]
\end{theorem}
\begin{proof}
  \leanok
  Work locally on the trimmed Amitsur product cover of \(X_{B \otimes_A (B \otimes_A B)}\)
  (Definition~\ref{pa-def:amitsurProductCover}). By the pullback expansion
  (Theorem~\ref{pa-thm:exists_glued_unitsTrivTwist}) each coface pullback \(w_{ij}^{\sharp}v_V\) is the
  coface pullback of the witness \(\theta\) times two insertion pullbacks of the trivialization; the
  three tensor-slot insertions \(v_1, v_2, v_3\) are common to numerator and denominator, so the
  trivialization factors telescope, \((T_{23}\,t_c\,t_b^{-1})(T_{12}\,t_b\,t_a^{-1}) =
  T_{23}T_{12}\,t_c\,t_a^{-1}\). The landed coherence (Theorem~\ref{pa-thm:normalizedComparison_coherent})
  supplies \((w_{23}^{\sharp}\theta)(w_{12}^{\sharp}\theta) = w_{13}^{\sharp}\theta\), so the two
  sides agree locally; by sheaf separation the identity holds globally. This is Kleiman's
  descent-cocycle identity \(v_{13}^{-1}v_{23}v_{12} = 1\).
\end{proof}

\begin{theorem}[The transported comparison unit and its descent hypotheses]
  \label{pa-thm:pieceComparisonTensorUnit}
  \lean{AlgebraicGeometry.Over.pieceComparisonTensorUnit,
    AlgebraicGeometry.Over.pieceComparisonTensorUnit_val,
    AlgebraicGeometry.Over.pieceComparisonTensorUnit_lmul,
    AlgebraicGeometry.Over.pieceComparisonTensorUnit_cocycle}
  \leanok
  \dcref{custom}

  \uses{pa-thm:pieceComparisonUnit, pa-thm:pieceRingEquiv_whiskerLeft_naturality,
    pa-thm:pieceComparisonUnit_diagonal_eq_one, pa-thm:pieceComparisonUnit_coherent, pa-def:pieceAlgEquiv,
    pa-def:pieceDescentEquiv}
  Let \(C\) be proper, geometrically irreducible and geometrically reduced. Through the section-ring
  identification \(\Gamma(X_{B \otimes_A B}, \mathrm{cgq}^{-1}V) \cong \Gamma(X_A, V) \otimes_A (B
  \otimes_A B)\) (Definition~\ref{pa-def:pieceAlgEquiv}), the comparison unit \(v_V\) is a unit
  \(\tilde v_V \in (\Gamma(X_A, V) \otimes_A (B \otimes_A B))^{\times}\) of the descent-cocycle ring.
  It is diagonal-normalized and Amitsur-coherent through the maps of the descent \(1\)-cocycle:
  \[
    (\id \otimes \mu)(\tilde v_V) = 1, \qquad
    (\id \otimes \partial_{23})(\tilde v_V)\,(\id \otimes \partial_{12})(\tilde v_V)
      = (\id \otimes \partial_{13})(\tilde v_V) .
  \]
\end{theorem}
\begin{proof}
  \leanok
  The whiskered maps \(\Delta\) and \(w_{ij}\) are the whiskerings of \(\Spec\) of \(\mu\) and the
  three cofaces \(\partial_{ij}\); so by the algebra-naturality of the identification
  (Theorem~\ref{pa-thm:pieceRingEquiv_whiskerLeft_naturality}) the map \(\id \otimes \mu\) applied to
  \(\tilde v_V\) equals the transport of the diagonal pullback \(\Delta^{\sharp}v_V = 1\)
  (Theorem~\ref{pa-thm:pieceComparisonUnit_diagonal_eq_one}), and likewise the three \(\id \otimes
  \partial_{ij}\) applied to \(\tilde v_V\) are the transports of the three coface pullbacks of
  \(v_V\). The cocycle identity of the coface pullbacks
  (Theorem~\ref{pa-thm:pieceComparisonUnit_coherent}) therefore transports to the stated identity of
  the three \(\id \otimes \partial_{ij}\) images (Definition~\ref{pa-def:pieceDescentEquiv}).
\end{proof}

\begin{theorem}[The descended per-piece class and its cover pullback]
  \label{pa-def:pieceDescentClass}
  \lean{AlgebraicGeometry.Over.pieceDescentClass,
    AlgebraicGeometry.Over.mapAlgebra_pieceDescentClass_eq_one}
  \leanok
  \dcref{custom}

  \uses{pa-thm:pieceComparisonTensorUnit, pa-thm:isDescentCocycle_comparisonDescentUnit,
    pa-def:comparisonDescentClass}
  Assume in addition \(A \to B\) is faithfully flat. The \emph{descended per-piece class}
  \[
    M_V \;\in\; \operatorname{Pic}\bigl(\Gamma(X_A, V)\bigr)
  \]
  is the Picard class of the descent \(1\)-cocycle of the transported comparison unit \(\tilde v_V\),
  whose diagonal-normalization and coherence are those of
  Theorem~\ref{pa-thm:pieceComparisonTensorUnit}. Its pullback to the cover-piece ring \(\Gamma(X_A, V)
  \otimes_A B \cong \Gamma(X_B, \mathrm{cg}^{-1}V)\) is trivial.
\end{theorem}
\begin{proof}
  \leanok
  Faithful flatness of \(A \to B\) passes to \(\Gamma(X_A, V) \to \Gamma(X_A, V) \otimes_A B\) by
  base change, so the coherent comparison unit \(\tilde v_V\)
  (Theorem~\ref{pa-thm:pieceComparisonTensorUnit}) is a descent \(1\)-cocycle
  (Theorem~\ref{pa-thm:isDescentCocycle_comparisonDescentUnit}) with a well-defined Picard class
  \(M_V\) (Definition~\ref{pa-def:comparisonDescentClass}); the same definition records that its base
  change to the cover-piece ring is the trivial class.
\end{proof}

\begin{theorem}[The piece restriction of the descending class is trivial]
  \label{pa-thm:cechPicMap_pieceI_eq_one}
  \lean{AlgebraicGeometry.Scheme.Hom.pullbackUnitsCocycle_isCohomologous_one,
    AlgebraicGeometry.Over.cechPicMap_pieceι_eq_one}
  \leanok
  \dcref{custom}

  \uses{pa-def:pieceTrivialization, pa-cor:cechPic_mk_eq_one}
  The class \(L = [\gamma]\) of the \((C2)\) setting, restricted to the open subscheme
  \(\mathrm{cg}^{-1}V\), is trivial: \(\iota^{*}L = 1\) for the open immersion \(\iota \colon
  \mathrm{cg}^{-1}V \hookrightarrow X_B\). Together with the cover-pullback triviality of
  \(M_V\) (Theorem~\ref{pa-def:pieceDescentClass}), this is the recovery identity \(\mathrm{cg}^{*}M_V
  = [L|_{\mathrm{cg}^{-1}V}]\), both sides trivial by construction.
\end{theorem}
\begin{proof}
  \leanok
  Over abstract schemes, if \(f^{-1}(DZ) = \top\) and a \(0\)-cochain \(t\) trivializes the pair
  values of \(\gamma\) on the \(DZ\)-trimmed opens, then the \(f\)-pullbacks of \(t\) cobound the
  pulled cocycle \(f^{\sharp}\gamma\) against \(1\), so \(f^{\sharp}\gamma\) is cohomologous to
  \(1\). Apply this to the open immersion \(\iota\) of the piece \(\mathrm{cg}^{-1}V\), whose
  self-preimage is \(\top\), with \(t = T.\mathrm{triv}\) the piece trivialization
  (Definition~\ref{pa-def:pieceTrivialization}): the pulled cocycle is cohomologous to \(1\), hence its
  class is trivial (\ref{pa-cor:cechPic_mk_eq_one}).
\end{proof}

\begin{theorem}[Base-change naturality of the descent extraction]
  \label{pa-thm:comparisonDescentClass_mapAlgebra}
  \lean{Module.comparisonDescentUnit_baseChange, Module.comparisonDescentClass_mapAlgebra}
  \leanok
  \dcref{custom}

  \uses{pa-def:comparisonDescentClass, pa-thm:isDescentCocycle_comparisonDescentUnit,
    pa-thm:module_descentCocycle_baseChange, pa-thm:module_picClass_baseChange,
    pa-def:module_tensorSqBaseChange, pa-def:pieceDescentEquiv}
  Let \(S \to S'\) be a map of \(A\)-algebras and \(v\) a coherent comparison unit over \(S\). The
  class descended over \(S'\) from the base-changed comparison \((\mathrm{algebraMap} \otimes \id)(v)\)
  is the \(\operatorname{Pic}\)-image of the class descended over \(S\):
  \[
    \mathrm{comparisonDescentClass}_{S'}\bigl((\mathrm{algebraMap} \otimes \id)(v)\bigr)
      \;=\; \operatorname{Pic.mapAlgebra}_{S \to S'}\bigl(\mathrm{comparisonDescentClass}_S(v)\bigr).
  \]
\end{theorem}
\begin{proof}
  \leanok
  First, the comparison descent unit of the base-changed comparison is the image of the descent unit
  of \(v\) under the tensor-square base change along \(S \to S'\), identified through the
  base-change cancellation on each tensor factor: this is an elementwise computation on pure
  tensors, using the descent-cocycle ring identification (Definition~\ref{pa-def:pieceDescentEquiv}) and
  the tensor-square base change (Definition~\ref{pa-def:module_tensorSqBaseChange}). Now chain the
  descent extraction (Theorem~\ref{pa-thm:isDescentCocycle_comparisonDescentUnit}): the \(S'\)-class is
  the class of the base-changed descent cocycle further mapped by the cancellation
  (\(\operatorname{picClass}\) is invariant under the mapping and under replacing a cocycle by a
  cohomologous one), which equals the class of the base-changed cocycle
  (Theorem~\ref{pa-thm:module_descentCocycle_baseChange}), which in turn is the
  \(\operatorname{mapAlgebra}\)-image of the \(S\)-class
  (Theorem~\ref{pa-thm:module_picClass_baseChange}).
\end{proof}

\begin{theorem}[Restriction canonicity of the descended per-piece classes]
  \label{pa-thm:pieceDescentClass_res}
  \lean{AlgebraicGeometry.Over.PieceTrivialization.res, AlgebraicGeometry.Over.pieceComparisonUnit_res,
    AlgebraicGeometry.Over.pieceComparisonTensorUnit_res,
    AlgebraicGeometry.Over.pieceDescentClass_res}
  \leanok
  \dcref{custom}

  \uses{pa-def:pieceDescentClass, pa-thm:comparisonDescentClass_mapAlgebra, pa-thm:pieceComparisonUnit,
    pa-thm:pieceComparisonTensorUnit, pa-thm:pieceRingEquiv_naturality}
  Assume \(A \to B\) faithfully flat and \(C\) proper, geometrically irreducible and geometrically
  reduced. For affine pieces \(V' \le V\) of \(X_A\) and a piece trivialization \(T\) on \(V\): the
  trivialization restricts to one on \(V'\), the glued comparison unit of the restricted datum is the
  restriction of the glued unit, and the descended classes are natural in the piece,
  \[
    M_{V'} \;=\; \operatorname{Pic.mapAlgebra}_{\Gamma(X_A, V) \to \Gamma(X_A, V')}\,(M_V) .
  \]
\end{theorem}
\begin{proof}
  \leanok
  The trivialization restricts by restricting each trimmed cochain value along \(\mathrm{cg}^{-1}V'
  \le \mathrm{cg}^{-1}V\), preserving the trivializing relation. Its comparison unit is the
  restriction of \(v_V\) by uniqueness of gluing (Theorem~\ref{pa-thm:pieceComparisonUnit}): the
  restriction of \(v_V\) has the defining restriction property of the restricted datum. Through the
  section-ring identifications, this restriction is \(\mathrm{algebraMap} \otimes \id\) on the
  descent-cocycle ring, by the open-variable restriction seam
  (Theorem~\ref{pa-thm:pieceRingEquiv_naturality}) applied to the transported unit
  (Theorem~\ref{pa-thm:pieceComparisonTensorUnit}). Hence \(\tilde v_{V'}\) is the base change of
  \(\tilde v_V\) along \(\Gamma(X_A, V) \to \Gamma(X_A, V')\), and \(M_{V'}\) is the class of the
  base-changed unit, which by the base-change naturality of the extraction
  (Theorem~\ref{pa-thm:comparisonDescentClass_mapAlgebra}) is
  \(\operatorname{mapAlgebra}(M_V)\). On an overlap of two pieces both descend from the same global
  \(\sigma\)-normalized comparison, so the descended classes agree through their common restriction:
  this is what the refinement-splice consumes.
\end{proof}

\section{Trivializing an invertible module away from finitely many primes}
\label{pa-sec:effectivityInvertibleAvoid}

The pure-algebra core of the \((C2)\) piece-selection moving lemma. An invertible module over a
commutative ring becomes free after inverting a single element chosen to avoid any prescribed finite
set of primes; and, for an integral extension, an element of the extension avoiding every prime over
a base prime divides the image of a base element avoiding that prime --- the \emph{algebraic tube}.
These are the two bricks that convert triviality of a class on a basic open of the cover into a
basic open of the base.

\begin{lemma}[A cyclic invertible module is free]
  \label{pa-lem:invertible_cyclic_free}
  \lean{Module.Invertible.bijective_toSpanSingleton_of_span_eq_top,
    Module.Invertible.free_of_span_singleton_eq_top}
  \leanok
  \dcref{custom}

  Let \(R\) be a commutative ring and \(N\) an invertible \(R\)-module generated by a single element
  \(m_0\), i.e.\ \(R \cdot m_0 = N\). Then the scalar map \(r \mapsto r \cdot m_0\) is an isomorphism
  \(R \xrightarrow{\sim} N\); in particular \(N\) is free of rank one.
\end{lemma}
\begin{proof}
  \leanok
  Surjectivity is the hypothesis \(R \cdot m_0 = N\). For injectivity, suppose \(r \cdot m_0 = 0\).
  Since \(m_0\) generates \(N\), \(r\) annihilates every element of \(N\), hence \(r \cdot \id_N = 0\)
  as an endomorphism, and therefore \(r\) annihilates every element of \(N^{\vee} \otimes_R N\) (the
  endomorphism acts through the right tensor factor). Invertibility of \(N\) means the contraction
  pairing \(N^{\vee} \otimes_R N \to R\) is surjective; choose \(w\) with \(\mathrm{contract}(w) = 1\).
  Then \(r = r \cdot \mathrm{contract}(w) = \mathrm{contract}(r \cdot w) = \mathrm{contract}(0) = 0\).
  So the scalar map is bijective, whence \(N \cong R\) is free.
\end{proof}

\begin{theorem}[The avoidance brick]
  \label{pa-thm:invertible_avoid_free}
  \lean{Module.Invertible.exists_notMem_isUnit_free}
  \leanok
  \dcref{custom}

  \uses{pa-lem:invertible_cyclic_free}
  Let \(R\) be a commutative ring, \(N\) an invertible \(R\)-module, and \((\mathfrak q_i)_{i \in I}\)
  a finite family of prime ideals of \(R\). Then there is an element \(f \in R\) with \(f \notin
  \mathfrak q_i\) for every \(i\) such that, for every \(R\)-algebra \(S\) in which \(f\) becomes a
  unit, the base change \(S \otimes_R N\) is a free \(S\)-module.
\end{theorem}
\begin{proof}
  \leanok
  If \(I = \varnothing\) there are no primes to avoid and the claim holds trivially; assume \(I\) is
  nonempty. Localize \(R\) at the multiplicative set \(S_0 := \bigcap_i (R \setminus \mathfrak q_i)\),
  the complement of \(\bigcup_i \mathfrak q_i\). Every maximal ideal of \(R_{S_0}\) is the
  localization of one of the \(\mathfrak q_i\): its contraction to \(R\) is a prime disjoint from
  \(S_0\), hence contained in \(\bigcup_i \mathfrak q_i\), hence --- by prime avoidance --- inside
  some \(\mathfrak q_i\), and maximality forces equality. Thus \(R_{S_0}\) is semilocal with finitely
  many maximal ideals, so its Picard group is trivial, and therefore \(S_0^{-1}N\) is a \emph{cyclic}
  \(R_{S_0}\)-module, generated by a single element \(z\).

  Fix finitely many generators \(g_1, \dots, g_n\) of \(N\) over \(R\), and write \(z = m_0 / s_0\)
  with \(m_0 \in N\), \(s_0 \in S_0\). Clearing denominators in the relations expressing each \(g_j\)
  as an \(R_{S_0}\)-multiple of \(z\) yields elements \(d_j \in S_0\) and \(r_j \in R\) with
  \(d_j \cdot g_j = r_j \cdot m_0\) in \(N\). Set \(f := \prod_j d_j\); each \(d_j \in S_0\) avoids
  every \(\mathfrak q_i\), and a prime containing a product contains a factor, so \(f\) avoids every
  \(\mathfrak q_i\).

  Now let \(S\) be an \(R\)-algebra in which \(f\) is a unit. In \(S \otimes_R N\) each \(1 \otimes
  g_j\) lies in the \(S\)-span of \(1 \otimes m_0\): from \((\prod_l d_l)\cdot g_j = \bigl(\prod_{l
  \ne j} d_l\bigr) r_j \cdot m_0\), inverting the unit \(f = \prod_l d_l\) expresses \(1 \otimes g_j\)
  as an \(S\)-multiple of \(1 \otimes m_0\). Since the \(g_j\) generate \(N\), every pure tensor
  \(1 \otimes x\) --- and hence all of \(S \otimes_R N\) --- lies in the \(S\)-span of \(1 \otimes
  m_0\). So \(S \otimes_R N\) is cyclic, and invertible because invertibility is preserved by base
  change; by Lemma~\ref{pa-lem:invertible_cyclic_free} it is free.
\end{proof}

\begin{lemma}[The algebraic tube]
  \label{pa-lem:algebraic_tube}
  \lean{Algebra.exists_notMem_algebraMap_eq_mul}
  \leanok
  \dcref{custom}

  Let \(S \to R\) be an integral extension of commutative rings, \(\mathfrak p\) a prime of \(S\),
  and \(f \in R\) an element avoiding every prime of \(R\) lying over \(\mathfrak p\) (that is,
  \(f \notin Q\) whenever \(Q\) is prime with \(Q \cap S = \mathfrak p\)). Then there are \(g \in S
  \setminus \mathfrak p\) and \(h \in R\) with \(g = f \cdot h\) in \(R\).
\end{lemma}
\begin{proof}
  \leanok
  Consider the contracted ideal \((fR) \cap S\) of \(S\). Were it contained in \(\mathfrak p\),
  going-up for the integral extension \(S \to R\) would produce a prime \(Q \supseteq fR\) of \(R\)
  with \(Q \cap S = \mathfrak p\); but then \(f \in Q\) is a prime over \(\mathfrak p\) containing
  \(f\), contradicting the hypothesis. Hence \((fR) \cap S \not\subseteq \mathfrak p\): some
  \(g \in (fR) \cap S\) lies outside \(\mathfrak p\). By construction \(g \in fR\), so \(g = f \cdot
  h\) for some \(h \in R\).
\end{proof}

\section{The moving lemma: pieces on which the class trivializes}
\label{pa-sec:effectivityMoving}

We now select, for a \emph{module-finite} faithfully flat cover \(A \to B\), the affine pieces of the
base curve \(X_A\) on which a prescribed \v Cech Picard class of the cover curve \(X_B\) becomes
trivial after pullback to the cover piece. The route is pure commutative algebra --- no integrality
of the curve, no divisors --- through the affine-open class and the two bricks of the previous
section. Throughout, \(\mathrm{cg} \colon X_B \to X_A\) is the cover inclusion.

\begin{definition}[The affine-open \v Cech--Picard class]
  \label{pa-def:cechPicClass_affineOpen}
  \lean{AlgebraicGeometry.Scheme.Opens.cechPicClass, AlgebraicGeometry.Scheme.Opens.ιTop,
    AlgebraicGeometry.Scheme.Opens.cechPicClass_of_le}
  \leanok
  \dcref{custom}

  \uses{pa-def:cechPic_toPic}
  Let \(Z\) be a scheme and \(O \subseteq Z\) an affine open. Pullback along the open immersion \(O
  \hookrightarrow Z\) identifies the ambient sections \(\Gamma(Z, O)\) with the global sections
  \(\Gamma(O, \top)\) of the open subscheme; write \(\iota_O^{\top}\) for this ring isomorphism. The
  \emph{affine-open class} of a \v Cech Picard class \(L\) on \(Z\) is
  \[
    L\langle O\rangle \;:=\; \bigl(\iota_O^{\top}\bigr)^{-1}_{*}\,
      \mathrm{toPic}_O\bigl(\iota^{*}L\bigr) \;\in\; \operatorname{Pic}\bigl(\Gamma(Z, O)\bigr),
  \]
  the invertible \(\Gamma(Z, O)\)-module obtained by transporting, through the affine \v
  Cech--Picard dictionary \(\mathrm{toPic}\) (Definition~\ref{pa-def:cechPic_toPic}), the restriction of
  \(L\) to \(O\). It is natural in the open: for affine opens \(O' \le O\) the class \(L\langle
  O'\rangle\) is the image of \(L\langle O\rangle\) under the restriction homomorphism
  \(\operatorname{Pic}(\Gamma(Z, O)) \to \operatorname{Pic}(\Gamma(Z, O'))\).
\end{definition}
\begin{proof}
  \leanok
  \uses{pa-thm:cechPic_toPic_map}
  The section identification \(\iota_O^{\top}\) is an isomorphism because \(O \hookrightarrow Z\) is
  an open immersion with self-preimage \(O\), so its pullback on sections over \(O\) is invertible;
  transporting the affine class \(\mathrm{toPic}_O(\iota^{*}L)\) back along it gives a well-defined
  element of \(\operatorname{Pic}(\Gamma(Z, O))\). Naturality is the commuting square of the two open
  immersions \(O' \le O \le Z\): the restriction map of ambient sections and the two section
  identifications intertwine the \(\mathrm{toPic}\)-images by naturality of the dictionary in the
  scheme (Theorem~\ref{pa-thm:cechPic_toPic_map}), so the class formed on \(O'\) is the restriction of
  the class formed on \(O\).
\end{proof}

\begin{theorem}[Triviality detection for the affine-open class]
  \label{pa-thm:cechPicClass_triv}
  \lean{AlgebraicGeometry.Scheme.Opens.cechPicMap_ι_eq_one_of_cechPicClass_eq_one,
    AlgebraicGeometry.Scheme.cechPicMap_ι_eq_one_of_le,
    AlgebraicGeometry.Scheme.Opens.cechPicClass_basicOpen_eq_one_of_free}
  \leanok
  \dcref{custom}

  \uses{pa-def:cechPicClass_affineOpen, pa-thm:cechPic_toPic_injective}
  Let \(Z\) be a scheme and \(L\) a \v Cech Picard class on \(Z\).
  \begin{enumerate}
    \item If the affine-open class \(L\langle O\rangle\) is trivial for an affine open \(O\), then
      the restriction \(\iota^{*}L\) of \(L\) to \(O\) is trivial.
    \item Restriction triviality is monotone: if \(\iota_O^{*}L = 1\) and \(O' \le O\), then
      \(\iota_{O'}^{*}L = 1\).
    \item If, for an affine open \(O\) and \(f \in \Gamma(Z, O)\), the invertible module \(L\langle
      O\rangle\) becomes free over every \(\Gamma(Z, O)\)-algebra in which \(f\) is a unit, then the
      affine-open class on the basic open \(D(f)\) is trivial, \(L\langle D(f)\rangle = 1\); hence
      \(L\) restricts trivially to \(D(f)\).
  \end{enumerate}
\end{theorem}
\begin{proof}
  \leanok
  (1) The dictionary \(\mathrm{toPic}\) is injective (Theorem~\ref{pa-thm:cechPic_toPic_injective}), so
  \(\iota^{*}L = 1\) follows from \(\mathrm{toPic}_O(\iota^{*}L) = 1\); the latter is the transport
  of \(L\langle O\rangle = 1\) through the section identification, which reflects triviality because
  it has a two-sided inverse. (2) Restriction along \(O' \le O\) factors the restriction to \(O'\)
  through the restriction to \(O\), and a group homomorphism sends \(1\) to \(1\). (3) The
  restriction homomorphism \(\Gamma(Z, O) \to \Gamma(Z, D(f))\) inverts \(f\), because \(f\) is a
  unit on its own basic open; so by hypothesis \(L\langle O\rangle\) base-changes to a free, hence
  trivial, class over \(\Gamma(Z, D(f))\), and by naturality of the affine-open class
  (Definition~\ref{pa-def:cechPicClass_affineOpen}) this base change is \(L\langle D(f)\rangle\), which
  is therefore \(1\). Triviality of \(L\) restricted to \(D(f)\) then follows from~(1).
\end{proof}

\begin{lemma}[Module-finiteness of the cover piece over the base piece]
  \label{pa-lem:finite_pieceCover}
  \lean{AlgebraicGeometry.Over.finite_pieceCover}
  \leanok
  \dcref{custom}

  \uses{pa-def:pieceAlgEquiv}
  Let \(A \to B\) be a module-finite homomorphism of \(k\)-algebras and \(V \subseteq X_A\) an affine
  open. Then the cover-piece ring \(\Gamma(X_B, \mathrm{cg}^{-1}V)\) is a module-finite \(\Gamma(X_A,
  V)\)-algebra.
\end{lemma}
\begin{proof}
  \leanok
  By the piece identification (Definition~\ref{pa-def:pieceAlgEquiv}) the cover-piece ring is \(\Gamma(
  X_A, V) \otimes_A B\) as a \(\Gamma(X_A, V)\)-algebra. Base change preserves module-finiteness: as
  \(B\) is module-finite over \(A\), the base change \(\Gamma(X_A, V) \otimes_A B\) is module-finite
  over \(\Gamma(X_A, V)\).
\end{proof}

\begin{theorem}[The moving lemma]
  \label{pa-thm:moving_lemma}
  \lean{AlgebraicGeometry.Over.exists_isAffineOpen_cechPicMap_preimage_eq_one}
  \leanok
  \dcref{custom}

  \uses{pa-thm:cechPicClass_triv, pa-thm:invertible_avoid_free, pa-lem:algebraic_tube, pa-lem:finite_pieceCover,
    pa-def:cechPicClass_affineOpen, pa-thm:exists_isAffineOpen_cg_le, pa-def:pieceAlgEquiv}
  Let \(A \to B\) be module-finite. For every \v Cech Picard class \(L\) on the cover curve \(X_B\)
  and every point \(x \in X_A\), there is an affine open \(V \ni x\) of the base curve such that \(L\)
  restricts trivially to the open subscheme \(\mathrm{cg}^{-1}V\) of \(X_B\).
\end{theorem}
\begin{proof}
  \leanok
  Choose a \(\mathrm{cg}\)-saturated affine open \(V_0 \ni x\) of \(X_A\)
  (Theorem~\ref{pa-thm:exists_isAffineOpen_cg_le}); the cover piece \(\mathrm{cg}^{-1}V_0\) is then
  affine with ring \(N := \Gamma(X_B, \mathrm{cg}^{-1}V_0)\), and \(L\langle \mathrm{cg}^{-1}V_0
  \rangle\) (Definition~\ref{pa-def:cechPicClass_affineOpen}) is an invertible \(N\)-module. By
  module-finiteness (Lemma~\ref{pa-lem:finite_pieceCover}) the piece map \(\Gamma(X_A, V_0) \to N\) is
  module-finite, hence integral and quasi-finite; therefore the fibre of \(\operatorname{Spec} N \to
  \operatorname{Spec}\Gamma(X_A, V_0)\) over the germ prime \(\mathfrak p_x\) of \(x\) is a
  \emph{finite} set of primes \(\mathfrak q_1, \dots, \mathfrak q_m\) of \(N\).

  Apply the avoidance brick (Theorem~\ref{pa-thm:invertible_avoid_free}) to the invertible module
  \(L\langle \mathrm{cg}^{-1}V_0\rangle\) and the finite family \((\mathfrak q_j)\): there is \(f \in
  N\) with \(f \notin \mathfrak q_j\) for all \(j\) such that \(L\langle \mathrm{cg}^{-1}V_0\rangle\)
  becomes free wherever \(f\) is inverted. Since \(f\) avoids every prime of \(N\) over \(\mathfrak
  p_x\), the algebraic tube (Lemma~\ref{pa-lem:algebraic_tube}) for the integral extension \(\Gamma(X_A,
  V_0) \to N\) produces \(g \notin \mathfrak p_x\) and \(h \in N\) with \(\mathrm{cg}^{\sharp}g = f
  \cdot h\). Set \(V := D(g) \subseteq X_A\). Then \(x \in V\), because \(g \notin \mathfrak p_x\)
  means the germ of \(g\) at \(x\) is a unit, and
  \[
    \mathrm{cg}^{-1}V = D(\mathrm{cg}^{\sharp}g) = D(f \cdot h) \le D(f) .
  \]
  Freeness of \(L\langle \mathrm{cg}^{-1}V_0\rangle\) after inverting \(f\) trivializes the class on
  \(D(f)\) (Theorem~\ref{pa-thm:cechPicClass_triv}(3)), and restriction triviality is monotone
  (Theorem~\ref{pa-thm:cechPicClass_triv}(2)), so \(L\) restricts trivially to \(\mathrm{cg}^{-1}V \le
  D(f)\).

  \emph{Module-finiteness is necessary.} For a merely quasi-finite (for instance \'etale) non-finite
  cover the conclusion fails. There the fibre of the piece map over \(\mathfrak p_x\) can be infinite,
  and the bare restricted class need not become trivial on any \(\mathrm{cg}\)-saturated neighbourhood
  of the fibre: take \(B = A_u \times A_v\), a product of two localizations of a two-dimensional
  normal local ring \(A\) with non-trivial class group, and a class whose restrictions to the two
  factors are non-trivial classes of the two localizations; no single saturated piece trivializes it.
  The general faithfully-flat effectivity therefore genuinely needs the per-piece descent datum of
  Section~\ref{pa-sec:effectivityComparisonUnit} rather than bare piece-triviality, and survives with
  piece-selection taken as a hypothesis (Theorem~\ref{pa-thm:effectivity_of_pieces}).
\end{proof}

\section{From subscheme triviality to a trimmed trivializing cochain}
\label{pa-sec:effectivityTrivialization}

The moving lemma delivers triviality of the class \emph{as a class} on the cover piece. The per-piece
descent of Section~\ref{pa-sec:effectivityComparisonUnit} instead consumes a trivializing \emph{cochain}
of the representing cocycle, trimmed to the piece and indexed by all ambient points. This section
converts the one to the other.

\begin{theorem}[The trimmed trivializing cochain of a class trivial on an open subscheme]
  \label{pa-thm:exists_trimmed_trivializing}
  \lean{AlgebraicGeometry.Scheme.exists_trimmed_trivializing_of_cechPicMap_ι_eq_one}
  \leanok
  \dcref{custom}

  \uses{pa-cor:cechPic_mk_eq_one, pa-def:unitsPreimageEquiv, pa-lem:units_gluing}
  Let \(Z\) be a scheme, \(\mathcal N\) a pointed cover of \(Z\), \(\gamma\) a unit \(1\)-cocycle on
  \(\mathcal N\), and \(D \subseteq Z\) an open subscheme. If the class \([\gamma]\) restricts
  trivially to \(D\) (that is \(\iota_D^{*}[\gamma] = 1\)), then there is a family of units
  \[
    t_b \in \Gamma\bigl(Z,\, \mathcal N(b) \cap D\bigr)^{\times}, \qquad b \in Z,
  \]
  indexed by \emph{all} ambient points, satisfying the one-sided trivializing relation \(t_b \cdot
  \gamma_{b b'} = t_{b'}\) on the \(D\)-trimmed pairwise overlaps.
\end{theorem}
\begin{proof}
  \leanok
  Triviality of the restricted class means the pulled-back cocycle \(\iota_D^{\sharp}\gamma\) on the
  cover \(\iota_D^{-1}\mathcal N\) of \(D\) is cohomologous to \(1\)
  (Corollary~\ref{pa-cor:cechPic_mk_eq_one}); as restriction to an open subscheme needs no further
  refinement, this yields a trivializing \(0\)-cochain \(\alpha\), indexed by points \(p\) of \(D\),
  with \(\alpha_p \cdot (\iota_D^{\sharp}\gamma)_{p p'} = \alpha_{p'}\). Transport each \(\alpha_p\)
  across the open immersion \(\iota_D\): on the preimage of the trimmed member \(\mathcal
  N(\iota_D p) \cap D\) the unit \(\alpha_p\) is the image of an ambient unit
  (Definition~\ref{pa-def:unitsPreimageEquiv}, unit sections across an open immersion), giving units
  \(t'_p\) on \(\mathcal N(\iota_D p) \cap D\) that satisfy the ambient trivializing relation on
  subscheme pairs. Finally re-index from points of \(D\) to points of \(Z\): for each ambient \(b\)
  the trimmed member \(\mathcal N(b) \cap D\) is covered by the domains of the \(t'_p\), the values
  \(t'_p \cdot \gamma_{(\iota_D p)\, b}\) agree on overlaps by the cocycle identity, so they glue
  (Lemma~\ref{pa-lem:units_gluing}) to a unit \(t_b\) on \(\mathcal N(b) \cap D\). The glued family
  satisfies \(t_b \cdot \gamma_{b b'} = t_{b'}\), again by the cocycle identity together with the
  local relations.
\end{proof}

\begin{theorem}[A piece trivialization from restriction triviality]
  \label{pa-thm:pieceTrivialization_of_triv}
  \lean{AlgebraicGeometry.Over.pieceTrivialization_of_cechPicMap_eq_one}
  \leanok
  \dcref{custom}

  \uses{pa-thm:exists_trimmed_trivializing, pa-def:pieceTrivialization}
  Let \(\gamma\) be a representing unit cocycle on a pointed cover \(\mathcal N\) of \(X_B\) and \(V
  \subseteq X_A\) a piece. If the class \([\gamma]\) restricts trivially to the cover piece
  \(\mathrm{cg}^{-1}V\), then \(V\) carries a piece trivialization of \(\gamma\)
  (Definition~\ref{pa-def:pieceTrivialization}).
\end{theorem}
\begin{proof}
  \leanok
  Instantiate Theorem~\ref{pa-thm:exists_trimmed_trivializing} at the open subscheme \(D = \mathrm{cg}
  ^{-1}V\): the trimmed trivializing cochain it produces is precisely the field data of a piece
  trivialization (Definition~\ref{pa-def:pieceTrivialization}).
\end{proof}

\section{Twisting a piece trivialization}
\label{pa-sec:effectivityTwist}

Two piece trivializations of the same representing cocycle on the same piece differ, member by member,
by the restrictions of a single global unit of the cover piece (this is proved in the next section).
Here we record how the per-piece comparison unit responds to such a twist, and the algebraic and
geometric seams that identify the twist with a descent coboundary --- the mechanism behind both the
trivialization-independence of the descended class and the flattening of a trivialization.

\begin{theorem}[Twisting a piece trivialization by a global unit]
  \label{pa-thm:pieceComparisonUnit_twist}
  \lean{AlgebraicGeometry.Over.PieceTrivialization.twist, AlgebraicGeometry.Over.geomCoboundary,
    AlgebraicGeometry.Over.pieceComparisonUnit_eq_of_triv_eq}
  \leanok
  \dcref{custom}

  \uses{pa-def:pieceTrivialization, pa-thm:pieceComparisonUnit}
  Let \(T\) be a piece trivialization of \(\gamma\) on \(V\) and \(e \in \Gamma(X_B, \mathrm{cg}^{-1}
  V)^{\times}\) a global unit of the cover piece. The \emph{twist} \(T \cdot e\) is the piece
  trivialization with members \(t_b \cdot e|_{\mathcal N(b) \cap \mathrm{cg}^{-1}V}\). Write the
  \emph{geometric descent coboundary} of \(e\),
  \[
    \partial e \;:=\; \frac{u_2^{\sharp} e}{u_1^{\sharp} e} \;\in\; \Gamma\bigl(X_{B \otimes_A B},\,
      \mathrm{cgq}^{-1}V\bigr)^{\times},
  \]
  the ratio of its two coprojection pullbacks to the double piece. If two piece trivializations \(T,
  T'\) differ member by member by \(e\), i.e.\ \(T' = T \cdot e\), then for any \(\sigma\)-normalized
  comparison \(W\) their per-piece comparison units (Theorem~\ref{pa-thm:pieceComparisonUnit}) differ by
  \(\partial e\):
  \[
    v_V(T') = v_V(T) \cdot \partial e .
  \]
\end{theorem}
\begin{proof}
  \leanok
  Multiplying every member of \(T\) by the restrictions of the single global unit \(e\) preserves the
  trivializing relation \(t_b \cdot \gamma_{b b'} = t_{b'}\): the two \(e\)-factors coincide on the
  overlap. For the comparison units, the per-piece comparison unit is characterized by its
  restriction to the trimmed comparison cover (Theorem~\ref{pa-thm:pieceComparisonUnit}), where it
  restricts to \(\theta_x \cdot (u_2^{\sharp}T)(x) / (u_1^{\sharp}T)(x)\). Replacing \(T\) by \(T
  \cdot e\) multiplies this local value by \((u_2^{\sharp}e)/(u_1^{\sharp}e)\) at every point; by
  uniqueness of the glued unit this multiplies \(v_V(T)\) by the global unit \(\partial e\) with that
  restriction.
\end{proof}

\begin{theorem}[The coboundary seams of the identification]
  \label{pa-thm:geomCoboundary_seams}
  \lean{Module.comparisonDescentUnit_coboundary,
    AlgebraicGeometry.Over.map_pieceEquiv_geomCoboundary}
  \leanok
  \dcref{custom}

  \uses{pa-thm:pieceComparisonUnit_twist, pa-def:pieceDescentEquiv,
    pa-thm:isDescentCocycle_comparisonDescentUnit, pa-thm:pieceRingEquiv_whiskerLeft_naturality}
  Fix an affine piece \(V\) and \(e \in \Gamma(X_B, \mathrm{cg}^{-1}V)^{\times}\); let \(\beta \in
  (\Gamma(X_A, V) \otimes_A B)^{\times}\) be its image under the piece identification
  (Definition~\ref{pa-def:pieceAlgEquiv}). Then:
  \begin{enumerate}
    \item \textup{(algebra seam)} the comparison descent unit of the base-changed coboundary of
      \(\beta\) --- the ratio of the two inclusions \(\beta \mapsto \beta \otimes 1\), \(\beta
      \mapsto 1 \otimes \beta\) of \(B \to B \otimes_A B\), base changed to \(\Gamma(X_A, V)\) ---
      equals the descent coboundary of \(\beta\) relative to \(\Gamma(X_A, V) \to \Gamma(X_A, V)
      \otimes_A B\);
    \item \textup{(geometric seam)} through the piece identification, the geometric descent coboundary
      \(\partial e\) (Theorem~\ref{pa-thm:pieceComparisonUnit_twist}) is carried to the base-changed
      coboundary of \(\beta\).
  \end{enumerate}
\end{theorem}
\begin{proof}
  \leanok
  (1) Transport both sides through the descent-cocycle ring identification
  (Definition~\ref{pa-def:pieceDescentEquiv}) to \((\Gamma(X_A, V) \otimes_A (B \otimes_A B))^{\times}\):
  there the identification carries a left pure tensor to the base change of the first inclusion and a
  right pure tensor to that of the second (checked on pure tensors, extended by bilinearity), so the
  comparison descent unit of the inclusion ratio (Theorem~\ref{pa-thm:isDescentCocycle_comparisonDescentUnit})
  is the descent coboundary of \(\beta\). (2) The two coprojections \(u_1, u_2\) are the whiskerings
  of the two inclusions \(B \to B \otimes_A B\), so by the algebra-naturality of the section-ring
  identification (Theorem~\ref{pa-thm:pieceRingEquiv_whiskerLeft_naturality}) each pullback \(u_i^{\sharp}
  e\) is carried to the corresponding inclusion applied to \(\beta\), and their ratio is the
  base-changed coboundary of \(\beta\).
\end{proof}

\section{Trivialization independence and flat trivializations}
\label{pa-sec:effectivityOverlap}

We deduce the \((B)\) deliverable of the \((C2)\) close --- the descended per-piece class is
independent of the chosen trivialization --- and its constructive upgrade: on a piece whose descended
class is trivial some trivialization is \emph{flat}, with per-piece comparison unit \(1\) on the nose.
The flat pieces produced here are what the refinement splice glues.

\begin{theorem}[Two piece trivializations differ by a glued global unit]
  \label{pa-thm:exists_ratio_unit}
  \lean{AlgebraicGeometry.Over.exists_ratio_unit}
  \leanok
  \dcref{custom}

  \uses{pa-def:pieceTrivialization, pa-lem:units_gluing}
  Any two piece trivializations \(T, T'\) of the same cocycle \(\gamma\) on the same piece \(V\)
  differ by a single global unit of the cover piece: there is \(e \in \Gamma(X_B, \mathrm{cg}^{-1}
  V)^{\times}\) with \(T' = T \cdot e\), i.e.\ \(t'_b = t_b \cdot e|_{\mathcal N(b) \cap \mathrm{cg}
  ^{-1}V}\) for every \(b\).
\end{theorem}
\begin{proof}
  \leanok
  The member-wise ratio \(t'_b / t_b\), a unit on the trimmed member \(\mathcal N(b) \cap \mathrm{cg}
  ^{-1}V\), is \v Cech-compatible: the two trivializing relations \(t_b \cdot \gamma_{b b'} =
  t_{b'}\) and \(t'_b \cdot \gamma_{b b'} = t'_{b'}\) cancel the cocycle \(\gamma\), so \(t'_b/t_b\)
  and \(t'_{b'}/t_{b'}\) agree on the overlap. The trimmed members cover \(\mathrm{cg}^{-1}V\), so the
  ratios glue (Lemma~\ref{pa-lem:units_gluing}) to a unit \(e\) with the stated restrictions.
\end{proof}

\begin{theorem}[Trivialization independence of the descended per-piece class]
  \label{pa-thm:pieceDescentClass_congr}
  \lean{AlgebraicGeometry.Over.comparisonDescentUnit_twist,
    AlgebraicGeometry.Over.pieceDescentClass_congr}
  \leanok
  \dcref{custom}

  \uses{pa-thm:exists_ratio_unit, pa-thm:pieceComparisonUnit_twist, pa-thm:geomCoboundary_seams,
    pa-def:pieceDescentClass, pa-def:module_picClass}
  Let \(C\) be proper, geometrically irreducible and geometrically reduced, and \(A \to B\)
  faithfully flat. For a fixed \(\sigma\)-normalized comparison \(W\) and an affine piece \(V\), the
  descended per-piece class \(M_V\) (Theorem~\ref{pa-def:pieceDescentClass}) does not depend on the
  chosen piece trivialization.
\end{theorem}
\begin{proof}
  \leanok
  Let \(T, T'\) be two piece trivializations; by Theorem~\ref{pa-thm:exists_ratio_unit} they differ by a
  global unit \(e\), \(T' = T \cdot e\). Transported to the descent-cocycle ring, the comparison unit
  of \(T'\) is that of \(T\) multiplied by the descent coboundary of \(\beta\), the piece
  identification of \(e\): this is the twist lemma (Theorem~\ref{pa-thm:pieceComparisonUnit_twist})
  combined with the two coboundary seams (Theorem~\ref{pa-thm:geomCoboundary_seams}). The Picard class
  of a descent \(1\)-cocycle is unchanged when the cocycle is multiplied by a descent coboundary,
  since a coboundary alters the descent datum by an isomorphism and not the isomorphism class of the
  descended module (Definition~\ref{pa-def:module_picClass}). Hence \(M_V(T') = M_V(T)\).
\end{proof}

\begin{theorem}[Flattening a trivialization along a trivial descended class]
  \label{pa-thm:exists_pieceComparisonUnit_eq_one}
  \lean{AlgebraicGeometry.Over.exists_pieceComparisonUnit_eq_one}
  \leanok
  \dcref{custom}

  \uses{pa-thm:pieceComparisonUnit_twist, pa-thm:geomCoboundary_seams, pa-def:pieceDescentClass,
    pa-def:module_picClass, pa-thm:pieceComparisonUnit}
  Let \(C\) be proper, geometrically irreducible and geometrically reduced, \(A \to B\) faithfully
  flat, \(W\) a \(\sigma\)-normalized comparison, and \(T\) a piece trivialization on \(V\) whose
  descended class \(M_V(T)\) is trivial. Then \(V\) carries a piece trivialization whose per-piece
  comparison unit is \(1\) on the nose.
\end{theorem}
\begin{proof}
  \leanok
  Triviality of \(M_V(T)\) means the descent \(1\)-cocycle of the transported comparison unit
  cobounds: there is \(\beta \in (\Gamma(X_A, V) \otimes_A B)^{\times}\) whose descent coboundary is
  that cocycle (Definition~\ref{pa-def:module_picClass}: the class is \(1\) exactly when the descended
  module is free). Pull \(\beta\) back through the piece identification to a global unit \(e \in
  \Gamma(X_B, \mathrm{cg}^{-1}V)^{\times}\). By the coboundary seams
  (Theorem~\ref{pa-thm:geomCoboundary_seams}) the per-piece comparison unit of \(T\) is the geometric
  descent coboundary \(\partial e\). Twisting \(T\) by \(e^{-1}\) multiplies the comparison unit by
  \(\partial(e^{-1}) = (\partial e)^{-1}\) (Theorem~\ref{pa-thm:pieceComparisonUnit_twist}), giving
  comparison unit \(v_V(T) \cdot (\partial e)^{-1} = 1\).
\end{proof}

\begin{theorem}[Localize and flatten]
  \label{pa-thm:exists_flat_pieceTrivialization}
  \lean{AlgebraicGeometry.Over.exists_flat_pieceTrivialization_of_le}
  \leanok
  \dcref{custom}

  \uses{pa-thm:exists_pieceComparisonUnit_eq_one, pa-thm:invertible_avoid_free, pa-thm:pieceDescentClass_res}
  Let \(C\) be proper, geometrically irreducible and geometrically reduced, \(A \to B\) faithfully
  flat, and \(W\) a \(\sigma\)-normalized comparison. Every point \(x\) of a trivialized affine piece
  \((V_1, T_1)\) lies in a basic affine subpiece \(V \le V_1\) carrying a \emph{flat} piece
  trivialization --- one whose per-piece comparison unit is \(1\).
\end{theorem}
\begin{proof}
  \leanok
  Apply the avoidance brick (Theorem~\ref{pa-thm:invertible_avoid_free}) to the invertible module
  \(M_{V_1}(T_1)\) and the single germ prime \(\mathfrak p_x\) of \(x\) on \(\Gamma(X_A, V_1)\): it
  yields \(f \notin \mathfrak p_x\) such that \(M_{V_1}(T_1)\) becomes free wherever \(f\) is
  inverted. Set \(V := D(f) \subseteq X_A\); then \(x \in V\) (the germ of \(f\) at \(x\) is a unit)
  and \(V \le V_1\). Restrict \(T_1\) to \(V\); by restriction canonicity of the descended classes
  (Theorem~\ref{pa-thm:pieceDescentClass_res}) the restricted descended class is the base change of
  \(M_{V_1}(T_1)\) along \(\Gamma(X_A, V_1) \to \Gamma(X_A, V)\), which inverts \(f\) and so is free,
  hence trivial. Flattening (Theorem~\ref{pa-thm:exists_pieceComparisonUnit_eq_one}) then produces a
  piece trivialization on \(V\) with comparison unit \(1\).
\end{proof}

\section{The refinement splice}
\label{pa-sec:effectivitySplice}

The assembly of the descended class from a covering family of flat pieces, run in the constructive
direction of the \((C1)\) kernel lemma's final assembly.

\begin{theorem}[The refinement splice]
  \label{pa-thm:refinement_splice}
  \lean{AlgebraicGeometry.Over.exists_cechPic_map_whiskerLeft_eq_of_flat}
  \leanok
  \dcref{custom}

  \uses{pa-thm:exists_ratio_unit, pa-thm:pieceComparisonUnit_twist, pa-thm:exists_unitsAppLE_eq,
    pa-thm:appLE_whiskerLeft_injective, pa-thm:isAffineOpen_inf, pa-def:pieceTrivialization,
    pa-cor:cechPic_mk_eq_one}
  Let \(C\) be proper, geometrically irreducible and geometrically reduced, and \(A \to B\)
  faithfully flat. Suppose every point of \(X_A\) lies in an affine piece carrying a \emph{flat}
  piece trivialization of \(\gamma\) (per-piece comparison unit \(1\)). Then the class \(L =
  [\gamma]\) descends along the cover inclusion: there is a \v Cech Picard class \(M\) on \(X_A\)
  with \(\mathrm{cg}^{*}M = L\).
\end{theorem}
\begin{proof}
  \leanok
  Fix a covering family \((V_x, T_x)\) of flat trivialized pieces. On the overlap \(V_x \cap V_y\)
  the two restricted trivializations differ by a glued unit \(c_{xy}\) of the overlap cover piece
  (Theorem~\ref{pa-thm:exists_ratio_unit}); both restrictions are still flat, so the twist lemma
  (Theorem~\ref{pa-thm:pieceComparisonUnit_twist}) forces the geometric descent coboundary of
  \(c_{xy}\) to be \(1\) --- the two coprojection pullbacks of \(c_{xy}\) to the double piece agree
  on the nose. Faithfully flat descent of unit sections along \(\mathrm{cg}\)
  (Theorem~\ref{pa-thm:exists_unitsAppLE_eq}) therefore descends each \(c_{xy}\) to a unit \(m_{xy}\) of
  the affine overlap \(V_x \cap V_y\) downstairs, the overlaps being affine by separatedness of the
  curve product (Theorem~\ref{pa-thm:isAffineOpen_inf}). The \(m_{xy}\) satisfy the \v Cech cocycle
  identity: their \(\mathrm{cg}\)-pullbacks telescope through the trivializations, and
  \(\mathrm{cg}\)-pullback is injective on sections (Theorem~\ref{pa-thm:appLE_whiskerLeft_injective}).
  Let \(M\) be the class of the resulting unit cocycle \(m\) on the pointed cover \((V_x)\). Then
  \(\mathrm{cg}^{*}M = L\) by \(\mathrm{mk}\)-calculus over the common refinement of the
  \(\mathrm{cg}\)-pullback of \((V_x)\) and \(\mathcal N\): the trivializations \(T_x\) themselves
  cobound the two pulled cocycles (Corollary~\ref{pa-cor:cechPic_mk_eq_one}).
\end{proof}

\section{The \((C2)\) effectivity theorem and the surjectivity close}
\label{pa-sec:effectivityClose}

The campaign close: the effectivity theorem in both forms, and the resulting bijectivity of the unit
of the \'etale plus construction over a section-admitting field test.

\begin{theorem}[Effectivity from a covering family of trivialized pieces]
  \label{pa-thm:effectivity_of_pieces}
  \lean{AlgebraicGeometry.Over.exists_cechPic_map_whiskerLeft_eq_of_pieces}
  \leanok
  \dcref{custom}

  \uses{pa-thm:exists_normalizedCechComparison, pa-thm:exists_flat_pieceTrivialization,
    pa-thm:refinement_splice, pa-def:isRigidified}
  Let \(C\) be proper, geometrically irreducible and geometrically reduced, and \(A \to B\)
  faithfully flat. Let \(\gamma\) be a representing cocycle on \(X_B\) whose class \(L = [\gamma]\) is
  rigidified along the base-changed section (Definition~\ref{pa-def:isRigidified}) and satisfies the
  on-the-nose descent equation \(u_1^{*}L = u_2^{*}L\). If every point of \(X_A\) lies in an affine
  piece on which \(\gamma\) admits a piece trivialization, then \(L\) descends along the cover
  inclusion: \(\exists M,\ \mathrm{cg}^{*}M = L\).
\end{theorem}
\begin{proof}
  \leanok
  The descent equation together with rigidification produces a \(\sigma\)-normalized comparison \(W\)
  of \(\gamma\), with no further geometric hypotheses
  (Theorem~\ref{pa-thm:exists_normalizedCechComparison}). At each point, localize-and-flatten
  (Theorem~\ref{pa-thm:exists_flat_pieceTrivialization}) turns the given trivialized piece into a flat
  one through the point. The resulting covering family of flat pieces feeds the refinement splice
  (Theorem~\ref{pa-thm:refinement_splice}), which descends \(L\).
\end{proof}

\begin{theorem}[The \((C2)\) effectivity theorem]
  \label{pa-thm:effectivity}
  \lean{AlgebraicGeometry.Over.exists_cechPic_map_whiskerLeft_eq}
  \leanok
  \source{kleiman-picard}

  \uses{pa-thm:effectivity_of_pieces, pa-thm:moving_lemma, pa-thm:pieceTrivialization_of_triv}
  \dcref{custom}
  Let \(C\) be proper, geometrically irreducible and geometrically reduced, and \(A \to B\) a
  \emph{module-finite} faithfully flat cover. Every \v Cech Picard class \(L\) on the cover curve
  \(X_B\) that is rigidified along the base-changed section and satisfies the on-the-nose descent
  equation \(u_1^{*}L = u_2^{*}L\) descends along the cover inclusion: \(\exists M,\ \mathrm{cg}^{*}M
  = L\).
\end{theorem}
\begin{proof}
  \leanok
  Represent \(L = [\gamma]\). Because the cover is module-finite, the moving lemma
  (Theorem~\ref{pa-thm:moving_lemma}) applies: every point of \(X_A\) has an affine neighbourhood on
  which \(L\) restricts trivially to the cover piece, and such a restriction triviality furnishes a
  piece trivialization (Theorem~\ref{pa-thm:pieceTrivialization_of_triv}). The covering family of
  trivialized pieces feeds the general effectivity Theorem~\ref{pa-thm:effectivity_of_pieces}. This is
  the cocycle recast of the descent step in Kleiman's Comparison Theorem, Part~2: the two projection
  pullbacks of the rigidified pair agree, so the automorphism \(v_{13}^{-1}v_{23}v_{12}\) of its
  first-projection pullback is trivial and the pair descends along the cover.
\end{proof}

\begin{theorem}[Surjectivity of the unit over a section-admitting field test]
  \label{pa-thm:unit_surjective_of_section}
  \lean{AlgebraicGeometry.PicEtAff.unit_surjective_of_section}
  \leanok
  \source{kleiman-picard}

  \uses{pa-thm:effectivity, pa-thm:etale_cover_field_cofinality, pa-def:isRigidified,
    pa-thm:PicEtAff_unit_eq_mk, pa-def:relPicAlgMap, pa-lem:PicEtAff_mk_descentMap, pa-def:descentClasses,
    pa-lem:PicEtAff_ind}
  \dcref{custom}
  Let \(C\) be a curve over \(k\) that is proper, geometrically irreducible and geometrically
  reduced, and let \(K\) be a field extension of \(k\) admitting a \(K\)-point \(\sigma \colon
  \operatorname{Spec} K \to C\). Then the unit \(\eta_K \colon \mathrm{relPic}(C, \operatorname{Spec}
  K) \to \mathrm{PicEtAff}(C, K)\) of the \'etale plus construction is surjective.
\end{theorem}
\begin{proof}
  \leanok
  A plus class is represented on an \'etale cover \(E\) of \(\operatorname{Spec} K\)
  (Lemma~\ref{pa-lem:PicEtAff_ind}). \'Etale covers of a field are cofinal in finite separable field
  extensions (Theorem~\ref{pa-thm:etale_cover_field_cofinality}): refine \(E\) to a finite separable
  field extension \(L/K\), giving a module-finite --- hence faithfully flat --- cover \(F =
  \operatorname{Spec} L\) and a descent class \(y\) over \(F\) representing the same plus class
  (Lemma~\ref{pa-lem:PicEtAff_mk_descentMap}). The rigidification layer produces a rigidified \v Cech
  representative \(L_c\) of \(y\) (Definition~\ref{pa-def:isRigidified}) whose two double-cover pullbacks
  agree on the nose. As \(K \to L\) is module-finite faithfully flat, the effectivity theorem
  (Theorem~\ref{pa-thm:effectivity}) descends \(L_c\) along the cover inclusion to a class \(M\) on
  \(X_K\) with \(\mathrm{cg}^{*}M = L_c\). The relative Picard class \([M]\) maps under \(\eta_K\) to
  the original plus class, by the identification of the unit with the tautological \(\mathrm{mk}\)
  (Theorem~\ref{pa-thm:PicEtAff_unit_eq_mk}) and \(\mathrm{mk}\)-calculus for the algebra map \(K \to L\)
  (Definition~\ref{pa-def:relPicAlgMap}, Definition~\ref{pa-def:descentClasses}). Hence \(\eta_K\) is
  surjective. This is the surjectivity direction of Kleiman's Comparison Theorem Part~2 for a base
  admitting a section --- over a field test the \'etale-local section is an actual point.
\end{proof}

\begin{theorem}[The unit is an isomorphism over a section-admitting field test]
  \label{pa-thm:unitEquiv_of_section}
  \lean{AlgebraicGeometry.PicEtAff.unitEquiv_of_section}
  \leanok
  \source{kleiman-picard}

  \uses{pa-thm:unit_surjective_of_section, pa-thm:PicEtAff_unit_injective}
  \dcref{custom}
  Under the hypotheses of Theorem~\ref{pa-thm:unit_surjective_of_section}, the unit is a group
  isomorphism
  \[
    \eta_K \colon \mathrm{relPic}(C, \operatorname{Spec} K) \;\xrightarrow{\ \sim\ }\;
      \mathrm{PicEtAff}(C, K).
  \]
\end{theorem}
\begin{proof}
  \leanok
  The unit is injective by \'etale separatedness (the \((C1)\) close,
  Theorem~\ref{pa-thm:PicEtAff_unit_injective}) and surjective by
  Theorem~\ref{pa-thm:unit_surjective_of_section}; a bijective group homomorphism is an isomorphism.
\end{proof}

\section{Divisor classes from local equations}

An effective Cartier divisor \(D\) on a scheme \(X\) is cut out, near each point, by a single
regular equation. Adapted to the pointed covers of the \v Cech Picard group \(\Pic(X)\) recalled
earlier in this chapter, this data --- one regular section \(f_x\) on each member
of a pointed cover, the sections agreeing up to a unit on overlaps --- is a trivialization of the
invertible sheaf \(\struct X(D)\): on the member \(\mathcal U(x)\) the sheaf is trivialized by
multiplication by \(f_x\). The transition units \(g_{xy} = f_x/f_y \in \struct
X(\mathcal U(x) \cap \mathcal U(y))^{\times}\) form a unit \v Cech \(1\)-cocycle, and its class in
\(\Pic(X)\) is \(\struct X(D)\). This section records that construction and its independence of
the chosen equations up to units; the next builds from it the class of a Weil divisor on a curve.

\begin{definition}[Local equations for an effective divisor]
  \label{pa-def:localEquations}
  \lean{AlgebraicGeometry.Scheme.LocalEquations}
  \leanok
  \dcref{custom}

  A system of \emph{local equations} on \(X\), adapted to a pointed cover \(\mathcal U\), consists
  of a section \(f_x \in \struct X(\mathcal U(x))\) on each member, subject to two conditions:
  \begin{itemize}
    \item (regularity) the germ of \(f_x\) at every point \(y \in \mathcal U(x)\) is a
      nonzerodivisor in the stalk \(\struct{X,y}\);
    \item (overlap) for every pair \(x, y\) the restrictions of \(f_x\) and \(f_y\) to the overlap
      \(\mathcal U(x) \cap \mathcal U(y)\) differ by a unit factor: there is a unit \(u \in
      \struct X(\mathcal U(x) \cap \mathcal U(y))^{\times}\) with \(f_x|_{\cap} = u \cdot
      f_y|_{\cap}\).
  \end{itemize}
\end{definition}

\begin{lemma}[The transition units]
  \label{pa-lem:localEquations_ratioUnit}
  \lean{AlgebraicGeometry.Scheme.LocalEquations.ratioUnit,
    AlgebraicGeometry.Scheme.LocalEquations.eqn_restrict_eq,
    AlgebraicGeometry.Scheme.LocalEquations.ratioUnit_unique}
  \leanok
  \dcref{custom}

  \uses{pa-def:localEquations}
  For a local-equation system and points \(x, y\) there is a \emph{unique} unit \(g_{xy} \in
  \struct X(\mathcal U(x) \cap \mathcal U(y))^{\times}\) with \(f_x|_{\cap} = g_{xy} \cdot
  f_y|_{\cap}\) on the overlap; write \(g_{xy} = f_x/f_y\) for the transition unit.
\end{lemma}
\begin{proof}
  \leanok
  Existence is the overlap condition of Definition~\ref{pa-def:localEquations}. For uniqueness, the
  restriction of \(f_y\) to the overlap is again a nonzerodivisor: a section of a sheaf is
  determined by its germs, so it is a nonzerodivisor as soon as each of its germs is, and the germs
  of \(f_y\) are nonzerodivisors by regularity. Hence in the equation \(f_x|_{\cap} = u \cdot
  f_y|_{\cap}\) the right factor \(f_y|_{\cap}\) can be cancelled on the right, determining the unit
  \(u\) uniquely.
\end{proof}

\begin{lemma}[The cocycle identity of the transition units]
  \label{pa-lem:localEquations_cocycle}
  \lean{AlgebraicGeometry.Scheme.LocalEquations.ratioUnit_trans}
  \leanok
  \dcref{custom}

  \uses{pa-lem:localEquations_ratioUnit}
  On the triple overlap \(\mathcal U(x) \cap \mathcal U(y) \cap \mathcal U(z)\) the transition
  units satisfy \(g_{xy} \cdot g_{yz} = g_{xz}\).
\end{lemma}
\begin{proof}
  \leanok
  Restrict everything to the triple overlap and multiply through by \(f_z\). By the defining
  equation (Lemma~\ref{pa-lem:localEquations_ratioUnit}) applied to the pairs \((y,z)\) and \((x,y)\),
  \[
    g_{xy}\, g_{yz}\, f_z = g_{xy}\, f_y = f_x ,
  \]
  while applied to \((x,z)\) it gives \(g_{xz}\, f_z = f_x\). The restriction of \(f_z\) to the
  triple overlap is a nonzerodivisor (regularity, as in
  Lemma~\ref{pa-lem:localEquations_ratioUnit}), so cancelling it yields \(g_{xy}\, g_{yz} = g_{xz}\).
\end{proof}

\begin{definition}[The Picard class of a local-equation system]
  \label{pa-def:localEquations_picClass}
  \lean{AlgebraicGeometry.Scheme.LocalEquations.unitsCocycle,
    AlgebraicGeometry.Scheme.LocalEquations.picClass}
  \leanok
  \dcref{custom}

  \uses{pa-def:localEquations, pa-lem:localEquations_ratioUnit, pa-lem:localEquations_cocycle}
  The transition units \(g_{xy}\) of a local-equation system \(d\) form a unit \v Cech
  \(1\)-cocycle on its cover \(\mathcal U\); its class
  \[
    \mathrm{picClass}(d) = [\mathcal U,\, g] \in \Pic(X)
  \]
  is the \emph{Picard class} \(\struct X(D)\) of the effective divisor \(D\) cut out by \(d\).
\end{definition}
\begin{proof}
  \leanok
  Each \(g_{xy}\) is a unit on the double overlap \(\mathcal U(x) \cap \mathcal U(y)\), and the
  cocycle identity \(g_{xy}\, g_{yz} = g_{xz}\) holds on triple overlaps
  (Lemma~\ref{pa-lem:localEquations_cocycle}); this is exactly the data of a unit \v Cech
  \(1\)-cocycle on \(\mathcal U\). Take its class in the colimit \(\Pic(X)\).
\end{proof}

\begin{lemma}[Independence of the class under unit rescaling]
  \label{pa-lem:localEquations_rescale}
  \lean{AlgebraicGeometry.Scheme.LocalEquations.rescale,
    AlgebraicGeometry.Scheme.LocalEquations.rescale_ratioUnit,
    AlgebraicGeometry.Scheme.LocalEquations.picClass_rescale}
  \leanok
  \dcref{custom}

  \uses{pa-def:localEquations, pa-lem:localEquations_ratioUnit, pa-def:localEquations_picClass}
  Let \(d\) be a local-equation system on \(\mathcal U\) and \((v_x)_x\) a family of units
  \(v_x \in \struct X(\mathcal U(x))^{\times}\). Rescaling each equation, \(f'_x = v_x \cdot f_x\),
  is again a local-equation system on \(\mathcal U\), and it has the same Picard class:
  \(\mathrm{picClass}(d\text{ rescaled by }v) = \mathrm{picClass}(d)\). In particular \(\struct
  X(D)\) depends only on the divisor, not on the chosen equations.
\end{lemma}
\begin{proof}
  \leanok
  Each \(v_x \cdot f_x\) is regular --- at every point its germ is the product of a unit germ and a
  nonzerodivisor germ, hence a nonzerodivisor --- and its overlap ratio is the conjugated unit
  \[
    g'_{xy} = (v_x|_{\cap}) \cdot g_{xy} \cdot (v_y|_{\cap})^{-1},
  \]
  so the rescaled data is a local-equation system on the same cover \(\mathcal U\) with transition
  units \(g'\). The rings \(\struct X(U)\) are commutative, so their units form an abelian group;
  hence \(g'_{xy} = (v_x|_{\cap})(v_y|_{\cap})^{-1}\, g_{xy}\), which exhibits \(g'\) as the product
  of \(g\) with the coboundary of the \(0\)-cochain \((v_x)_x\). Cohomologous cocycles on the same
  cover have the same class in \(\Pic(X)\), so \(\mathrm{picClass}\) is unchanged.
\end{proof}

\section{Pullback of a local-equation system}

A morphism \(f \colon Y \to X\) pulls a local-equation system on \(X\) back to one on \(Y\),
provided the pulled-back equations remain regular; on Picard classes this realizes the pullback of
the associated invertible sheaf, \(f^{*}\struct X(D) = \struct Y(f^{*}D)\).

\begin{definition}[Pullback of a local-equation system]
  \label{pa-def:localEquations_pullback}
  \lean{AlgebraicGeometry.Scheme.LocalEquations.pullback,
    AlgebraicGeometry.Scheme.LocalEquations.pullbackEqn}
  \leanok
  \dcref{custom}

  \uses{pa-def:localEquations, pa-lem:localEquations_ratioUnit}
  Let \(f \colon Y \to X\) be a morphism of schemes and \(E\) a local-equation system on \(X\)
  (\ref{pa-def:localEquations}), adapted to a pointed cover \(\mathcal U\). Pull the cover back
  pointwise: the pulled-back pointed cover \(f^{*}\mathcal U\) assigns to \(y \in Y\) the open
  \(f^{-1}\mathcal U(f(y))\). The \emph{pulled-back equations} are the \(f\)-pullbacks
  \(f^{*}\bigl(f_{f(y)}\bigr) \in \struct Y\bigl(f^{-1}\mathcal U(f(y))\bigr)\) of the equations of
  \(E\).

  These form a local-equation system \(f^{*}E\) on \(f^{*}\mathcal U\), \emph{under the additional
  hypothesis that the pulled-back equations are regular}: that the germ of \(f^{*}(f_{f(y)})\) at
  every point of its member is a nonzerodivisor. This regularity is not automatic --- pulling a
  nonzerodivisor back along a non-flat morphism can produce a zerodivisor, so the naive claim
  ``pull the equations back and stay a local-equation system'' is false in general, and the
  hypothesis is genuinely required. (Flatness of \(f\) would be a stronger sufficient condition,
  since flat pullback preserves nonzerodivisors; it is not assumed here.)
\end{definition}
\begin{proof}
  \leanok
  Regularity of the pulled-back equations is the added hypothesis. For the overlap condition, fix
  \(y, y'\) and restrict the defining equation
  \(f_{f(y)}|_{\cap} = g_{f(y)f(y')}\cdot f_{f(y')}|_{\cap}\) of \(E\)
  (\ref{pa-lem:localEquations_ratioUnit}) along \(f\). Since pullback of sections is a ring
  homomorphism commuting with restriction, the two pulled-back equations differ on the overlap
  \(f^{-1}\mathcal U(f(y)) \cap f^{-1}\mathcal U(f(y'))\) by the factor
  \(f^{*}\bigl(g_{f(y)f(y')}\bigr)\), the \(f\)-pullback of the transition unit of \(E\). The
  pullback of a unit section is again a unit, so this factor is a unit and the overlap condition
  holds.
\end{proof}

\begin{lemma}[Transition data of a pullback]
  \label{pa-lem:localEquations_pullback_ratioUnit}
  \lean{AlgebraicGeometry.Scheme.LocalEquations.pullback_ratioUnit,
    AlgebraicGeometry.Scheme.LocalEquations.pullback_unitsCocycle}
  \leanok
  \dcref{custom}

  \uses{pa-def:localEquations_pullback, pa-lem:localEquations_ratioUnit, pa-def:localEquations_picClass}
  The transition units of the pulled-back system \(f^{*}E\) are the \(f\)-pullbacks of those of
  \(E\): for all \(y, y'\),
  \[
    g^{f^{*}E}_{y y'} = f^{*}\bigl(g_{f(y)f(y')}\bigr).
  \]
  Consequently the unit \v Cech \(1\)-cocycle of \(f^{*}E\) is the \(f\)-pullback of the cocycle of
  \(E\).
\end{lemma}
\begin{proof}
  \leanok
  The overlap ratio arising in the construction of \(f^{*}E\) (\ref{pa-def:localEquations_pullback})
  is \(f^{*}(g_{f(y)f(y')})\), and the transition unit of \(f^{*}E\) is the unique unit realizing
  that ratio (\ref{pa-lem:localEquations_ratioUnit}); the two therefore agree. Assembling over all
  pairs, the cocycle of \(f^{*}E\) (\ref{pa-def:localEquations_picClass}) is the entry-by-entry
  \(f\)-pullback of the cocycle of \(E\).
\end{proof}

\begin{theorem}[The Picard class of a pullback is the pullback of the class]
  \label{pa-thm:localEquations_picClass_pullback}
  \lean{AlgebraicGeometry.Scheme.LocalEquations.picClass_pullback}
  \leanok
  \dcref{custom}

  \uses{pa-def:localEquations_pullback, pa-def:localEquations_picClass,
    pa-lem:localEquations_pullback_ratioUnit}
  For a morphism \(f \colon Y \to X\) and a local-equation system \(E\) on \(X\) with regular
  pullback, the Picard class of \(f^{*}E\) is the pullback of the class of \(E\):
  \[
    \mathrm{picClass}(f^{*}E) = f^{*}\,\mathrm{picClass}(E) \quad\text{in } \Pic(Y),
  \]
  that is, \(\struct Y(f^{*}D) = f^{*}\struct X(D)\).
\end{theorem}
\begin{proof}
  \leanok
  Both sides are represented on the pulled-back cover \(f^{*}\mathcal U\). The class of \(f^{*}E\)
  is the class of its cocycle (\ref{pa-def:localEquations_picClass}), and that cocycle is the
  \(f\)-pullback of the cocycle of \(E\) (\ref{pa-lem:localEquations_pullback_ratioUnit}); pullback of
  Picard classes is, by definition, induced by pulling cocycles back along \(f\). Hence the class
  of the pulled-back cocycle is the pullback \(f^{*}\) of \(\mathrm{picClass}(E)\).
\end{proof}

\section{The point divisor and the class of a Weil divisor}

We return to the curve bundle \(X\) of Chapter~\ref{pa-chap:RiemannRoch}: integral, smooth of relative
dimension one and quasi-compact over the field \(K\), with generic point \(\eta\). To each closed
point \(x\) we attach the local equations of the effective divisor \(1 \cdot x\), whose only zero
is a simple one at \(x\); passing to Picard classes and taking products to the multiplicities sends
each Weil divisor \(D \in \Div(X)\) to its class \(\struct X(D) \in \Pic(X)\), a homomorphism
extending the point classes.

\begin{lemma}[A uniformizer isolated to a single zero]
  \label{pa-lem:exists_pointLocalEquation}
  \lean{AlgebraicGeometry.Scheme.exists_pointLocalEquation}
  \leanok
  \dcref{custom}

  \uses{pa-lem:uniformizer, pa-lem:exists_stalk_of_ord_le_one, pa-lem:germ_generic_eq_algebraMap_germ,
    pa-thm:ordZ_support_finite, pa-thm:closedpoint_isClosed, pa-lem:ord_eq_one_of_mem_basicOpen,
    pa-def:scheme_ordZ}
  For a closed point \(x\) of the curve there are an open neighbourhood \(V \ni x\) and a section
  \(s \in \struct X(V)\) whose germ at the generic point \(\eta\) is the chosen uniformizer \(t\) at
  \(x\), whose germ at every point \(y \in V\) is a nonzerodivisor of \(\struct{X,y}\), and is
  moreover a unit whenever \(y \neq x\). Thus \(s\) is a uniformizer at \(x\) whose only zero on
  \(V\) is the point \(x\) itself, and whose realization as a rational function is recorded to be
  \(t\).
\end{lemma}
\begin{proof}
  \leanok
  Let \(t \in K(X)\) be a uniformizer at \(x\), so \(\mathrm{ord}_x t = 1\)
  (Lemma~\ref{pa-lem:uniformizer}); it is a nonzero rational function. Having nonnegative order at
  \(x\), it is integral there: \(t\) is the image of an element of the stalk \(\struct{X,x}\)
  (Lemma~\ref{pa-lem:exists_stalk_of_ord_le_one}). Represent that stalk element by a section \(s_0\)
  over some open neighbourhood \(W_0 \ni x\). Its germ at the generic point is the image of its
  germ at \(x\) under \(\struct{X,x} \to K(X)\) (Lemma~\ref{pa-lem:germ_generic_eq_algebraMap_germ}),
  so \(\mathrm{germ}_\eta(s_0) = t\).

  Now \(t\), being a nonzero rational function, has nonzero order at only finitely many closed
  points (Theorem~\ref{pa-thm:ordZ_support_finite}); let \(\mathrm{bad}\) be the finite set of those
  points \emph{other than} \(x\). Each closed point is a closed singleton
  (Theorem~\ref{pa-thm:closedpoint_isClosed}), so \(\mathrm{bad}\) is a finite closed set; removing it
  from \(W_0\) leaves an open set
  \[
    V := W_0 \setminus \mathrm{bad}, \qquad x \in V
  \]
  (indeed \(x \notin \mathrm{bad}\)). Restrict \(s_0\) to \(V\) and call the result \(s\); its germ
  at \(y \in V\) still maps to \(t\) under \(\struct{X,y} \to K(X)\). In particular at the generic
  point the germ is exactly \(\mathrm{germ}_\eta(s) = t\), the tracked realization.

  \emph{Regularity.} For \(y \in V\), the germ \(\mathrm{germ}_y(s)\) maps to \(t \neq 0\) in the
  function field under the localization \(\struct{X,y} \to K(X)\); since that map is injective on
  the domain stalk, \(\mathrm{germ}_y(s) \neq 0\), i.e.\ it is a nonzerodivisor.

  \emph{Unit away from \(x\).} Let \(y \in V\) with \(y \neq x\). If \(y = \eta\), then
  \(\mathrm{germ}_\eta(s) = t \neq 0\) is a unit of the field \(K(X)\). Otherwise \(y\) is a closed
  point; suppose \(\mathrm{germ}_y(s)\) were not a unit of \(\struct{X,y}\). A section invertible at
  a closed point has order of vanishing \(0\) there (Lemma~\ref{pa-lem:ord_eq_one_of_mem_basicOpen});
  contrapositively, non-invertibility of \(s\) at \(y\) forces \(\mathrm{ord}_y t \neq 0\)
  (recall \(\mathrm{germ}_\eta(s) = t\), and \(\mathrm{ord}_y\) is the order homomorphism
  \ref{pa-def:scheme_ordZ}). Then \(y\) is a closed point where \(t\) has nonzero order, distinct from
  \(x\), so \(y \in \mathrm{bad}\) --- contradicting \(y \in V = W_0 \setminus \mathrm{bad}\). Hence
  \(\mathrm{germ}_y(s)\) is a unit.
\end{proof}

\begin{definition}[The point divisor \(1 \cdot x\)]
  \label{pa-def:pointDivisor}
  \lean{AlgebraicGeometry.Scheme.pointDivisor,
    AlgebraicGeometry.Scheme.germGeneric_pointDivisorEqn_self}
  \leanok
  \dcref{custom}

  \uses{pa-def:localEquations, pa-lem:exists_pointLocalEquation, pa-thm:closedpoint_isClosed}
  For a closed point \(x\), the \emph{point divisor} \(1 \cdot x\) is the local-equation system on
  the pointed cover
  \[
    \mathcal U(y) = \begin{cases} V & y = x, \\ X \setminus \{x\} & y \neq x, \end{cases}
  \]
  with \(V\) the neighbourhood of Lemma~\ref{pa-lem:exists_pointLocalEquation}, whose equation on the
  distinguished member \(V\) is the uniformizer-section \(s\) and whose equation on every other
  member \(X \setminus \{x\}\) is the constant \(1\). Its defining data is fully accessible: in
  particular the germ at the generic point \(\eta\) of the equation on the member at \(x\) is the
  chosen uniformizer \(t\), which is what pins the Weil divisor these equations cut out.
\end{definition}
\begin{proof}
  \leanok
  The complement \(X \setminus \{x\}\) is open because \(\{x\}\) is closed
  (Theorem~\ref{pa-thm:closedpoint_isClosed}), and it contains every \(y \neq x\), so \(\mathcal U\) is
  a pointed cover. Regularity holds member by member: on \(V\) the germs of \(s\) are
  nonzerodivisors (Lemma~\ref{pa-lem:exists_pointLocalEquation}), and on \(X \setminus \{x\}\) the
  germs of the constant \(1\) are units, hence nonzerodivisors. For the overlap ratios, consider a
  pair \((a, b)\). If neither or both of \(a, b\) equal \(x\), the two equations on the overlap are
  equal (both \(1\), or both restrictions of \(s\) on \(V \cap V = V\)), so the ratio is the unit
  \(1\). If exactly one of \(a, b\) equals \(x\), the overlap is contained in \(V \cap (X \setminus
  \{x\}) = V \setminus \{x\}\), where \(s\) is a unit (Lemma~\ref{pa-lem:exists_pointLocalEquation});
  the ratio of the two equations \(s\) and \(1\) is then \(s\) or \(s^{-1}\), a unit there. Thus the
  data is a local-equation system. The equation on the member at \(x\) is the section \(s\), whose
  germ at \(\eta\) is the tracked uniformizer \(t\) (Lemma~\ref{pa-lem:exists_pointLocalEquation}).
\end{proof}

\begin{definition}[The class of a Weil divisor]
  \label{pa-def:divisorClass}
  \lean{AlgebraicGeometry.Scheme.divisorClass}
  \leanok
  \dcref{custom}

  \uses{pa-def:pointDivisor, pa-def:localEquations_picClass, pa-def:curveDivisor}
  For a Weil divisor \(D \in \Div(X)\) (\ref{pa-def:curveDivisor}), its \emph{Picard class} is the
  finite product
  \[
    \struct X(D) = \prod_{x} \struct X(x)^{D_x} \in \Pic(X),
  \]
  over the finite support of \(D\), where \(\struct X(x) := \mathrm{picClass}(1 \cdot x) \in
  \Pic(X)\) is the class of the point divisor (\ref{pa-def:pointDivisor}, \ref{pa-def:localEquations_picClass}).
  A negative coefficient \(D_x\) contributes the inverse class, so the exponent is an integer power
  in the abelian group \(\Pic(X)\); the assignment is defined on all of \(\Div(X)\), with no
  effective/anti-effective case split. The product is finite because \(D\) is finitely supported
  and \(\struct X(x)^0 = 1\).
\end{definition}

\begin{theorem}[Additivity of the divisor class]
  \label{pa-thm:divisorClass_add}
  \lean{AlgebraicGeometry.Scheme.divisorClass_add}
  \leanok
  \dcref{custom}

  \uses{pa-def:divisorClass}
  The class map \(D \mapsto \struct X(D)\) is a homomorphism \(\Div(X) \to \Pic(X)\): for all
  \(D, D' \in \Div(X)\),
  \[
    \struct X(D + D') = \struct X(D) \cdot \struct X(D') .
  \]
\end{theorem}
\begin{proof}
  \leanok
  The class is the finitely supported product \(D \mapsto \prod_x \struct X(x)^{D_x}\). At each
  point \(x\) the exponentiation \(n \mapsto \struct X(x)^n\) is a homomorphism \(\Z \to \Pic(X)\):
  it sends \(0\) to the identity and satisfies \(\struct X(x)^{m+n} = \struct X(x)^m \cdot \struct
  X(x)^n\). Hence, taking the product over the pointwise sum \((D + D')_x = D_x + D'_x\) in the
  abelian group \(\Pic(X)\), the summands separate and the product over \(D + D'\) factors as the
  product over \(D\) times the product over \(D'\), giving \(\struct X(D + D') = \struct X(D) \cdot
  \struct X(D')\).
\end{proof}

\begin{theorem}[Normalization at a one-point divisor]
  \label{pa-thm:divisorClass_single}
  \lean{AlgebraicGeometry.Scheme.divisorClass_single_eq_pointDivisor}
  \leanok
  \dcref{custom}

  \uses{pa-def:divisorClass, pa-def:curveDivisor_single, pa-def:pointDivisor, pa-def:localEquations_picClass}
  For a closed point \(x\), the class of the one-point divisor \(1 \cdot x\)
  (\ref{pa-def:curveDivisor_single}) is the point-divisor class:
  \[
    \struct X(1 \cdot x) = \mathrm{picClass}(1 \cdot x) = \struct X(x) .
  \]
\end{theorem}
\begin{proof}
  \leanok
  The one-point divisor \(1 \cdot x\) is supported at the single point \(x\) with coefficient \(1\)
  (\ref{pa-def:curveDivisor_single}), so the defining product of \(\struct X(1 \cdot x)\)
  (\ref{pa-def:divisorClass}) reduces to the single factor \(\struct X(x)^1 = \struct X(x)\).
\end{proof}

\section{Meromorphic presentations of \v Cech Picard classes}

The remainder of this chapter carries out, at the level of \v Cech cocycles, the classical
identification of the group of Weil divisors modulo principal equivalence with the Picard group
of an integral curve. The device that replaces every sheaf-theoretic gluing argument is the
\emph{generic-point embedding}: on an integral scheme every nonempty open contains the generic
point \(\eta\), and taking germs at \(\eta\) embeds the sections over such an open into the
function field \(K(X)\). A \v Cech Picard class can therefore be \emph{presented meromorphically}
--- by a pointed cover, a unit cocycle on it, and one rational function per member --- after
which every cochain comparison downstream is a literal equality of elements of the field
\(K(X)\), never a gluing of sheaf sections.

\begin{definition}[The generic-point embedding of unit sections]
  \label{pa-def:germGenericUnits}
  \lean{AlgebraicGeometry.Scheme.germGenericUnits,
    AlgebraicGeometry.Scheme.germGenericUnits_unitsRestrict,
    AlgebraicGeometry.Scheme.germGenericUnits_injective}
  \leanok
  \dcref{custom}

  Let \(X\) be an integral scheme with generic point \(\eta\) and \(U\) an open with
  \(\eta \in U\). Taking the germ at \(\eta\) is a homomorphism of groups
  \[
    \mathrm{germ}_\eta \colon \Gamma(X,U)^{\times} \longrightarrow K(X)^{\times},
  \]
  the \emph{generic-point embedding} of unit sections into the units of the function field. It is
  restriction-invariant --- for \(\eta \in W \le U\) it agrees with restriction
  \(\Gamma(X,U)^{\times} \to \Gamma(X,W)^{\times}\) followed by the embedding on \(W\) --- and, on
  an integral scheme, injective.
\end{definition}
\begin{proof}
  \leanok
  The germ at \(\eta\) is a ring homomorphism \(\Gamma(X,U) \to K(X)\) carrying units to units, so
  it restricts to a group homomorphism on the unit groups. Restriction-invariance is the
  functoriality of germs along \(W \le U\). For injectivity, on an integral scheme a section over
  a nonempty open is determined by its germ at \(\eta\) (which lies in every nonempty open); hence
  two unit sections with the same germ at \(\eta\) are equal.
\end{proof}

\begin{lemma}[Pointed covers contain the generic point]
  \label{pa-lem:pointedCover_genericPoint_mem}
  \lean{AlgebraicGeometry.Scheme.PointedCover.genericPoint_mem_opens,
    AlgebraicGeometry.Scheme.PointedCover.genericPoint_mem_inf}
  \leanok
  \dcref{custom}

  On an irreducible scheme \(X\), every member \(\mathcal U(x)\) of a pointed cover contains the
  generic point \(\eta\), and hence so does every pairwise overlap \(\mathcal U(x) \cap \mathcal
  U(y)\).
\end{lemma}
\begin{proof}
  \leanok
  A pointed cover assigns to each \(x\) an open \(\mathcal U(x) \ni x\), in particular nonempty;
  on an irreducible space the generic point lies in every nonempty open, so \(\eta \in \mathcal
  U(x)\). Intersecting two such memberships gives \(\eta \in \mathcal U(x) \cap \mathcal U(y)\).
\end{proof}

\begin{definition}[Meromorphic presentation]
  \label{pa-def:meromorphicPresentation}
  \lean{AlgebraicGeometry.Scheme.MeromorphicPresentation,
    AlgebraicGeometry.Scheme.MeromorphicPresentation.picClass}
  \leanok
  \dcref{custom}

  \uses{pa-def:germGenericUnits, pa-lem:pointedCover_genericPoint_mem}
  A \emph{meromorphic presentation} of a \v Cech Picard class on an integral scheme \(X\) consists
  of a pointed cover \(\mathcal U\), a unit \v Cech \(1\)-cocycle \(\gamma\) on \(\mathcal U\), and
  a family \((f_x)_x\) of units \(f_x \in K(X)^{\times}\), one per member, satisfying the
  \emph{ratio law}
  \[
    f_x \cdot f_y^{-1} = \mathrm{germ}_\eta(\gamma_{xy}) \qquad \text{for every pair } x, y,
  \]
  where \(\gamma_{xy} \in \Gamma(X, \mathcal U(x) \cap \mathcal U(y))^{\times}\) is the cocycle
  value. Its \emph{Picard class} is
  \[
    \mathrm{picClass}(P) = [\mathcal U, \gamma] \in \Pic(X).
  \]
\end{definition}
\begin{proof}
  \leanok
  Each overlap \(\mathcal U(x) \cap \mathcal U(y)\) contains \(\eta\)
  (\ref{pa-lem:pointedCover_genericPoint_mem}), so the germ \(\mathrm{germ}_\eta(\gamma_{xy})\)
  (\ref{pa-def:germGenericUnits}) is defined and the ratio law is a condition in \(K(X)^{\times}\).
  The Picard class is the class in the colimit \(\Pic(X)\) of the unit cocycle \(\gamma\) on
  \(\mathcal U\).
\end{proof}

\begin{theorem}[Existence of meromorphic presentations]
  \label{pa-thm:exists_meromorphicPresentation}
  \lean{AlgebraicGeometry.Scheme.exists_meromorphicPresentation}
  \leanok
  \dcref{custom}

  \uses{pa-def:meromorphicPresentation, pa-def:germGenericUnits, pa-lem:pointedCover_genericPoint_mem}
  Every \v Cech Picard class \(L \in \Pic(X)\) of an integral scheme is the class of a meromorphic
  presentation.
\end{theorem}
\begin{proof}
  \leanok
  Represent \(L\) by a unit \(1\)-cocycle \(\gamma\) on a pointed cover \(\mathcal U\). Use the
  generic point itself as base index: set
  \[
    f_x := \mathrm{germ}_\eta(\gamma_{x\eta}) \in K(X)^{\times},
  \]
  the germ at \(\eta\) of the cocycle value on the pair \((x, \eta)\), whose overlap contains
  \(\eta\) (\ref{pa-lem:pointedCover_genericPoint_mem}). The cocycle identity for the triple
  \((x, y, \eta)\) reads \(\gamma_{xy}\, \gamma_{y\eta} = \gamma_{x\eta}\) on the triple overlap,
  which contains \(\eta\); applying the generic-point embedding, whose value is unchanged by the
  restriction to the triple overlap (\ref{pa-def:germGenericUnits}), and rearranging gives
  \[
    \mathrm{germ}_\eta(\gamma_{xy}) = \mathrm{germ}_\eta(\gamma_{x\eta}) \cdot
      \mathrm{germ}_\eta(\gamma_{y\eta})^{-1} = f_x \cdot f_y^{-1},
  \]
  the ratio law. Thus \((\mathcal U, \gamma, f)\) is a meromorphic presentation
  (\ref{pa-def:meromorphicPresentation}) with class \([\mathcal U, \gamma] = L\).
\end{proof}

\begin{lemma}[The presentation of a local-equation system]
  \label{pa-lem:localEquations_presentation}
  \lean{AlgebraicGeometry.Scheme.LocalEquations.presentation,
    AlgebraicGeometry.Scheme.LocalEquations.presentation_picClass}
  \leanok
  \dcref{custom}

  \uses{pa-def:meromorphicPresentation, pa-def:localEquations, pa-lem:localEquations_ratioUnit,
    pa-def:localEquations_picClass}
  Let \(E\) be a local-equation system on \(X\) (\ref{pa-def:localEquations}). The germs at \(\eta\)
  of its regular equations, \(f_x := \mathrm{germ}_\eta(E.\mathrm{eqn}_x) \in K(X)^{\times}\), form
  a meromorphic presentation whose cocycle is the transition-unit cocycle of \(E\) and whose class
  is the Picard class \(\mathrm{picClass}(E)\) of \(E\) (\ref{pa-def:localEquations_picClass}).
\end{lemma}
\begin{proof}
  \leanok
  Each equation \(E.\mathrm{eqn}_x\) is regular, so its germ at \(\eta\) is a nonzerodivisor of the
  field \(K(X)\), that is, a nonzero element, hence a unit \(f_x \in K(X)^{\times}\). The defining
  equation \(E.\mathrm{eqn}_x = g_{xy} \cdot E.\mathrm{eqn}_y\) on the overlap
  (\ref{pa-lem:localEquations_ratioUnit}), pushed to the generic point, gives \(f_x = \mathrm{germ}_\eta(g_{xy}) \cdot f_y\),
  i.e.\ the ratio law \(f_x \cdot f_y^{-1} = \mathrm{germ}_\eta(g_{xy})\) with the transition-unit
  cocycle \(g\) of \(E\). The presented class is \([\mathcal U, g] = \mathrm{picClass}(E)\)
  (\ref{pa-def:localEquations_picClass}).
\end{proof}

\section{The divisor of a meromorphic presentation}

We now specialize to the curve bundle \(X\) of Chapter~\ref{pa-chap:RiemannRoch}: integral, smooth of
relative dimension one and quasi-compact over the field \(K\), with generic point \(\eta\). The
trivializing elements of a meromorphic presentation cut out a Weil divisor, read off point by
point from their orders of vanishing.

\begin{theorem}[Unit sections have trivial order]
  \label{pa-thm:ordZ_germGenericUnits}
  \lean{AlgebraicGeometry.Scheme.ordZ_germGenericUnits}
  \leanok
  \dcref{custom}

  \uses{pa-def:germGenericUnits, pa-def:scheme_ordZ, pa-lem:ord_eq_one_of_mem_basicOpen}
  For an open \(U \ni \eta\), a unit section \(u \in \Gamma(X,U)^{\times}\), and a closed point
  \(x \in U\), the germ at \(\eta\) of \(u\) has trivial order at \(x\):
  \[
    \mathrm{ord}_x\bigl(\mathrm{germ}_\eta(u)\bigr) = 0.
  \]
\end{theorem}
\begin{proof}
  \leanok
  The section \(u\) is invertible at \(x\), so \(x\) lies in the basic open set where the
  underlying section is invertible. A rational function that is the germ at \(\eta\) of a section
  invertible at \(x\) has order \(0\) there (\ref{pa-lem:ord_eq_one_of_mem_basicOpen}).
\end{proof}

\begin{theorem}[Piece-independence of the order]
  \label{pa-thm:ordZ_elem_eq}
  \lean{AlgebraicGeometry.Scheme.MeromorphicPresentation.ordZ_elem_eq}
  \leanok
  \dcref{custom}

  \uses{pa-def:meromorphicPresentation, pa-thm:ordZ_germGenericUnits, pa-def:scheme_ordZ}
  Let \(P\) be a meromorphic presentation and \(x\) a closed point lying in a member
  \(\mathcal U(y)\) of its cover. Then the trivializing element \(f_y\) of that piece has the same
  order at \(x\) as the element \(f_x\) indexed by \(x\) itself:
  \[
    \mathrm{ord}_x(f_y) = \mathrm{ord}_x(f_x).
  \]
\end{theorem}
\begin{proof}
  \leanok
  By the ratio law, \(f_x \cdot f_y^{-1} = \mathrm{germ}_\eta(\gamma_{xy})\) with \(\gamma_{xy}\) a
  unit section over \(\mathcal U(x) \cap \mathcal U(y)\), an overlap containing the closed point
  \(x\) (which lies in \(\mathcal U(x)\) and, by hypothesis, in \(\mathcal U(y)\)). Its germ at
  \(\eta\) has order \(0\) at \(x\) (\ref{pa-thm:ordZ_germGenericUnits}). Since \(\mathrm{ord}_x\) is a
  homomorphism (\ref{pa-def:scheme_ordZ}), \(\mathrm{ord}_x(f_x) - \mathrm{ord}_x(f_y) =
  \mathrm{ord}_x(f_x \cdot f_y^{-1}) = 0\).
\end{proof}

\begin{theorem}[Finite support of the orders]
  \label{pa-thm:ordZ_elem_support_finite}
  \lean{AlgebraicGeometry.Scheme.MeromorphicPresentation.ordZ_elem_support_finite}
  \leanok
  \dcref{custom}

  \uses{pa-def:meromorphicPresentation, pa-thm:ordZ_elem_eq, pa-thm:ordZ_support_finite}
  For a meromorphic presentation \(P\) on the quasi-compact curve, the set of closed points \(x\)
  at which the piece-indexed element \(f_x\) has nonzero order is finite.
\end{theorem}
\begin{proof}
  \leanok
  The curve is quasi-compact, so the members \(\mathcal U(i)\) of the cover admit a finite
  subcover indexed by a finite set \(T\). If a closed point \(x\) has \(\mathrm{ord}_x(f_x) \neq
  0\), choose \(i \in T\) with \(x \in \mathcal U(i)\); by piece-independence
  (\ref{pa-thm:ordZ_elem_eq}), \(\mathrm{ord}_x(f_x) = \mathrm{ord}_x(f_i) \neq 0\), so \(x\) lies in
  the order support of the single rational function \(f_i\), which is finite
  (\ref{pa-thm:ordZ_support_finite}). Hence the whole support is contained in the finite union over
  \(i \in T\) of these finite sets.
\end{proof}

\begin{definition}[The divisor of a presentation]
  \label{pa-def:presentationDivisor}
  \lean{AlgebraicGeometry.Scheme.presentationDivisor,
    AlgebraicGeometry.Scheme.coeffAt_presentationDivisor}
  \leanok
  \dcref{custom}

  \uses{pa-def:meromorphicPresentation, pa-thm:ordZ_elem_eq, pa-thm:ordZ_elem_support_finite,
    pa-def:curveDivisor, pa-def:scheme_ordZ}
  For a meromorphic presentation \(P\) on the curve, its \emph{divisor} is the Weil divisor
  \[
    \mathrm{presentationDivisor}(P) = \sum_x \mathrm{ord}_x(f_x)\, x \in \Div(X),
  \]
  whose coefficient at each closed point \(x\) is the order at \(x\) of the piece-indexed
  trivializing element \(f_x\).
\end{definition}
\begin{proof}
  \leanok
  The assignment \(x \mapsto \mathrm{ord}_x(f_x)\) has finite support
  (\ref{pa-thm:ordZ_elem_support_finite}), so it defines an element of \(\Div(X)\)
  (\ref{pa-def:curveDivisor}); by piece-independence (\ref{pa-thm:ordZ_elem_eq}) the coefficient is the
  order of the local trivialization near \(x\), whichever piece is used to read it off.
\end{proof}

\section{The calculus of meromorphic presentations}

Meromorphic presentations multiply, invert, and have a unit --- pointwise on the trivializing
elements, and, on the cocycles, by restriction to the common refinement. No group laws are
claimed (the covers of \((P \cdot Q) \cdot R\) and \(P \cdot (Q \cdot R)\) differ by
associativity of intersection), and none are needed: every identification is made at the level of
the Picard class and of the divisor.

\begin{definition}[The principal presentation]
  \label{pa-def:meromorphicPresentation_principal}
  \lean{AlgebraicGeometry.Scheme.MeromorphicPresentation.principal}
  \leanok
  \dcref{custom}

  \uses{pa-def:meromorphicPresentation}
  For \(g \in K(X)^{\times}\), the \emph{principal presentation} of \(g\) is the trivial cocycle on
  the one-member cover \(\{X\}\), trivialized by the constant \(g\) (so \(f_x = g\) for all
  \(x\)).
\end{definition}
\begin{proof}
  \leanok
  The cocycle is trivial, and the ratio law holds because \(f_x \cdot f_y^{-1} = g \cdot g^{-1} =
  1 = \mathrm{germ}_\eta(1)\).
\end{proof}

\begin{theorem}[Multiplicativity of the presented class]
  \label{pa-thm:meromorphicPresentation_picClass_hom}
  \lean{AlgebraicGeometry.Scheme.MeromorphicPresentation.picClass_one,
    AlgebraicGeometry.Scheme.MeromorphicPresentation.picClass_mul,
    AlgebraicGeometry.Scheme.MeromorphicPresentation.picClass_inv,
    AlgebraicGeometry.Scheme.MeromorphicPresentation.picClass_principal}
  \leanok
  \dcref{custom}

  \uses{pa-def:meromorphicPresentation, pa-def:meromorphicPresentation_principal}
  Write \(1\) for the \emph{unit presentation} (the trivial cocycle on \(\{X\}\) with all elements
  \(1\)), \(P \cdot Q\) for the \emph{product} on the common refinement \(\mathcal U_P \cap
  \mathcal U_Q\) (trivializing elements multiplied pointwise, cocycles restricted and multiplied),
  and \(P^{-1}\) for the \emph{inverse} (trivializing elements and cocycle inverted). Then the
  presented class is multiplicative:
  \[
    \mathrm{picClass}(1) = 1, \quad \mathrm{picClass}(P \cdot Q) = \mathrm{picClass}(P) \cdot
      \mathrm{picClass}(Q), \quad \mathrm{picClass}(P^{-1}) = \mathrm{picClass}(P)^{-1},
  \]
  and \(\mathrm{picClass}(\mathrm{principal}\, g) = 1\).
\end{theorem}
\begin{proof}
  \leanok
  The class of a presentation is the class of its cocycle. The trivial cocycle has trivial class,
  giving \(\mathrm{picClass}(1) = 1\); the principal presentation
  (\ref{pa-def:meromorphicPresentation_principal}) also carries the trivial cocycle, so
  \(\mathrm{picClass}(\mathrm{principal}\, g) = 1\). Inverting a cocycle inverts its class,
  giving \(\mathrm{picClass}(P^{-1}) = \mathrm{picClass}(P)^{-1}\). For the product, restricting
  each cocycle to the common refinement leaves its class unchanged, and the class of a product of
  cocycles is the product of the classes, so \(\mathrm{picClass}(P \cdot Q) = \mathrm{picClass}(P)
  \cdot \mathrm{picClass}(Q)\).
\end{proof}

\begin{theorem}[Additivity of the presentation divisor]
  \label{pa-thm:presentationDivisor_hom}
  \lean{AlgebraicGeometry.Scheme.MeromorphicPresentation.presentationDivisor_one,
    AlgebraicGeometry.Scheme.MeromorphicPresentation.presentationDivisor_mul,
    AlgebraicGeometry.Scheme.MeromorphicPresentation.presentationDivisor_inv,
    AlgebraicGeometry.Scheme.MeromorphicPresentation.presentationDivisor_principal}
  \leanok
  \dcref{custom}

  \uses{pa-def:presentationDivisor, pa-def:meromorphicPresentation_principal, pa-def:divOf}
  With the operations of \ref{pa-thm:meromorphicPresentation_picClass_hom}, the divisor is additive:
  \[
    \mathrm{presentationDivisor}(1) = 0, \quad \mathrm{presentationDivisor}(P \cdot Q) =
      \mathrm{presentationDivisor}(P) + \mathrm{presentationDivisor}(Q),
  \]
  \(\mathrm{presentationDivisor}(P^{-1}) = -\mathrm{presentationDivisor}(P)\), and the divisor of
  the principal presentation of \(g\) is the principal divisor:
  \[
    \mathrm{presentationDivisor}(\mathrm{principal}\, g) = \mathrm{div}(g).
  \]
\end{theorem}
\begin{proof}
  \leanok
  Compare coefficients at each closed point \(x\) (\ref{pa-def:presentationDivisor}). The element of
  \(1\) at \(x\) is \(1\), of order \(0\), so \(\mathrm{presentationDivisor}(1) = 0\). The element
  of \(P \cdot Q\) at \(x\) is \(f_x \cdot f'_x\), and \(\mathrm{ord}_x\) is a homomorphism
  (\ref{pa-def:scheme_ordZ}), giving additivity; inversion negates the order. For the principal
  presentation the element at \(x\) is \(g\), so the coefficient is \(\mathrm{ord}_x(g)\), the
  coefficient of \(\mathrm{div}(g)\) (\ref{pa-def:divOf}).
\end{proof}

\section{The class law: equal divisors give equal classes}

The heart of the meromorphic bridge is that the \v Cech Picard class of a presentation depends
only on its divisor. Its local input is that a rational function of trivial order is the germ of a
unit section; its global mechanism is that the cohomology relating two presentations of equal
divisor can be verified entirely inside \(K(X)^{\times}\), through the injective generic-point
embedding --- on an integral scheme no gluing of sheaf sections is ever required.

\begin{lemma}[Order matching produces a unit section]
  \label{pa-lem:exists_germGenericUnits_ordMatch}
  \lean{AlgebraicGeometry.Scheme.exists_germGenericUnits_eq_of_forall_ordZ_eq_one,
    AlgebraicGeometry.Scheme.exists_germGenericUnits_eq_mul_inv_of_forall_ordZ_eq}
  \leanok
  \dcref{custom}

  \uses{pa-def:germGenericUnits, pa-def:scheme_ordZ, pa-def:divisorSections, pa-lem:exists_section_germ_eq}
  Let \(V\) be a nonempty open (\(\eta \in V\)). A rational function \(q \in K(X)^{\times}\) whose
  order vanishes at every closed point of \(V\) is the germ at \(\eta\) of a unit section over
  \(V\): there is \(u \in \Gamma(X,V)^{\times}\) with \(\mathrm{germ}_\eta(u) = q\). Consequently
  two rational functions \(g, h \in K(X)^{\times}\) with equal order at every closed point of
  \(V\) differ by the germ of a unit section: \(\mathrm{germ}_\eta(u) = g \cdot h^{-1}\) for some
  \(u \in \Gamma(X,V)^{\times}\).
\end{lemma}
\begin{proof}
  \leanok
  Since \(\mathrm{ord}_x(q) = 0\) at every closed point \(x \in V\), both \(q\) and \(q^{-1}\) have
  poles bounded by the zero divisor on \(V\), so both lie in \(\mathcal O(0)(V)\)
  (\ref{pa-def:divisorSections}). Each is therefore the germ at \(\eta\) of a section over \(V\)
  (\ref{pa-lem:exists_section_germ_eq}); let \(s, s'\) be sections with \(\mathrm{germ}_\eta(s) = q\)
  and \(\mathrm{germ}_\eta(s') = q^{-1}\). Their product has \(\mathrm{germ}_\eta(s s') = q q^{-1}
  = 1\), and on the integral scheme a section is determined by its germ at \(\eta\), so \(s s' =
  1\); thus \(s\) is a unit section with \(\mathrm{germ}_\eta(s) = q\). The two-element form is
  this applied to \(q = g \cdot h^{-1}\), whose order at each closed point of \(V\) is
  \(\mathrm{ord}_x(g) - \mathrm{ord}_x(h) = 0\).
\end{proof}

\begin{theorem}[The class law]
  \label{pa-thm:picClass_eq_of_presentationDivisor_eq}
  \lean{AlgebraicGeometry.Scheme.MeromorphicPresentation.picClass_eq_of_presentationDivisor_eq}
  \leanok
  \dcref{custom}

  \uses{pa-def:presentationDivisor, pa-thm:ordZ_elem_eq, pa-lem:exists_germGenericUnits_ordMatch,
    pa-def:meromorphicPresentation}
  Two meromorphic presentations with equal divisors present equal \v Cech Picard classes: if
  \(\mathrm{presentationDivisor}(P) = \mathrm{presentationDivisor}(Q)\) then \(\mathrm{picClass}(P)
  = \mathrm{picClass}(Q)\).
\end{theorem}
\begin{proof}
  \leanok
  Equal divisors mean \(\mathrm{ord}_z(f_z) = \mathrm{ord}_z(f'_z)\) at every closed point \(z\),
  comparing coefficients (\ref{pa-def:presentationDivisor}). Fix a member indexed by \(x\) of the
  common refinement \(\mathcal W := \mathcal U_P \cap \mathcal U_Q\) and a closed point \(z \in
  \mathcal W(x)\). Piece-independence (\ref{pa-thm:ordZ_elem_eq}) gives \(\mathrm{ord}_z(f_x) =
  \mathrm{ord}_z(f_z)\) and \(\mathrm{ord}_z(f'_x) = \mathrm{ord}_z(f'_z)\), so \(\mathrm{ord}_z(f_x)
  = \mathrm{ord}_z(f'_x)\): the elements \(f_x\) and \(f'_x\) have equal order at every closed
  point of \(\mathcal W(x)\). By order matching (\ref{pa-lem:exists_germGenericUnits_ordMatch}) there
  is a unit section \(u_x \in \Gamma(X, \mathcal W(x))^{\times}\) with \(\mathrm{germ}_\eta(u_x) =
  f_x \cdot (f'_x)^{-1}\).

  The \(0\)-cochain \((u_x^{-1})_x\) exhibits the restricted cocycles \(\gamma|_{\mathcal W}\) and
  \(\gamma'|_{\mathcal W}\) as cohomologous. The required coboundary identity is, on each overlap
  \(\mathcal W(x) \cap \mathcal W(y)\), an equality of unit sections; applying the injective
  generic-point embedding (\ref{pa-def:germGenericUnits}) turns it into the identity
  \[
    \mathrm{germ}_\eta(u_x)^{-1} \cdot \bigl(f_x f_y^{-1}\bigr) = \bigl(f'_x (f'_y)^{-1}\bigr) \cdot
      \mathrm{germ}_\eta(u_y)^{-1}
  \]
  in the abelian group \(K(X)^{\times}\), where \(f_x f_y^{-1}\) and \(f'_x (f'_y)^{-1}\) are the
  germs of the restricted cocycle values by the ratio law; substituting \(\mathrm{germ}_\eta(u_x) =
  f_x (f'_x)^{-1}\) makes both sides equal. Cohomologous cocycles on \(\mathcal W\) have the same
  class, and restricting a cocycle to a refinement preserves its \v Cech class, so
  \(\mathrm{picClass}(P) = [\mathcal W, \gamma|_{\mathcal W}] = [\mathcal W, \gamma'|_{\mathcal W}]
  = \mathrm{picClass}(Q)\).
\end{proof}

\section{Extraction: equal classes differ by a principal divisor}

\begin{theorem}[Extraction and the meromorphic characterization]
  \label{pa-thm:meromorphicPresentation_extraction}
  \lean{AlgebraicGeometry.Scheme.MeromorphicPresentation.exists_divOf_of_picClass_eq,
    AlgebraicGeometry.Scheme.MeromorphicPresentation.picClass_eq_iff_exists_divOf}
  \leanok
  \dcref{custom}

  \uses{pa-thm:picClass_eq_of_presentationDivisor_eq, pa-thm:meromorphicPresentation_picClass_hom,
    pa-thm:presentationDivisor_hom, pa-def:presentationDivisor, pa-thm:ordZ_germGenericUnits,
    pa-def:meromorphicPresentation_principal, pa-def:divOf}
  Two meromorphic presentations \(P, Q\) present the same \v Cech Picard class if and only if their
  divisors differ by a principal divisor:
  \[
    \mathrm{picClass}(P) = \mathrm{picClass}(Q) \iff \mathrm{presentationDivisor}(P) -
      \mathrm{presentationDivisor}(Q) = \mathrm{div}(g) \text{ for some } g \in K(X)^{\times}.
  \]
\end{theorem}
\begin{proof}
  \leanok
  \((\Rightarrow)\) Equal classes are, on a common refinement \(\mathcal W\), cohomologous
  cocycles: there is a \(0\)-cochain \((\alpha_x)\) of unit sections \(\alpha_x \in \Gamma(X,
  \mathcal W(x))^{\times}\) with \(\alpha_x \cdot \gamma_{xy}|_{\mathcal W} = \gamma'_{xy}|_{\mathcal
  W} \cdot \alpha_y\) on each overlap. Put
  \[
    c_x := f_x \cdot (f'_x)^{-1} \cdot \mathrm{germ}_\eta(\alpha_x) \in K(X)^{\times}.
  \]
  Applying the generic-point embedding to the coboundary identity and rearranging in the abelian
  group \(K(X)^{\times}\) shows \(c_x = c_y\) for all \(x, y\); as the members cover \(X\), the
  \(c_x\) are a single rational function \(g\). At a closed point \(z\), choosing \(x\) with \(z
  \in \mathcal W(x)\), the germ \(\mathrm{germ}_\eta(\alpha_x)\) has order \(0\) at \(z\)
  (\ref{pa-thm:ordZ_germGenericUnits}), so
  \[
    \mathrm{ord}_z(g) = \mathrm{ord}_z(f_z) - \mathrm{ord}_z(f'_z) =
      \mathrm{presentationDivisor}(P)_z - \mathrm{presentationDivisor}(Q)_z .
  \]
  Hence \(\mathrm{div}(g) = \mathrm{presentationDivisor}(P) - \mathrm{presentationDivisor}(Q)\).

  \((\Leftarrow)\) If \(\mathrm{presentationDivisor}(P) - \mathrm{presentationDivisor}(Q) =
  \mathrm{div}(g)\), then by additivity (\ref{pa-thm:presentationDivisor_hom}) the presentation \(Q
  \cdot \mathrm{principal}\, g\) has divisor \(\mathrm{presentationDivisor}(Q) + \mathrm{div}(g) =
  \mathrm{presentationDivisor}(P)\); the class law (\ref{pa-thm:picClass_eq_of_presentationDivisor_eq})
  gives \(\mathrm{picClass}(Q \cdot \mathrm{principal}\, g) = \mathrm{picClass}(P)\), while
  \(\mathrm{picClass}(Q \cdot \mathrm{principal}\, g) = \mathrm{picClass}(Q)\) by multiplicativity
  and \(\mathrm{picClass}(\mathrm{principal}\, g) = 1\)
  (\ref{pa-thm:meromorphicPresentation_picClass_hom}). Hence \(\mathrm{picClass}(P) =
  \mathrm{picClass}(Q)\).
\end{proof}

\section{The point presentation}

To feed the class map with concrete divisors we realize each one-point divisor \(1 \cdot x\) by a
meromorphic presentation whose defining data is tracked: its equation at \(x\) is a uniformizer
whose germ at \(\eta\) is recorded.

\begin{lemma}[A tracked point-uniformizer]
  \label{pa-lem:pointUniformizerData}
  \lean{AlgebraicGeometry.Scheme.PointUniformizerData,
    AlgebraicGeometry.Scheme.nonempty_pointUniformizerData,
    AlgebraicGeometry.Scheme.pointUniformizerData}
  \leanok
  \dcref{custom}

  \uses{pa-lem:uniformizer, pa-lem:exists_stalk_of_ord_le_one, pa-lem:germ_generic_eq_algebraMap_germ,
    pa-thm:ordZ_support_finite, pa-thm:closedpoint_isClosed, pa-def:scheme_ordZ,
    pa-lem:ord_eq_one_of_mem_basicOpen}
  For a closed point \(x\) of the curve, a \emph{tracked point-uniformizer} at \(x\) is an open
  neighbourhood \(V \ni x\) together with a section \(s \in \Gamma(X,V)\) that is regular (its
  germs are nonzerodivisors), whose germ at \(\eta\) is a chosen uniformizer \(t\) at \(x\)
  (\(\mathrm{ord}_x t = 1\)), and which is a unit at every point of \(V\) other than \(x\). Such
  data exists.
\end{lemma}
\begin{proof}
  \leanok
  Choose a uniformizer \(t \in K(X)\) at \(x\) (\ref{pa-lem:uniformizer}), a nonzero rational function
  with \(\mathrm{ord}_x t = 1\). Having nonnegative order at \(x\), \(t\) is integral there, so it
  is the germ at \(\eta\) of a section \(s_0\) over a neighbourhood \(W_0 \ni x\)
  (\ref{pa-lem:exists_stalk_of_ord_le_one}, \ref{pa-lem:germ_generic_eq_algebraMap_germ}). Remove from
  \(W_0\) the finitely many closed points other than \(x\) at which \(t\) has nonzero order
  (finite by \ref{pa-thm:ordZ_support_finite}, and a finite union of closed points is closed by
  \ref{pa-thm:closedpoint_isClosed}); the result \(V := W_0 \setminus \mathrm{bad}\) is an open
  neighbourhood of \(x\), and the restriction \(s\) of \(s_0\) to \(V\) still has germ \(t\) at
  \(\eta\).

  \emph{Regularity.} For \(y \in V\), the germ of \(s\) maps to \(t \neq 0\) under the injective
  localization \(\struct{X,y} \to K(X)\), hence is a nonzerodivisor.

  \emph{Unit away from \(x\).} Let \(y \in V\), \(y \neq x\). If \(y = \eta\), then
  \(\mathrm{germ}_\eta(s) = t\) is a unit of the field. If \(y\) is a closed point and the germ of
  \(s\) at \(y\) were not a unit, then \(t\) would have nonzero order at \(y\) (a section invertible
  at a closed point has order \(0\) there, \ref{pa-lem:ord_eq_one_of_mem_basicOpen}, contrapositively,
  reading order through \ref{pa-def:scheme_ordZ}); but then \(y\) is a removed point, contradicting
  \(y \in V\). So \(s\) is a unit at \(y\).
\end{proof}

\begin{definition}[The local equations of the point divisor]
  \label{pa-def:pointEquations}
  \lean{AlgebraicGeometry.Scheme.pointCover, AlgebraicGeometry.Scheme.pointEquations,
    AlgebraicGeometry.Scheme.germGeneric_pointEqn_self}
  \leanok
  \dcref{custom}

  \uses{pa-lem:pointUniformizerData, pa-def:localEquations, pa-thm:closedpoint_isClosed}
  From a tracked point-uniformizer \(d\) at \(x\) (\ref{pa-lem:pointUniformizerData}), the
  \emph{point equations} of \(1 \cdot x\) are the local-equation system on the pointed cover
  \[
    \mathcal U(z) = \begin{cases} V & z = x, \\ X \setminus \{x\} & z \neq x, \end{cases}
  \]
  whose equation on the member at \(x\) is the tracked section \(s\) and whose equation on every
  other member is the constant \(1\); the germ at \(\eta\) of the equation at \(x\) is the chosen
  uniformizer.
\end{definition}
\begin{proof}
  \leanok
  The complement \(X \setminus \{x\}\) is open (\ref{pa-thm:closedpoint_isClosed}) and contains every
  \(z \neq x\), so \(\mathcal U\) is a pointed cover. Regularity holds member by member: on \(V\)
  the germs of \(s\) are nonzerodivisors (\ref{pa-lem:pointUniformizerData}), and the constant \(1\)
  is a unit. For overlap ratios, a pair both at \(x\) or both away from \(x\) has equal equations,
  ratio \(1\); a mixed pair has overlap inside \(V \setminus \{x\}\), where \(s\) is a unit
  (\ref{pa-lem:pointUniformizerData}), so the ratio of \(s\) and \(1\) is a unit. Thus the data is a
  local-equation system (\ref{pa-def:localEquations}), and the germ at \(\eta\) of the equation at
  \(x\) is \(\mathrm{germ}_\eta(s) = t\).
\end{proof}

\begin{theorem}[The divisor of the point presentation]
  \label{pa-thm:presentationDivisor_pointEquations}
  \lean{AlgebraicGeometry.Scheme.presentationDivisor_pointEquations}
  \leanok
  \dcref{custom}

  \uses{pa-def:pointEquations, pa-lem:localEquations_presentation, pa-def:presentationDivisor,
    pa-def:curveDivisor_single, pa-lem:uniformizer}
  The divisor of the meromorphic presentation of the point equations at a closed point \(x\) is
  the one-point divisor:
  \[
    \mathrm{presentationDivisor}\bigl((\text{point equations at } x)\text{'s presentation}\bigr) =
      1 \cdot x .
  \]
\end{theorem}
\begin{proof}
  \leanok
  The presentation of the point equations (\ref{pa-lem:localEquations_presentation}) has, at each
  closed point \(z\), trivializing element the germ at \(\eta\) of the equation at \(z\). At \(z =
  x\) this is the uniformizer \(t\), of order \(\mathrm{ord}_x t = 1\) (\ref{pa-lem:uniformizer}), so
  the coefficient at \(x\) is \(1\); at \(z \neq x\) the equation is the constant \(1\), of order
  \(0\), so the coefficient vanishes. Hence the divisor is \(1 \cdot x\)
  (\ref{pa-def:curveDivisor_single}, \ref{pa-def:presentationDivisor}).
\end{proof}

\section{The \v Cech Picard class of a Weil divisor}

Assembling the realization of divisors, the class law, and extraction yields the total class map
\(D \mapsto \struct X(D)\) from Weil divisors to \v Cech Picard classes, and its kernel: the class
map descends to an isomorphism of \(\Div(X)\) modulo principal divisors onto \(\Pic(X)\).

\begin{theorem}[Realization of Weil divisors]
  \label{pa-thm:exists_presentationDivisor_eq}
  \lean{AlgebraicGeometry.Scheme.exists_presentationDivisor_eq}
  \leanok
  \dcref{custom}

  \uses{pa-thm:presentationDivisor_pointEquations, pa-thm:presentationDivisor_hom,
    pa-lem:pointUniformizerData, pa-def:curveDivisor_single, pa-lem:curveDivisor_single_arith}
  Every Weil divisor \(D \in \Div(X)\) on the curve is the divisor of a meromorphic presentation.
\end{theorem}
\begin{proof}
  \leanok
  First realize the one-point divisors \(n \cdot x\). For \(n \ge 0\), induct on \(n\): the divisor
  \(0 \cdot x = 0\) is realized by the unit presentation (\(\mathrm{presentationDivisor}(1) = 0\),
  \ref{pa-thm:presentationDivisor_hom}), and \((n+1) \cdot x = n \cdot x + 1 \cdot x\) by the product
  of a realization of \(n \cdot x\) with the point presentation at \(x\), using additivity
  (\ref{pa-thm:presentationDivisor_hom}), the point-divisor anchor
  (\ref{pa-thm:presentationDivisor_pointEquations}, \ref{pa-lem:pointUniformizerData}), and the one-point
  arithmetic (\ref{pa-lem:curveDivisor_single_arith}). Negative multiplicities use the inverse
  presentation, whose divisor is negated (\ref{pa-thm:presentationDivisor_hom}). A general \(D\) is a
  finite sum of one-point divisors (\ref{pa-def:curveDivisor_single}), realized by the product of
  their presentations, again by additivity.
\end{proof}

\begin{definition}[The class map and its surjectivity]
  \label{pa-def:curveDivisor_picClass}
  \lean{AlgebraicGeometry.Scheme.CurveDivisor.picClass,
    AlgebraicGeometry.Scheme.CurveDivisor.picClass_presentationDivisor,
    AlgebraicGeometry.Scheme.CurveDivisor.exists_picClass_eq,
    AlgebraicGeometry.Scheme.CurveDivisor.picClass_surjective}
  \leanok
  \dcref{custom}

  \uses{pa-thm:exists_presentationDivisor_eq, pa-thm:picClass_eq_of_presentationDivisor_eq,
    pa-thm:exists_meromorphicPresentation, pa-def:curveDivisor}
  The \emph{\v Cech Picard class} \(\struct X(D) = \mathrm{picClass}(D) \in \Pic(X)\) of a Weil
  divisor \(D\) is the class of any meromorphic presentation whose divisor is \(D\). It is
  well-defined, satisfies \(\mathrm{picClass}(\mathrm{presentationDivisor}(P)) = \mathrm{picClass}(P)\)
  for every presentation \(P\), and the map \(D \mapsto \struct X(D)\) is surjective onto
  \(\Pic(X)\).
\end{definition}
\begin{proof}
  \leanok
  A realizing presentation exists (\ref{pa-thm:exists_presentationDivisor_eq}). If \(P, P'\) both have
  divisor \(D\), the class law (\ref{pa-thm:picClass_eq_of_presentationDivisor_eq}) gives
  \(\mathrm{picClass}(P) = \mathrm{picClass}(P')\), so \(\mathrm{picClass}(D)\) is well-defined;
  taking \(D = \mathrm{presentationDivisor}(P)\) with the realization \(P\) itself gives
  \(\mathrm{picClass}(\mathrm{presentationDivisor}(P)) = \mathrm{picClass}(P)\). For surjectivity,
  every class \(L\) is presented by some \(P\) (\ref{pa-thm:exists_meromorphicPresentation}), and then
  \(\struct X(\mathrm{presentationDivisor}(P)) = \mathrm{picClass}(P) = L\).
\end{proof}

\begin{theorem}[The class map is a homomorphism trivial on principal divisors]
  \label{pa-thm:curveDivisor_picClass_hom}
  \lean{AlgebraicGeometry.Scheme.CurveDivisor.picClass_zero,
    AlgebraicGeometry.Scheme.CurveDivisor.picClass_add,
    AlgebraicGeometry.Scheme.CurveDivisor.picClass_divOf}
  \leanok
  \dcref{custom}

  \uses{pa-def:curveDivisor_picClass, pa-thm:presentationDivisor_hom,
    pa-thm:meromorphicPresentation_picClass_hom, pa-def:divOf}
  The class map is a group homomorphism \(\Div(X) \to \Pic(X)\), trivial on principal divisors:
  \[
    \struct X(0) = 1, \quad \struct X(D + D') = \struct X(D) \cdot \struct X(D'), \quad \struct
    X(\mathrm{div}\, g) = 1 \ (g \in K(X)^{\times}).
  \]
\end{theorem}
\begin{proof}
  \leanok
  Realizing divisors by presentations and using
  \(\mathrm{picClass}(\mathrm{presentationDivisor}(P)) = \mathrm{picClass}(P)\)
  (\ref{pa-def:curveDivisor_picClass}): \(\struct X(0) = \mathrm{picClass}(1) = 1\) since
  \(\mathrm{presentationDivisor}(1) = 0\) (\ref{pa-thm:presentationDivisor_hom},
  \ref{pa-thm:meromorphicPresentation_picClass_hom}). If \(P, P'\) realize \(D, D'\), then \(P \cdot
  P'\) realizes \(D + D'\) (\ref{pa-thm:presentationDivisor_hom}), so \(\struct X(D + D') =
  \mathrm{picClass}(P \cdot P') = \mathrm{picClass}(P) \cdot \mathrm{picClass}(P') = \struct X(D)
  \cdot \struct X(D')\) (\ref{pa-thm:meromorphicPresentation_picClass_hom}). Finally
  \(\mathrm{div}\, g = \mathrm{presentationDivisor}(\mathrm{principal}\, g)\)
  (\ref{pa-thm:presentationDivisor_hom}), so \(\struct X(\mathrm{div}\, g) =
  \mathrm{picClass}(\mathrm{principal}\, g) = 1\) (\ref{pa-thm:meromorphicPresentation_picClass_hom}).
\end{proof}

\begin{theorem}[The kernel is the principal divisors]
  \label{pa-thm:curveDivisor_picClass_eq_iff}
  \lean{AlgebraicGeometry.Scheme.CurveDivisor.exists_divOf_of_picClass_eq,
    AlgebraicGeometry.Scheme.CurveDivisor.picClass_eq_iff}
  \leanok
  \dcref{custom}

  \uses{pa-def:curveDivisor_picClass, pa-thm:meromorphicPresentation_extraction,
    pa-thm:curveDivisor_picClass_hom, pa-def:divOf}
  Two Weil divisors have the same \v Cech Picard class if and only if they differ by a principal
  divisor:
  \[
    \struct X(D) = \struct X(D') \iff D - D' = \mathrm{div}(g) \text{ for some } g \in
      K(X)^{\times}.
  \]
  Together with surjectivity (\ref{pa-def:curveDivisor_picClass}) this identifies \(\Pic(X)\) with the
  group of Weil divisors modulo principal divisors, the classical isomorphism
  \(\mathrm{CaCl}(X) = \Pic(X)\) on an integral curve.
\end{theorem}
\begin{proof}
  \leanok
  \((\Rightarrow)\) Realize \(D, D'\) by presentations \(P, P'\), so \(\mathrm{picClass}(P) =
  \struct X(D) = \struct X(D') = \mathrm{picClass}(P')\); extraction
  (\ref{pa-thm:meromorphicPresentation_extraction}) yields \(g\) with
  \(\mathrm{presentationDivisor}(P) - \mathrm{presentationDivisor}(P') = \mathrm{div}(g)\), that is
  \(D - D' = \mathrm{div}(g)\). \((\Leftarrow)\) If \(D - D' = \mathrm{div}(g)\), then \(D = D' +
  \mathrm{div}(g)\), so \(\struct X(D) = \struct X(D') \cdot \struct X(\mathrm{div}\, g) = \struct
  X(D')\) by additivity and triviality on principal divisors
  (\ref{pa-thm:curveDivisor_picClass_hom}). Since the class map is onto \(\Pic(X)\) with fibres the
  cosets of the subgroup of principal divisors, it descends to an isomorphism of \(\Div(X)\) modulo
  principal divisors onto \(\Pic(X)\).
\end{proof}

\begin{theorem}[Normalization at a one-point divisor]
  \label{pa-thm:curveDivisor_picClass_single}
  \lean{AlgebraicGeometry.Scheme.CurveDivisor.picClass_single}
  \leanok
  \dcref{custom}

  \uses{pa-def:curveDivisor_picClass, pa-thm:presentationDivisor_pointEquations,
    pa-lem:localEquations_presentation, pa-def:curveDivisor_single}
  For a closed point \(x\), the class of the one-point divisor \(1 \cdot x\) is the class of the
  tracked point equations at \(x\):
  \[
    \struct X(1 \cdot x) = \mathrm{picClass}(\text{point equations at } x).
  \]
\end{theorem}
\begin{proof}
  \leanok
  The presentation of the point equations at \(x\) has divisor \(1 \cdot x\)
  (\ref{pa-thm:presentationDivisor_pointEquations}), so \(\struct X(1 \cdot x)\) is its class
  (\ref{pa-def:curveDivisor_picClass}), which is the class of the point equations
  (\ref{pa-lem:localEquations_presentation}, \ref{pa-def:curveDivisor_single}).
\end{proof}

\section{The point-divisor class map and the meromorphic class map coincide}

The chapter carries two constructions of the divisor-to-class assignment \(D \mapsto \struct
X(D)\). The first, \(\mathrm{divisorClass}\) (\ref{pa-def:divisorClass}), is the finitely supported
product of the point-divisor classes \(\struct X(x) = \mathrm{picClass}(1 \cdot x)\) read off the
point local equations (\ref{pa-def:pointDivisor}); the second, the anchored
\(\mathrm{picClass}\) (\ref{pa-def:curveDivisor_picClass}), is the class of any meromorphic
presentation realizing \(D\). This closing section identifies them. The bridge is that the point
local equations, whose equation at \(x\) has generic-point germ the uniformizer
(\ref{pa-def:pointDivisor}), present the one-point divisor \(1 \cdot x\); their two classes agree, and
the agreement propagates from one-point divisors to all divisors by homomorphism extensionality.

\begin{theorem}[The divisor of the point-divisor presentation]
  \label{pa-thm:presentationDivisor_pointDivisor}
  \lean{AlgebraicGeometry.Scheme.presentationDivisor_pointDivisor}
  \leanok
  \dcref{custom}

  \uses{pa-def:pointDivisor, pa-lem:localEquations_presentation, pa-def:presentationDivisor,
    pa-def:curveDivisor_single, pa-lem:uniformizer}
  The divisor of the meromorphic presentation of the point-divisor local equations at a closed
  point \(x\) is the one-point divisor:
  \[
    \mathrm{presentationDivisor}\bigl((1 \cdot x)\text{'s presentation}\bigr) = 1 \cdot x .
  \]
\end{theorem}
\begin{proof}
  \leanok
  Compare coefficients at each closed point \(z\) (\ref{pa-def:presentationDivisor}). The presentation
  of the point-divisor equations (\ref{pa-lem:localEquations_presentation}) has, at \(z\), trivializing
  element the germ at \(\eta\) of the equation at \(z\). At \(z = x\) that germ is the chosen
  uniformizer \(t\) (\ref{pa-def:pointDivisor}), of order \(\mathrm{ord}_x t = 1\)
  (\ref{pa-lem:uniformizer}), so the coefficient at \(x\) is \(1\); at \(z \neq x\) the equation is the
  constant \(1\), of order \(0\), so the coefficient vanishes. Hence the divisor is \(1 \cdot x\)
  (\ref{pa-def:curveDivisor_single}).
\end{proof}

\begin{theorem}[The point-divisor class is the anchored one-point class]
  \label{pa-thm:pointDivisor_picClass_eq}
  \lean{AlgebraicGeometry.Scheme.pointDivisor_picClass_eq}
  \leanok
  \dcref{custom}

  \uses{pa-def:pointDivisor, pa-def:curveDivisor_single, pa-thm:presentationDivisor_pointDivisor,
    pa-def:curveDivisor_picClass, pa-lem:localEquations_presentation}
  For a closed point \(x\), the Picard class of the point-divisor local equations at \(x\) equals
  the anchored class of the one-point divisor \(1 \cdot x\):
  \[
    \mathrm{picClass}(1 \cdot x) \;=\; \struct X(1 \cdot x).
  \]
\end{theorem}
\begin{proof}
  \leanok
  The presentation of the point-divisor equations has divisor \(1 \cdot x\)
  (\ref{pa-thm:presentationDivisor_pointDivisor}), so the anchored class \(\struct X(1 \cdot x)\) is
  the class of that presentation (\ref{pa-def:curveDivisor_picClass}); and the class of the
  presentation of a local-equation system is the Picard class of the system itself
  (\ref{pa-lem:localEquations_presentation}), which is \(\mathrm{picClass}(1 \cdot x)\)
  (\ref{pa-def:pointDivisor}). Chaining the two equalities gives \(\mathrm{picClass}(1 \cdot x) =
  \struct X(1 \cdot x)\).
\end{proof}

\begin{theorem}[The two divisor-class maps coincide]
  \label{pa-thm:divisorClass_eq_picClass}
  \lean{AlgebraicGeometry.Scheme.divisorClass_eq_picClass}
  \leanok
  \dcref{custom}

  \uses{pa-def:divisorClass, pa-def:curveDivisor_picClass, pa-thm:divisorClass_add,
    pa-thm:curveDivisor_picClass_hom, pa-thm:pointDivisor_picClass_eq, pa-def:curveDivisor_single}
  The point-product class map (\ref{pa-def:divisorClass}) and the anchored meromorphic class map
  (\ref{pa-def:curveDivisor_picClass}) agree on every Weil divisor:
  \[
    \mathrm{divisorClass}(D) \;=\; \struct X(D), \qquad D \in \Div(X).
  \]
\end{theorem}
\begin{proof}
  \leanok
  Both assignments are homomorphisms \(\Div(X) \to \Pic(X)\) from the additive group of divisors to
  the multiplicative Picard group: \(\mathrm{divisorClass}\) by additivity
  (\ref{pa-thm:divisorClass_add}) and the anchored class map by
  Theorem~\ref{pa-thm:curveDivisor_picClass_hom}. Since \(\Div(X)\) is freely generated by the
  one-point divisors (\ref{pa-def:curveDivisor_single}), it suffices to check agreement on each
  \(n \cdot x\). On the point-product side the defining product (\ref{pa-def:divisorClass}) collapses
  to the single factor \(\mathrm{divisorClass}(n \cdot x) = \mathrm{picClass}(1 \cdot x)^{n}\). On
  the anchored side, additivity (\ref{pa-thm:curveDivisor_picClass_hom}) gives \(\struct X(n \cdot x) =
  \struct X(1 \cdot x)^{n}\). The two base classes are equal, \(\mathrm{picClass}(1 \cdot x) =
  \struct X(1 \cdot x)\) (\ref{pa-thm:pointDivisor_picClass_eq}), so their \(n\)-th powers agree and
  \(\mathrm{divisorClass}(n \cdot x) = \struct X(n \cdot x)\). By homomorphism extensionality on the
  generators, the two maps coincide on all of \(\Div(X)\).
\end{proof}

\section{The unit of the \'etale sheafification}
\label{pa-sec:picEtUnit}

The honest relative Picard functor maps into its \'etale-sheafified Layer-2 vehicle
\(\mathrm{picEt}\) by a natural transformation — the unit of the sheafification. On an affine
test it is the unit of the plus construction; on a general test it is assembled over the
affine opens.

\begin{definition}[Restriction of plus classes along an explicit algebra map]
  \label{pa-def:PicEtAff_mapAlg}
  \lean{AlgebraicGeometry.PicEtAff.mapAlg, AlgebraicGeometry.PicEtAff.mapAlg_id,
    AlgebraicGeometry.PicEtAff.mapAlg_comp, AlgebraicGeometry.PicEtAff.mapAlg_unit}
  \leanok
  \dcref{custom}

  \uses{pa-def:PicEtAff_map, pa-thm:PicEtAff_map_id, pa-thm:PicEtAff_map_map, pa-thm:PicEtAff_map_unit,
    pa-def:PicEtAff_unit, pa-def:relPicAlgMap}
  For an explicit homomorphism of \(k\)-algebras \(\varphi \colon A \to A'\), restriction of
  plus classes \(\varphi^{*} \colon \mathrm{PicEtAff}(C,A) \to \mathrm{PicEtAff}(C,A')\):
  regard \(A'\) as an \(A\)-algebra through \(\varphi\) and apply the base-extension map of
  Definition~\ref{pa-def:PicEtAff_map}. It satisfies \((\id_A)^{*} = \id\) and
  \((\psi \circ \varphi)^{*} = \psi^{*} \circ \varphi^{*}\), and it is natural on units of
  the plus construction:
  \[
    \varphi^{*}\bigl(\mathrm{unit}(x)\bigr) = \mathrm{unit}\bigl(\varphi^{*} x\bigr),
    \qquad x \in \mathrm{relPic}(C,A).
  \]
\end{definition}
\begin{proof}
  \leanok
  The scalar tower \(k \to A \to A'\) determined by \(\varphi\) makes
  Definition~\ref{pa-def:PicEtAff_map} applicable; the two functor laws are
  Theorems~\ref{pa-thm:PicEtAff_map_id} and~\ref{pa-thm:PicEtAff_map_map} read through this tower,
  and naturality on units is Theorem~\ref{pa-thm:PicEtAff_map_unit}.
\end{proof}

\begin{definition}[The affine-open test objects of a test object]
  \label{pa-def:fromSpecAffine}
  \lean{AlgebraicGeometry.Over.fromSpecAffine, AlgebraicGeometry.Over.fromSpecAffine_naturality,
    AlgebraicGeometry.Over.fromSpecAffine_resAlgHom,
    AlgebraicGeometry.Over.overSpecMap_ΓTop_fromSpecAffine_top}
  \leanok
  \dcref{custom}

  \uses{pa-def:overSpecMap, pa-def:overSpecGammaTop}
  For a test object \(T\) and an affine open \(U \subseteq T_{\mathrm{left}}\), the spectrum
  of the section ring \(\Spec \Gamma(T, U)\) is a test object over \(T\), via the canonical
  morphism \(\Spec \Gamma(T,U) \to T\) of the affine open. This assignment is coherent:
  \begin{enumerate}
  \item (restriction) for \(U \le V\), the map of \(U\) factors through that of \(V\) along
    \(\Spec\) of the section restriction \(\Gamma(T,V) \to \Gamma(T,U)\);
  \item (naturality) for a morphism of tests \(f \colon T' \to T\) and affine opens
    \(W \le f^{-1}V\), the square formed by \(\Spec\) of the section pullback
    \(\Gamma(T,V) \to \Gamma(T',W)\) and the two affine-open test maps commutes;
  \item (top of an affine test) for \(T = \Spec A\), the affine-open test object of the top
    open, composed with \(\Spec\) of the identification \(\Gamma(\Spec A, \top) \cong A\)
    (Definition~\ref{pa-def:overSpecGammaTop}), is the identity of \(\Spec A\).
  \end{enumerate}
\end{definition}
\begin{proof}
  \leanok
  The structure morphism to \(\Spec k\) is \(\Spec\) of the \(k\)-algebra structure of the
  section ring, so the canonical morphism of the affine open lies over \(\Spec k\); this is
  the compatibility of the affine-open inclusion with the structure morphisms. Coherences (i)
  and (ii) are both instances of the standard identity
  \(\Spec(f^{\sharp}_{V,W}) \circ \iota_W = \iota_V \circ f\) for the canonical maps of
  affine opens, applied to the identity of \(T\) and to \(f\) respectively; (iii) is the
  fact that the canonical map of the top affine open of an affine scheme is the isomorphism
  \(\Spec \Gamma(\Spec A, \top) \cong \Spec A\), inverse to \(\Spec\) of the section
  identification.
\end{proof}

\begin{definition}[The component of the unit at a test object]
  \label{pa-def:relPicToPicEt}
  \lean{AlgebraicGeometry.relPicToPicEt, AlgebraicGeometry.relPicToPicEt_val}
  \leanok
  \dcref{custom}

  \uses{pa-def:picEt, pa-def:PicEtAff_unit, pa-def:fromSpecAffine, pa-def:relPicMap, pa-def:PicEtAff_mapAlg}
  For a test object \(T\), the homomorphism
  \(\mathrm{relPic}(C,T) \to \mathrm{picEt}(C,T)\) sending a relative Picard class \(z\) to
  the family whose value at an affine open \(U \subseteq T_{\mathrm{left}}\) is the plus unit
  of the restriction of \(z\) along the affine-open test object of \(U\):
  \[
    z \;\longmapsto\;
    \Bigl( U \mapsto \mathrm{unit}\bigl(\iota_U^{*} z\bigr) \Bigr)_{U} .
  \]
\end{definition}
\begin{proof}
  \leanok
  The family is compatible under restriction of affine opens, as required by
  Definition~\ref{pa-def:picEt}: for \(U \le V\), the restriction coherence of the affine-open
  test objects (Definition~\ref{pa-def:fromSpecAffine}) gives \(\iota_U = \iota_V \circ
  \Spec(\mathrm{res}_{V,U})\), so functoriality of restriction (Definition~\ref{pa-def:relPicMap})
  and naturality of the unit on explicit algebra maps (Definition~\ref{pa-def:PicEtAff_mapAlg})
  identify \(\mathrm{res}_{V,U}^{*}\,\mathrm{unit}(\iota_V^{*}z) =
  \mathrm{unit}(\iota_U^{*}z)\). The assignment is a homomorphism because each constituent
  map is.
\end{proof}

\begin{theorem}[Affine consistency of the unit]
  \label{pa-thm:picEtAffineEquiv_relPicToPicEt}
  \lean{AlgebraicGeometry.picEtAffineEquiv_relPicToPicEt}
  \leanok
  \dcref{custom}

  \uses{pa-def:relPicToPicEt, pa-def:picEtAffineEquiv, pa-def:PicEtAff_mapAlg, pa-def:fromSpecAffine,
    pa-def:relPicMap, pa-def:PicEtAff_unit}
  On an affine test \(\Spec A\), the component of the unit collapses through the affine
  comparison \(\mathrm{picEt}(C, \Spec A) \cong \mathrm{PicEtAff}(C, A)\)
  (Definition~\ref{pa-def:picEtAffineEquiv}) to the unit of the plus construction:
  \[
    \mathrm{picEtAffineEquiv}\bigl(\mathrm{relPicToPicEt}(z)\bigr)
      = \mathrm{unit}(z), \qquad z \in \mathrm{relPic}(C,A).
  \]
\end{theorem}
\begin{proof}
  \leanok
  The affine comparison evaluates the family at the top affine open and transports along the
  section identification \(\Gamma(\Spec A, \top) \cong A\). The value there is
  \(\mathrm{unit}(\iota_\top^{*} z)\); transporting along the identification and using
  naturality of the unit (Definition~\ref{pa-def:PicEtAff_mapAlg}) yields the unit of
  \(z\) restricted along \(\Spec(\Gamma(\Spec A,\top) \cong A) \circ \iota_\top\), which is
  the identity of \(\Spec A\) by coherence (iii) of Definition~\ref{pa-def:fromSpecAffine}; so
  the value is \(\mathrm{unit}(z)\) by functoriality (Definition~\ref{pa-def:relPicMap}).
\end{proof}

\begin{theorem}[Naturality of the unit]
  \label{pa-thm:picEtMap_relPicToPicEt}
  \lean{AlgebraicGeometry.picEtMap_relPicToPicEt}
  \leanok
  \dcref{custom}

  \uses{pa-def:relPicToPicEt, pa-def:picEtMap, pa-def:picEt_IsPullbackValue,
    pa-thm:picEt_pullbackValue_unique, pa-def:fromSpecAffine, pa-def:PicEtAff_mapAlg, pa-def:relPicMap}
  Let the curve \(C\) be proper, geometrically irreducible and geometrically reduced. For a
  morphism of test objects \(f \colon T' \to T\) and \(z \in \mathrm{relPic}(C,T)\),
  \[
    f^{*}\bigl(\mathrm{relPicToPicEt}(z)\bigr)
      = \mathrm{relPicToPicEt}\bigl(f^{*} z\bigr)
    \quad\text{in } \mathrm{picEt}(C,T').
  \]
\end{theorem}
\begin{proof}
  \leanok
  The restriction map of \(\mathrm{picEt}\) is characterised value-by-value by the
  \(\exists!\)-property of Definition~\ref{pa-def:picEt_IsPullbackValue} together with its
  uniqueness (Theorem~\ref{pa-thm:picEt_pullbackValue_unique}); it therefore suffices to verify
  that the right-hand family has the characterising property of the pullback of the
  left-hand family. Unwinding, this is the statement that for affine opens
  \(W_0 \le W \subseteq T'_{\mathrm{left}}\) with \(W_0 \le f^{-1}V\),
  \[
    \mathrm{res}^{*}\,\mathrm{unit}\bigl(\iota_W^{*} f^{*} z\bigr)
      = (f^{\sharp}_{V,W_0})^{*}\,\mathrm{unit}\bigl(\iota_V^{*} z\bigr).
  \]
  Both sides are units of restrictions of \(z\) (naturality of the unit,
  Definition~\ref{pa-def:PicEtAff_mapAlg}), and the two restrictions agree because the two
  composites of affine-open test maps agree: \(\iota_{W_0} \circ (\text{res}) = \iota_W\) and
  \(f \circ \iota_{W_0} = \iota_V \circ \Spec(f^{\sharp}_{V,W_0})\), the two coherences of
  Definition~\ref{pa-def:fromSpecAffine}.
\end{proof}

\begin{definition}[The unit of the \'etale sheafification]
  \label{pa-def:picEtUnit}
  \lean{AlgebraicGeometry.picEtUnit, AlgebraicGeometry.picEtUnit_app,
    AlgebraicGeometry.picEtAffineEquiv_picEtUnit_app}
  \leanok
  \dcref{custom}

  \uses{pa-def:relPicToPicEt, pa-thm:picEtMap_relPicToPicEt, pa-def:picEtFunctor, pa-def:relPicMap,
    pa-thm:picEtAffineEquiv_relPicToPicEt}
  Let the curve \(C\) be proper, geometrically irreducible and geometrically reduced. The
  components of Definition~\ref{pa-def:relPicToPicEt} assemble into a natural transformation
  \[
    \mathrm{relPicFunctor}(C) \;\longrightarrow\; \mathrm{picEtFunctor}(C)
  \]
  of presheaves of commutative groups on \(\Over(\Spec k)\), whose affine components collapse
  to the unit of the plus construction.
\end{definition}
\begin{proof}
  \leanok
  Naturality of the components is Theorem~\ref{pa-thm:picEtMap_relPicToPicEt}; the affine
  collapse is Theorem~\ref{pa-thm:picEtAffineEquiv_relPicToPicEt}.
\end{proof}

\section{The degree of a plus class}
\label{pa-sec:degAff}

The relative degree \(\mathrm{relPicDeg}\) of Definition~\ref{pa-def:relPicDeg} reads an honest
integer off a relative Picard class over a \emph{field}. This section descends it to the plus
construction: a plus class over \(K\) is represented on an \'etale cover, every cover is
refined by a single finite separable field extension (cofinality,
Theorem~\ref{pa-thm:etale_cover_field_cofinality}), and the \emph{zigzag} shows the field
reading is independent of every choice.

\begin{lemma}[The relative degree is additive]
  \label{pa-lem:relPicDeg_hom}
  \lean{AlgebraicGeometry.relPicDeg_mul, AlgebraicGeometry.relPicDeg_one}
  \leanok
  \dcref{custom}

  \uses{pa-def:relPicDeg}
  For a field extension \(K\) of \(k\) and relative Picard classes
  \(x, y \in \mathrm{relPic}(C, K)\),
  \[
    \mathrm{relPicDeg}_K(x \cdot y) = \mathrm{relPicDeg}_K(x) + \mathrm{relPicDeg}_K(y),
    \qquad \mathrm{relPicDeg}_K(1) = 0 .
  \]
\end{lemma}
\begin{proof}
  \leanok
  \(\mathrm{relPicDeg}_K\) is a homomorphism from the multiplicative relative Picard group to
  \(\Z\) by construction (Definition~\ref{pa-def:relPicDeg}); these are its two homomorphism
  laws.
\end{proof}

\begin{theorem}[The zigzag, leg (i): one degree through any two field readings]
  \label{pa-thm:relPicDeg_relPicAlgMap_congr}
  \lean{AlgebraicGeometry.relPicDeg_relPicAlgMap_congr}
  \leanok
  \dcref{custom}

  \uses{pa-def:relPicDeg, pa-thm:relPicDeg_relPicAlgMap, pa-thm:relPicAlgMap_congr,
    pa-thm:etale_cover_field_cofinality, pa-def:etale_cover_prod, pa-def:etale_cover_prod_coproj,
    pa-def:etale_cover_ofField, pa-def:descentClasses, pa-lem:relPicAlgMap_functor}
  Let \(K\) be a field extension of \(k\), \(E\) a presented \'etale cover of \(K\), and
  \(x \in \mathrm{relPic}(C, B_E)\) a descent class on \(E\). Let \(L_1, L_2\) be finite
  separable field extensions of \(K\) and \(j_1 \colon B_E \to L_1\),
  \(j_2 \colon B_E \to L_2\) two \(K\)-algebra maps. Then
  \[
    \mathrm{relPicDeg}_{L_1}\bigl(j_1^{*} x\bigr)
      = \mathrm{relPicDeg}_{L_2}\bigl(j_2^{*} x\bigr).
  \]
\end{theorem}
\begin{proof}
  \leanok
  Refine the product of the two field covers \(\mathrm{ofField}(L_1) \times
  \mathrm{ofField}(L_2)\) (Definitions~\ref{pa-def:etale_cover_ofField},
  \ref{pa-def:etale_cover_prod}) to a single finite separable field extension \(M\) of \(K\)
  (cofinality, Theorem~\ref{pa-thm:etale_cover_field_cofinality}); composing with the
  coprojections (Definition~\ref{pa-def:etale_cover_prod_coproj}) yields \(K\)-algebra maps
  \(\theta_1 \colon L_1 \to M\) and \(\theta_2 \colon L_2 \to M\). Then
  \[
    \mathrm{relPicDeg}_{L_1}(j_1^{*} x)
      = \mathrm{relPicDeg}_{M}(\theta_1^{*} j_1^{*} x)
      = \mathrm{relPicDeg}_{M}\bigl((\theta_1 j_1)^{*} x\bigr),
  \]
  using invariance of the relative degree under base-field transitions (E-iv-alg descended,
  Theorem~\ref{pa-thm:relPicDeg_relPicAlgMap}) and functoriality
  (Lemma~\ref{pa-lem:relPicAlgMap_functor}); likewise for the second leg. Now \(\theta_1 j_1\)
  and \(\theta_2 j_2\) are two \(K\)-algebra maps \(B_E \to M\) out of the cover carrier, and
  \(x\) is a descent class, so the transport-independence keystone of the plus construction
  (Theorem~\ref{pa-thm:relPicAlgMap_congr}) gives \((\theta_1 j_1)^{*} x = (\theta_2 j_2)^{*} x\).
  The two degrees over \(M\) therefore agree, hence so do the original readings.
\end{proof}

\begin{theorem}[The zigzag: well-definedness of the plus-class degree]
  \label{pa-thm:relPicDeg_eq_of_descentMap_eq}
  \lean{AlgebraicGeometry.relPicDeg_eq_of_descentMap_eq}
  \leanok
  \dcref{custom}

  \uses{pa-thm:relPicDeg_relPicAlgMap_congr, pa-def:descentMap, pa-thm:etale_cover_field_cofinality,
    pa-lem:relPicAlgMap_functor, pa-def:descentClasses}
  Let \(E_1, E_2, G\) be presented \'etale covers of \(K\), \(x_1, x_2\) descent classes on
  \(E_1, E_2\), and \(f \colon B_{E_1} \to B_G\), \(g \colon B_{E_2} \to B_G\) \(K\)-algebra
  maps with \(f_{*} x_1 = g_{*} x_2\) as descent classes on \(G\) (the identification of the
  plus-construction setoid). Then for any two field readings — finite separable field
  extensions \(L_1, L_2\) of \(K\) with \(K\)-algebra maps \(j_1 \colon B_{E_1} \to L_1\) and
  \(j_2 \colon B_{E_2} \to L_2\) —
  \[
    \mathrm{relPicDeg}_{L_1}\bigl(j_1^{*} x_1\bigr)
      = \mathrm{relPicDeg}_{L_2}\bigl(j_2^{*} x_2\bigr).
  \]
  No Galois theory appears: Galois twisting of the transport is a special case of the
  transport-independence keystone.
\end{theorem}
\begin{proof}
  \leanok
  Refine \(G\) to a single finite separable field extension \(N\) of \(K\) by cofinality
  (Theorem~\ref{pa-thm:etale_cover_field_cofinality}), with refinement map
  \(\ell \colon B_G \to N\). By leg (i) of the zigzag
  (Theorem~\ref{pa-thm:relPicDeg_relPicAlgMap_congr}) applied to the descent class \(x_1\) on
  \(E_1\) and the two readings \(j_1\) and \(\ell f\),
  \[
    \mathrm{relPicDeg}_{L_1}(j_1^{*} x_1) = \mathrm{relPicDeg}_{N}\bigl((\ell f)^{*} x_1\bigr)
      = \mathrm{relPicDeg}_{N}\bigl(\ell^{*} f^{*} x_1\bigr)
  \]
  (functoriality, Lemma~\ref{pa-lem:relPicAlgMap_functor}). Cross over along the identification
  \(f^{*} x_1 = g^{*} x_2\) (the underlying relative classes of \(f_{*}x_1 = g_{*}x_2\),
  Definition~\ref{pa-def:descentMap}) and run the same two steps backwards on the second leg:
  \(\mathrm{relPicDeg}_{N}(\ell^{*} g^{*} x_2) = \mathrm{relPicDeg}_{N}((\ell g)^{*} x_2) =
  \mathrm{relPicDeg}_{L_2}(j_2^{*} x_2)\), again by leg (i).
\end{proof}

\begin{definition}[The degree of a plus class]
  \label{pa-def:degAff}
  \lean{AlgebraicGeometry.PicEtAff.degAff, AlgebraicGeometry.PicEtAff.degAff_mk,
    AlgebraicGeometry.PicEtAff.degAff_mk_ofField}
  \leanok
  \dcref{custom}

  \uses{pa-def:PicEtAff, pa-def:PicEtAff_mk, pa-thm:relPicDeg_eq_of_descentMap_eq,
    pa-thm:relPicDeg_relPicAlgMap_congr, pa-thm:etale_cover_field_cofinality, pa-def:relPicDeg,
    pa-def:etale_cover_ofField}
  The \emph{degree} \(\deg^{+}_K \colon \mathrm{PicEtAff}(C, K) \to \Z\) of an
  \'etale-locally trivial relative Picard class: represent the class by a descent class
  \(x\) on a cover \(E\), refine \(E\) to a finite separable field extension \(L/K\) by
  cofinality (Theorem~\ref{pa-thm:etale_cover_field_cofinality}), and read
  \(\mathrm{relPicDeg}_L\) of the transported class. The value is independent of the
  representative and of the refinement, and it is pinned by the \emph{master equality}: for
  every descent-class representative \((E, x)\), every finite separable field extension
  \(L/K\) and every \(K\)-algebra map \(j \colon B_E \to L\),
  \[
    \deg^{+}_K [E, x] = \mathrm{relPicDeg}_L\bigl(j^{*} x\bigr).
  \]
  In particular, on a field-cover representative \((\mathrm{ofField}(L), y)\) the degree is
  \(\mathrm{relPicDeg}_L\) of \(y\) transported through the carrier identification.
\end{definition}
\begin{proof}
  \leanok
  Existence of a field refinement of each cover is cofinality
  (Theorem~\ref{pa-thm:etale_cover_field_cofinality}). Independence of the representative is
  exactly the zigzag (Theorem~\ref{pa-thm:relPicDeg_eq_of_descentMap_eq}): two representatives
  of one plus class are identified after transport to a common refinement
  (Definition~\ref{pa-def:PicEtAff_mk}), and the zigzag equates any two field readings of such
  a pair — in particular the readings through the chosen refinements. The master equality is
  leg (i) of the zigzag (Theorem~\ref{pa-thm:relPicDeg_relPicAlgMap_congr}), comparing the
  chosen refinement map with the given \(j\); the field-cover normalization is the master
  equality at the carrier identification of \(\mathrm{ofField}(L)\)
  (Definition~\ref{pa-def:etale_cover_ofField}).
\end{proof}

\begin{theorem}[Additivity of the plus-class degree]
  \label{pa-thm:degAff_hom}
  \lean{AlgebraicGeometry.PicEtAff.degAff_one, AlgebraicGeometry.PicEtAff.degAff_mul,
    AlgebraicGeometry.PicEtAff.degAff_inv}
  \leanok
  \dcref{custom}

  \uses{pa-def:degAff, pa-def:PicEtAff_one, pa-thm:PicEtAff_mk_mul_mk, pa-lem:PicEtAff_ind,
    pa-def:etale_cover_prod, pa-def:etale_cover_prod_coproj, pa-thm:etale_cover_field_cofinality,
    pa-lem:relPicDeg_hom, pa-def:etale_cover_self, pa-lem:relPicAlgMap_functor}
  \(\deg^{+}_K\) is a homomorphism: \(\deg^{+}_K(1) = 0\),
  \(\deg^{+}_K(a b) = \deg^{+}_K(a) + \deg^{+}_K(b)\) and
  \(\deg^{+}_K(a^{-1}) = -\deg^{+}_K(a)\).
\end{theorem}
\begin{proof}
  \leanok
  \emph{Unit.} The trivial plus class is represented by the trivial descent class on the
  trivial cover (Definitions~\ref{pa-def:PicEtAff_one}, \ref{pa-def:etale_cover_self}); its field
  reading is \(\mathrm{relPicDeg}_K(1) = 0\) (Lemma~\ref{pa-lem:relPicDeg_hom}), and the master
  equality (Definition~\ref{pa-def:degAff}) transfers this to \(\deg^{+}\).

  \emph{Additivity.} Represent \(a = [E,x]\) and \(b = [F,y]\)
  (Lemma~\ref{pa-lem:PicEtAff_ind}); the product is represented on the product cover
  \(E \times F\) by the product of the transported classes
  (Theorem~\ref{pa-thm:PicEtAff_mk_mul_mk}). Refine \(E \times F\) to a single finite separable
  field \(L\) (Theorem~\ref{pa-thm:etale_cover_field_cofinality}) with refinement map \(\ell\),
  and read all three degrees there through the master equality: the reading of \(ab\) is
  \(\mathrm{relPicDeg}_L\) of the product of the two transported classes (transport is a
  homomorphism, Lemma~\ref{pa-lem:relPicAlgMap_functor}), which splits as the sum by
  Lemma~\ref{pa-lem:relPicDeg_hom}; and the two summands are the readings of \(a\) and \(b\)
  through \(\ell\) composed with the coprojections
  (Definition~\ref{pa-def:etale_cover_prod_coproj}).

  \emph{Inverses.} From \(\deg^{+}(a^{-1}) + \deg^{+}(a) = \deg^{+}(1) = 0\).
\end{proof}

\begin{theorem}[The unit normalization of the plus-class degree]
  \label{pa-thm:degAff_unit}
  \lean{AlgebraicGeometry.PicEtAff.degAff_unit}
  \leanok
  \dcref{custom}

  \uses{pa-def:degAff, pa-def:PicEtAff_unit, pa-def:etale_cover_self,
    pa-lem:PicEtAff_tautological_mem_descentClasses, pa-lem:relPicAlgMap_functor}
  On a class in the image of the unit of the plus construction, the degree is the relative
  degree over the base field itself:
  \[
    \deg^{+}_K\bigl(\mathrm{unit}(z)\bigr) = \mathrm{relPicDeg}_K(z),
    \qquad z \in \mathrm{relPic}(C,K).
  \]
\end{theorem}
\begin{proof}
  \leanok
  The unit represents its value by the tautological descent class of \(z\) on the trivial
  cover (Definition~\ref{pa-def:PicEtAff_unit},
  Lemma~\ref{pa-lem:PicEtAff_tautological_mem_descentClasses}), and the trivial cover is
  refined by the field cover of \(K\) itself. The master equality
  (Definition~\ref{pa-def:degAff}) reads \(\mathrm{relPicDeg}_K\) of \(z\) transported along the
  composite \(K \to B_{\mathrm{self}} \to K\), which is the identity
  (Lemma~\ref{pa-lem:relPicAlgMap_functor}).
\end{proof}

\section{The degree at field points and the degree-zero Picard functor}
\label{pa-sec:pic0}

\begin{definition}[The degree of an \'etale Picard class at a field point]
  \label{pa-def:degAt}
  \lean{AlgebraicGeometry.degAt}
  \leanok
  \dcref{custom}

  \uses{pa-def:degAff, pa-def:picEtAffineEquiv, pa-def:picEtMap}
  Let \(T\) be a test object, \(\lambda \in \mathrm{picEt}(C,T)\), \(K\) a field extension of
  \(k\) and \(t \colon \Spec K \to T\) a \emph{field point} of \(T\). The \emph{degree of
  \(\lambda\) at \(t\)} is the plus-class degree of the restriction:
  \[
    \deg_t(\lambda) \;:=\; \deg^{+}_K\bigl(\mathrm{picEtAffineEquiv}(t^{*}\lambda)\bigr) \in \Z ,
  \]
  the restriction along \(t\) (Definition~\ref{pa-def:picEtMap}) collapsed to a plus class by
  the affine comparison (Definition~\ref{pa-def:picEtAffineEquiv}).
\end{definition}

\begin{lemma}[The degree at a field point is additive]
  \label{pa-lem:degAt_hom}
  \lean{AlgebraicGeometry.degAt_one, AlgebraicGeometry.degAt_mul, AlgebraicGeometry.degAt_inv}
  \leanok
  \dcref{custom}

  \uses{pa-def:degAt, pa-thm:degAff_hom}
  For classes \(\lambda, \mu \in \mathrm{picEt}(C,T)\) and a field point \(t\),
  \[
    \deg_t(1) = 0, \qquad
    \deg_t(\lambda\mu) = \deg_t(\lambda) + \deg_t(\mu), \qquad
    \deg_t(\lambda^{-1}) = -\deg_t(\lambda).
  \]
\end{lemma}
\begin{proof}
  \leanok
  Restriction and the affine comparison are homomorphisms, and \(\deg^{+}\) is additive
  (Theorem~\ref{pa-thm:degAff_hom}).
\end{proof}

\begin{theorem}[Functoriality of the degree at field points]
  \label{pa-thm:degAt_picEtMap}
  \lean{AlgebraicGeometry.degAt_picEtMap, AlgebraicGeometry.degAt_picEtFunctor_map}
  \leanok
  \dcref{custom}

  \uses{pa-def:degAt, pa-thm:picEtMap_functor_laws}
  For \(f \colon T' \to T\), \(\lambda \in \mathrm{picEt}(C,T)\) and a field point
  \(t' \colon \Spec K \to T'\),
  \[
    \deg_{t'}\bigl(f^{*}\lambda\bigr) = \deg_{f \circ t'}(\lambda).
  \]
\end{theorem}
\begin{proof}
  \leanok
  Both sides collapse the restriction of \(\lambda\) along the same composite field point:
  \(t'^{*}(f^{*}\lambda) = (f \circ t')^{*}\lambda\) by functoriality of restriction
  (Theorem~\ref{pa-thm:picEtMap_functor_laws}).
\end{proof}

\begin{definition}[The degree-zero subgroup]
  \label{pa-def:pic0Subgroup}
  \lean{AlgebraicGeometry.pic0Subgroup, AlgebraicGeometry.mem_pic0Subgroup_iff}
  \leanok
  \dcref{custom}

  \uses{pa-def:degAt, pa-lem:degAt_hom}
  For a test object \(T\), the subgroup
  \[
    \mathrm{pic}^0(C,T) \;:=\;
    \bigl\{\, \lambda \in \mathrm{picEt}(C,T) \;:\; \deg_t(\lambda) = 0
      \text{ for every field point } t \text{ of } T \,\bigr\}
  \]
  of classes of degree zero at \emph{every} field point of \(T\).
\end{definition}
\begin{proof}
  \leanok
  Closure under unit, product and inverse is Lemma~\ref{pa-lem:degAt_hom}, pointwise in the
  field point.
\end{proof}

\begin{definition}[The degree-zero relative Picard functor]
  \label{pa-def:pic0Functor}
  \lean{AlgebraicGeometry.pic0Map, AlgebraicGeometry.pic0Map_coe,
    AlgebraicGeometry.pic0Functor, AlgebraicGeometry.pic0Functor_obj,
    AlgebraicGeometry.pic0Functor_map, AlgebraicGeometry.pic0Inclusion}
  \leanok
  \dcref{custom}

  \uses{pa-def:pic0Subgroup, pa-thm:degAt_picEtMap, pa-def:picEtFunctor, pa-thm:picEtMap_functor_laws}
  The degree-zero condition is stable under restriction: for \(f \colon T' \to T\), the
  restriction map of \(\mathrm{picEt}\) carries \(\mathrm{pic}^0(C,T)\) into
  \(\mathrm{pic}^0(C,T')\). The functor
  \[
    \mathrm{Pic}^0_{C/k} \colon \Over(\Spec k)^{\mathrm{op}} \longrightarrow \mathrm{CommGrp},
    \qquad T \mapsto \mathrm{pic}^0(C,T),
  \]
  with these restrictions, is the \emph{degree-zero relative Picard functor} — a subgroup
  subfunctor of \(\mathrm{picEtFunctor}(C)\), with inclusion natural transformation
  \(\mathrm{Pic}^0_{C/k} \to \mathrm{picEtFunctor}(C)\). This is the carrier of the
  representability statement: the Jacobian represents \emph{this} functor.
\end{definition}
\begin{proof}
  \leanok
  Stability: for \(\lambda\) of degree zero at every field point of \(T\) and \(t'\) a field
  point of \(T'\), functoriality (Theorem~\ref{pa-thm:degAt_picEtMap}) gives
  \(\deg_{t'}(f^{*}\lambda) = \deg_{f \circ t'}(\lambda) = 0\). The functor laws are
  inherited from those of \(\mathrm{picEtFunctor}\)
  (Theorem~\ref{pa-thm:picEtMap_functor_laws}) on underlying classes, and the inclusion is
  natural because the restriction of a degree-zero class \emph{is} the restriction of its
  underlying class.
\end{proof}

\section{Functoriality of the \'etale-sheafified Picard functor in the curve}

Throughout this section \(k\) is a field and \(D, E, F\) are objects of
\(\Over(\Spec k)\); for the statements involving restriction along arbitrary morphisms
of test objects the curves are in addition proper, geometrically irreducible and
geometrically reduced over \(k\).

\begin{definition}[Curve-direction transport of the limit]
  \label{pa-def:picEtPullback}
  \lean{AlgebraicGeometry.picEtPullback, AlgebraicGeometry.picEtPullback_val}
  \leanok
  \dcref{custom}

  \uses{pa-def:picEt, pa-def:PicEtAff_curveMap, pa-def:PicEtAff_mapAlg}
  Let \(g \colon D \to E\) be a morphism of \(\Over(\Spec k)\) and \(T\) a test object.
  For a compatible family \(s = (s_U)_U \in \mathrm{picEt}(E, T)\) over the affine opens
  of \(T\), the componentwise transported family
  \[
    g^{*} s := \bigl( g^{*}(s_U) \bigr)_{U},
  \]
  each component transported by the affine curve transport
  \(g^{*} \colon \mathrm{PicEtAff}(E, \Gamma(T, U)) \to
  \mathrm{PicEtAff}(D, \Gamma(T, U))\) of Definition~\ref{pa-def:PicEtAff_curveMap}, is
  again compatible, and the assignment
  \(g^{*} \colon \mathrm{picEt}(E, T) \to \mathrm{picEt}(D, T)\) is a group
  homomorphism.
\end{definition}
\begin{proof}
  \leanok
  \uses{pa-thm:PicEtAff_curveMap_laws}
  Compatibility: for affine opens \(U \subseteq V\) of \(T\), the restriction of the
  transported value is
  \(\mathrm{res}^V_U\bigl(g^{*}(s_V)\bigr) = g^{*}\bigl(\mathrm{res}^V_U(s_V)\bigr)
  = g^{*}(s_U)\), by the bilateral restriction square of
  Theorem~\ref{pa-thm:PicEtAff_curveMap_laws} applied to the \(k\)-algebra map
  \(\mathrm{res}^V_U\). The homomorphism property holds componentwise.
\end{proof}

\begin{theorem}[Functor laws and affine compatibility]
  \label{pa-thm:picEtPullback_laws}
  \lean{AlgebraicGeometry.picEtPullback_id, AlgebraicGeometry.picEtPullback_comp,
    AlgebraicGeometry.picEtAffineEquiv_picEtPullback,
    AlgebraicGeometry.picEtPullback_picEtAffineEquiv_symm}
  \leanok
  \dcref{custom}

  \uses{pa-def:picEtPullback, pa-def:picEtAffineEquiv}
  The curve-direction transport of the limit is contravariantly functorial in the curve:
  \(\id_E^{*} = \id\), and \((g' \circ g)^{*} = g^{*} \circ g'^{*}\) for
  \(g \colon D \to E\) and \(g' \colon E \to F\). At an affine test \(\Spec A\) it
  agrees with the affine curve transport under the affine comparison of
  Definition~\ref{pa-def:picEtAffineEquiv}:
  \[
    \mathrm{picEtAffineEquiv}_{D,A} \circ g^{*}
      = g^{*} \circ \mathrm{picEtAffineEquiv}_{E,A},
  \]
  and likewise with the inverse comparisons.
\end{theorem}
\begin{proof}
  \leanok
  \uses{pa-thm:PicEtAff_curveMap_laws}
  The functor laws hold componentwise by Theorem~\ref{pa-thm:PicEtAff_curveMap_laws}. For
  the affine compatibility: both sides evaluate the family at the top affine open and
  transport, and the \(\Gamma\)--\(\Spec\) identification of the test algebra is a
  \(k\)-algebra map, so it commutes with the curve transport by the bilateral
  restriction square. The inverse form follows because both comparisons are
  isomorphisms.
\end{proof}

\begin{theorem}[Naturality of the curve transport in the test object]
  \label{pa-thm:picEtMap_picEtPullback}
  \lean{AlgebraicGeometry.picEtMap_picEtPullback}
  \leanok
  \dcref{custom}

  \uses{pa-def:picEtPullback, pa-def:picEtMap, pa-def:picEt_IsPullbackValue,
    pa-thm:picEt_pullbackValue_unique}
  Let \(D, E\) be proper, geometrically irreducible and geometrically reduced over
  \(k\), let \(g \colon D \to E\), and let \(f \colon T' \to T\) be an arbitrary
  morphism of test objects. Then restriction commutes with the curve transport:
  \[
    f^{*} \circ g^{*} = g^{*} \circ f^{*}
    \colon \mathrm{picEt}(E, T) \longrightarrow \mathrm{picEt}(D, T').
  \]
\end{theorem}
\begin{proof}
  \leanok
  \uses{pa-thm:PicEtAff_curveMap_laws}
  The value of the restricted family \(f^{*}\) at each affine open of \(T'\) is
  characterized uniquely by its pullback property
  (Theorem~\ref{pa-thm:picEt_pullbackValue_unique}), so it suffices to show that the curve
  transport of the glued value for \(s\) has the pullback property for the transported
  family \(g^{*}s\). On every affine sub-open lying in the preimage of an affine open
  \(V\) of \(T\), the required identity is obtained from the pullback property of the
  glued value for \(s\) by applying \(g^{*}\) and commuting it past the two restriction
  \(k\)-algebra maps by the bilateral square of
  Theorem~\ref{pa-thm:PicEtAff_curveMap_laws}.
\end{proof}

\begin{definition}[The bundled curve-direction natural transformation]
  \label{pa-def:picEtPullbackNat}
  \lean{AlgebraicGeometry.picEtPullbackNat, AlgebraicGeometry.picEtPullbackNat_app,
    AlgebraicGeometry.picEtPullbackNat_id, AlgebraicGeometry.picEtPullbackNat_comp}
  \leanok
  \dcref{custom}

  \uses{pa-def:picEtFunctor, pa-def:picEtPullback, pa-thm:picEtMap_picEtPullback}
  For \(D, E\) proper, geometrically irreducible and geometrically reduced over \(k\)
  and \(g \colon D \to E\), the componentwise transports assemble into a natural
  transformation
  \[
    g^{*} \colon \mathrm{picEtFunctor}(E) \longrightarrow \mathrm{picEtFunctor}(D)
  \]
  of \(\mathrm{CommGrp}\)-valued functors on the test category, with
  \(\id_E^{*}\) the identity transformation and
  \((g' \circ g)^{*} = g^{*} \circ g'^{*}\).
\end{definition}
\begin{proof}
  \leanok
  \uses{pa-thm:picEtPullback_laws}
  Naturality is Theorem~\ref{pa-thm:picEtMap_picEtPullback}; the identity and composition
  laws are the pointwise laws of Theorem~\ref{pa-thm:picEtPullback_laws} at each test
  object.
\end{proof}

\begin{theorem}[Degree of a pulled-back class at a field point]
  \label{pa-thm:degAt_picEtPullback}
  \lean{AlgebraicGeometry.degAt_picEtPullback}
  \leanok
  \dcref{custom}

  \uses{pa-def:degAt, pa-thm:picEtMap_picEtPullback, pa-thm:picEtPullback_laws}
  Let \(D\) be smooth of relative dimension \(1\), proper and geometrically irreducible
  over \(k\), let \(E\) be proper, geometrically irreducible and geometrically reduced
  over \(k\), let \(g \colon D \to E\), let \(\lambda \in \mathrm{picEt}(E, T)\), and
  let \(t \colon \Spec K \to T\) be a field point of the test object. Then the degree
  of the pulled-back class at \(t\) is the degree of the affine curve transport of the
  collapsed restriction of \(\lambda\):
  \[
    \deg_{t}\bigl(g^{*}\lambda\bigr)
      = \deg\Bigl( g^{*}\bigl(\mathrm{picEtAffineEquiv}_{E,K}(t^{*}\lambda)\bigr) \Bigr).
  \]
\end{theorem}
\begin{proof}
  \leanok
  By definition, \(\deg_t(g^{*}\lambda)\) is the degree of the plus class obtained by
  restricting \(g^{*}\lambda\) along \(t\) and collapsing by the affine comparison.
  Restriction along \(t\) commutes with the curve transport
  (Theorem~\ref{pa-thm:picEtMap_picEtPullback}), and at the affine test \(\Spec K\) the
  transport is the affine curve transport under the comparison
  (Theorem~\ref{pa-thm:picEtPullback_laws}).
\end{proof}

\section{Degree-zero preservation under the curve transport}

Throughout this section \(k\) is a field, \(D\) and \(E\) (and, where a composite
appears, \(F\)) are curve bundles over \(k\) — objects of \(\Over(\Spec k)\), proper,
smooth of relative dimension one and geometrically irreducible over \(k\) — and
\(g \colon D \to E\) is an arbitrary morphism of \(\Over(\Spec k)\). Pullback along
\(g\) preserves degree zero. The degree control is obtained \emph{over each field
independently}: over a field extension \(L\) of \(k\) the base-changed morphism is
finite surjective or constant, and in either case a degree-zero class pulls back to a
degree-zero class — the multiplier of the finite case is never compared across
different fields.

\begin{theorem}[Pullback preserves degree zero on every fibre]
  \label{pa-thm:classDeg_cechPicMap_whiskerRight_eq_zero}
  \lean{AlgebraicGeometry.classDeg_cechPicMap_whiskerRight_eq_zero}
  \leanok
  \dcref{custom}

  \uses{pa-def:classDeg, pa-thm:baseChangeCurve_bundle}
  Let \(L\) be a field extension of \(k\) and let \(\Lambda\) be a \v Cech Picard class
  on the fibre \((E \otimes \Spec L)_{\mathrm{left}}\) with
  \(\deg_{\Pic}(\Lambda) = 0\). Then the pullback of \(\Lambda\) along the whisker
  \((g \rhd \Spec L) \colon (D \otimes \Spec L)_{\mathrm{left}} \to
  (E \otimes \Spec L)_{\mathrm{left}}\) satisfies
  \(\deg_{\Pic}\bigl((g \rhd \Spec L)^{*}\Lambda\bigr) = 0\).
\end{theorem}
\begin{proof}
  \leanok
  \uses{pa-thm:curveHom_isFinite_or_constant, pa-thm:classDeg_cechPicMap_eq_finrank_mul,
    pa-cor:cechPicMap_eq_one_of_range_eq_singleton, pa-thm:overHom_isProper,
    pa-lem:genericPoint_eq_of_surjective}
  Both fibres \(D_L := (D \otimes \Spec L)_{\mathrm{left}}\) and \(E_L\) are proper,
  smooth of relative dimension one and geometrically irreducible over \(L\) via their
  second projections (Theorem~\ref{pa-thm:baseChangeCurve_bundle}), and the whisker
  commutes with the second projections, so it is a morphism of curve bundles over
  \(L\). By the dichotomy (Theorem~\ref{pa-thm:curveHom_isFinite_or_constant}) it is
  either finite and surjective, or constant with a single closed point as set-range.
  In the finite surjective case the generic points match
  (Lemma~\ref{pa-lem:genericPoint_eq_of_surjective}) and degree multiplicativity
  (Theorem~\ref{pa-thm:classDeg_cechPicMap_eq_finrank_mul}) gives
  \(\deg_{\Pic}\bigl((g \rhd \Spec L)^{*}\Lambda\bigr)
    = \bigl[L(D_L) : L(E_L)\bigr] \cdot \deg_{\Pic}(\Lambda) = 0\).
  In the constant case the whisker is proper (Theorem~\ref{pa-thm:overHom_isProper}), in
  particular quasi-compact, so the pulled-back class is already trivial
  (Corollary~\ref{pa-cor:cechPicMap_eq_one_of_range_eq_singleton}), of degree zero.
\end{proof}

\begin{theorem}[Relative-degree-zero preservation]
  \label{pa-thm:relPicDeg_relPicMapCurveApp_eq_zero}
  \lean{AlgebraicGeometry.relPicDeg_relPicMapCurveApp_eq_zero}
  \leanok
  \dcref{custom}

  \uses{pa-def:relPicDeg, pa-lem:relPicMapCurve}
  Let \(L\) be a field extension of \(k\) and \(y \in \mathrm{relPic}(E, \Spec L)\) a
  relative Picard class with \(\mathrm{relPicDeg}_L(y) = 0\). Then the curve transport
  satisfies \(\mathrm{relPicDeg}_L\bigl(g^{*}y\bigr) = 0\).
\end{theorem}
\begin{proof}
  \leanok
  \uses{pa-thm:classDeg_cechPicMap_whiskerRight_eq_zero}
  Choose a representative \v Cech class \(\Lambda\) of \(y\); then
  \((g \rhd \Spec L)^{*}\Lambda\) represents \(g^{*}y\)
  (Lemma~\ref{pa-lem:relPicMapCurve}), and the relative degree reads the class degree of
  any representative (Definition~\ref{pa-def:relPicDeg}). Conclude by
  Theorem~\ref{pa-thm:classDeg_cechPicMap_whiskerRight_eq_zero}.
\end{proof}

\begin{theorem}[The affine curve transport preserves degree zero]
  \label{pa-thm:degAff_curveMap_eq_zero}
  \lean{AlgebraicGeometry.PicEtAff.degAff_curveMap_eq_zero}
  \leanok
  \dcref{custom}

  \uses{pa-def:degAff, pa-def:PicEtAff_curveMap}
  Let \(K\) be a field extension of \(k\) and \(a \in \mathrm{PicEtAff}(E, K)\) a plus
  class with \(\deg^{+}_K(a) = 0\). Then \(\deg^{+}_K\bigl(g^{*}a\bigr) = 0\).
\end{theorem}
\begin{proof}
  \leanok
  \uses{pa-thm:etale_cover_field_cofinality, pa-thm:PicEtAff_curveMap_laws,
    pa-lem:relPicMapCurve, pa-thm:relPicDeg_relPicMapCurveApp_eq_zero}
  Write \(a\) as the class of a descent-class representative \(x\) on an \'etale cover
  \(U\) of \(K\), so that \(g^{*}a\) is the class of the cover-wise transport
  \(g^{*}x\) on the same cover (Theorem~\ref{pa-thm:PicEtAff_curveMap_laws}). Choose a
  finite separable field refinement \(j \colon B_U \to L\) of the cover
  (Theorem~\ref{pa-thm:etale_cover_field_cofinality}). By the master equality of the plus
  degree (Definition~\ref{pa-def:degAff}),
  \(\deg^{+}_K(a) = \mathrm{relPicDeg}_L\bigl(j^{*}x\bigr)\) and
  \(\deg^{+}_K(g^{*}a) = \mathrm{relPicDeg}_L\bigl(j^{*}(g^{*}x)\bigr)\). The
  transports along \(j\) and along the curve morphism commute
  (Lemma~\ref{pa-lem:relPicMapCurve}), so the latter equals
  \(\mathrm{relPicDeg}_L\bigl(g^{*}(j^{*}x)\bigr)\), which vanishes by
  Theorem~\ref{pa-thm:relPicDeg_relPicMapCurveApp_eq_zero} applied over \(L\) to the class
  \(j^{*}x\), of relative degree \(\deg^{+}_K(a) = 0\).
\end{proof}

\begin{definition}[Curve transport of the degree-zero functor]
  \label{pa-def:pic0Pullback}
  \lean{AlgebraicGeometry.pic0Pullback, AlgebraicGeometry.pic0Pullback_coe}
  \leanok
  \dcref{custom}

  \uses{pa-def:pic0Subgroup, pa-def:picEtPullback, pa-def:degAt}
  For a test object \(T\), the curve transport of Definition~\ref{pa-def:picEtPullback}
  restricts to the degree-zero subgroups: a group homomorphism
  \(g^{*} \colon \mathrm{pic}^0(E, T) \to \mathrm{pic}^0(D, T)\) whose underlying map
  of \'etale Picard classes is the curve transport.
\end{definition}
\begin{proof}
  \leanok
  \uses{pa-thm:degAt_picEtPullback, pa-thm:degAff_curveMap_eq_zero}
  Let \(\lambda \in \mathrm{pic}^0(E,T)\) and let \(t \colon \Spec K \to T\) be a field
  point of the test object. The degree of \(g^{*}\lambda\) at \(t\) is the plus degree
  of the affine curve transport of the collapsed restriction \(t^{*}\lambda\)
  (Theorem~\ref{pa-thm:degAt_picEtPullback}); the collapsed restriction has plus degree
  \(\deg_t(\lambda) = 0\) (Definition~\ref{pa-def:degAt}), so the transport has degree
  zero (Theorem~\ref{pa-thm:degAff_curveMap_eq_zero}, applied over the field \(K\) of the
  point itself). Hence \(g^{*}\lambda\) has degree zero at every field point of \(T\).
\end{proof}

\begin{theorem}[Naturality and functor laws of the degree-zero transport]
  \label{pa-thm:pic0Pullback_laws}
  \lean{AlgebraicGeometry.pic0Map_pic0Pullback, AlgebraicGeometry.pic0Pullback_id,
    AlgebraicGeometry.pic0Pullback_comp}
  \leanok
  \dcref{custom}

  \uses{pa-def:pic0Pullback, pa-def:pic0Functor}
  The degree-zero curve transport commutes with restriction along an arbitrary
  morphism of test objects \(f \colon T' \to T\),
  \(f^{*} \circ g^{*} = g^{*} \circ f^{*}\); it is the identity for \(g = \id_E\); and
  \((g' \circ g)^{*} = g^{*} \circ g'^{*}\) for composable curve morphisms
  \(g \colon D \to E\), \(g' \colon E \to F\).
\end{theorem}
\begin{proof}
  \leanok
  \uses{pa-thm:picEtMap_picEtPullback, pa-thm:picEtPullback_laws}
  Membership in the degree-zero subgroup is a property of the underlying class, and
  all three equalities hold on underlying classes: naturality is
  Theorem~\ref{pa-thm:picEtMap_picEtPullback}, and the functor laws are
  Theorem~\ref{pa-thm:picEtPullback_laws}.
\end{proof}

\begin{definition}[The bundled degree-zero curve transport]
  \label{pa-def:pic0PullbackNat}
  \lean{AlgebraicGeometry.pic0PullbackNat, AlgebraicGeometry.pic0PullbackNat_app,
    AlgebraicGeometry.pic0PullbackNat_id, AlgebraicGeometry.pic0PullbackNat_comp,
    AlgebraicGeometry.pic0PullbackNat_inclusion}
  \leanok
  \dcref{custom}

  \uses{pa-def:pic0Functor, pa-def:pic0Pullback, pa-def:picEtPullbackNat}
  The degree-zero curve transports assemble into a natural transformation
  \[
    g^{*} \colon \mathrm{Pic}^0_{E/k} \longrightarrow \mathrm{Pic}^0_{D/k}
  \]
  of \(\mathrm{CommGrp}\)-valued functors on the test category, with \(\id_E^{*}\) the
  identity transformation and \((g' \circ g)^{*} = g^{*} \circ g'^{*}\), compatibly
  with the inclusions into the \'etale-sheafified relative Picard functors and their
  curve transport (Definition~\ref{pa-def:picEtPullbackNat}).
\end{definition}
\begin{proof}
  \leanok
  \uses{pa-thm:pic0Pullback_laws}
  Naturality and the two functor laws are Theorem~\ref{pa-thm:pic0Pullback_laws} at each
  test object; the inclusion square commutes because the transport of a degree-zero
  class is, by construction, the transport of its underlying class.
\end{proof}

\section{The degree seam at field points}
\label{pa-sec:degreeSeam}

For a \emph{cocycle-presented} plus class on an affine test — one which is the unit image of
an honest \v Cech Picard class — the degree at a field point collapses to the curve-level
class degree \(\deg_{\Pic}\) of the fibre class. This seam is where the campaign's coverage
and stratification arguments spend degree-zero-ness.

\begin{theorem}[The degree seam, unit form]
  \label{pa-thm:degAt_of_affineEquiv_eq_unit}
  \lean{AlgebraicGeometry.degAt_of_affineEquiv_eq_unit}
  \leanok
  \dcref{custom}

  \uses{pa-def:degAt, pa-thm:picEtAffineEquiv_naturality, pa-def:PicEtAff_mapAlg, pa-thm:degAff_unit,
    pa-def:relPicAlgMap}
  Let \(A\) be a \(k\)-algebra, \(\lambda \in \mathrm{picEt}(C, \Spec A)\) a class whose
  affine collapse is the plus unit of a relative Picard class \(z \in \mathrm{relPic}(C,A)\),
  and \(\varphi \colon A \to K\) a \(k\)-algebra map into a field. Then the degree of
  \(\lambda\) at the field point \(\Spec \varphi\) is the relative degree of the fibre class:
  \[
    \deg_{\Spec\varphi}(\lambda) = \mathrm{relPicDeg}_K\bigl(\varphi^{*} z\bigr).
  \]
\end{theorem}
\begin{proof}
  \leanok
  The affine comparison intertwines restriction along \(\Spec\varphi\) with the explicit
  transport \(\varphi^{*}\) of plus classes (Theorem~\ref{pa-thm:picEtAffineEquiv_naturality}),
  which is natural on units (Definition~\ref{pa-def:PicEtAff_mapAlg}); so the collapse of
  \(t^{*}\lambda\) is \(\mathrm{unit}(\varphi^{*} z)\), whose degree is
  \(\mathrm{relPicDeg}_K(\varphi^{*} z)\) by the unit normalization
  (Theorem~\ref{pa-thm:degAff_unit}).
\end{proof}

\begin{theorem}[The degree seam: cocycle-presented classes]
  \label{pa-thm:degAt_of_affineEquiv_eq_unit_relPicMk}
  \lean{AlgebraicGeometry.degAt_of_affineEquiv_eq_unit_relPicMk}
  \leanok
  \dcref{custom}

  \uses{pa-thm:degAt_of_affineEquiv_eq_unit, pa-def:relPicMap, pa-def:relPic, pa-def:relPicDeg}
  In the situation of Theorem~\ref{pa-thm:degAt_of_affineEquiv_eq_unit}, suppose
  \(z = \pi_{\Spec A}(L)\) for an honest \v Cech Picard class \(L\) on \(C \otimes \Spec A\).
  Then the degree at the field point is the class degree of the fibre \v Cech class:
  \[
    \deg_{\Spec\varphi}(\lambda)
      = \deg_{\Pic}\Bigl( (C \lhd \Spec\varphi)^{*} L \Bigr),
  \]
  the pullback of \(L\) along the point.
\end{theorem}
\begin{proof}
  \leanok
  By Theorem~\ref{pa-thm:degAt_of_affineEquiv_eq_unit} the degree is
  \(\mathrm{relPicDeg}_K(\varphi^{*}\pi_{\Spec A}(L))\); the restriction of a projected class
  is the projection of the pullback (Definition~\ref{pa-def:relPicMap}), and on projected
  classes over a field the relative degree reads the class degree of the representative
  (Definition~\ref{pa-def:relPicDeg}).
\end{proof}

\begin{theorem}[Base-field constancy of the degree]
  \label{pa-thm:degAt_baseChange}
  \lean{AlgebraicGeometry.degAt_of_affineEquiv_eq_unit_baseChange,
    AlgebraicGeometry.degAt_eq_classDeg_of_affineEquiv_eq_unit_baseChange}
  \leanok
  \dcref{custom}

  \uses{pa-thm:degAt_of_affineEquiv_eq_unit, pa-lem:relPicAlgMap_functor,
    pa-thm:relPicDeg_relPicAlgMap, pa-def:relPicDeg, pa-def:relPic}
  In the situation of Theorem~\ref{pa-thm:degAt_of_affineEquiv_eq_unit}, suppose \(z\) is pulled
  back from the base: \(z = \iota_A^{*} w\) for a class \(w \in \mathrm{relPic}(C,k)\), with
  \(\iota_A \colon k \to A\) the structure map. Then the degree of \(\lambda\) is
  \(\mathrm{relPicDeg}_k(w)\) at \emph{every} field point of \(\Spec A\) — independent of the
  point. If moreover \(w = \pi(L_0)\) for a \v Cech class \(L_0\) on \(C \otimes \Spec k\),
  the constant is \(\deg_{\Pic}(L_0)\).
\end{theorem}
\begin{proof}
  \leanok
  At a field point \(\Spec\varphi\), Theorem~\ref{pa-thm:degAt_of_affineEquiv_eq_unit} reads
  \(\mathrm{relPicDeg}_K(\varphi^{*}\iota_A^{*}w) =
  \mathrm{relPicDeg}_K\bigl((\varphi\iota_A)^{*}w\bigr)\) (functoriality,
  Lemma~\ref{pa-lem:relPicAlgMap_functor}); invariance of the relative degree under base-field
  transitions (Theorem~\ref{pa-thm:relPicDeg_relPicAlgMap}) collapses this to
  \(\mathrm{relPicDeg}_k(w)\), with no dependence on \(\varphi\). The \v Cech form reads the
  representative through Definition~\ref{pa-def:relPicDeg}.
\end{proof}

\begin{theorem}[The fibre degree of a degree-zero twist]
  \label{pa-thm:degAt_pic0_mul_pow}
  \lean{AlgebraicGeometry.degAt_pow, AlgebraicGeometry.degAt_pic0_mul_pow}
  \leanok
  \dcref{custom}

  \uses{pa-lem:degAt_hom, pa-def:pic0Subgroup}
  For any class \(\theta \in \mathrm{picEt}(C,T)\), \(m \in \N\), and a field point \(t\),
  \(\deg_t(\theta^m) = m \cdot \deg_t(\theta)\); and for
  \(\lambda \in \mathrm{pic}^0(C,T)\),
  \[
    \deg_t\bigl(\lambda \cdot \theta^{m}\bigr) = m \cdot \deg_t(\theta)
  \]
  on the nose — the degree-zero factor contributes nothing.
\end{theorem}
\begin{proof}
  \leanok
  The power law is additivity (Lemma~\ref{pa-lem:degAt_hom}) by induction on \(m\); the twist
  law adds \(\deg_t(\lambda) = 0\) (Definition~\ref{pa-def:pic0Subgroup}).
\end{proof}

\section{Field points across a base-field extension}
\label{pa-sec:fieldPointMatching}

Let \(k \to L\) be a field extension, \(\sigma : \Spec L \to \Spec k\) the induced
morphism, and \(T\) a test object of \(\Over(\Spec L)\). Write \(\sigma_{!}T\) for
\(T\) regarded as a test object of \(\Over(\Spec k)\) through \(\sigma\): the same
underlying scheme, with structure morphism the structure morphism of \(T\) followed by
\(\sigma\); and write \(\Spec K/\Spec F\) for the affine test of an \(F\)-algebra
\(K\), as in Chapter~\ref{pa-chap:Cohomology}. This section matches the
\(k\)-field-points of \(\sigma_{!}T\) with the \(L\)-field-points of \(T\): a
\(k\)-morphism \(\Spec K/\Spec k \to \sigma_{!}T\) from the affine test of a field
\(K\) is the same thing as an \(L\)-algebra structure on \(K\) completing the tower
\(k \to L \to K\) together with an \(L\)-morphism \(\Spec K/\Spec L \to T\), both
directions preserving the underlying morphism of schemes. The matching then exchanges
the quantifier defining the degree-zero subgroup on \(\sigma_{!}T\) for a quantifier
over \(L\)-field-points of \(T\). For the last two statements, \(C\) is a curve bundle
over \(k\) --- proper, smooth of relative dimension one, geometrically irreducible ---
with base change \(C_L := \sigma^{*}C\) over \(L\) (\ref{pa-def:baseChange}).

\begin{definition}[Pushing a field point forward]
  \label{pa-def:pushFieldPoint}
  \lean{AlgebraicGeometry.pushFieldPoint, AlgebraicGeometry.pushFieldPoint_left}
  \leanok
  \dcref{custom}

  \uses{pa-def:mapOverSpecIso}
  Let \(K\) be a field in a scalar tower \(k \to L \to K\) and
  \(s : \Spec K/\Spec L \to T\) an \(L\)-field-point of \(T\). The \emph{pushed field
  point} is the \(k\)-field-point
  \[
    \mathrm{push}(s) : \Spec K/\Spec k \longrightarrow \sigma_{!}T
  \]
  obtained by composing the inverse of the affine matching
  \(\sigma_{!}(\Spec K/\Spec L) \cong \Spec K/\Spec k\) (\ref{pa-def:mapOverSpecIso}) with
  \(s\) pushed forward along \(\sigma\). Its underlying morphism of schemes is that of
  \(s\).
\end{definition}
\begin{proof}
  \leanok
  The affine matching has underlying morphism the identity of \(\Spec K\)
  (\ref{pa-def:mapOverSpecIso}), and pushing forward along \(\sigma\) leaves underlying
  morphisms unchanged; so the composite is a morphism of \(\Over(\Spec k)\) whose
  underlying morphism is that of \(s\).
\end{proof}

\begin{definition}[The algebra structure unpacked from a field point]
  \label{pa-def:fieldPointAlgebra}
  \lean{AlgebraicGeometry.fieldPointAlgebra,
    AlgebraicGeometry.fieldPointAlgebra_algebraMap}
  \leanok
  \dcref{custom}

  Let \(K\) be a field extension of \(k\) and
  \(t : \Spec K/\Spec k \to \sigma_{!}T\) a \(k\)-field-point of the pushed test. The
  composite of the underlying morphism of \(t\) with the structure morphism of \(T\)
  is a morphism of schemes \(\Spec K \to \Spec L\), hence \(\Spec\) of a unique ring
  homomorphism \(L \to K\), by full faithfulness of \(\Spec\). This homomorphism makes
  \(K\) an \(L\)-algebra: the \emph{field-point algebra structure} of \(t\), whose
  algebra map is that unique homomorphism.
\end{definition}
\begin{proof}
  \leanok
  \(\Spec\) is a fully faithful contravariant functor from commutative rings to
  schemes, so every morphism \(\Spec K \to \Spec L\) is \(\Spec\varphi\) for a unique
  ring homomorphism \(\varphi : L \to K\), and a ring homomorphism into \(K\) equips
  \(K\) with the structure of an \(L\)-algebra.
\end{proof}

\begin{lemma}[The unpacked structure completes the tower]
  \label{pa-lem:isScalarTower_fieldPointAlgebra}
  \lean{AlgebraicGeometry.isScalarTower_fieldPointAlgebra}
  \leanok
  \dcref{custom}

  \uses{pa-def:fieldPointAlgebra}
  In the situation of Definition~\ref{pa-def:fieldPointAlgebra}, write
  \(\varphi : L \to K\) for the unpacked algebra map. Then
  \(\varphi \circ (k \to L) = (k \to K)\): the field-point algebra structure, the
  given \(k\)-structure of \(L\) and the given \(k\)-structure of \(K\) form a scalar
  tower \(k \to L \to K\).
\end{lemma}
\begin{proof}
  \leanok
  The point \(t\) lies over \(\Spec k\): the underlying morphism of \(t\) followed by
  the structure morphism of \(\sigma_{!}T\) is the structure morphism of the affine
  test \(\Spec K/\Spec k\), which is \(\Spec\) of \(k \to K\). The structure morphism
  of \(\sigma_{!}T\) is that of \(T\) followed by \(\sigma\), so the left-hand side is
  \(\Spec\varphi\) followed by \(\Spec(k \to L)\), which by contravariant
  functoriality is \(\Spec\bigl(\varphi \circ (k \to L)\bigr)\). Faithfulness of
  \(\Spec\) turns \(\Spec\bigl(\varphi \circ (k \to L)\bigr) = \Spec(k \to K)\) into
  the tower identity.
\end{proof}

\begin{definition}[Lifting a field point of the pushed test]
  \label{pa-def:liftFieldPoint}
  \lean{AlgebraicGeometry.liftFieldPoint, AlgebraicGeometry.liftFieldPoint_left}
  \leanok
  \dcref{custom}

  \uses{pa-def:fieldPointAlgebra}
  In the situation of Definition~\ref{pa-def:fieldPointAlgebra}, endow \(K\) with the
  field-point algebra structure of \(t\). Then the underlying morphism of \(t\) is an
  \(L\)-field-point
  \[
    \mathrm{lift}(t) : \Spec K/\Spec L \longrightarrow T
  \]
  with the same underlying morphism of schemes as \(t\).
\end{definition}
\begin{proof}
  \leanok
  The required triangle over \(\Spec L\) --- the underlying morphism of \(t\) followed
  by the structure morphism of \(T\) equals \(\Spec\) of the algebra map \(L \to K\)
  --- holds by the very definition of the field-point algebra structure
  (\ref{pa-def:fieldPointAlgebra}).
\end{proof}

\begin{theorem}[The field-point matching round trips]
  \label{pa-thm:fieldPoint_roundTrip}
  \lean{AlgebraicGeometry.pushFieldPoint_liftFieldPoint,
    AlgebraicGeometry.fieldPointAlgebra_pushFieldPoint,
    AlgebraicGeometry.liftFieldPoint_pushFieldPoint_left}
  \leanok
  \dcref{custom}

  \uses{pa-def:pushFieldPoint, pa-def:fieldPointAlgebra,
    pa-lem:isScalarTower_fieldPointAlgebra, pa-def:liftFieldPoint}
  The pack and unpack constructions are mutually inverse:
  (i) for a \(k\)-field-point \(t\) of \(\sigma_{!}T\), pushing \(\mathrm{lift}(t)\)
  forward --- over the field-point algebra structure of \(t\) and its tower ---
  recovers \(t\): \(\mathrm{push}(\mathrm{lift}(t)) = t\);
  (ii) for a field \(K\) in a scalar tower \(k \to L \to K\) and an \(L\)-field-point
  \(s\) of \(T\), the field-point algebra structure of \(\mathrm{push}(s)\) is the
  given \(L\)-structure of \(K\);
  (iii) and \(\mathrm{lift}(\mathrm{push}(s))\) has the same underlying morphism of
  schemes as \(s\).
\end{theorem}
\begin{proof}
  \leanok
  (i) and (iii): push and lift both preserve the underlying morphism of schemes
  (\ref{pa-def:pushFieldPoint}, \ref{pa-def:liftFieldPoint}), and a morphism of a slice
  category is determined by its underlying morphism; the tower of
  Lemma~\ref{pa-lem:isScalarTower_fieldPointAlgebra} guarantees that source and target
  objects agree on both sides. (ii): the underlying morphism of \(\mathrm{push}(s)\)
  followed by the structure morphism of \(T\) is the underlying morphism of \(s\) so
  followed, which is \(\Spec\) of the given algebra map \(L \to K\) --- the defining
  triangle of \(s\). By uniqueness in Definition~\ref{pa-def:fieldPointAlgebra}, the
  unpacked algebra map is the given one, so the two \(L\)-algebra structures coincide.
\end{proof}

\begin{theorem}[The quantifier exchange for degree zero on a pushed test]
  \label{pa-thm:mem_pic0Subgroup_iff_forall_pushFieldPoint}
  \lean{AlgebraicGeometry.mem_pic0Subgroup_iff_forall_pushFieldPoint}
  \leanok
  \dcref{custom}

  \uses{pa-def:pic0Subgroup, pa-def:degAt, pa-def:pushFieldPoint, pa-def:fieldPointAlgebra,
    pa-lem:isScalarTower_fieldPointAlgebra, pa-def:liftFieldPoint, pa-thm:fieldPoint_roundTrip}
  For a class \(\lambda \in \mathrm{picEt}(C, \sigma_{!}T)\), the following are
  equivalent: \(\lambda \in \mathrm{pic}^0(C, \sigma_{!}T)\); and for every field
  \(K\) in a scalar tower \(k \to L \to K\) and every \(L\)-field-point \(s\) of
  \(T\),
  \[
    \deg_{\mathrm{push}(s)}(\lambda) = 0 .
  \]
\end{theorem}
\begin{proof}
  \leanok
  Membership in \(\mathrm{pic}^0(C, \sigma_{!}T)\) is vanishing of the degree at every
  \(k\)-field-point of \(\sigma_{!}T\) (\ref{pa-def:pic0Subgroup}). Forward: each pushed
  point \(\mathrm{push}(s)\) is such a field point (\ref{pa-def:pushFieldPoint}).
  Backward: let \(t : \Spec K \to \sigma_{!}T\) be any \(k\)-field-point. Endow \(K\)
  with the field-point algebra structure of \(t\); it completes the tower
  \(k \to L \to K\) (\ref{pa-def:fieldPointAlgebra},
  \ref{pa-lem:isScalarTower_fieldPointAlgebra}), and over it
  \(t = \mathrm{push}(\mathrm{lift}(t))\) (\ref{pa-def:liftFieldPoint},
  \ref{pa-thm:fieldPoint_roundTrip}). The hypothesis at this tower and at the point
  \(\mathrm{lift}(t)\) gives
  \(\deg_t(\lambda) = \deg_{\mathrm{push}(\mathrm{lift}(t))}(\lambda) = 0\).
\end{proof}

\begin{theorem}[Classes matching in degree at pushed points are simultaneously
    degree zero]
  \label{pa-thm:mem_pic0Subgroup_iff_of_degAt_pushFieldPoint_eq}
  \lean{AlgebraicGeometry.mem_pic0Subgroup_iff_of_degAt_pushFieldPoint_eq}
  \leanok
  \dcref{custom}

  \uses{pa-thm:mem_pic0Subgroup_iff_forall_pushFieldPoint, pa-def:pic0Subgroup, pa-def:degAt,
    pa-def:pushFieldPoint, pa-def:picEt, pa-def:baseChange}
  Let \(\lambda \in \mathrm{picEt}(C, \sigma_{!}T)\) and
  \(\lambda_L \in \mathrm{picEt}(C_L, T)\) --- the \'etale Picard group of the
  base-changed curve \(C_L\) at the test \(T\), formed over the base field \(L\)
  exactly as in Definition~\ref{pa-def:picEt}. Suppose the two classes match in degree at
  every pushed field-point pair: for every field \(K\) in a scalar tower
  \(k \to L \to K\) and every \(L\)-field-point \(s\) of \(T\),
  \[
    \deg_{\mathrm{push}(s)}(\lambda) \;=\; \deg_{s}(\lambda_L).
  \]
  Then \(\lambda \in \mathrm{pic}^0(C, \sigma_{!}T)\) if and only if
  \(\lambda_L \in \mathrm{pic}^0(C_L, T)\).
\end{theorem}
\begin{proof}
  \leanok
  By the quantifier exchange
  (\ref{pa-thm:mem_pic0Subgroup_iff_forall_pushFieldPoint}), membership of \(\lambda\) is
  equivalent to the vanishing of \(\deg_{\mathrm{push}(s)}(\lambda)\) over all towers
  \(K\) and \(L\)-field-points \(s\) of \(T\).

  Suppose this vanishing holds, and let \(s : \Spec K \to T\) be an \(L\)-field-point
  with \(K\) a field extension of \(L\). Restricting the \(L\)-structure of \(K\)
  through \(k \to L\) makes \(K\) a field in a scalar tower \(k \to L \to K\), so the
  matching hypothesis applies:
  \(\deg_s(\lambda_L) = \deg_{\mathrm{push}(s)}(\lambda) = 0\). As \(s\) was
  arbitrary, \(\lambda_L \in \mathrm{pic}^0(C_L, T)\) (\ref{pa-def:pic0Subgroup}, over
  the base field \(L\)).

  Conversely, if \(\lambda_L\) is of degree zero at every \(L\)-field-point of \(T\),
  then for every tower \(K\) and every \(L\)-field-point \(s\),
  \(\deg_{\mathrm{push}(s)}(\lambda) = \deg_s(\lambda_L) = 0\); conclude by the
  exchange.
\end{proof}

\section{The base-field comparison of the degree-zero Picard functor}
\label{pa-sec:pic0Theta}

Fix a field embedding \(k \to L\), write \(\sigma \colon \Spec L \to \Spec k\) for the
induced morphism and \(\sigma_{!}\) for the postcomposition functor
\(\Over(\Spec L) \to \Over(\Spec k)\), and let \(C_L\) denote the base change of the
curve \(C\) along \(k \to L\).  The base-field shuffle of the plus construction
(Definition~\ref{pa-def:PicEtAff_baseFieldShuffle}) extends from affine tests to all test
objects, respects the degree, and assembles into a natural isomorphism \(\theta\)
between the degree-zero Picard functor of \(C_L\) and the pullback of that of \(C\)
along \(\sigma_{!}\).

\begin{lemma}[The section-ring scalar tower of a pushed test]
  \label{pa-lem:isScalarTower_sections_map}
  \lean{AlgebraicGeometry.Over.ofHom_algebraMap_sections_map,
    AlgebraicGeometry.Over.isScalarTower_sections_map}
  \leanok
  \dcref{custom}

  \uses{pa-def:baseChange}
  Let \(T\) be a test object of \(\Over(\Spec L)\) and \(U \subseteq T_{\mathrm{left}}\)
  an open.  The section ring \(\Gamma(T, U)\) carries an \(L\)-algebra structure
  through the structure morphism of \(T\), and a \(k\)-algebra structure through the
  structure morphism of the pushed test \(\sigma_{!}T\) (the same scheme, viewed over
  \(\Spec k\) through \(\sigma\)).  These form a scalar tower: the \(k\)-structure map
  is the composite of \(k \to L\) with the \(L\)-structure map.
\end{lemma}
\begin{proof}
  \leanok
  The structure morphism of \(\sigma_{!}T\) is the composite of that of \(T\) with
  \(\sigma\), so its action on sections over \(U\) factors as
  \(\Gamma(\Spec k, \top) \to \Gamma(\Spec L, \top) \to \Gamma(T, U)\).  Under the
  canonical identifications \(\Gamma(\Spec k, \top) \cong k\) and
  \(\Gamma(\Spec L, \top) \cong L\), the first map is the embedding \(k \to L\) and
  the second is the \(L\)-structure map of \(\Gamma(T, U)\).
\end{proof}

\begin{definition}[The componentwise lift of the base-field shuffle]
  \label{pa-def:picEtCrossBase}
  \lean{AlgebraicGeometry.sectionShuffle, AlgebraicGeometry.mapAlg_sectionShuffle,
    AlgebraicGeometry.mapAlg_sectionShuffle_symm,
    AlgebraicGeometry.Over.resAlgHom_map_restrictScalars,
    AlgebraicGeometry.picEtCrossBase, AlgebraicGeometry.picEtCrossBase_val,
    AlgebraicGeometry.picEtCrossBaseInv, AlgebraicGeometry.picEtCrossBaseInv_val,
    AlgebraicGeometry.picEtCrossBaseEquiv,
    AlgebraicGeometry.picEtCrossBaseEquiv_apply,
    AlgebraicGeometry.picEtCrossBaseEquiv_symm_apply}
  \leanok
  \dcref{custom}

  \uses{pa-def:picEt, pa-def:PicEtAff_baseFieldShuffle, pa-lem:isScalarTower_sections_map,
    pa-def:baseChange}
  Let \(T\) be a test object of \(\Over(\Spec L)\).  The underlying schemes of \(T\)
  and \(\sigma_{!}T\) coincide, so their affine opens and section rings are literally
  shared, and by the scalar tower of
  Lemma~\ref{pa-lem:isScalarTower_sections_map} each section ring \(\Gamma(T, U)\) is a
  legal test algebra for the base-field shuffle.  Applying the shuffle
  \(\mathrm{Sh}_{\Gamma(T,U)}\) of Definition~\ref{pa-def:PicEtAff_baseFieldShuffle}
  componentwise to a compatible family defines a group isomorphism
  \[
    \mathrm{picEt}(C, \sigma_{!}T) \;\cong\; \mathrm{picEt}(C_L, T),
  \]
  the \emph{componentwise lift} of the shuffle.
\end{definition}
\begin{proof}
  \leanok
  A family transported componentwise remains compatible: the shuffle commutes with
  restriction of sections, since restriction along \(U \subseteq V\) is the same ring
  homomorphism \(\Gamma(T, V) \to \Gamma(T, U)\) for the two base structures, and the
  shuffle satisfies the restriction square of
  Theorem~\ref{pa-thm:PicEtAff_baseFieldShuffle_laws} at it.  Componentwise application
  of a group isomorphism preserves products, and the two directions are componentwise
  inverse, so the lift is a group isomorphism.
  \uses{pa-thm:PicEtAff_baseFieldShuffle_laws}
\end{proof}

\begin{theorem}[Naturality of the componentwise lift in the test object]
  \label{pa-thm:picEtMap_picEtCrossBase}
  \lean{AlgebraicGeometry.Over.appLEAlgHom_map_restrictScalars,
    AlgebraicGeometry.appLE_sectionShuffle,
    AlgebraicGeometry.picEtMap_picEtCrossBase,
    AlgebraicGeometry.picEtMap_picEtCrossBaseInv}
  \leanok
  \dcref{custom}

  \uses{pa-def:picEtCrossBase, pa-def:picEtMap}
  For a morphism \(f \colon T' \to T\) of test objects of \(\Over(\Spec L)\),
  restriction along \(f\) commutes with the componentwise lift:
  restricting a class of \(\mathrm{picEt}(C, \sigma_{!}T)\) along
  \(\sigma_{!}f\) and lifting equals lifting and restricting along \(f\).
\end{theorem}
\begin{proof}
  \leanok
  The value at an affine open \(W \subseteq T'\) of a restricted family is
  characterized by the pullback property of Definition~\ref{pa-def:picEtMap}: on every
  affine sub-open of \(W\) contained in the preimage of an affine open \(V\) of the
  target, it restricts to the pullback of the family value at \(V\).  The lift of the
  \(\sigma_{!}f\)-restriction satisfies this property for \(f\): pullback of sections
  along \(f\) and along \(\sigma_{!}f\) are the same ring homomorphism, and the
  shuffle commutes with both restriction and pullback of sections
  (Theorem~\ref{pa-thm:PicEtAff_baseFieldShuffle_laws} at the section rings, through the
  tower of Lemma~\ref{pa-lem:isScalarTower_sections_map}).  Uniqueness of the
  characterized value concludes.
  \uses{pa-thm:PicEtAff_baseFieldShuffle_laws, pa-lem:isScalarTower_sections_map}
\end{proof}

\begin{theorem}[The affine face of the componentwise lift]
  \label{pa-thm:picEtAffineEquiv_picEtCrossBase}
  \lean{AlgebraicGeometry.overSpecΓTop_map_restrictScalars,
    AlgebraicGeometry.picEtAffineEquiv_picEtCrossBase}
  \leanok
  \dcref{custom}

  \uses{pa-def:picEtCrossBase, pa-def:picEtAffineEquiv, pa-def:mapOverSpecIso, pa-def:picEtMap}
  Let \(A\) be an algebra in a scalar tower \(k \to L \to A\).  Under the affine
  comparisons of Definition~\ref{pa-def:picEtAffineEquiv} on the two sides, the
  componentwise lift at the affine test \(\Spec A\) over \(L\) is the base-field
  shuffle \(\mathrm{Sh}_A\) itself, the pushed side being read through the affine
  matching \(\sigma_{!}(\Spec A / L) \cong (\Spec A / k)\) of
  Definition~\ref{pa-def:mapOverSpecIso}.
\end{theorem}
\begin{proof}
  \leanok
  Evaluate both sides at the top affine open.  On the lifted side the value is the
  shuffle of the value of the original family at the top open, transported into \(A\)
  by the \(L\)-comparison \(\Gamma(\Spec A, \top) \cong A\); by the restriction square
  of the shuffle this is \(\mathrm{Sh}_A\) applied to the transport of the original
  value along the \(k\)-restriction of that comparison.  On the pushed side, the
  restriction along the affine matching has value the pullback of the original value
  along the identity carrier, transported by the \(k\)-comparison; the two composite
  transports \(\Gamma((\Spec A/L)_{\mathrm{left}}, \top) \to A\) agree, both being the
  canonical identification.
  \uses{pa-thm:PicEtAff_baseFieldShuffle_laws}
\end{proof}

\begin{theorem}[Degree invariance of the base-field shuffle at a field]
  \label{pa-thm:degAff_baseFieldShuffle}
  \lean{AlgebraicGeometry.degAff_baseFieldShuffle}
  \leanok
  \dcref{custom}

  \uses{pa-def:degAff, pa-def:PicEtAff_baseFieldShuffle, pa-def:baseChange}
  Let \(K\) be a field in a scalar tower \(k \to L \to K\).  For every plus class
  \(x \in \mathrm{PicEtAff}(C, K)\),
  \[
    \deg_K\bigl(\mathrm{Sh}_K(x)\bigr) \;=\; \deg_K(x),
  \]
  where the left degree is that of the base-changed curve \(C_L\) over the base field
  \(L\) and the right one that of \(C\) over \(k\).
\end{theorem}
\begin{proof}
  \leanok
  Represent \(x\) by a descent class \(y\) on an \'etale cover \(E\) of \(\Spec K\);
  then \(\mathrm{Sh}_K(x)\) is represented on the same cover by the relative Picard
  shuffle of \(y\) (Theorem~\ref{pa-thm:PicEtAff_baseFieldShuffle_relPic}).  Choose a
  finite separable field extension \(M/K\) refining \(E\) through some
  \(j \colon E \to M\), and endow \(M\) with the composite \(k\)- and \(L\)-structures
  of the tower \(k \to L \to K \to M\).  By the master equality of the plus-class
  degree (Definition~\ref{pa-def:degAff}) both sides are read at \((M, j)\): the left
  side is \(\mathrm{relPicDeg}_M\) of the transport along \(j\) over \(L\) of the
  shuffled class, which by the naturality square of the relative Picard shuffle
  (Theorem~\ref{pa-thm:relPicCrossBase_relPicAlgMap}) is the shuffle at \(M\) of the
  transport along \(j\) over \(k\) of \(y\); by the degree invariance of the relative
  Picard shuffle at the field \(M\)
  (Theorem~\ref{pa-thm:relPicDeg_relPicCrossBase}) this reads the same integer as the
  transport itself, which is the right side.
  \uses{pa-thm:PicEtAff_baseFieldShuffle_relPic, pa-thm:relPicCrossBase_relPicAlgMap,
    pa-thm:relPicDeg_relPicCrossBase}
\end{proof}

\begin{theorem}[The degree matching at pushed field points]
  \label{pa-thm:degAt_pushFieldPoint_picEtCrossBase}
  \lean{AlgebraicGeometry.degAt_pushFieldPoint_picEtCrossBase}
  \leanok
  \dcref{custom}

  \uses{pa-def:degAt, pa-def:pushFieldPoint, pa-def:picEtCrossBase}
  Let \(T\) be a test object of \(\Over(\Spec L)\) and
  \(\lambda \in \mathrm{picEt}(C, \sigma_{!}T)\).  For every field \(K\) in a scalar
  tower \(k \to L \to K\) and every \(L\)-field-point \(s\) of \(T\),
  \[
    \deg_{\mathrm{push}(s)}(\lambda)
      \;=\; \deg_{s}\bigl(\mathrm{lift}(\lambda)\bigr),
  \]
  where \(\mathrm{lift}(\lambda) \in \mathrm{picEt}(C_L, T)\) is the componentwise
  lift of \(\lambda\).
\end{theorem}
\begin{proof}
  \leanok
  The pushed point factors as the affine matching followed by \(\sigma_{!}s\), so the
  left side is the degree at \(K\) of the restriction of \(\lambda\) along
  \(\sigma_{!}s\), read through the affine matching and the \(k\)-side affine
  comparison.  By naturality of the lift
  (Theorem~\ref{pa-thm:picEtMap_picEtCrossBase}) the restriction of
  \(\mathrm{lift}(\lambda)\) along \(s\) is the lift of the restriction of \(\lambda\)
  along \(\sigma_{!}s\); by the affine face
  (Theorem~\ref{pa-thm:picEtAffineEquiv_picEtCrossBase}) its reading under the
  \(L\)-side affine comparison is the shuffle \(\mathrm{Sh}_K\) of the \(k\)-side
  reading.  The shuffle preserves the degree at the field \(K\)
  (Theorem~\ref{pa-thm:degAff_baseFieldShuffle}), so the two sides agree.
  \uses{pa-thm:picEtMap_picEtCrossBase, pa-thm:picEtAffineEquiv_picEtCrossBase,
    pa-thm:degAff_baseFieldShuffle, pa-def:picEtAffineEquiv}
\end{proof}

\begin{theorem}[The componentwise lift restricts to the degree-zero subgroups]
  \label{pa-thm:mem_pic0Subgroup_picEtCrossBase_iff}
  \lean{AlgebraicGeometry.mem_pic0Subgroup_picEtCrossBase_iff}
  \leanok
  \dcref{custom}

  \uses{pa-def:pic0Subgroup, pa-def:picEtCrossBase,
    pa-thm:mem_pic0Subgroup_iff_of_degAt_pushFieldPoint_eq}
  A class \(\lambda \in \mathrm{picEt}(C, \sigma_{!}T)\) is of degree zero at every
  field point if and only if its componentwise lift
  \(\mathrm{lift}(\lambda) \in \mathrm{picEt}(C_L, T)\) is.
\end{theorem}
\begin{proof}
  \leanok
  The pair \((\lambda, \mathrm{lift}(\lambda))\) matches in degree at every pushed
  field-point pair (Theorem~\ref{pa-thm:degAt_pushFieldPoint_picEtCrossBase}), so the
  simultaneous-vanishing criterion
  (Theorem~\ref{pa-thm:mem_pic0Subgroup_iff_of_degAt_pushFieldPoint_eq}) applies.
  \uses{pa-thm:degAt_pushFieldPoint_picEtCrossBase}
\end{proof}

\begin{definition}[\(\theta\): the base-field comparison of the degree-zero Picard
    functor]
  \label{pa-def:pic0Theta}
  \lean{AlgebraicGeometry.pic0CrossBaseEquiv,
    AlgebraicGeometry.pic0CrossBaseEquiv_apply_coe,
    AlgebraicGeometry.pic0CrossBaseEquiv_symm_apply_coe,
    AlgebraicGeometry.pic0Theta, AlgebraicGeometry.pic0Theta_hom_app,
    AlgebraicGeometry.pic0Theta_inv_app}
  \leanok
  \dcref{custom}

  \uses{pa-def:pic0Functor, pa-def:picEtCrossBase,
    pa-thm:mem_pic0Subgroup_picEtCrossBase_iff}
  The componentwise lift restricts to the degree-zero subgroups
  (Theorem~\ref{pa-thm:mem_pic0Subgroup_picEtCrossBase_iff}) and so induces, for each
  test object \(T\) of \(\Over(\Spec L)\), a group isomorphism
  \[
    \theta_T \colon \mathrm{pic}^0(C_L, T) \;\cong\; \mathrm{pic}^0(C, \sigma_{!}T).
  \]
  These are the components of a natural isomorphism of commutative-group-valued
  functors
  \[
    \theta \colon \mathrm{pic}^0_{C_L/L}
      \;\cong\; \sigma_{!}^{\mathrm{op}} \circ \mathrm{pic}^0_{C/k}
  \]
  on \(\Over(\Spec L)^{\mathrm{op}}\).
\end{definition}
\begin{proof}
  \leanok
  Naturality in \(T\) is the naturality of the componentwise lift in the test object
  (Theorem~\ref{pa-thm:picEtMap_picEtCrossBase}), restricted to the degree-zero
  subgroups, whose restriction maps are those of the ambient \'etale Picard groups.
  \uses{pa-thm:picEtMap_picEtCrossBase}
\end{proof}

\begin{definition}[The set-valued form of \(\theta\)]
  \label{pa-def:pic0ThetaType}
  \lean{AlgebraicGeometry.pic0ThetaType, AlgebraicGeometry.pic0ThetaType_hom_app,
    AlgebraicGeometry.pic0ThetaType_inv_app}
  \leanok
  \dcref{custom}

  \uses{pa-def:pic0Theta, pa-def:pic0SigmaFunctor}
  Composing \(\theta\) with the forgetful functor from commutative groups to sets
  yields the natural isomorphism
  \(\mathrm{pic}^0_{\mathrm{Set}}(C_L) \cong
    \sigma_{!}^{\mathrm{op}} \circ \mathrm{pic}^0_{\mathrm{Set}}(C)\)
  of the set-valued degree-zero functors of
  Definition~\ref{pa-def:pic0SigmaFunctor}, with the same components as \(\theta\).
\end{definition}

\section{The \texorpdfstring{\(\theta\)}{theta}-family and the multiplication shift}
\label{pa-sec:thetaShift}

Fix a \v Cech Picard class \(L_0\) on the base-changed curve
\(C_k := (C \otimes \Spec k)_{\mathrm{left}}\). Pulling its plus class back from the base
test along the canonical structure maps produces a family \(\theta_T\) of \'etale Picard
classes, natural in \(T\) and of constant degree \(d_1 = \deg_{\Pic}(L_0)\) at every field
point; multiplication by \(\theta_T^{m}\) then shifts the degree-zero functor onto the
degree-\(m d_1\) layer.

\begin{definition}[The canonical map to the base test]
  \label{pa-def:toBaseTest}
  \lean{AlgebraicGeometry.toBaseTest, AlgebraicGeometry.toBaseTest_left,
    AlgebraicGeometry.toBaseTest_naturality, AlgebraicGeometry.toBaseTest_overSpec}
  \leanok
  \dcref{custom}

  \uses{pa-def:overSpecMap}
  The base test \(\Spec k\) (the affine test of the \(k\)-algebra \(k\)) is terminal in
  \(\Over(\Spec k)\) up to the identification of its structure morphism with the identity:
  every test object \(T\) has a canonical morphism \(T \to \Spec k\) given by its structure
  morphism, natural in \(T\) (\(f \circ (\text{canonical}) = \text{canonical}\)); on an
  affine test \(\Spec A\) it is \(\Spec\) of the structure map \(k \to A\).
\end{definition}
\begin{proof}
  \leanok
  The structure morphism of the base test is \(\Spec\) of \(\mathrm{id}_k\), the identity;
  hence the structure morphism of \(T\) is a morphism of \(\Over(\Spec k)\) into it.
  Naturality is the defining triangle of morphisms of \(\Over(\Spec k)\); the affine case is
  the compatibility of \(\Spec\) with structure maps (Definition~\ref{pa-def:overSpecMap}).
\end{proof}

\begin{definition}[The \(\theta\)-family]
  \label{pa-def:thetaFamily}
  \lean{AlgebraicGeometry.thetaBase, AlgebraicGeometry.thetaFamily,
    AlgebraicGeometry.thetaFamily_natural}
  \leanok
  \dcref{custom}

  \uses{pa-def:toBaseTest, pa-def:picEtAffineEquiv, pa-def:PicEtAff_unit, pa-def:relPic, pa-def:picEtMap,
    pa-thm:picEtMap_functor_laws}
  For a \v Cech class \(L_0\) on \(C_k\), let
  \(\theta_{\mathrm{base}} \in \mathrm{picEt}(C, \Spec k)\) be the affine-comparison preimage
  of the plus unit of \(\pi(L_0)\), and for every test \(T\) set
  \[
    \theta_T \;:=\; (\text{canonical } T \to \Spec k)^{*}\,\theta_{\mathrm{base}}
      \;\in\; \mathrm{picEt}(C,T).
  \]
  The family is natural: \(f^{*}\theta_T = \theta_{T'}\) for every \(f \colon T' \to T\) —
  \(\theta\) pulls back to \(\theta\), so twisting by \(\theta_T^{m}\) is uniform in the
  test.
\end{definition}
\begin{proof}
  \leanok
  Naturality: \(f^{*}\theta_T\) is the pullback of \(\theta_{\mathrm{base}}\) along the
  composite \(f \circ (\text{canonical})\), which is the canonical map of \(T'\)
  (Definition~\ref{pa-def:toBaseTest}); conclude by functoriality of restriction
  (Theorem~\ref{pa-thm:picEtMap_functor_laws}).
\end{proof}

\begin{theorem}[The degree of the \(\theta\)-family at every field point]
  \label{pa-thm:degAt_thetaFamily}
  \lean{AlgebraicGeometry.thetaFamily_overSpec_affineEquiv,
    AlgebraicGeometry.degAt_thetaFamily, AlgebraicGeometry.degAt_thetaFamily_pow}
  \leanok
  \dcref{custom}

  \uses{pa-def:thetaFamily, pa-thm:degAt_baseChange, pa-thm:picEtAffineEquiv_naturality,
    pa-def:PicEtAff_mapAlg, pa-thm:degAt_picEtMap, pa-thm:degAt_pic0_mul_pow, pa-def:toBaseTest,
    pa-def:relPicAlgMap}
  For every test \(T\), every field point \(t \colon \Spec K \to T\) and every \(m \in \N\),
  \[
    \deg_t(\theta_T) = \deg_{\Pic}(L_0) =: d_1,
    \qquad
    \deg_t\bigl(\theta_T^{\,m}\bigr) = m \, d_1 ,
  \]
  a base-field-invariant constant.
\end{theorem}
\begin{proof}
  \leanok
  On an affine test \(\Spec A\), the affine collapse of \(\theta_{\Spec A}\) is the plus unit
  of the base-changed class \(\iota_A^{*}\pi(L_0)\): the canonical map is \(\Spec\iota_A\)
  (Definition~\ref{pa-def:toBaseTest}), the affine comparison intertwines restriction with the
  explicit transport (Theorem~\ref{pa-thm:picEtAffineEquiv_naturality}), and the transport is
  natural on units (Definition~\ref{pa-def:PicEtAff_mapAlg}). Base-field constancy
  (Theorem~\ref{pa-thm:degAt_baseChange}) then reads \(\deg = \deg_{\Pic}(L_0)\) at every field
  point of \(\Spec A\). For a general test \(T\) and field point \(t\), naturality of the
  family (Definition~\ref{pa-def:thetaFamily}) plus functoriality of the degree
  (Theorem~\ref{pa-thm:degAt_picEtMap}) reduce \(\deg_t(\theta_T)\) to the degree of
  \(\theta_{\Spec K}\) at the identity point of \(\Spec K\), covered by the affine case.
  The power law is Theorem~\ref{pa-thm:degAt_pic0_mul_pow}.
\end{proof}

\begin{definition}[The degree-\(d\) layer functor]
  \label{pa-def:picDegLayerFunctor}
  \lean{AlgebraicGeometry.picDegLayer, AlgebraicGeometry.picDegLayerFunctor}
  \leanok
  \dcref{custom}

  \uses{pa-def:degAt, pa-def:picEtFunctor, pa-thm:degAt_picEtMap, pa-thm:picEtMap_functor_laws}
  For \(d \in \Z\), the \emph{degree-\(d\) layer}
  \[
    \mathrm{pic}^{d}(C,T) \;:=\;
    \bigl\{\, \lambda \in \mathrm{picEt}(C,T) \;:\; \deg_t(\lambda) = d
      \text{ for every field point } t \,\bigr\},
  \]
  a \emph{set}-valued subfunctor \(\Over(\Spec k)^{\mathrm{op}} \to \mathrm{Set}\) of the
  \'etale Picard functor. For \(d \neq 0\) the layer is a coset of the degree-zero subgroup,
  not a subgroup: it carries no group structure.
\end{definition}
\begin{proof}
  \leanok
  Stability of the degree-\(d\) condition under restriction is functoriality of the degree
  (Theorem~\ref{pa-thm:degAt_picEtMap}); the functor laws are inherited from
  \(\mathrm{picEtFunctor}\) (Theorem~\ref{pa-thm:picEtMap_functor_laws}).
\end{proof}

\begin{theorem}[The multiplication-by-\(\theta^m\) natural isomorphism and the shift of a representation]
  \label{pa-thm:mulThetaPowNatIso}
  \lean{AlgebraicGeometry.mulThetaPowEquiv, AlgebraicGeometry.mulThetaPowNatIso,
    AlgebraicGeometry.representableByOfShift}
  \leanok
  \dcref{custom}

  \uses{pa-def:picDegLayerFunctor, pa-def:pic0Functor, pa-thm:degAt_pic0_mul_pow,
    pa-thm:degAt_thetaFamily, pa-def:thetaFamily, pa-lem:degAt_hom}
  For every \(m \in \N\), multiplication by \(\theta_T^{m}\) is a bijection
  \[
    \mathrm{pic}^0(C,T) \;\xrightarrow{\ \sim\ }\; \mathrm{pic}^{m d_1}(C,T),
    \qquad \lambda \mapsto \lambda \cdot \theta_T^{m},
  \]
  natural in \(T\): an isomorphism of set-valued functors from (the underlying-set functor
  of) \(\mathrm{Pic}^0_{C/k}\) onto the degree-\(m d_1\) layer functor. Consequently a
  representation of the shifted layer functor by an object \(J\) transports to a
  representation of \(\mathrm{Pic}^0_{C/k}\) (as a set-valued functor) by \(J\).
\end{theorem}
\begin{proof}
  \leanok
  Forward: \(\deg_t(\lambda\theta_T^m) = m\,\deg_t(\theta_T) = m d_1\) at every field point
  (Theorem~\ref{pa-thm:degAt_pic0_mul_pow}, Theorem~\ref{pa-thm:degAt_thetaFamily}). Backward,
  \(\mu \mapsto \mu \cdot (\theta_T^{m})^{-1}\) lands in degree zero by additivity
  (Lemma~\ref{pa-lem:degAt_hom}); the two maps are mutually inverse since the ambient group is
  commutative. Naturality is naturality of the \(\theta\)-family
  (Definition~\ref{pa-def:thetaFamily}): restriction is a homomorphism carrying
  \(\theta_T^{m}\) to \(\theta_{T'}^{m}\). Transport of a representation along a natural
  isomorphism of set-valued functors is formal.
\end{proof}

\begin{definition}[The pinned \(\Theta\)-class]
  \label{pa-def:thetaCechClass}
  \lean{AlgebraicGeometry.baseCurveObj, AlgebraicGeometry.thetaP1,
    AlgebraicGeometry.isFinite_thetaP1, AlgebraicGeometry.isDominant_thetaP1,
    AlgebraicGeometry.thetaCechClass, AlgebraicGeometry.one_le_classDeg_thetaCechClass}
  \leanok
  \dcref{custom}

  \uses{pa-thm:exists_finite_dominant_toP1, pa-thm:baseChangeCurve_bundle, pa-def:fiberTwist,
    pa-thm:classDeg_fiberTwist_one}
  The base-changed curve \(C_k = (C \otimes \Spec k)_{\mathrm{left}}\) is again a curve over
  \(k\) through the second projection (Theorem~\ref{pa-thm:baseChangeCurve_bundle}), so it
  carries a finite dominant morphism \(\pi \colon C_k \to \PP^1\)
  (Theorem~\ref{pa-thm:exists_finite_dominant_toP1}). The \emph{pinned \(\Theta\)-class} is the
  unit fiber twist
  \[
    \Theta \;:=\; \Theta_1(\pi) = \mathrm{fiberTwist}(\pi, 1) \in \Pic(C_k),
  \]
  and it satisfies \(1 \le \deg_{\Pic}(\Theta)\). It is the canonical choice of \(L_0\) for
  the \(\theta\)-family; its degree \(d_1 := \deg_{\Pic}(\Theta) \ge 1\) is the shift
  constant.
\end{definition}
\begin{proof}
  \leanok
  Positivity is the \(\Theta\)-positivity theorem
  (Theorem~\ref{pa-thm:classDeg_fiberTwist_one}) applied to the finite dominant \(\pi\) on the
  curve \(C_k\).
\end{proof}

\section{The Zariski sheaf property of the degree-zero functor on the slice}
\label{pa-sec:pic0ZariskiSheaf}

The functor \(\mathrm{picEt}\) satisfies the Zariski sheaf axioms for arbitrary open covers
of arbitrary tests — separation was established as Theorem~\ref{pa-thm:picEt_ext_of_le_cover};
here we add gluing, and the S-lemma making the degree-zero subfunctor a subsheaf.

\begin{definition}[Compatible local data on an open cover]
  \label{pa-def:picEt_LocalData}
  \lean{AlgebraicGeometry.picEt.LocalData, AlgebraicGeometry.picEt.IsGlueValue}
  \leanok
  \dcref{custom}

  \uses{pa-def:picEt, pa-def:PicEtAff_mapAlg}
  Let \(T\) be a test object and \((O_i)_{i \in \iota}\) an open cover of
  \(T_{\mathrm{left}}\). \emph{Compatible local data} on the cover is a family assigning to
  every index \(i\) and every affine open \(W \le O_i\) a plus class
  \(v_{i,W} \in \mathrm{PicEtAff}(C, \Gamma(T,W))\), such that
  (\emph{res}) within a member, \(v_{i,W}\) restricts to \(v_{i,W'}\) for \(W' \le W \le O_i\),
  and (\emph{glue}) across members, \(v_{i,W} = v_{j,W}\) whenever \(W \le O_i\) and
  \(W \le O_j\). A plus class \(z\) at an affine open \(W\) is a \emph{glue value} of the
  data when its restriction to every affine sub-open \(W_0 \le W\) contained in some member
  \(O_i\) equals \(v_{i,W_0}\).
\end{definition}

\begin{theorem}[Uniqueness of glue values]
  \label{pa-thm:picEt_glueValue_unique}
  \lean{AlgebraicGeometry.picEt.glueValue_unique}
  \leanok
  \dcref{custom}

  \uses{pa-def:picEt_LocalData, pa-thm:exists_basic_subcover, pa-thm:picEt_eq_of_basic_eq}
  If the \(O_i\) cover \(T_{\mathrm{left}}\), then any two glue values of the same data at
  the same affine open agree.
\end{theorem}
\begin{proof}
  \leanok
  Choose a finite family of basic opens of \(W\) covering \(W\), each contained in some
  member (Theorem~\ref{pa-thm:exists_basic_subcover}, with the subordination predicate
  ``contained in some \(O_i\)''). Two glue values restrict, on each such basic open, to the
  same local datum; classes agreeing on a spanning family of basic opens agree
  (Theorem~\ref{pa-thm:picEt_eq_of_basic_eq}).
\end{proof}

\begin{theorem}[Existence of glue values]
  \label{pa-thm:picEt_exists_isGlueValue}
  \lean{AlgebraicGeometry.picEt.exists_isGlueValue}
  \leanok
  \dcref{custom}

  \uses{pa-def:picEt_LocalData, pa-thm:PicEtAff_exists_mapAlg_eq_of_compat,
    pa-thm:exists_basic_subcover, pa-thm:picEt_eq_of_basic_eq, pa-def:PicEtAff_mapAlg}
  Compatible local data on an open cover of \(T_{\mathrm{left}}\) has a glue value at every
  affine open of \(T_{\mathrm{left}}\).
\end{theorem}
\begin{proof}
  \leanok
  Fix an affine open \(W\) and choose a finite basic-open refinement \((D(r))_{r \in s}\) of
  \(W\) subordinate to the cover, with \(\sum_{r \in s}(r) = (1)\)
  (Theorem~\ref{pa-thm:exists_basic_subcover}). The local data restricted to the \(D(r)\) is a
  compatible family in the sense of the affine Zariski gluing: on an overlap
  \(D(rr') = D(r) \cap D(r')\) the two restrictions agree by (\emph{res}) within each
  chosen member followed by (\emph{glue}) across the two members. By the affine gluing
  theorem (Theorem~\ref{pa-thm:PicEtAff_exists_mapAlg_eq_of_compat}) there is a plus class
  \(z\) at \(W\) restricting to the data on each \(D(r)\).

  \(z\) is a glue value: given an affine sub-open \(W_0 \le W\) with \(W_0 \le O_i\), refine
  \(W_0\) by basic opens of \(W_0\) each of the form \(D(q) \subseteq D(r)\) for some
  \(r \in s\) (Theorem~\ref{pa-thm:exists_basic_subcover} again). On each such \(D(q)\), the
  restriction of \(z\) computes through \(D(r)\) as the local datum at \(D(q)\) by
  (\emph{res}), and the local datum at \(D(q)\) read through member \(i\) agrees with that
  read through the member of \(r\) by (\emph{glue}). So the restrictions of \(z\) and of
  \(v_{i,W_0}\) to a spanning family of basic opens of \(W_0\) agree, and
  Theorem~\ref{pa-thm:picEt_eq_of_basic_eq} concludes.
\end{proof}

\begin{theorem}[The gluing half of the Zariski sheaf property of \(\mathrm{picEt}\)]
  \label{pa-thm:picEt_existsUnique_glue_of_le_cover}
  \lean{AlgebraicGeometry.picEt.glueValue, AlgebraicGeometry.picEt.glueValue_spec,
    AlgebraicGeometry.picEt.glueValue_eq_of, AlgebraicGeometry.picEt.glueSection,
    AlgebraicGeometry.picEt.glueSection_eq,
    AlgebraicGeometry.picEt.existsUnique_glue_of_le_cover}
  \leanok
  \dcref{custom}

  \uses{pa-thm:picEt_glueValue_unique, pa-thm:picEt_exists_isGlueValue, pa-thm:picEt_ext_of_le_cover,
    pa-def:picEt, pa-def:PicEtAff_mapAlg}
  Compatible local data on an open cover of \(T_{\mathrm{left}}\) glues to a \emph{unique}
  section \(s \in \mathrm{picEt}(C,T)\) whose value at every affine open subordinate to a
  member is the local datum there.
\end{theorem}
\begin{proof}
  \leanok
  Take \(s\) to be the family of glue values (existence:
  Theorem~\ref{pa-thm:picEt_exists_isGlueValue}; they are unique by
  Theorem~\ref{pa-thm:picEt_glueValue_unique}). The family is compatible under restriction of
  affine opens — the restriction of the glue value at \(V\) to \(U \le V\) is itself a glue
  value at \(U\), since glue-value restrictions compose — so \(s\) is a section of
  \(\mathrm{picEt}(C,T)\) (Definition~\ref{pa-def:picEt}). At an affine open \(W \le O_i\), the
  glue value restricted along \(W \le W\) is the local datum, so \(s\) extends the data.
  Uniqueness: two such sections agree at every affine open subordinate to a member, hence
  agree by separation (Theorem~\ref{pa-thm:picEt_ext_of_le_cover}).
\end{proof}

\begin{theorem}[Separation of the degree-zero subfunctor]
  \label{pa-thm:pic0Subgroup_ext_of_le_cover}
  \lean{AlgebraicGeometry.pic0Subgroup_ext_of_le_cover}
  \leanok
  \dcref{custom}

  \uses{pa-def:pic0Subgroup, pa-thm:picEt_ext_of_le_cover}
  Two degree-zero classes on \(T\) whose values agree at every affine open subordinate to an
  open cover of \(T_{\mathrm{left}}\) are equal.
\end{theorem}
\begin{proof}
  \leanok
  The underlying \(\mathrm{picEt}\) classes agree by separation
  (Theorem~\ref{pa-thm:picEt_ext_of_le_cover}); membership in the subgroup carries no further
  data.
\end{proof}

\begin{theorem}[The S-lemma: degree-zero membership is Zariski-local]
  \label{pa-thm:mem_pic0Subgroup_of_cover}
  \lean{AlgebraicGeometry.mem_pic0Subgroup_of_cover}
  \leanok
  \dcref{custom}

  \uses{pa-def:pic0Subgroup, pa-thm:degAt_picEtMap, pa-def:picEtMap}
  Let \((f_i \colon T_i \to T)\) be a family of test morphisms whose left legs are open
  immersions with jointly surjective ranges, and \(\lambda \in \mathrm{picEt}(C,T)\) a class
  whose restriction \(f_i^{*}\lambda\) is degree-zero for every \(i\). Then \(\lambda\) is
  degree-zero. This is what makes the degree-zero subfunctor a \emph{subsheaf}, not merely a
  subfunctor, of the Zariski sheaf \(\mathrm{picEt}\).
\end{theorem}
\begin{proof}
  \leanok
  Let \(t \colon \Spec K \to T\) be a field point. Its underlying space is a single point,
  so its topological image lies in the range of some \(f_i\); as that left leg is an open
  immersion, \(t\) factors through it: \(t = f_i \circ t'\) for a field point \(t'\) of
  \(T_i\) (lift along an open immersion containing the image). Then
  \[
    \deg_t(\lambda) = \deg_{t'}\bigl(f_i^{*}\lambda\bigr) = 0
  \]
  by functoriality of the degree (Theorem~\ref{pa-thm:degAt_picEtMap}) and the hypothesis. As
  \(t\) was arbitrary, \(\lambda \in \mathrm{pic}^0(C,T)\).
\end{proof}

\section{The morphism-form cover bridge}
\label{pa-sec:coverBridge}

The sheaf statements of the previous section speak \emph{topological} covers — opens of
\(T_{\mathrm{left}}\) and data at subordinate affine opens. The big-site packaging of the
next section speaks \emph{morphism-form} covers: families \(f_i \colon T_i \to T\) of test
morphisms with open-immersion left legs carrying sections. This section is the bridge.

\begin{lemma}[The section-ring isomorphism of an open immersion of tests]
  \label{pa-lem:appLEAlgEquiv}
  \lean{AlgebraicGeometry.Over.appLEAlgHom_congr_hom, AlgebraicGeometry.Over.bijective_appLEAlgHom,
    AlgebraicGeometry.Over.appLEAlgEquiv,
    AlgebraicGeometry.Over.appLEAlgHom_comp_appLEAlgEquiv_symm,
    AlgebraicGeometry.PicEtAff.mapAlg_injective}
  \leanok
  \dcref{custom}

  \uses{pa-def:PicEtAff_mapAlg}
  Let \(f \colon T' \to T\) be a test morphism whose left leg is an open immersion, and
  \(V \subseteq T_{\mathrm{left}}\) an open contained in its range. Then pullback of sections
  \(\Gamma(T,V) \to \Gamma(T', f^{-1}V)\) is a \(k\)-algebra \emph{isomorphism}. Moreover
  transport of plus classes along any bijective \(k\)-algebra map is injective.
\end{lemma}
\begin{proof}
  \leanok
  An open immersion restricts to an isomorphism from \(f^{-1}V\) onto \(V\) when
  \(V\) lies in the range, and an isomorphism of schemes induces an isomorphism on sections;
  the algebra structures match because \(f\) lies over \(\Spec k\). For the second claim,
  transport along the inverse algebra isomorphism is a two-sided inverse up to the functor
  laws of Definition~\ref{pa-def:PicEtAff_mapAlg}, so the transport is injective.
\end{proof}

\begin{definition}[The evaluation of a morphism-form family into local data]
  \label{pa-def:picEt_coverValue}
  \lean{AlgebraicGeometry.picEt.coverValue, AlgebraicGeometry.picEt.mapAlg_coverValue,
    AlgebraicGeometry.picEt.coverValue_res, AlgebraicGeometry.picEt.overlapTest,
    AlgebraicGeometry.picEt.overlapInclusion, AlgebraicGeometry.picEt.coverValue_glue}
  \leanok
  \dcref{custom}

  \uses{pa-lem:appLEAlgEquiv, pa-def:picEt_LocalData, pa-def:picEt, pa-def:picEtMap,
    pa-def:PicEtAff_mapAlg, pa-thm:picEt_mapAlg_appLE_eq}
  Let \((f_i \colon T_i \to T)\) have open-immersion left legs and let
  \(x_i \in \mathrm{picEt}(C, T_i)\) be a family of sections, \emph{compatible} in the
  morphism form: for every pair \(i, j\), every test \(Z\) and every pair of morphisms
  \(g_i \colon Z \to T_i\), \(g_j \colon Z \to T_j\) with \(f_i g_i = f_j g_j\), the two
  restrictions of the family to \(Z\) agree. The \emph{cover value} of the family at an
  affine open \(W\) contained in the range of \(f_i\) is the value of \(x_i\) at the
  preimage affine open \(f_i^{-1}W\), transported back along the section-ring isomorphism
  (Lemma~\ref{pa-lem:appLEAlgEquiv}). The cover values are compatible local data
  (Definition~\ref{pa-def:picEt_LocalData}) on the ranges \(O_i := \mathrm{range}(f_i)\).
\end{definition}
\begin{proof}
  \leanok
  (\emph{res}) Restriction within a member commutes with the section-ring transports and
  with the internal compatibility of the section \(x_i\), by the restriction calculus of the
  vehicle (Theorem~\ref{pa-thm:picEt_mapAlg_appLE_eq}) and injectivity of transport along the
  bijective pullback (Lemma~\ref{pa-lem:appLEAlgEquiv}).

  (\emph{glue}) Fix a shared affine open \(W \le O_i \cap O_j\) and form the mediating test:
  the open subscheme of \(T\) at \(W\), over \(\Spec k\) through the inclusion. The
  inclusion lifts through both members (an open immersion whose range contains \(W\) admits
  a lift of the inclusion of \(W\)), producing \(g_i, g_j\) with
  \(f_i g_i = f_j g_j = (\text{inclusion of } W)\). The compatibility hypothesis identifies
  the two restricted sections on the mediating test; evaluating both at its top affine open
  and cancelling the common (bijective) evaluation map
  (Lemma~\ref{pa-lem:appLEAlgEquiv}) identifies the two cover values.
\end{proof}

\begin{theorem}[The sheaf property of \(\mathrm{picEt}\) in morphism-cover form]
  \label{pa-thm:picEt_existsUnique_of_cover}
  \lean{AlgebraicGeometry.picEt.ext_of_cover, AlgebraicGeometry.picEt.existsUnique_of_cover}
  \leanok
  \dcref{custom}

  \uses{pa-def:picEt_coverValue, pa-thm:picEt_existsUnique_glue_of_le_cover,
    pa-thm:picEt_ext_of_le_cover, pa-thm:picEt_glueValue_unique, pa-lem:appLEAlgEquiv}
  Let \((f_i \colon T_i \to T)\) have open-immersion left legs with jointly surjective
  ranges. Then two sections of \(\mathrm{picEt}(C,T)\) agreeing after restriction to every
  member agree; and every compatible family \((x_i)\) as in
  Definition~\ref{pa-def:picEt_coverValue} glues to a unique section \(s\) with
  \(f_i^{*} s = x_i\) for all \(i\).
\end{theorem}
\begin{proof}
  \leanok
  \emph{Separation.} At an affine open subordinate to a member the values of the two
  sections compute through the member's restriction, which agree; conclude by topological
  separation (Theorem~\ref{pa-thm:picEt_ext_of_le_cover}) after cancelling the bijective
  transports (Lemma~\ref{pa-lem:appLEAlgEquiv}).

  \emph{Gluing.} The cover values of the family are compatible local data on the ranges
  (Definition~\ref{pa-def:picEt_coverValue}), which glue to a unique section \(s\)
  (Theorem~\ref{pa-thm:picEt_existsUnique_glue_of_le_cover}). That \(f_i^{*} s = x_i\) is
  checked at every affine open \(U'\) of \(T_i\): the value of \(f_i^{*}s\) there computes
  from the value of \(s\) at the image affine open \(f_i(U')\), which is the cover value,
  i.e.\ the value of \(x_i\) at \(f_i^{-1}f_i(U') \supseteq U'\), and the internal
  compatibility of \(x_i\) restricts it to \(U'\). Uniqueness follows from separation.
\end{proof}

\begin{theorem}[The sheaf property of the degree-zero functor in morphism-cover form]
  \label{pa-thm:pic0Subgroup_existsUnique_of_cover}
  \lean{AlgebraicGeometry.pic0Subgroup_ext_of_cover,
    AlgebraicGeometry.pic0Subgroup_existsUnique_of_cover}
  \leanok
  \dcref{custom}

  \uses{pa-thm:picEt_existsUnique_of_cover, pa-thm:mem_pic0Subgroup_of_cover, pa-def:pic0Functor,
    pa-def:pic0Subgroup}
  In the situation of Theorem~\ref{pa-thm:picEt_existsUnique_of_cover}, a compatible family of
  \emph{degree-zero} classes \(x_i \in \mathrm{pic}^0(C, T_i)\) glues to a unique degree-zero
  class \(s \in \mathrm{pic}^0(C,T)\) with \(f_i^{*} s = x_i\); and two degree-zero classes
  agreeing on every member agree.
\end{theorem}
\begin{proof}
  \leanok
  Glue the underlying \(\mathrm{picEt}\) classes
  (Theorem~\ref{pa-thm:picEt_existsUnique_of_cover}); the glued class is degree-zero by the
  S-lemma (Theorem~\ref{pa-thm:mem_pic0Subgroup_of_cover}), since each restriction is. Both
  uniqueness claims reduce to the underlying classes.
\end{proof}

\section{The \texorpdfstring{\(\Sigma\)}{Sigma}-extension and the big-site sheaf}
\label{pa-sec:sigmaExtension}

A presheaf on the slice \(\Over(S)\) extends to a presheaf on the ambient category by
summing over structure morphisms. For \(S = \Spec k\) and the degree-zero Picard functor
this \emph{slice trick} produces the big-site carrier consumed by the local-representability
engine, and a representation of the extension descends to one of the original functor.

\begin{definition}[The canonical slice morphism of an equality]
  \label{pa-def:mkCongr}
  \lean{CategoryTheory.Over.mkCongr, CategoryTheory.Over.mkCongr_left,
    CategoryTheory.Over.mkCongr_rfl, CategoryTheory.Over.mkCongr_comp_mkCongr,
    CategoryTheory.Over.mkCongr_comp_homMk}
  \leanok
  \dcref{custom}

  Let \(\mathcal C\) be a category, \(S\) an object, and \(a = b\) two equal morphisms
  \(T \to S\). The identity of \(T\) is a morphism of slice objects
  \((T, a) \to (T, b)\) in \(\Over(S)\), the \emph{canonical morphism of the equality}. It
  satisfies: the canonical morphism of a reflexive equality is the identity; canonical
  morphisms compose to the canonical morphism of the transitive equality; and every slice
  morphism \(g \colon X \to (T, b)\) factors as the canonical morphism of its own
  compatibility equation followed by the slice morphism induced by its underlying map.
\end{definition}
\begin{proof}
  \leanok
  The identity of \(T\) satisfies \(\id_T \circ b = a\) precisely because \(a = b\); the
  three laws are immediate from composition of identities.
\end{proof}

\begin{definition}[The \(\Sigma\)-extension of a slice presheaf]
  \label{pa-def:sigmaExtension}
  \lean{CategoryTheory.Over.sigmaExtensionMap, CategoryTheory.Over.sigmaExtension,
    CategoryTheory.Over.sigmaExtension_obj, CategoryTheory.Over.sigmaExtension_map_mk,
    CategoryTheory.Over.sigmaExtension_map_fst, CategoryTheory.Over.sigmaExtension_ext,
    CategoryTheory.Over.sigmaExtension_snd_eq}
  \leanok
  \dcref{custom}

  \uses{pa-def:mkCongr}
  For a presheaf \(F\) on the slice \(\Over(S)\), the \emph{\(\Sigma\)-extension}
  \(\widetilde F\) is the presheaf on the ambient category with
  \[
    \widetilde F(T) \;:=\; \coprod_{a \colon T \to S} F(T, a) ,
  \]
  an element being a pair \(\langle a, \xi\rangle\) of a structure morphism and a fibre
  element. Restriction along \(g \colon T' \to T\) composes on the structure component and
  applies \(F\) of the induced slice morphism on the fibre:
  \(g^{*}\langle a, \xi \rangle = \langle g \circ a,\ F(g)^{*}\xi\rangle\). Two
  \(\Sigma\)-elements over equal structure morphisms are equal iff their fibres correspond
  under \(F\) of the canonical morphism of the equality (Definition~\ref{pa-def:mkCongr}).
\end{definition}
\begin{proof}
  \leanok
  Functoriality: restriction along the identity composes with \(a\) to \(a\), and the
  induced slice morphism is the identity, so \(F\) fixes the fibre. For a composite, the two
  structure components agree by associativity, and the two fibre components correspond under
  \(F\) of the canonical morphism, since the induced slice morphisms compose accordingly —
  the transport calculus of Definition~\ref{pa-def:mkCongr} mediates without any cast. The
  final characterisation of equality of \(\Sigma\)-elements is the elimination rule of the
  coproduct.
\end{proof}

\begin{lemma}[The amalgamation seam and the compatibility descent]
  \label{pa-lem:sigmaExtension_seams}
  \lean{CategoryTheory.Over.sigmaExtension_map_mk_eq_iff,
    CategoryTheory.Over.sigmaExtension_map_left_mk,
    CategoryTheory.Over.map_eq_of_sigmaExtension_map_eq}
  \leanok
  \dcref{custom}

  \uses{pa-def:sigmaExtension, pa-def:mkCongr}
  Let \(F\) be a presheaf on \(\Over(S)\) and \(\widetilde F\) its \(\Sigma\)-extension.
  \begin{enumerate}
  \item (\emph{Amalgamation seam}) For \(g \colon T' \to T\), \(a \colon T \to S\),
    \(b \colon T' \to S\) with \(g \circ a = b\), and elements \(s \in F(T,a)\),
    \(y \in F(T',b)\):
    \(g^{*}\langle a, s\rangle = \langle b, y\rangle\) in \(\widetilde F(T')\) iff
    \(s\) restricts to \(y\) along the slice morphism \((T',b) \to (T,a)\) induced by \(g\).
  \item (\emph{Left-leg rule}) For a slice morphism \(g \colon Z \to (T,a)\) and
    \(x \in F(T,a)\): restriction of \(\langle a, x\rangle\) along the underlying map of
    \(g\) is \(\langle Z_{\mathrm{struct}},\, F(g)^{*}x \rangle\), the honest slice
    restriction on the fibre.
  \item (\emph{Compatibility descent}) If two \(\Sigma\)-elements
    \(\langle a, x\rangle, \langle b, y\rangle\) have equal restrictions along the
    underlying maps of a pair of slice morphisms \(g_a \colon Z \to (T,a)\),
    \(g_b \colon Z' = Z \to (T',b)\), then \(F(g_a)^{*}x = F(g_b)^{*}y\).
  \end{enumerate}
\end{lemma}
\begin{proof}
  \leanok
  (i) After substituting the equality \(g \circ a = b\), both sides live over the same
  structure morphism and the characterisation of equal \(\Sigma\)-elements
  (Definition~\ref{pa-def:sigmaExtension}) reduces the claim to the fibres, where the canonical
  morphism of the reflexive equality acts as the identity. (ii) is (i) applied to the
  compatibility equation of \(g\), using that the induced slice morphism of the left leg
  \emph{is} \(g\) (Definition~\ref{pa-def:mkCongr}, factorisation law). (iii) Combine (ii) on
  both sides: both restrictions are \(\Sigma\)-elements over the structure morphism of
  \(Z\), and equality of \(\Sigma\)-elements over one structure morphism gives equality of
  fibres.
\end{proof}

\begin{theorem}[The \(\Sigma\)-descent of a representation]
  \label{pa-thm:representableBy_overSlice}
  \lean{CategoryTheory.Functor.RepresentableBy.overSlice}
  \leanok
  \dcref{custom}

  \uses{pa-def:sigmaExtension, pa-lem:sigmaExtension_seams}
  Let \(F\) be a presheaf on \(\Over(S)\) and suppose the \(\Sigma\)-extension
  \(\widetilde F\) is representable by an object \(J\) of the ambient category, with
  universal element \(\langle u_1, u_2 \rangle \in \widetilde F(J)\) (the image of
  \(\id_J\)). Then \(F\) is representable on the slice by \((J, u_1)\), with universal
  correspondence sending a slice morphism \(g \colon X \to (J,u_1)\) to the slice
  restriction \(F(g)^{*} u_2\).
\end{theorem}
\begin{proof}
  \leanok
  \emph{Injectivity.} If \(F(g)^{*}u_2 = F(g')^{*}u_2\) then, by the left-leg rule
  (Lemma~\ref{pa-lem:sigmaExtension_seams}(ii)), the restrictions of the universal element
  along the underlying maps of \(g\) and \(g'\) are the same \(\Sigma\)-element of
  \(\widetilde F(X_{\mathrm{left}})\); by representability of \(\widetilde F\) the
  underlying maps agree, hence \(g = g'\) (slice morphisms are determined by their
  underlying maps).

  \emph{Surjectivity.} Given \(\xi \in F(X)\), the \(\Sigma\)-element
  \(\langle X_{\mathrm{struct}}, \xi\rangle\) corresponds under the representation to a map
  \(h \colon X_{\mathrm{left}} \to J\) with
  \(h^{*}\langle u_1, u_2\rangle = \langle X_{\mathrm{struct}}, \xi\rangle\). Comparing
  structure components, \(h \circ u_1 = X_{\mathrm{struct}}\), so \(h\) underlies a slice
  morphism \(g \colon X \to (J, u_1)\); comparing fibres through the amalgamation seam
  (Lemma~\ref{pa-lem:sigmaExtension_seams}(i)) gives \(F(g)^{*}u_2 = \xi\).

  Compatibility with composition is functoriality of \(F\).
\end{proof}

\begin{definition}[The big-site carrier of the degree-zero functor]
  \label{pa-def:pic0SigmaFunctor}
  \lean{AlgebraicGeometry.pic0TypeFunctor, AlgebraicGeometry.pic0TypeFunctor_obj,
    AlgebraicGeometry.pic0TypeFunctor_map_apply, AlgebraicGeometry.pic0SigmaFunctor}
  \leanok
  \dcref{custom}

  \uses{pa-def:pic0Functor, pa-def:sigmaExtension}
  Write \(\mathrm{pic}^0_{\mathrm{Set}}(C)\) for the set-valued functor underlying the
  degree-zero relative Picard functor (Definition~\ref{pa-def:pic0Functor}). Its
  \(\Sigma\)-extension along the slice \(\Over(\Spec k)\),
  \[
    T \;\longmapsto\; \coprod_{a \colon T \to \Spec k} \mathrm{pic}^0\bigl(C, (T,a)\bigr),
    \qquad T \in \mathrm{Sch},
  \]
  is a presheaf of sets on the category of schemes — the big-site carrier of the slice
  trick.
\end{definition}

\begin{theorem}[The big-site Zariski sheaf certificate]
  \label{pa-thm:pic0SigmaFunctor_isSheaf}
  \lean{AlgebraicGeometry.pic0SigmaFunctor_isSheaf, AlgebraicGeometry.pic0SigmaSheaf}
  \leanok
  \dcref{custom}

  \uses{pa-def:pic0SigmaFunctor, pa-thm:pic0Subgroup_existsUnique_of_cover,
    pa-lem:sigmaExtension_seams, pa-def:sigmaExtension}
  The \(\Sigma\)-extension of the degree-zero Picard functor is a sheaf for the big Zariski
  topology on schemes.
\end{theorem}
\begin{proof}
  \leanok
  The Zariski pretopology is generated by open covers and is stable under base change, so
  the sheaf condition reduces to elementwise unique amalgamation over an honest open cover
  \((\iota_i \colon U_i \to Y)\) carrying a compatible family
  \(x_i = \langle a_i, \xi_i \rangle\) of \(\Sigma\)-elements.

  \emph{Structure components.} Compatibility on pullbacks forces the \(a_i\) to agree on
  overlaps, so they glue to a unique \(a_Y \colon Y \to \Spec k\) with
  \(\iota_i \circ a_Y = a_i\) (gluing of scheme morphisms along an open cover).

  \emph{Fibre components.} The cover members induce slice morphisms
  \((U_i, a_i) \to (Y, a_Y)\) with open-immersion left legs and jointly surjective ranges.
  The fibres \(\xi_i\) are a compatible family on this slice cover: by the compatibility
  descent (Lemma~\ref{pa-lem:sigmaExtension_seams}(iii)), the \(\Sigma\)-level compatibility of
  the \(x_i\) descends to slice-level compatibility of the \(\xi_i\). By the morphism-form
  sheaf property of the degree-zero functor
  (Theorem~\ref{pa-thm:pic0Subgroup_existsUnique_of_cover}) the \(\xi_i\) glue to a unique
  degree-zero class \(s\) on \((Y, a_Y)\) restricting to each.

  \emph{Stitching.} The \(\Sigma\)-element \(\langle a_Y, s\rangle\) restricts to each
  \(x_i\) by the amalgamation seam (Lemma~\ref{pa-lem:sigmaExtension_seams}(i)). Uniqueness:
  the structure component of any amalgamation is determined by gluing of morphisms, and its
  fibre by the separation half of
  Theorem~\ref{pa-thm:pic0Subgroup_existsUnique_of_cover}, again through the seam.
\end{proof}

\begin{definition}[Local representability descends to the slice]
  \label{pa-def:pic0RepresentableByOfCharts}
  \lean{AlgebraicGeometry.pic0RepresentableByOfCharts}
  \leanok
  \dcref{custom}

  \uses{pa-thm:pic0SigmaFunctor_isSheaf, pa-thm:representableBy_overSlice, pa-def:pic0SigmaFunctor}
  Suppose given a family of schemes \(X_i\) with maps
  \(f_i \colon \mathrm{yoneda}(X_i) \to \widetilde{\mathrm{pic}^0}(C)\) of big-site sheaves
  which are relatively representable open immersions, jointly Zariski-locally surjective.
  Then the local-representability theorem glues the \(X_i\) into a scheme \(J\)
  representing the \(\Sigma\)-extension, and by \(\Sigma\)-descent
  (Theorem~\ref{pa-thm:representableBy_overSlice}) the set-valued degree-zero Picard functor is
  representable on the slice by \(J\) placed over \(\Spec k\) via the structure component of
  the universal element.
\end{definition}
\begin{proof}
  \leanok
  The local-representability theorem for Zariski sheaves on schemes (mathlib's
  01JJ engine) applies to the sheaf of Theorem~\ref{pa-thm:pic0SigmaFunctor_isSheaf} and the
  given chart family, producing a representation of the \(\Sigma\)-extension by the glued
  scheme; compose with Theorem~\ref{pa-thm:representableBy_overSlice}.
\end{proof}

This remains a valid conditional local-representability theorem. Spelled out, its hypothesis
asks for a family of schemes \(X_i\) together
with maps \(\mathrm{yoneda}(X_i) \to \widetilde{\mathrm{pic}^0}(C)\) such that

\begin{itemize}
  \item each \(f_i\) is \emph{relatively representable by an open immersion}: for every
    scheme \(T\) and every \(\Sigma\)-element \(\xi \in \widetilde{\mathrm{pic}^0}(C)(T)\),
    the fibre product \(\mathrm{yoneda}(X_i) \times_{\widetilde{\mathrm{pic}^0}(C)} T\) is
    representable by an open subscheme of \(T\) --- so each \(f_i\) exhibits \(X_i\) as an
    \emph{open piece} of the functor, with the openness verified at every test, not merely
    at points;
  \item the family is \emph{jointly Zariski-locally surjective}: every element of
    \(\widetilde{\mathrm{pic}^0}(C)(T)\) is, after restriction to the members of some open
    cover of \(T\), in the image of some \(f_i\).
\end{itemize}

These hypotheses are not supplied by the raw admissible Abel map.  Its fibres are complete
linear systems and can have positive dimension, so that map is neither an open immersion nor
an open atlas.  Section~\ref{pa-sec:rankOneAbelDescent} instead restricts the degree-\(g(C)\) Abel
map to the locus where \(H^1\) vanishes and the pushforward has rank one.  The canonical
evaluation divisor identifies that divisor open with the corresponding Picard open.  Over a
separably closed field, its translates cover every field-valued Picard class and therefore
supply the chart family above; finite Galois descent then returns the arbitrary-field
representer together with its representation certificate.  The larger admissible divisor
source remains available for arbitrary-degree divisor geometry, but its positive-dimensional
linear systems do not enter this open-atlas construction.

\section{The Abel element}
\label{pa-sec:abelElement}

For a rational point \(P\) of the curve, the class of the diagonal minus the constant graph
is a degree-zero family on the test \(T = C\) — the element whose image under the universal
correspondence of the representation is the Abel map.

\begin{theorem}[The unit seam at arbitrary tests]
  \label{pa-thm:degAt_relPicToPicEt}
  \lean{AlgebraicGeometry.degAt_relPicToPicEt}
  \leanok
  \dcref{custom}

  \uses{pa-def:degAt, pa-thm:picEtMap_relPicToPicEt, pa-thm:picEtAffineEquiv_relPicToPicEt,
    pa-thm:degAff_unit, pa-def:relPicToPicEt, pa-def:relPicMap}
  For a test object \(T\), a relative Picard class \(z \in \mathrm{relPic}(C,T)\) and a
  field point \(t \colon \Spec K \to T\), the degree at \(t\) of the unit image of \(z\) is
  the relative degree of the restriction:
  \[
    \deg_t\bigl(\mathrm{relPicToPicEt}(z)\bigr) = \mathrm{relPicDeg}_K\bigl(t^{*} z\bigr).
  \]
\end{theorem}
\begin{proof}
  \leanok
  Naturality of the unit (Theorem~\ref{pa-thm:picEtMap_relPicToPicEt}) collapses the
  restriction along \(t\) to the unit image of \(t^{*}z\) over \(\Spec K\); affine
  consistency (Theorem~\ref{pa-thm:picEtAffineEquiv_relPicToPicEt}) identifies its affine
  collapse with the plus unit of \(t^{*}z\); the unit normalization of the degree
  (Theorem~\ref{pa-thm:degAff_unit}) reads \(\mathrm{relPicDeg}_K\).
\end{proof}

\begin{definition}[The Abel \v Cech class]
  \label{pa-def:abelCechClass}
  \lean{AlgebraicGeometry.abelCechClass, AlgebraicGeometry.abelPicEt}
  \leanok
  \source{kleiman-picard, abelian-varieties:page-0107}

  \uses{pa-def:graphLocalEquations, pa-def:relPicToPicEt, pa-def:relPic}
  \dcref{custom}

  Let \(P \colon \mathbf{1} \to C\) be a rational point of the curve (a morphism from the
  unit test). The \emph{Abel \v Cech class} on the test \(T = C\) is
  \[
    \mathcal A_P \;:=\; \struct{C \otimes C}(\Delta) \cdot \struct{C \otimes C}(P \times C)^{-1}
      \;\in\; \Pic(C \otimes C),
  \]
  the graph class of the identity \(\id_C\) (the diagonal) times the inverse graph class of
  the constant point \(C \to \mathbf 1 \xrightarrow{P} C\) — both factors through the graph
  divisor of Definition~\ref{pa-def:graphLocalEquations}, so one degree certificate serves
  both. The \emph{Abel class} \(\alpha_P \in \mathrm{picEt}(C, C)\) is the unit image of its
  relative class. Milne's cited section assumes \(g(C)>0\); this graph-class construction
  uses only the rational point \(P\) and therefore also applies in genus zero.
\end{definition}

\begin{theorem}[The functor-point formula, \v Cech level]
  \label{pa-thm:cechPicMap_abelCechClass}
  \lean{AlgebraicGeometry.cechPicMap_abelCechClass}
  \leanok
  \dcref{custom}

  \uses{pa-def:abelCechClass, pa-thm:graphLocalEquations_base_change}
  For any test \(T\) and any point \(f \colon T \to C\), the pullback of the Abel \v Cech
  class along \(C \lhd f\) is the graph-minus-constant-graph class of \(f\):
  \[
    (C \lhd f)^{*} \mathcal A_P
      = \struct{C \otimes T}(\Gamma_f)\cdot\struct{C \otimes T}(\Gamma_{P_T})^{-1},
  \]
  where \(P_T \colon T \to \mathbf 1 \xrightarrow{P} C\) is the constant point of \(T\).
\end{theorem}
\begin{proof}
  \leanok
  Base change of the graph class (Theorem~\ref{pa-thm:graphLocalEquations_base_change}) sends
  \(\struct{}(\Gamma_{\id_C})\) to \(\struct{}(\Gamma_{f \circ \id}) = \struct{}(\Gamma_f)\)
  and \(\struct{}(\Gamma_{P_C})\) to \(\struct{}(\Gamma_{f \circ P_C})\); and
  \(f \circ P_C = P_T\) by terminality of the unit test (both factor through
  \(\mathbf 1\)). Pullback is a homomorphism, so the product formula follows.
\end{proof}

\begin{theorem}[The degree-zero certificate of the Abel class]
  \label{pa-thm:degAt_abelPicEt}
  \lean{AlgebraicGeometry.degAt_abelPicEt}
  \leanok
  \dcref{custom}

  \uses{pa-thm:degAt_relPicToPicEt, pa-thm:cechPicMap_abelCechClass, pa-thm:classDeg_graphPicClass,
    pa-def:relPicMap, pa-def:relPicDeg, pa-lem:degAt_hom}
  At every field point \(t \colon \Spec K \to C\), the Abel class has degree zero:
  \(\deg_t(\alpha_P) = 0\). The certificate is a proved equality at every point — no
  choice is involved.
\end{theorem}
\begin{proof}
  \leanok
  By the unit seam (Theorem~\ref{pa-thm:degAt_relPicToPicEt}) the degree is
  \(\mathrm{relPicDeg}_K\) of the restriction of the relative Abel class along \(t\), which
  the relative degree reads as \(\deg_{\Pic}\) of the fibre \v Cech class
  (Definition~\ref{pa-def:relPicDeg} on the projected representative,
  Definition~\ref{pa-def:relPicMap}). By the functor-point formula
  (Theorem~\ref{pa-thm:cechPicMap_abelCechClass}) at \(f := t\), the fibre class is
  \(\struct{}(\Gamma_t) \cdot \struct{}(\Gamma_{P_{\Spec K}})^{-1}\) — the ratio of two
  graph classes of \(K\)-points. Each has degree one by the degree-one certificate
  (Theorem~\ref{pa-thm:classDeg_graphPicClass}), so the degree is \(1 - 1 = 0\).
\end{proof}

\begin{definition}[The Abel element]
  \label{pa-def:abelElement}
  \lean{AlgebraicGeometry.abelElement, AlgebraicGeometry.abelElement_coe}
  \leanok
  \source{kleiman-picard, abelian-varieties:page-0107}

  \uses{pa-def:pic0Subgroup, pa-thm:degAt_abelPicEt, pa-def:abelCechClass}
  \dcref{custom}

  The \emph{Abel element} is the Abel class as an element of the degree-zero subgroup,
  \[
    \alpha_P \in \mathrm{pic}^0(C, C),
  \]
  membership being the certificate of Theorem~\ref{pa-thm:degAt_abelPicEt}. On points it is the
  family \(t \mapsto [\,\Gamma_t - \Gamma_P\,]\), the divisor-class form \(t \mapsto [t - P]\)
  of the Abel map with base point \(P\). As above, Milne states this formula under the
  standing hypothesis \(g(C)>0\), whereas the formal element is defined in genus zero as
  well.
\end{definition}

\begin{theorem}[The functor-point formula of the Abel element]
  \label{pa-thm:abelElement_map}
  \lean{AlgebraicGeometry.abelPicEt_map, AlgebraicGeometry.abelElement_map}
  \leanok
  \dcref{custom}

  \uses{pa-def:abelElement, pa-thm:picEtMap_relPicToPicEt, pa-thm:cechPicMap_abelCechClass,
    pa-def:pic0Functor, pa-def:relPicMap}
  For any point \(f \colon T \to C\), the restriction of the Abel element along \(f\) is the
  unit image of the graph-minus-constant-graph class of \(f\):
  \[
    f^{*}\alpha_P = \mathrm{relPicToPicEt}\Bigl(\pi\bigl(
      \struct{C \otimes T}(\Gamma_f)\cdot\struct{C\otimes T}(\Gamma_{P_T})^{-1}\bigr)\Bigr)
    \quad\text{in } \mathrm{pic}^0(C,T).
  \]
\end{theorem}
\begin{proof}
  \leanok
  Naturality of the unit (Theorem~\ref{pa-thm:picEtMap_relPicToPicEt}) commutes the restriction
  into the relative class, where it is the projection of the pullback of the Abel \v Cech
  class (Definition~\ref{pa-def:relPicMap}); conclude by the \v Cech-level functor-point
  formula (Theorem~\ref{pa-thm:cechPicMap_abelCechClass}).
\end{proof}

\begin{theorem}[The pointing law]
  \label{pa-thm:abelElement_map_point}
  \lean{AlgebraicGeometry.abelElement_map_point}
  \leanok
  \dcref{custom}

  \uses{pa-thm:abelElement_map, pa-def:pic0Functor}
  Restriction of the Abel element along the rational point \(P\) itself is trivial:
  \[
    P^{*} \alpha_P = 1 \quad\text{in } \mathrm{pic}^0(C, \mathbf 1).
  \]
\end{theorem}
\begin{proof}
  \leanok
  By the functor-point formula (Theorem~\ref{pa-thm:abelElement_map}) at \(f := P\), the
  restricted class is the unit image of
  \(\struct{}(\Gamma_P)\cdot\struct{}(\Gamma_{P_{\mathbf 1}})^{-1}\) on
  \(C \otimes \mathbf 1\). The constant point of the unit test is \(P\) itself
  (terminality: the canonical map \(\mathbf 1 \to \mathbf 1\) is the identity), so the two
  graph factors coincide and the \v Cech class is already \(1\); its unit image is \(1\).
\end{proof}
